\documentclass[aps,prx,onecolumn,10pt,notitlepage,nofootinbib,longbibliography]{revtex4-2}
\usepackage[T1]{fontenc}
\usepackage[utf8]{inputenc}
\usepackage{lmodern}
\usepackage{amsmath,amssymb,mathtools,bm,mathrsfs}
\usepackage{array,booktabs,longtable,multirow}
\makeatletter
\def\LT@bot{\nobreak\copy\LT@foot\vss}
\makeatother
\usepackage{microtype}
\usepackage{xurl}
\usepackage[hidelinks]{hyperref}
\pdfstringdefDisableCommands{\def\\{ }}
\hypersetup{pdftitle={Supplemental Material for Physics Without an Outside Observer: A Finite Geometry, Broken and Completed},pdfauthor={Yongzi CHEN},pdfsubject={Finite comparison, continuing histories, physical sources, calculations and proofs}}
\setlength{\textwidth}{6.95in}
\setlength{\textheight}{9.1in}
\setlength{\oddsidemargin}{-.225in}
\setlength{\evensidemargin}{-.225in}
\setlength{\topmargin}{-.45in}
\setlength{\parskip}{2pt}
\setlength{\parindent}{1.2em}
\setlength{\emergencystretch}{3em}
\allowdisplaybreaks[2]
\renewcommand{\thesection}{S\arabic{section}}
\numberwithin{equation}{section}
\numberwithin{table}{section}
\newcommand{\Id}{\mathbb I}
\newcommand{\CC}{\mathbb C}
\newcommand{\RR}{\mathbb R}
\newcommand{\EE}{\mathbb E}
\newcommand{\Tr}{\operatorname{Tr}}
\newcommand{\tr}{\operatorname{tr}}
\newcommand{\Var}{\operatorname{Var}}
\newcommand{\Cov}{\operatorname{Cov}}
\newcommand{\diag}{\operatorname{diag}}
\newcommand{\rank}{\operatorname{rank}}
\newcommand{\sgn}{\operatorname{sgn}}
\newcommand{\dd}{\mathrm d}
\newcommand{\ii}{\mathrm i}
\newcommand{\MeV}{\,\mathrm{MeV}}
\newcommand{\GeV}{\,\mathrm{GeV}}
\newcommand{\eV}{\,\mathrm{eV}}
\newcommand{\fm}{\,\mathrm{fm}}
\newcommand{\code}[1]{\texttt{#1}}
\newcommand{\MR}{M_R}
\newcommand{\MP}{M_{\mathrm{Pl}}}
\newcommand{\lamP}{\lambda_p}
\newcommand{\Fret}{F_{\mathrm{ret}}}
\newcommand{\hover}{h_\nu}
\newcommand{\Phis}{\Phi_{\rm ret}}
\newcommand{\trn}{\operatorname{tr}_0}
\begin{document}
\title{Supplemental Material for\\
Physics Without an Outside Observer:\\
A Finite Geometry, Broken and Completed}
\author{Yongzi CHEN}
\email{23103646d@connect.polyu.hk}
\affiliation{Department of Physics and Materials, The Hong Kong Polytechnic University, Hung Hom, Kowloon, Hong Kong, China}
\date[]{}
\begin{abstract}
The finite comparison $L$ and the remembered step $B$ have one determinant law. Part I derives the geometry and follows the eight physical steps in the order of the article. Part II expresses their physical quantities in the one spectral variable. Part III gives the current choices, several-line sources and conventional evolution, with their inputs. Part IV gives the formula and comparison registers and the routes to the programs and proofs.
\end{abstract}
\maketitle
\raggedbottom
\begingroup\small\tableofcontents\endgroup
\clearpage
\part{Detailed derivations}\label{part:derivations}

\section{The method, the postulates and the geometry}
\label{sec:geometry}
Describe nothing that cannot be observed; regard only what is universal as law. These two rules give the equal standing of internal comparisons and the existence of a particular continuing history. The five postulates express those requirements.

\subsection{The five postulates and their mathematical expression}
An elementary comparison includes its participants, their local frames and the history that brings them together.

\paragraph*{P1. There is no outside observer.}
Every participant, including the observer, is represented within the comparison. In the mathematical formulation used here, an elementary system has a finite complex Hilbert space $H$; two systems have $H\otimes H$. Pure states are rays, mixed states are positive operators $\rho$ with $\Tr\rho=1$, and a complete set of outcomes consists of positive effects $F_r$ with $\sum_rF_r=\Id$. The outcome probability is $\Tr(\rho F_r)$. Complex amplitudes, tensor composition, the adjoint, and this probability pairing are the mathematical representation of P1.

Let $\Theta={\mathsf J_\Theta}{\mathsf C_{\rm conj}}$ be the antiunitary pairing of the two roles, with ${\mathsf C_{\rm conj}}$ complex conjugation and ${\mathsf J_\Theta}$ its unitary linear part.

\paragraph*{P2. All observers inside have equal standing.}
The comparison is invariant under exchange of its participants, a common unitary frame change preserving the structure of P3, and symplectic transposition on either operator factor. Its linear involution on operators is
\begin{equation}
 X^\sharp=\Theta X^\dagger\Theta^{-1}={\mathsf J_\Theta}X^{\mathsf T}{\mathsf J_\Theta}^{-1}.
 \label{eq:symplectictranspose}
\end{equation}
One's own earlier state is a participant in this sense.

\paragraph*{P3. Observing and being observed have equal standing.}
Each side holds half of a distinction. The state conjugation is antiunitary, $\Theta^2=-\Id$. The whole nonzero real space of distinctions is
\begin{equation}
 \mathscr Q=\{a=a^\dagger:\Tr a=0,\ a\Theta=\Theta a\},\qquad
 a^2=\trn(a^2)\Id\quad(a\in\mathscr Q),\qquad
 \trn=\frac{\Tr}{\dim_{\CC}H}.
 \label{eq:fullspacepostulate}
\end{equation}
The scalar-square condition applies to every member of $\mathscr Q$.

\paragraph*{P4. A happening is the same for every observer.}
A closed comparison is a fact. Transport compares separately chosen frames; its unitary, reciprocal form is
\begin{equation}
 U_{yx}=U_{xy}^{-1},\qquad U_{xy}\longmapsto g_xU_{xy}g_y^{-1}.
 \label{eq:reciprocaltransport}
\end{equation}
A physical amplitude contracts this transport with its preparation and detection states, or takes the trace around a closed history. The result is invariant under the independent frame changes.

\paragraph*{P5. Existence is a world line that never immediately undoes its last step.}
At each step one unit distinction from a local orthonormal frame acts on one member of the pair. The continuing state consists of the relation and its last direction. The next direction differs from the last. The product of the intervening steps compares two moments of the line; the continuation rule retains this one-step memory.

The step in an orthonormal five-frame is
\begin{equation}
 (\psi,s_{\rm last})\longmapsto
 ((\gamma_s\otimes\Id)\psi,s),\qquad s\ne s_{\rm last}.
 \label{eq:postulatestep}
\end{equation}
The distinction acts on one participant and the remembered label belongs to the continuing relation.

\subsection{Equal sides and an explicit representation}
A distinction compares equal sides within the same elementary system. The antiunitary map $\Theta$ pairs each vector with a perpendicular vector of the same length, since $\Theta^2=-\Id$. Thus $H\simeq\mathbb H^n$, of complex dimension $2n$. A distinction is any traceless quaternionic Hermitian operator. The dimension of this entire real vector space is $2n^2-n-1$. For $n=1$ it is zero. For $n\geq3$ it contains $\diag(1,-1,0,\ldots)$, whose square has unequal diagonal entries. For $n=2$ every member is
\begin{equation}
 a=\begin{pmatrix}r&q\\\bar q&-r\end{pmatrix},\qquad
 a^2=(r^2+|q|^2)\Id,\qquad r\in\RR,\ q\in\mathbb H.
 \label{eq:quaternionic}
\end{equation}
The scalar-square condition on the whole space of distinctions therefore fixes $H=\CC^4$ and five real components of a distinction. Polarizing the square gives $ab+ba=2\trn(ab)\Id$, where $\trn=\Tr/4$. An orthonormal basis obeys $\{\gamma_s,\gamma_t\}=2\delta_{st}\Id$.

An explicit model makes every finite calculation reproducible. With the Pauli matrices $\sigma_x,\sigma_y,\sigma_z$ and complex conjugation ${\mathsf C_{\rm conj}}$, take
\begin{align}
 \Theta&=(\Id_2\otimes\ii\sigma_y){\mathsf C_{\rm conj}},\notag\\
 (\gamma_1,\gamma_2,\gamma_3,\gamma_4,\gamma_5)
 &=(\sigma_x\otimes\Id_2,\sigma_z\otimes\Id_2,
 \sigma_y\otimes\sigma_x,\sigma_y\otimes\sigma_y,\sigma_y\otimes\sigma_z).
 \label{eq:explicitclifford}
\end{align}
Their product is $\Id$. Every allowed quaternionic orthonormal frame gives the same representation up to $\mathrm{Sp}(2)$ conjugation. For unit $n,m\in\RR^5$, the sides $Q_n=(\Id+\gamma(n))/2$ obey
\begin{equation}
 Q_nQ_mQ_n=\frac{1+n\cdot m}{2}Q_n.\label{eq:sideoverlap}
\end{equation}
The probability of remaining on a side depends only on its angle to the other side. The possible sides form $\mathbb HP^1=S^4$.

In any complex orthonormal basis $u_1,\ldots,u_4$, put
\begin{equation}
 \Omega=\frac12\sum_{j=1}^4u_j\otimes\Theta u_j,\qquad
 \Pi_0=|\Omega\rangle\langle\Omega|,\qquad E_\Omega=4\Pi_0,
 \qquad P(v\otimes w)=w\otimes v.
 \label{eq:agreementstate}
\end{equation}
Then $P\Omega=-\Omega$, $E_\Omega^2=4E_\Omega$, and $PE_\Omega=E_\Omega P=-E_\Omega$. The common $\mathrm{Sp}(2)$ action on the pair has three irreducible spaces of dimensions $1,10,5$, so its commutant is spanned by $\Id,P,E_\Omega$. Symplectic transposition on one factor interchanges $P$ and $E_\Omega$. The equality of the two participants and equal unit weights give
\begin{equation}
 L=\Id+\sum_t\gamma_t\otimes\gamma_t=2(P+E_\Omega),\qquad
 \Delta=\frac12\sum_t(\gamma_t\otimes\Id-\Id\otimes\gamma_t)^2.
 \label{eq:paircomparison}
\end{equation}
The even Clifford products $\lvert S\rangle=(\gamma_S\otimes\Id)\Omega$ form an orthonormal basis, with $L\lvert S\rangle=(6-2|S|)\lvert S\rangle$. The grades $|S|=0,2,4$ are the singlet, ten symmetric states and five vector states. The individual distinctions anticommute, but their paired agreements commute: interchanging $\gamma_s\otimes\gamma_s$ and $\gamma_t\otimes\gamma_t$ introduces two minus signs, which cancel. Complete agreement can therefore be jointly exact in $\Omega$ although neither participant separately has definite values for all five distinctions. For a product preparation, the two expectation vectors of the $\gamma_t$ each have length at most one; therefore $\langle L\rangle\leq2$ and $\langle\Delta\rangle\geq4$. The same inequality survives convex mixing. Agreement is a property of the relation between the participants.

The coefficient matrix of $\Omega$ is ${\mathsf J_\Theta}^{\mathsf T}/2$, with ${\mathsf J_\Theta}$ unitary. Consequently its reduced density is $\Id_4/4$. A unitary change on either participant preserves this density:
\begin{equation}
 \operatorname{Tr}_2\!\left[|(U\otimes\Id)\Omega\rangle
 \langle(U\otimes\Id)\Omega|\right]
 =U\frac{\Id_4}{4}U^\dagger=\frac{\Id_4}{4}.
 \label{eq:pairmarginal}
\end{equation}

\subsection{Moving a matrix across the agreement state}
\label{sec:crossingidentity}

Write the antiunitary reciprocity as $\Theta={\mathsf J_\Theta}{\mathsf C_{\rm conj}}$ in the orthonormal
basis of Eq.~\eqref{eq:agreementstate}, where ${\mathsf C_{\rm conj}}$ is complex conjugation.
Its matrix obeys ${\mathsf J_\Theta}^\dagger {\mathsf J_\Theta}=\Id$ and ${\mathsf J_\Theta}^{\mathsf T}=-{\mathsf J_\Theta}$.
Symplectic transposition is therefore
$X^\sharp={\mathsf J_\Theta}X^{\mathsf T}{\mathsf J_\Theta}^{-1}$. For every complex linear operator $X$
on the four amplitudes,
\begin{equation}
 (\Id\otimes X)\Omega=(X^\sharp\otimes\Id)\Omega.
 \label{eq:crossingidentity}
\end{equation}
Indeed, the coefficient matrix of
$\Omega=\tfrac12\sum_j u_j\otimes\Theta u_j$ is
${\mathsf J_\Theta}^{\mathsf T}/2$. Acting on its second factor gives
${\mathsf J_\Theta}^{\mathsf T}X^{\mathsf T}/2$; acting on its first factor gives
$X^\sharp {\mathsf J_\Theta}^{\mathsf T}/2$. The two matrices agree because
${\mathsf J_\Theta}^{\mathsf T}=-{\mathsf J_\Theta}$.

\subsection{The six-dimensional form}
An occurrence separates directions within a relation that previously admitted all of them equally. Its signed geometry and its positive probability norm come from the same amplitudes. The antisymmetric pair has six complex components. Write $\Omega_t=(\gamma_t\otimes\Id)\Omega$. Its normalized Pfaffian is
\begin{equation}
 p\left(\xi_0\Omega+\sum_t \xi_t\Omega_t\right)=\xi_0^2-\sum_t \xi_t^2.
 \label{eq:pfaffiansi}
\end{equation}
For a direct verification, let $X=\xi_0\Id+\sum_t \xi_t\gamma_t$ and $\Theta={\mathsf J_\Theta}{\mathsf C_{\rm conj}}$. The antisymmetric coefficient matrix is proportional to $X{\mathsf J_\Theta}$; $\det X=(\xi_0^2-\sum_t \xi_t^2)^2$. The Pfaffian normalized to one at $X=\Id$ is the displayed polynomial. On the real span of $\Omega,\ii\Omega_t$ it is positive. On the real span of $\Omega,\Omega_t$ it has signature $(1,5)$.

Let $P_D$ project onto one part of an orthogonal division of the Euclidean six-dimensional space. Give that part a quarter turn,
\begin{equation}
 q_D=\Id-P_D+\ii P_D,\qquad
 q_D^\dagger q_D=\Id,\qquad
 q_D^{\mathsf T}q_D=q_D^2=\Id-2P_D.
 \label{eq:divisionroot}
\end{equation}
The Hermitian pairing conjugates one factor and preserves the positive norm. The bilinear pairing of the same factors gives the signed form. Write $d|6-d$ with the divided part of dimension $d\leq3$. The operation $q_D$ acts on that part; interchanging the parts changes its phase even though it gives the same unordered division. There are four divisions:
\begin{center}
\begin{tabular}{cccl}
\toprule
Division&Preserved rotations&Mixed rotations&Spin real form\\\midrule
$0|6$&15&0&$SU(4)$\\
$1|5$&10&5&$SU^*(4)$\\
$2|4$&7&8&$SU(2,2)$\\
$3|3$&6&9&$SL(4,\RR)$\\\bottomrule
\end{tabular}
\end{center}
For signature $(p,q)$ with both signs present, a projective null ray is a pair of unit vectors in $S^{p-1}\times S^{q-1}$ modulo simultaneous sign. Its conformal tangent signature is $(p-1,q-1)$. The $2|4$ division therefore has a four-dimensional Lorentzian conformal null geometry. Its compact $SO(2)$ rotation is the time coordinate of the conformal cylinder; ordinary time translations in a Minkowski chart are noncompact conformal transformations.

\subsection{Two encounters with a common participant}
When the same system meets two others, its two comparisons share its state. Their order can matter. Let $a,b,c$ denote three participants, and let $A,B$ compare their local five-frames. Write
\begin{equation}
 L_{ab}(A)=\Id+\sum_{s,t}A_{st}\gamma_s^{(a)}\gamma_t^{(b)}.
\end{equation}
The Clifford relation on the common participant gives
\begin{equation}
 \frac12\{L_{ab}(A),L_{ac}(B)\}
 =L_{ab}(A)+L_{ac}(B)+L_{bc}(A^{\mathsf T}B)-2\Id.
 \label{eq:threeparticipants}
\end{equation}
Indeed, the product of the nonconstant parts has symmetric contribution
$\sum_{s,t,v}A_{st}B_{sv}\gamma_t^{(b)}\gamma_v^{(c)}$.
The induced comparison of $b$ with $c$ therefore uses their transported relative frame $A^{\mathsf T}B$. The antisymmetric contribution contains $[\gamma_s,\gamma_u]^{(a)}$ and records the order of the two encounters. A common participant thus relates both the states and the histories of those it meets.

A coherent comparison of two orders keeps their relative phase available. Let two unitary transformations be $A,B$, and let $|0\rangle,|1\rangle$ be the two alternatives of an accessible control. The preparation retains no further record that distinguishes the orders. For a normalized input $|\psi\rangle$, the joint state is
\begin{equation}
 |\Psi\rangle=\frac{|0\rangle BA|\psi\rangle+|1\rangle AB|\psi\rangle}{\sqrt2}.
 \label{eq:orderstate}
\end{equation}
Comparing the order in the basis $|\pm\rangle=(|0\rangle\pm|1\rangle)/\sqrt2$ gives
\begin{equation}
 p_\pm=\frac14\|(BA\pm AB)|\psi\rangle\|^2.
 \label{eq:orderprobability}
\end{equation}
For the task in which the two transformations are known to commute or anticommute, precisely one probability is one. This is the ideal quantum-switch success. The experimental comparison for randomly chosen commuting and anticommuting pairs is recorded in H03.\cite{Procopio} The order is itself a physical degree of freedom; comparing its superpositions reveals the relative sign between the histories. If a further system retains orthogonal records of the orders, tracing that record removes the cross term and gives $p_+=p_-=1/2$ for either kind of pair. Which physical relations remain available determines which comparison the observer can make.

\section{One step}
\label{sec:step-one}
\label{sec:finite}

\subsection{The same comparison in its state and history bases}
The sixteen basis states describe complete agreement, ten independent comparisons of pairs of directions, and five comparisons of a single direction with the reference. We label them by the even subsets $\mathcal V$ of $\{1,\ldots,5\}$. If $g_t$ is the complement of $\{t\}$, one change takes $x$ to $x\mathbin{\triangle}g_t$. These adjacencies form the Clebsch graph. In its vertex basis,
\begin{equation}
 A_{y x}=\sum_{t=1}^5[\,y=x\mathbin{\triangle}g_t\,],
 \qquad A^2=3\Id-2A+2J_{16},
 \qquad \operatorname{spec}A=\{5^{(1)},1^{(10)},-3^{(5)}\}.
 \label{eq:adj}
\end{equation}
Here brackets denote an indicator and $J_{16}$ is the $16\times16$ all-ones matrix. The remembered state is $(x,t)$ and its next state is $(x\mathbin{\triangle}g_s,s)$, $s\ne t$. With target index first,
\begin{equation}
 B_{(y,s),(x,t)}=[s\ne t][y=x\mathbin{\triangle}g_s].
 \label{eq:B}
\end{equation}
Thus $B$ is an $80\times80$ matrix with four entries in every row and column; it has 320 directed transitions. The transition probability matrix $\mathsf P=B/4$ has stationary mass $1/80$ at each state.



\subsection{Closed histories and their determinant}
For a carried connection, put its transport on each directed transition and form $B[U]$. Its closed law is
\begin{equation}
 Z(u;U)=\det(\Id-uB[U])^{-1},\qquad
 \log Z=\sum_{r\geq1}\frac{u^r}{r}\Tr B[U]^r,
 \quad |u|<\tfrac14.
 \label{eq:closedlaw}
\end{equation}
A closed history has no preferred starting moment; the factor $1/r$ removes that choice. The finite determinant also defines the continuation beyond the convergent history series. A reciprocal rank-$d$ edge connection, $U_{yx}=U_{xy}^{-1}$, has the Bass prefactor $(1-u^2)^{24d}$; the scalar factors in Eq.~\eqref{eq:bass} belong to the free connection. For a reciprocal connection, a step carries the comparison from one local
frame to the next. A product round a closed reduced history transforms by
conjugation. Its trace is therefore available to every participant. For a
primitive oriented cycle $p$, write its length as $|p|$ and its transport as
$U_p$. The repeated traversals give
\begin{align}
 \log Z(u;U)
 &=\sum_p\sum_{r\geq1}\frac{u^{r|p|}}{r}\Tr U_p^r,
 &Z(u;U)&=\prod_p\det(\Id-u^{|p|}U_p)^{-1}.
 \label{eq:eulerproduct}
\end{align}

\subsection{The state resolution of the determinant}
The character matrix $\mathscr F$ and the map $\mathscr H_0$ from coordinate kets to pair states are
\begin{equation}
 \mathscr F_{xS}=\frac14(-1)^{|x\cap S|},\qquad
 \mathscr H_0\lvert S\rangle_{\rm coord}=(\gamma_S\otimes\Id)\Omega,\qquad
 T=\mathscr H_0\mathscr F^{\mathsf T}.
\end{equation}
They give $\mathscr F^{\mathsf T}A\mathscr F=\diag(5-2|S|)$ and
\begin{equation}
 T^\dagger\Delta T=5\Id-A.
\end{equation}
Thus the graph remembers precisely how two systems disagree. Returning histories propagate that disagreement. Put $\chi(u)=u^{-1}+4u-5=(1-u)(1-4u)/u$. Eliminating the remembered direction gives the Ihara--Bass identity\cite{Ihara,Bass}
\begin{align}
 D(u)&=\det(\Id-uB)
 =(1-u^2)^{24}\det[(1+4u^2)\Id-uA]\notag\\
 &=(1-u^2)^{24}u^{16}\det[\Delta+\chi(u)\Id]\notag\\
 &=(1-u^2)^{24}(1-u)(1-4u)
 (1-u+4u^2)^{10}(1+3u+4u^2)^5.
 \label{eq:bass}
\end{align}
The middle expression extends polynomially to $u=0$. The only zero in $|u|<1/2$ is the constant-mode zero $u_c=1/4$. The two quadratic factors have their zeros on $|u|=1/2$; the full determinant also contains the zero $u=1$ and the prefactor zeros $u=\pm1$. The displacement obeys $\chi(u)=\chi(1/(4u))$. On $|u|=1/2$ this involution sends $u$ to $\bar u$; its fixed points, and the zeros of $\chi'(u)$, are $u=\pm1/2$.

\subsection{The encounter resolved into its five channels}
\label{sec:encounterfactor}

The character transform resolves the vertex comparison before the next
direction is chosen. For the even set $S$ of grade $N=|S|$, put
$\epsilon_s=(-1)^{[s\in S]}$ and
$E_S=\diag(\epsilon_1,\ldots,\epsilon_5)$.
There are $a_+=5-N$ positive channels and $a_-=N$ negative channels.
The directed step on this five-dimensional character space is
\begin{equation}
 (B_S)_{st}=\epsilon_s(1-\delta_{st}),\qquad
 B_S=E_S(J_5-\Id_5),\qquad
 \Id_5-uB_S=\Id_5+uE_S-uE_S\mathbf1\mathbf1^{\mathsf T}.
 \label{eq:encounterchannelmatrix}
\end{equation}
The last term has rank one. For $u\ne\pm1$, the determinant lemma gives
\begin{align}
 \det(\Id_5-uB_S)
 &=(1+u)^{a_+}(1-u)^{a_-}
 \left[1-u\left(\frac{a_+}{1+u}-\frac{a_-}{1-u}\right)\right]\notag\\
 &=(1+u)^{a_+-1}(1-u)^{a_--1}
 \bigl[1+4u^2-u(a_+-a_-)\bigr].
 \label{eq:encounterfactor}
\end{align}
The coefficient four is $a_++a_--1$: one of the five directions is the
arrival direction. At grade zero the apparent factor $(1-u)^{-1}$
cancels the factor $1-u$ in the bracket. Every expression is therefore a
polynomial, and the identity also holds at $u=\pm1$ by continuation.

There is one block of grade zero, ten of grade two and five of grade four.
The two edge exponents add to
\begin{equation}
 \sum_S(a_+-1)=4+10\cdot2+5\cdot0=24,\qquad
 \sum_S(a_--1)=-1+10\cdot1+5\cdot3=24.
 \label{eq:encounteredgecounts}
\end{equation}
Consequently their product is
\begin{equation}
 \det(\Id-uB)
 =(1-u^2)^{24}k_0(u)k_2(u)^{10}k_4(u)^5,
 \qquad k_N(u)=1+4u^2-u(5-2N).
 \label{eq:encounterproduct}
\end{equation}
Eliminating the arrival first gives the Bass kernel. Resolving the
comparison first gives these five-channel blocks. Both operations retain
the same eighty directed states and give this one determinant.

The formal channel identity is
\path{Comparison.Law.pattern_factor}; its grade-zero cancellation is
\path{Comparison.Law.singlet_factor}. The numerical index identifies the
corresponding five-channel determinant calculation.
\subsection{The return source}
The inverse and its elementary source obey
\begin{equation}
 K_{80}(u,U)\mathcal G(u,U)=\Id,\qquad
 K_{80}=\Id-uB[U],\qquad
 \partial_j\log Z=-\Tr(K_{80}^{-1}(K_{80})_j).
 \label{eq:lawinverse}
\end{equation}
The formal power series determines the finite polynomial. Where the series
converges, it also gives its analytic reciprocal.

The graph constant vector is $\mathbf1/4$. The character transform sends it to
the paired singlet $\Omega$. The graph return vertex $\lvert0\rangle$ has weight $1/16$
in each character, so its spectral weights in the three comparison grades
are $1/16,10/16,5/16$. In particular,
\begin{equation}
 \langle0|(\zeta\Id-A)^{-1}|0\rangle
 =\frac1{16}\left(\frac1{\zeta-5}
 +\frac{10}{\zeta-1}+\frac5{\zeta+3}\right).
 \label{eq:returnvertex}
\end{equation}
At the leading zero the mode itself is the uniform vector. The return at a
specified vertex samples all three grades with the displayed weights.

The first matrices already give the complete history polynomial:
\begin{equation}
 \det(\lambda\Id-B)
 =(\lambda^2-1)^{24}(\lambda-4)(\lambda-1)
 (\lambda^2-\lambda+4)^{10}
 (\lambda^2+3\lambda+4)^5.
 \label{eq:Bspectrum}
\end{equation}
Every row and column of $B$ sums to four. Direct multiplication gives
$10\leq(B^5)_{xy}\leq24$, so the positive stationary mode is simple.

\subsection{The comparison in one spectral coordinate}
The history and state kernels carry the same finite comparison:
\begin{equation}
 K_{80}(u)=\Id-uB,\qquad
 K_{16}(u)=(1+4u^2)\Id-u(L-\Id)
          =u[\Delta+\chi(u)\Id],
 \quad \Delta=6\Id-L,\quad \chi=u^{-1}+4u-5.
 \label{eq:onekernel}
\end{equation}
The variable includes the weight and phase of one continuation. The factor
four counts the directions available after the previous change is remembered.
It fixes the critical modulus $|u|=1/4$ before a physical unit is assigned.

\section{Two steps}
\label{sec:step-two}

\subsection{The direction and the plane}
A continuing world line distinguishes its present, its preceding direction $h$, and the plane $\{h,k\}$ of its preceding two steps. Their relations are $\Omega$, $\Omega_h$, and $\pm \lvert\{h,k\}\rangle$. A new step $s\ne h$ changes that plane to $\{h,s\}$. The next participant brings another such particular frame. Unitary transport compares the two frames along the encounter.

Two clock generators associated with the same local history are $\ii\gamma_h$ and $\gamma_h\gamma_k$. Together with $\ii\gamma_k$ they form a quaternionic triple. The unitary
\begin{equation}
 H_k=\frac{\Id+\ii\gamma_k}{\sqrt2}
\end{equation}
conjugates one clock generator into the other, with the sign set by orientation. The continued comparison carries the clock by this conjugation. Transport of the selected central phase is path-independent precisely when its holonomy around every closed path is trivial.

\subsection{The central phase and the null cone}

For a division $d|(6-d)$ the preserving Lie algebra is
$\mathfrak{so}(d)\oplus\mathfrak{so}(6-d)$. Its centre has a one-dimensional
continuous factor precisely for $d=2$ or $d=4$. It is the rotation of the
retained two-plane. If $q_{xy}\in U(1)$ carries its phase, a common phase
assignment has the form
\begin{equation}
 z_x=q_{xy}z_y,
 \qquad
 \prod_{(x,y)\in\gamma}q_{xy}=1
 \quad\hbox{for every closed history }\gamma.
 \label{eq:commonphase}
\end{equation}
Necessity follows by multiplying the equations around the history. For
sufficiency, fix one included reference state, transport its phase along a
path to each state, and compare two such paths. Their concatenation is a
closed history, so both transports agree. The construction uses a reference
inside the comparison.

For the $2|4$ real form, a nonzero null vector has equal squared lengths in
the two parts. Quotienting by its common scale and sign gives
\begin{equation}
 \mathcal N_{2,4}=(S^1\times S^3)/\{(x,y)\sim(-x,-y)\},
 \qquad \operatorname{sign}T\mathcal N_{2,4}=(1,3).
 \label{eq:nullquotient}
\end{equation}
On a local two-spinor side, the same null relation is
\begin{equation}
 \det(t\Id_2+\bm x\cdot\bm\sigma)=t^2-|\bm x|^2.
 \label{eq:localcone}
\end{equation}
The clock phase and the three transverse coordinates belong to this one quadratic form. Let $\Delta t$ be the duration of one step and $\Delta x$ its spatial extent. Their ratio is $c=\Delta x/\Delta t$.

\section{All directions}
\label{sec:step-three}

\subsection{Adding a direction and rotating a frame}

On the full exterior space of the real six, let $\varepsilon_a$ add direction
$a$ and let $\iota_a$ remove it. With $a\ne b$ in the last identity,
\begin{equation}
 J_a=\varepsilon_a+\iota_a,\qquad
 J_aJ_b+J_bJ_a=2\delta_{ab}\Id,
 \qquad
 T_{ab}=\sum_k\Pi_kJ_aJ_b\Pi_k
       =\varepsilon_a\iota_b-\varepsilon_b\iota_a.
 \label{eq:verticalparallel}
\end{equation}
A product of two changes contains the degree-preserving rotation. The same
even products act on the original four amplitudes. Choose the reference
direction $n$ perpendicular to the five directions $t=1,\ldots,5$. Their
representation on $H=\CC^4$ is
\begin{equation}
 \rho_4(J_nJ_t)=-\ii\gamma_t,\qquad
 \rho_4(J_sJ_t)=\gamma_s\gamma_t\quad(s\ne t),\qquad
 \gamma_t=\ii\rho_4(J_nJ_t).
 \label{eq:verticalfourspinor}
\end{equation}
Indeed, the five $J_nJ_t$ square to $-\Id$ and anticommute, just as the five
$-\ii\gamma_t$ do. The Clifford relations therefore extend this assignment
to the even algebra. Since
$J_sJ_t=-(J_nJ_s)(J_nJ_t)$, the second identity follows. In the orientation
$w=J_nJ_1\cdots J_5$, Eq.~\eqref{eq:explicitclifford} gives
$\rho_4(w)=-\ii\Id$. The image contains $\Id$, $\ii\Id$, the five
$\gamma_t$ and their ten twofold products, and hence all of
$\operatorname{End}_{\CC}(\CC^4)$. Its real dimension is 32, the dimension
of the even algebra. A step along $t$ on the four amplitudes is thus the
two vertical changes along $n$ and $t$, with their quarter phase.\cite{Repository}

The spin lift of the rotation is
\begin{equation}
 R_{ab}(\theta)=\exp(\theta J_aJ_b/2),\qquad
 R_{ab}(2\pi)=-\Id,\qquad R_{ab}(4\pi)=\Id.
 \label{eq:spinreturn}
\end{equation}
The shortest perpendicular return is
$J_aJ_bJ_aJ_b=-\Id$. The spin sign and the rotation subsequently transported
between frames have therefore appeared in the vertical calculation.

\subsection{The spectral phase of a divided set}
The history of a line and the relation of two participants express the same comparison. Eliminating the remembered direction turns one into the other. The scalar displacement in Eq.~\eqref{eq:onekernel} says how far the continuing history stands from reproducing the singlet. Solving for its two history weights gives
\begin{equation}
 u_\pm(\chi)=\frac{\chi+5\pm\sqrt{(\chi+5)^2-16}}8,
 \qquad u_+(\chi)u_-(\chi)=\frac14.
 \label{eq:reciprocalroots}
\end{equation}
Interchanging the roots leaves the comparison fixed. At $\chi=0$ they are $1/4$ and $1$; the branch beginning at $1/4$ is the zero used by the long-range propagation.

The four equally weighted continuations fix the physical magnitude $|u|=1/4$. A divided set contributes a quarter-turn for each of its divided directions, so $u=\ii^{|D|}/4$. For a continuously carried phase $u=e^{\ii\theta}/4$,
\begin{equation}
 \chi(\theta)=5(\cos\theta-1)-3\ii\sin\theta,
 \qquad \frac{(\Re\chi+5)^2}{25}+\frac{(\Im\chi)^2}{9}=1.
 \label{eq:ellipse}
\end{equation}
This ellipse has foci $-1,-9$. The four phase values of the same comparison are
\begin{equation}
\begin{array}{c|rrrr}
 u&1/4&\ii/4&-1/4&-\ii/4\\\hline
 \chi&0&-5-3\ii&-10&-5+3\ii.
\end{array}
\label{eq:fourquarters}
\end{equation}
On the reciprocal circle $|u|=1/2$, the ellipse collapses to the interval $[-9,-1]$. This circle contains the roots of the two nonconstant quadratic factors in Eq.~\eqref{eq:bass}. Analytic continuation in $u$ therefore records both the physical phase and the distance of a zero from the continuing branch.

\subsection{The signed return and the divided sets}
\label{sec:signedreturn}
In the paired basis, the Clifford-signed adjacency is the matrix
\begin{equation}
 (A_-)_{yx}=\left\langle y\left|\sum_{s=1}^5(\gamma_s\otimes\Id)\right|x\right\rangle.
 \label{eq:signedadjacencydefinition}
\end{equation}
Anticommutation gives $A_-^2=5\Id$; its trace is zero, so each sign has multiplicity eight. Consequently
\begin{equation}
 D_-(u)=(1-u^2)^{24}(1+3u^2+16u^4)^8.
 \label{eq:signeddet}
\end{equation}


Eliminating the remembered arrival gives the diagonal sum of reduced signed walks,
\begin{equation}
 T_-(u)=\frac{(1-u^2)(1+4u^2)}{(1+4u^2)^2-5u^2}.
\end{equation}
The first direction has probability $1/5$ and each continuation has probability $1/4$. Thus the completed return, including its empty history, is
\begin{equation}
 R_-(z)=1+\frac45[T_-(z/4)-1],\qquad
 R_-(z)-1=-\frac{16u^4}{1+3u^2+16u^4},\quad z=4u.
\end{equation}
A nonempty return is its first return followed by a completed return. Hence $R_-=1+\Phis R_-$ and
\begin{equation}
 \Phis(z)=1-\frac1{R_-(z)}=-\frac{z^4}{16+3z^2}.
 \label{eq:signedrenewal}
\end{equation}
For each of the 64 divided sets $S$ in the six directions, let $r$ be the number of its five nonreference directions. Put
\begin{align}
 z_S&=\ii^{|S|},&p_S&=(-1)^{|S|},&L_S&=5-r+rp_S,\notag\\
 w_S&=z_S/L_S,&
 \Psi_S&=-\frac{\Phis(4w_S)}{16w_S^2}
 =\frac{p_S}{L_S^2+3p_S},\notag\\
 r_S&=1-2|w_S|,&t_S&=2|w_S|,&
 N(S)&=|S|(6-|S|)+[|S|\text{ odd}].
 \label{eq:support}
\end{align}
Their complete classification is
\phantomsection\label{sec:supportclassification}
\begin{center}
\begin{tabular}{rrrrrr}
\toprule
$p_S$&$|L_S|$&Number&$\Psi_S$&$r_S$&$t_S$\\\midrule
$+1$&5&32&$1/28$&$3/5$&$2/5$\\
$-1$&5&2&$-1/22$&$3/5$&$2/5$\\
$-1$&3&10&$-1/6$&$1/3$&$2/3$\\
$-1$&1&20&$1/2$&$-1$&2\\\bottomrule
\end{tabular}
\end{center}
These signed amplitudes depend on the size and parity of the divided set $S$. The even two-direction set has $\Psi_+=1/28$; adjoining the reference direction gives the odd set used below, with $\Psi_-=1/2$. The singular values of $B$ are $4$ with multiplicity 16 and $1$ with multiplicity 64. Their polar ratios give $\tanh(\log2)=3/5$, $\tanh(\tfrac12\log2)=1/3$, and $\tanh(\tfrac14\log2)=3-2\sqrt2$.

\subsection{The four phase evaluations}
The phase of a divided support is $z_S=\ii^{|S|}$ and its step weight is $u_S=z_S/4$. For the four unordered divisions, $|S|=d\leq3$ is the divided side. The three grade factors $k_N(u)=1+4u^2-u(5-2N)$, $N=0,2,4$, and the signed return give:
\begin{center}\small
\setlength{\tabcolsep}{5pt}
\begin{tabular}{c|cccc}
\toprule
$u$&$1/4$&$\ii/4$&$-1/4$&$-\ii/4$\\\midrule
$k_0$&$0$&$(3-5\ii)/4$&$5/2$&$(3+5\ii)/4$\\
$k_2$&$1$&$(3-\ii)/4$&$3/2$&$(3+\ii)/4$\\
$k_4$&$2$&$3(1+\ii)/4$&$1/2$&$3(1-\ii)/4$\\
$1-u^2$&$15/16$&$17/16$&$15/16$&$17/16$\\
$\Phis(4u)$&$-1/19$&$-1/13$&$-1/19$&$-1/13$\\
$u/(1+4u^2)$&$1/5$&$\ii/3$&$-1/5$&$-\ii/3$\\\bottomrule
\end{tabular}
\end{center}
The parity of the support determines the power of its amplitude: an even
support retains a Hermitian comparison and has power one; an odd support
retains the paired antilinear comparison and has power one half. Its number
of mixed rotations is $|S|(6-|S|)$, with one further antilinear channel for
odd $|S|$. This gives the last term in $N(S)$ in Eq.~\eqref{eq:support}.

\subsection{The amplitude of one continuation}
\label{sec:stepamplitude}

A direction can be compared with the direction just taken, or with one of the four alternatives. Equal standing makes all alternatives equivalent. Lossless comparison preserves the amplitude norm. These requirements fix the real five-direction amplitude directly.

\paragraph*{Theorem: the unique transmitting real amplitude.}
Let $X$ be a real matrix commuting with every permutation of the five directions and satisfying $X^{\mathsf T}X=\Id$. Permutation symmetry makes all its diagonal entries equal and all its off-diagonal entries equal, so $X=p\Id+qJ_5$. Its eigenvalue on the common direction is $p+5q$, and its four equal contrast eigenvalues are $p$. Orthogonality therefore requires
\begin{equation}
 p^2=1,\qquad (p+5q)^2=1.
\end{equation}
For nonzero transmission, $q\ne0$, the two eigenvalues have opposite signs. Choosing the common direction with eigenvalue $+1$ gives
\begin{equation}
 \boxed{\sigma_5=\frac25J_5-\Id,\qquad
 (\sigma_5)_{tt}=-\frac35,\qquad
 (\sigma_5)_{st}=\frac25\quad(s\ne t).}
 \label{eq:stepreflection}
\end{equation}
The only other transmitting solution is its overall negative. The remaining solutions $\pm\Id$ have zero amplitude between distinct directions.

One incoming direction thus has reflected probability $9/25$ and total transmitted probability $16/25$:
\begin{equation}
 \left(\frac35\right)^2+4\left(\frac25\right)^2=1.
 \label{eq:stepnorm}
\end{equation}
P5 describes an actual continuation, whose new direction differs from its previous direction. Conditional on those four alternatives, each has probability $(2/5)^2/[4(2/5)^2]=1/4$. The unitary comparison amplitude and the normalized continuation weights are both present: $t_0=2/5$ and $r_0=3/5$ enter endpoint transfer, while $u=1/4$ counts the four continuing histories. The identities $\tanh(\log2)=3/5$ and $\operatorname{sech}(\log2)=4/5$ express the same norm as the two sides of the light-clock triangle.

The three grade factors and their quarter values are
\begin{equation}
 k_N(u)=u(2N+\chi(u)),\qquad
 (k_0,k_2,k_4)(-1/4)=\left(\frac52,\frac32,\frac12\right).
\end{equation}
They supply the transmission coefficient $t_0$ and the complementary coefficient $r_0=1-t_0$, with corresponding coefficients for one retained vector direction (subscript $\mathrm v$):
\begin{equation}
 t_0=\frac25,\quad r_0=\frac35,\qquad
 t_{\rm v}=\frac23,\quad r_{\rm v}=\frac13.
 \label{eq:scattering}
\end{equation}
In the five-direction scattering matrix $t_0J_5-\Id$, the reflected amplitude is $t_0-1=-r_0$. With $R$ the reversal of an oriented graph edge, the corresponding unitary step is $\mathcal U=(2/5)B-(3/5)R$. The reversal term belongs to this amplitude; the four directed transitions counted in $B$ remain the continuations in Eq.~\eqref{eq:B}.

\section{The critical step}
\label{sec:step-four}

\subsection{The stopped return and its conditional visits}
We first ask what happens between leaving the singlet and reaching it again. All five initial directions are equivalent, so the first has probability $1/5$; each subsequent allowed direction has probability $1/4$. We stop at the first return. Its length is $N_{\rm ret}$, and $\varepsilon=\pm1$ records the Clifford sign carried around it. Length and sign are recorded by
\begin{align}
 \Fret(z)&=\EE z^{N_{\rm ret}}=\frac{z^4(1+z)}{16-8z-3z^2-3z^3},\notag\\
 \Phis(z)&=\EE(\varepsilon z^{N_{\rm ret}})=-\frac{z^4}{16+3z^2}.
 \label{eq:returns}
\end{align}
They give
\begin{equation}
 \Fret(1)=1,\quad \EE N_{\rm ret}=16,\quad \Var N_{\rm ret}=156,
 \quad \Pr(\varepsilon=+1)=\frac9{19},\quad
 \Pr(\varepsilon=-1)=\frac{10}{19}.
\end{equation}
The same sign function at the quarter turns gives $\Phis(\pm1)=-1/19$ and $\Phis(\pm\ii)=-1/13$.

A history also specifies which comparison sectors it passes. Let $n_L$ count arrivals among the five vector states and $n_F$ arrivals among the ten symmetric states, before the final singlet return. They have Clifford grades four and two respectively, and $N_{\rm ret}=n_L+n_F+1$. The four classes of remembered states outside the singlet have sizes $5,20,30,20$. Their continuation equations are
\begin{align}
 G_1&=zyG_2,&G_2&=\frac z4(xG_4+3yG_3),\notag\\
 G_3&=\frac z2(xG_4+yG_3),&G_4&=\frac z4(1+3yG_2),\notag\\
 \Fret(z;x,y)&=zxG_1
 =\frac{z^4x^2y(1+zy)}{16-8zy-3z^2xy-3z^3xy^2}.
 \label{eq:jointreturn}
\end{align}
Here $x$ marks a vector-state arrival and $y$ a symmetric-state arrival. Differentiation at $(z,x,y)=(1,1,1)$ gives
\begin{equation}
 \begin{array}{c|ccccc}
 &\EE n_F&\EE n_L&\Var n_F&\Var n_L&\Cov(n_F,n_L)\\\hline
 &10&5&84&12&30
 \end{array}.
 \label{eq:visitmoments}
\end{equation}
The conditional law is particularly useful:
\begin{align}
 \Pr(n_L=q)&=\frac14\left(\frac34\right)^{q-2},\quad q\geq2,\notag\\
 \EE[y^{n_F}\mid n_L=q]
 &=\left[\frac{y(1+y)}{2(2-y)}\right]^{q-1},\notag\\
 \EE[n_F\mid n_L=q]&=\frac52(q-1),\qquad
 \Var(n_F\mid n_L=q)=\frac94(q-1).
 \label{eq:conditional}
\end{align}
Conditioning on the number of vector-state arrivals leaves a conditional variance in the symmetric-state arrivals. Averaged over $n_L$, it is exactly
\begin{equation}
 V_C:=\EE\Var(n_F\mid n_L)=9
 =84-\frac{30^2}{12}.
 \label{eq:VC}
\end{equation}
The affine conditional mean makes this conditional variance equal to the regression residual. The value nine will become the first quark mixing angle after the return length and grade weight supply its normalization.

\subsection{The four transient classes and their signed return}

Set $D(z;x,y)=16-8zy-3z^2xy-3z^3xy^2$. Solving the four class equations in
Eq.~\eqref{eq:jointreturn} gives
\begin{align}
 G_1&=\frac{z^3xy(1+zy)}{D},&
 G_2&=\frac{z^2x(1+zy)}{D},\notag\\
 G_3&=\frac{2z^2x}{D},&
 G_4&=\frac{4z-2z^2y}{D}.
 \label{eq:classsolutions}
\end{align}
Multiplication by the first vector arrival $zx$ gives the numerator of
$\Fret(z;x,y)$. Substitution into each class equation proves the rational
identity without a limiting argument.

The signed first four-step return contributes $-z^4x^2y/16$. Every further
pair supplies one vector and one symmetric arrival with factor
$-3z^2xy/16$. Hence
\begin{equation}
 \Phis(z;x,y)=-\frac{z^4x^2y}{16+3z^2xy},\qquad
 {\Fret}_\pm(z;x,y)=\frac{\Fret(z;x,y)\pm\Phis(z;x,y)}2.
 \label{eq:signedjointreturn}
\end{equation}
At unit markers the two signs have probabilities $9/19$ and $10/19$.
The latter divided by the four continuations is $a=5/38$.

For a fixed encounter let $A$ of the other four directions agree. With
$q=z/4$, elimination of these four directions gives
\begin{equation}
 c_A(q)=\frac{A}{1+q}-\frac{4-A}{1-q},\qquad
 g_A(z)=\frac{q^2c_A(q)}{1-qc_A(q)}.
 \label{eq:retaineddirectionreturn}
\end{equation}
Its five values are
\begin{align}
 g_4&=\frac{z^2}{4-3z},&
 g_3&=\frac{z^2(2-z)}{16-8z+3z^2},&
 g_2&=-\frac{z^3}{16+3z^2},\notag\\
 g_1&=-\frac{z^2(2+z)}{16+8z+3z^2},&
 g_0&=-\frac{z^2}{4+3z}.&&
 \label{eq:fiveclockreturns}
\end{align}
In particular $g_4(z)=\Fret(1;z,1)$.

\subsection{Conditional moments and the quarter filter}

The conditional generating function in Eq.~\eqref{eq:conditional} factors as
\begin{equation}
 H_F(y)=\frac{y(1+y)}{2(2-y)}
     =y\,\frac{1+y}{2}\,\frac1{2-y}.
 \label{eq:conditionalfactor}
\end{equation}
The factors describe one compulsory arrival, a Bernoulli contribution of
probability $1/2$, and a geometric contribution of probabilities $2^{-r-1}$.
Their means sum to $5/2$ and their variances sum to $9/4$. The marginal
$x^2/(4-3x)$ then gives
\begin{equation}
 \EE n_F=10,\quad
 \Var n_F=\frac94\EE(n_L-1)+\frac{25}{4}\Var n_L=9+75=84,
 \quad\Cov(n_F,n_L)=30.
 \label{eq:totalvisitvariance}
\end{equation}
Since $N_{\rm ret}=n_F+n_L+1$, its variance is $84+12+60=156$. The independent
remainder is nine, and its normalized association is
\begin{equation}
 \frac13\left(1-\frac{\Cov(n_F,n_L)^2}{\Var n_F\Var n_L}\right)=\frac1{28}.
 \label{eq:associationfrommoments}
\end{equation}

The four-point Fourier filter separates return lengths modulo four:
\begin{equation}
 p_r=\frac14\sum_{j=0}^3\ii^{-jr}\Fret(\ii^j),\qquad
 s_r=\frac14\sum_{j=0}^3\ii^{-jr}\Phis(\ii^j).
 \label{eq:quarterfilter}
\end{equation}
The values are
\begin{center}
\begin{tabular}{c|rrrr}
\toprule $r$&0&1&2&3\\\midrule
$p_r$&$207/772$&$217/772$&$179/772$&$169/772$\\
$s_r$&$-16/247$&0&$3/247$&0\\\bottomrule
\end{tabular}
\end{center}
Their joint length-and-sign probabilities are $(p_r\pm s_r)/2$.

\subsection{The rest phase and a moving clock}
\label{sec:finiteclock}

A clock at rest must change while retaining its identity. Its internal direction supplies a phase $\mu$ even when there is no spatial phase. Motion adds a turn about a perpendicular Clifford direction. Let $\gamma_h^2=\gamma_j^2=\Id$ and $\{\gamma_h,\gamma_j\}=0$, with both generators Hermitian. The ordered finite step is
\begin{equation}
 U(k,\mu)=e^{-\ii k\gamma_j}e^{-\ii\mu\gamma_h}.
 \label{eq:finiteclockunitary}
\end{equation}
Expanding the two factors gives
\[
 U=\cos k\cos\mu\,\Id
 -\ii\sin k\cos\mu\,\gamma_j
 -\ii\cos k\sin\mu\,\gamma_h
 -\sin k\sin\mu\,\gamma_j\gamma_h.
\]
The last three terms are anti-Hermitian. Hence $U+U^\dagger=2\cos k\cos\mu\,\Id$. A unitary eigenvalue $e^{-\ii\omega}$ therefore satisfies the exact equation
\begin{equation}
 \boxed{\cos\omega=\cos\mu\cos k.}
 \label{eq:finiteclockdispersion}
\end{equation}
This is the cosine law of a right spherical triangle: the rest phase and spatial phase are its perpendicular legs. At rest, $k=0$ and the branch through the origin has $\omega=|\mu|$. With $\mu=0$, its principal massless branch has $\omega=|k|$. The internal phase that remains at $k=0$ is the rest energy, $mc^2=\hbar|\mu|/\Delta t$.

\paragraph*{Theorem: the same unitary clock shared by the pair.}
The crossing identity in Eq.~\eqref{eq:crossingidentity} gives
$(\gamma_h\otimes\Id)\Omega=(\Id\otimes\gamma_h)\Omega$, because $\gamma_h$ is Hermitian and commutes with the state conjugation $\Theta$. For $U_h(\mu)=e^{-\ii\mu\gamma_h}$ and every real $\mu$,
\begin{equation}
 (U_h(\mu)\otimes\Id)\Omega
 =[U_h(\mu/2)\otimes U_h(\mu/2)]\Omega
 =\cos\mu\,\Omega-\ii\sin\mu\,(\gamma_h\otimes\Id)\Omega.
 \label{eq:unitaryclocksharing}
\end{equation}
The two tensor generators commute; on $\Omega$ their sum acts as twice either one. Their half-angle exponentials therefore give the same full turn as a turn on one participant. The norm stays one throughout. This identity shares the phase of the agreement state. 

A wave packet reveals what is timed along the moving trajectory. Its phase at lattice position $x$ and step number $n$ is $kx-\omega n$. Stationary phase gives $x/n=\partial_k\omega$. Thus the exact group velocity and phase accumulated along the packet are
\begin{equation}
 \frac{v_g}{c}=\partial_k\omega
 =\frac{\cos\mu\sin k}{\sin\omega},
 \qquad
 \omega_{\rm packet}=\omega-k\,\partial_k\omega,
 \qquad c=\frac{\Delta x}{\Delta t}.
 \label{eq:packetclock}
\end{equation}
The spherical equation fixes both quantities before a continuum limit is taken.

For small $\mu$ and $k$, expanding both cosines to fourth order gives
\begin{equation}
 \omega^2=\mu^2+k^2-\frac13\mu^2k^2
 +O\!\left((|\mu|+|k|)^6\right).
 \label{eq:clockexpansion}
\end{equation}
Its leading generator is $H=\mu\gamma_h+k\gamma_j$, whose square is exactly $(\mu^2+k^2)\Id$. In the long-wave limit, $E=\hbar\omega/\Delta t$ and $p=\hbar k/\Delta x$ give
\begin{equation}
 E^2=p^2c^2+m^2c^4,\qquad
 \omega_{\rm packet}\longrightarrow\frac{\mu^2}{\omega},
 \qquad
 \frac{\dd\tau}{\dd t}=\frac{|\mu|}{\omega}
 =\frac{mc^2}{E}.
 \label{eq:clocklimit}
\end{equation}
The last ratio compares the phase along the packet with the clock's phase at rest. It is the light-clock time dilation. The same internal direction supplied the rest leg, the positive mass term and the phase with which that dilation is read.

\subsection{The exact spatial coordinate}
\label{sec:spectralcoordinate}
To resolve a step spatially, let $\mathbf e_t$ be the standard coordinate basis of $\RR^5$, set $v_t=\mathbf e_t-\frac15\sum_j\mathbf e_j$, and take spacing $\epsilon$. The two reciprocal sublattices are $\epsilon A_4$ and $\epsilon(A_4+v_h)$ in $\sum_jx_j=0$, with $A_4=\{z\in\mathbb Z^5:\sum_jz_j=0\}$. A momentum $p$ is dual to this four-dimensional coordinate space. The link sum is $a_\epsilon(p)=\sum_te^{\ii\epsilon p\cdot v_t}$, with $a_\epsilon(p)a_\epsilon(-p)=25-5\epsilon^2p^2+O(\epsilon^4p^4)$. Eliminating one sublattice from $\left(\begin{smallmatrix}\chi+5&-a_\epsilon(p)\\-a_\epsilon(-p)&\chi+5\end{smallmatrix}\right)$ gives the exact squared gap
\begin{equation}
 M_\epsilon^2=\frac{2\chi+\chi^2/5}{\epsilon^2}.
 \label{eq:exactgap}
\end{equation}
Near $u=1/4$, putting $u=1/4-\epsilon^2m^2/24$ gives $\chi=\epsilon^2m^2/2+O(\epsilon^4)$ and $M_\epsilon\to m$. The finite coordinate and its long-wave mass parameter remain explicitly related.

\subsection{The exact inverse range and phase mass}

The reciprocal spatial chart assigns inverse range $M_\epsilon$ to its
static zero. The positive branch of the original history variable is
\begin{equation}
 u(\epsilon,M_\epsilon)
 =\frac{\sqrt{25+5\epsilon^2M_\epsilon^2}
        -\sqrt{9+5\epsilon^2M_\epsilon^2}}8,
 \quad
 \chi=\sqrt{25+5\epsilon^2M_\epsilon^2}-5.
 \label{eq:exactradialbranch}
\end{equation}
The product of the two conjugate square-root combinations is sixteen.
It follows directly that $u^{-1}+4u=\chi+5$ and
$2\chi+\chi^2/5=\epsilon^2M_\epsilon^2$.

Continuing the chart momentum along the retained simplex direction $v_h$ gives the rest phase. In inverse-length units, its exact conversion on this direction is
\begin{align}
 m_{\rm phase}
 &=\frac4{\sqrt5\,\epsilon}
 \operatorname{arsinh}\!\left(\frac{\sqrt5\,\epsilon M_\epsilon}{4}\right),
 \notag\\
 M_\epsilon(m_{\rm phase})
 &=\frac4{\sqrt5\,\epsilon}
 \sinh\!\left(\frac{\sqrt5\,\epsilon m_{\rm phase}}4\right).
 \label{eq:finite_mass_conversion}
\end{align}
Both tend to the same mass in the tangent limit. At finite spacing their
conversion is part of the exact dispersion relation. The physical rest
energy is $\hbar c\,m_{\rm phase}$.

\subsection{Propagation in the resolved spatial chart}
The spatial Schur symbol defined above is
\begin{equation}
 Q_\epsilon(p)=\frac{(\chi+5)^2-a_\epsilon(p)a_\epsilon(-p)}{5\epsilon^2}.
 \label{eq:spatialsymbol}
\end{equation}
It has $Q_\epsilon=M_\epsilon^2+p^2+O(\epsilon^2p^4)$ and the bound
\begin{equation}
 0\leq p^2-\frac{25-a_\epsilon(p)a_\epsilon(-p)}{5\epsilon^2}
 \leq\frac{\epsilon^2p^4}{6}.
\end{equation}
Choosing $u=1/4-\epsilon^2m^2/24$ gives $M_\epsilon\to m$. The static projection perpendicular to the held direction is $p_t=v_t+v_h/4$ for $t\ne h$. Its four vectors sum to zero and obey $p_s\cdot p_t=\delta_{st}-1/4$. Any three are independent; their Gram determinant is $1/4$, giving cell volume $1/2$. Its symbol retains the held contribution:
\begin{equation}
 25-\left|1+\sum_{t\ne h}e^{\ii k\cdot p_t}\right|^2.
\end{equation}
The cell volume fixes the normalization of the spatial source. The continuum inverse is
\begin{equation}
 \mathcal G_{M_{\rm inv}}(r)=\int\frac{\dd^3p}{(2\pi)^3}\frac{e^{\ii\bm p\cdot\bm r}}{p^2+M_{\rm inv}^2}
 =\frac1{2\pi^2r}\int_0^\infty\frac{p\sin(pr)}{p^2+M_{\rm inv}^2}\,\dd p
 =\frac{e^{-M_{\rm inv}r}}{4\pi r},\qquad r>0.
 \label{eq:continuumgreen}
\end{equation}
For $M_{\rm inv}>0$ the last integral follows by closing the contour at the pole $p=\ii M_{\rm inv}$; $M_{\rm inv}=0$ is its limit. The massless disturbance spreads over spheres and gives $1/r$, while a nonzero rest frequency adds the attenuation $e^{-M_{\rm inv}r}$. End sources normalized by $\sqrt{4\pi}$ give the usual $1/r$ coefficient.

With Lorentz signature $(-+++)$ and $\{\Gamma^\mu,\Gamma^\nu\}=-2g^{\mu\nu}$, the free continuation has $(\Box-m^2)\psi=0$. Its first-order form is $(\ii\Gamma^\mu\partial_\mu-m)\psi=0$. Their common square gives the dispersion already derived from the finite phase. The transported square will include the curvature of the carried frame.

\subsection{One inverse propagator for interaction and self-interaction}
\label{sec:common}

A particle persists when the amplitude that leaves its present state returns with the phase required to reproduce it. Another particle can receive that amplitude along the way. Self-interaction and interaction consequently contain the same internal histories. At rest their sum reproduces the state. A current records how that state changes when another participant acts on it. Both questions require the time spent in the whole return.

For a specified state $X$, let ${\mathcal K}_X$ be its retained inverse propagator, including self-interaction. Persistence at rest means that this source reproduces the state:
\begin{equation}
 \det {\mathcal K}_X(E_X,0)=0,\qquad m_Xc^2=E_X\geq0.
 \label{eq:masszero}
\end{equation}


We use $\hbar=c=1$ for energies and masses unless a conversion is displayed. The physical energy variable is $E$ and $J$ is an applied source. Let ${\mathcal K}(E,J)$ be analytic near the frequency under consideration. Split the degrees of freedom into the part retained at the endpoints and the internal part explored between them:
\begin{equation}
 {\mathcal K}(E,J)=\begin{pmatrix}{\mathsf A}&{\mathsf B}\\{\mathsf C}&{\mathsf D}\end{pmatrix},\qquad
 S(E,J)={\mathsf A}(E,J)-\Sigma(E,J),\qquad
 \Sigma(E,J)={\mathsf B}(E,J){\mathsf D}(E,J)^{-1}{\mathsf C}(E,J).
 \label{eq:selfenergy}
\end{equation}
The self-energy $\Sigma$ sums departures from an endpoint, all intervening internal propagation and the return. The same term appears when the returning amplitude is absorbed by a second endpoint. For invertible ${\mathsf D}$, define
\begin{equation}
 {\mathcal E_R}=\begin{pmatrix}\Id\\-{\mathsf D}^{-1}{\mathsf C}\end{pmatrix},\quad
 {\mathcal E_L}=(\Id,-{\mathsf B}{\mathsf D}^{-1}),\quad
 {\mathcal Q_i}=\begin{pmatrix}0&0\\0&{\mathsf D}^{-1}\end{pmatrix}.
\end{equation}
\paragraph*{Theorem: exact elimination.}
For invertible ${\mathsf D}$, direct multiplication gives the determinant identity below. On the regular set where $S$ is also invertible, it gives the inverse:
\begin{equation}
 \det {\mathcal K}=\det {\mathsf D}\det S,\qquad {\mathcal K}^{-1}={\mathcal E_R}S^{-1}{\mathcal E_L}+{\mathcal Q_i}.
 \label{eq:schur}
\end{equation}
Thus the Schur complement retains exactly the amplitudes available to the endpoints. It also retains their responses. For source coordinates $J_a,J_b$,
\begin{align}
 S_a&={\mathcal E_L}{\mathcal K}_a{\mathcal E_R},\notag\\
 S_{ab}&={\mathcal E_L}({\mathcal K}_{ab}-{\mathcal K}_a{\mathcal Q_i}{\mathcal K}_b-{\mathcal K}_b{\mathcal Q_i}{\mathcal K}_a){\mathcal E_R},\label{eq:schurderivative}\\
 \partial_a\partial_b(-\log\det {\mathcal K})
 &=\Tr({\mathcal K}^{-1}{\mathcal K}_b{\mathcal K}^{-1}{\mathcal K}_a)-\Tr({\mathcal K}^{-1}{\mathcal K}_{ab}).
 \label{eq:twoinsertions}
\end{align}
The two orders of the encounter and its direct contact are therefore part of one response.

\subsection{The common clock derivative}
Near a continuing state, changing the energy changes the phase accumulated between departure and return. That change is the clock against which both the state's own displacement and a transmitted disturbance are measured. Suppose $S(E_0,0)r=0$ and $l^\dagger S(E_0,0)=0$, with a simple zero: $\rank S=n-1$ and $l^\dagger S_Er\ne0$. Expanding $S$ in $E-E_0$ on its one-dimensional null space gives
\begin{equation}
 S^{-1}(E,0)=\frac{rl^\dagger}{(E-E_0)\,l^\dagger S_Er}+O(1).
\end{equation}
Differentiating $S(E_0(J),J)r(J)=0$ and multiplying on the left by $l^\dagger$ then gives
\begin{equation}
 \operatorname*{Res}_{E=E_0}S^{-1}
 =\frac{rl^\dagger}{l^\dagger S_Er},\qquad
 \frac{\dd E_0}{\dd J}
 =-\frac{l^\dagger S_Jr}{l^\dagger S_Er}.
 \label{eq:residue}
\end{equation}
For preparation $X$ and detection $Y$, the amplitude passed through this state is
\begin{equation}
 \lim_{E\to E_0}(E-E_0)YS^{-1}X
 =\frac{(Yr)(l^\dagger X)}{l^\dagger S_Er}.
 \label{eq:sharedamplitude}
\end{equation}
The denominator measures the change of the continuing state's phase with energy. It normalizes the amplitude transmitted through the state and the shift of its rest frequency under an applied source. The mass itself is the zero $E_0/c^2$ in Eq.~\eqref{eq:masszero}; its source derivative is the measured response in Eq.~\eqref{eq:residue}. Energy-dependent internal propagation contributes through $S_E={\mathsf A}_E-\Sigma_E$, and thus contributes to both the zero's position and its residue.

\subsection{The norm carried by the internal return}
\label{sec:restnecessity}

An endpoint keeps only part of a continuing state. The amplitude inside its return still occupies the world and still accumulates phase. Eliminating it changes the apparent size of the endpoint's state, while its full energy derivative keeps the time spent inside. Differentiating the complete return gives
\begin{equation}
 \Sigma_E={\mathsf B}_E{\mathsf D}^{-1}{\mathsf C}+{\mathsf B}{\mathsf D}^{-1}{\mathsf C}_E-{\mathsf B}{\mathsf D}^{-1}{\mathsf D}_E{\mathsf D}^{-1}{\mathsf C}.
 \label{eq:selfenergyderivative}
\end{equation}
The first two terms change departure and arrival; the last changes the phase spent inside the eliminated state.

\paragraph*{Theorem: the clock is the full-state norm.}
Let the complete Hermitian rest equation have positive clock form $\mathcal C_{\rm clk}$,
\begin{equation}
 {\mathcal K}(E)=E\mathcal C_{\rm clk}-\mathsf H,\qquad \mathsf H=\mathsf H^\dagger,\qquad
 \mathcal C_{\rm clk}=\mathcal C_{\rm clk}^\dagger>0,
 \label{eq:fullrestform}
\end{equation}
with $\mathsf H$ and $\mathcal C_{\rm clk}$ independent of $E$. At a real energy where the eliminated block is invertible, ${\mathcal E_L}={\mathcal E_R}^\dagger$ in Eq.~\eqref{eq:schur}. For a nonzero retained vector $r$, the derivative identity in Eq.~\eqref{eq:schurderivative} therefore gives
\begin{equation}
 S_E={\mathcal E_R}^\dagger\mathcal C_{\rm clk} {\mathcal E_R},\qquad
 \psi={\mathcal E_R}r,\qquad
 r^\dagger S_Er=\psi^\dagger\mathcal C_{\rm clk}\psi>0.
 \label{eq:fullstatenorm}
\end{equation}
If $S(E)r=0$, the vector $\psi$ is precisely the full null state of ${\mathcal K}(E)$. Thus the denominator measured at the endpoint is the norm of the state everywhere it has gone.

In canonical coordinates this statement can be seen term by term. Write
\begin{equation}
 {\mathcal K}(E)=E\Id-
 \begin{pmatrix}\mathsf H_C&V\\V^\dagger&\mathsf H_D\end{pmatrix}.
\end{equation}
Its retained operator and clock are
\begin{align}
 S(E)&=E\Id-\mathsf H_C-V(E\Id-\mathsf H_D)^{-1}V^\dagger,\notag\\
 S_E&=\Id+V(E\Id-\mathsf H_D)^{-2}V^\dagger>0,
 \label{eq:positiveclock}
\end{align}
and the continuing state is
\begin{equation}
 \psi=\begin{pmatrix}r\\(E\Id-\mathsf H_D)^{-1}V^\dagger r\end{pmatrix},
 \qquad
 r^\dagger S_Er=\|r\|^2+
 \|(E\Id-\mathsf H_D)^{-1}V^\dagger r\|^2.
 \label{eq:returnnorm}
\end{equation}
For $\|r\|=1$, the inverse of this derivative is the endpoint's share of the normalized complete state. The pole residue and the transmitted current carry this same share. The retained diagonal term can vanish: $\mathsf H_C=0$ leaves the rest equation $Er=\Sigma(E)r$. A finite rest phase can therefore be sustained entirely by departure, internal propagation and return. 

Further elimination can make every retained coefficient depend on energy. Equation~\eqref{eq:selfenergyderivative} carries that dependence. For $\widehat {\mathcal K}(E)=E\Id-\mathsf H_0-\widehat\Sigma(E)$ and a normalized simple rest state $v$, the clock denominator is
\begin{equation}
 v^\dagger\bigl[\Id-\partial_E\widehat\Sigma(E_0)\bigr]v.
 \label{eq:nonlinearclock}
\end{equation}
The rest equation is $[\mathsf H_0+\widehat\Sigma(E_0)]v=E_0v$; Eq.~\eqref{eq:nonlinearclock} gives the weight of that same pole. In coordinates $S(E)=EC(E)-\mathsf H(E)$, the complete derivative is $S_E=C+EC_E-\mathsf H_E$. Each internal phase contributes through this derivative.

An analytic change of retained coordinates also preserves the response. If $\widetilde S=T_LST_R$, the derivatives of $T_L,T_R$ multiply $S$ and disappear between its null vectors. The surviving energy and source derivatives acquire the same factors, so their ratio in Eq.~\eqref{eq:residue} is unchanged. This is the precise sense in which two descriptions use one clock.

A symmetry can preserve a multiplet at a single energy. Let the columns of $V_0$ span its null space and put $C_0=V_0^\dagger S_EV_0>0$. Projection of the perturbed equation onto that space gives its first-order energy shifts as the eigenvalues of
\begin{equation}
 -C_0^{-1/2}V_0^\dagger S_JV_0C_0^{-1/2}\,J.
 \label{eq:multipletresponse}
\end{equation}
This is the same clock normalization for every member of the multiplet.

The finite walk supplies a concrete example. At $u_c=1/4$, $K_{80}=\Id-uB$ has the normalized uniform null vector $r=l=\mathbf1/\sqrt{80}$ and
\begin{equation}
 l^\dagger (K_{80})_ur=-4.
\end{equation}
An attenuation $e^{-J}$ on all departures in a particular direction has
$l^\dagger (K_{80})_Jr=1/5$ and $\dd u_c/\dd J=1/20$. A phase $e^{-\ii\vartheta}$ on one allowed transition gives $l^\dagger (K_{80})_\vartheta r=\ii/320$, hence $\dd\log(4u_c)/\dd\vartheta=\ii/320$ at zero. In the sixteen-state operator $\zeta\Id-A$, the pole at $\zeta=5$ has residue $J_{16}/16$. Relating it to the history variable carries the derivative of $\zeta=u^{-1}+4u$ and the prefactor $u$ in Eq.~\eqref{eq:bass}.

\subsection{The probability of a complete encounter}
An encounter can reach a detector through interfering histories. Their amplitudes add before the complete detection probability is formed. For endpoint matrices $X,Y$, introduce
\begin{equation}
 \mathcal K_{\rm pair}(s,t)=\begin{pmatrix}{\mathcal K}&-sXX^\dagger\\-tYY^\dagger&{\mathcal K}^\dagger\end{pmatrix}.
\end{equation}
Its mixed response satisfies
\begin{equation}
 \partial_s\partial_t[-\log\det\mathcal K_{\rm pair}]_{s=t=0}
 =\|Y^\dagger {\mathcal K}^{-1}X\|_F^2\geq0.
\end{equation}
Normalization over a complete set of detection alternatives uses the probability pairing of the first postulate. For a positive real quadratic response, $W[J]=\tfrac12J^{\mathsf T}S^{-1}J$ and its Legendre transform $\Gamma[\varphi]=\tfrac12\varphi^{\mathsf T}S\varphi$ carry the source and the measured field together through $\varphi=\partial W/\partial J$.

\subsection{One insertion, exterior completion and repeated return}
Let ${\mathcal K}$ be the completed comparison and $X={\mathcal K}^{-1}J$ its source insertion.
One determinant family describes the three operations used below:
\begin{equation}
 \mathcal Z_{\varsigma,g}(X)=\det(\Id-\varsigma X/g)^{-\varsigma g},
 \quad
 \log\mathcal Z_{\varsigma,g}
 =\sum_{r\geq1}\frac{\varsigma^{r+1}g^{1-r}}{r}\Tr X^r,
 \qquad \varsigma=\pm1,\quad g>0.
 \label{eq:sourceoperations}
\end{equation}
The logarithm is the branch through $X=0$, where the series converges for
$\|X\|<g$. For $\varsigma=1$, the limit $g\to\infty$ retains
$\exp\Tr X$. At $(\varsigma,g)=(-1,1)$ it gives $\det(\Id+X)$;
at $(1,1)$ it gives the repeated return $\det(\Id-X)^{-1}$.

For diagonal markers $J=\diag(j_i)$, multilinearity gives
\begin{equation}
 \det(\Id+JT)=\sum_S\left(\prod_{i\in S}j_i\right)\det T[S,S].
 \label{eq:exteriorsource}
\end{equation}
Choose the $JT$ term exactly on the rows in $S$. All surviving permutations
are permutations of $S$, so their signed sum is the displayed principal minor.
The full coefficient on a $d$-dimensional space is $\Lambda^dT$. The same
completion will give the charged exterior grades, the retained neutral
determinant and the positive volume of the whole.

For $J=jX+kY$ and $G={\mathcal K}^{-1}$, the connected return gives
\begin{equation}
 \partial_j[-\log\det(\Id-GJ)]_0=\Tr GX,\qquad
 \partial_j\partial_k[-\log\det(\Id-GJ)]_0=\Tr GXGY.
 \label{eq:connectedinsertions}
\end{equation}
The two cyclic orders cancel the factor $1/2$ in the logarithm's quadratic
term. A derivative of the logarithm is a connected insertion; a determinant
coefficient is a completed exterior volume. Both act on the same propagator.

\section{The plane}
\label{sec:step-five}
\label{sec:sources}

\subsection{Comparison in carried frames}
Let $P_D(x)$ select the division at $x$. Transport compares that selected state with the one at the neighbouring point:
\begin{equation}
 P_D(x)-U_{xy}P_D(y)U_{yx}\longmapsto
 g_x[P_D(x)-U_{xy}P_D(y)U_{yx}]g_x^{-1}.
 \label{eq:transportedsubtraction}
\end{equation}
The degree-preserving rotations constructed in the preceding step carry the axes of this comparison. For a step $\epsilon\,dx^\mu$, write $U=\exp(-\epsilon\mathcal A_\mu dx^\mu)$. The continuum derivative and an oriented small loop are
\begin{equation}
 D=d+\mathcal A,\qquad
 \mathcal F=d\mathcal A+\mathcal A\wedge\mathcal A,
 \qquad U_\square=\Id-\epsilon^2\mathcal F_{\mu\nu}+O(\epsilon^3).
 \label{eq:parallelcurvature}
\end{equation}
A rotation $X$ changes a selected division at the rate $[X,P_D]$. The squared change is the positive form
\begin{equation}
 \mathcal M_D(X,Y)=\Tr([X,P_D]^\dagger[Y,P_D]),\qquad
 \mathcal M_D(X,X)=\|[X,P_D]\|_F^2.
 \label{eq:positivemassform}
\end{equation}
Rotations within a part have zero cost. A unit antisymmetric generator $E_{ab}-E_{ba}$ mixing the two parts has squared cost two. This is how a particular division distinguishes massive directions from its surviving rotations.

\subsection{The common internal frame and the matter fields}
Each member of the pair carries five possible distinctions. Comparing both frames allows rotations among ten real directions. An explicit realization on $H\otimes H\otimes\CC^2$ is
\begin{equation}
 \Gamma_a=\gamma_a\otimes\Id\otimes\sigma_x,
 \qquad \Gamma_{5+a}=\Id\otimes\gamma_a\otimes\sigma_y,
 \qquad a=1,\ldots,5.
 \label{eq:tenclifford}
\end{equation}
These ten matrices obey the Euclidean Clifford relations. Their product is $\ii\Id\otimes\Id\otimes\sigma_z$, so each chiral subspace has dimension sixteen. On the $\sigma_z=+1$ subspace,
\begin{equation}
 -\ii\sum_{a=1}^5\Gamma_a\Gamma_{5+a}=L-\Id.
 \label{eq:paircomplexstructure}
\end{equation}
The bilinear rotations give the chiral spin representation of $\mathrm{Spin}(10)$. The orthogonal generator corresponding to the displayed spinor operator pairs the two real five-frames into a complex structure on $\RR^{10}$. An orthogonal transformation preserving this structure is a unitary transformation on $\CC^5$, so its centralizer is $U(5)$.

Exterior multiplication and contraction on these five complex directions realize the same ten Clifford generators. The even space $\Lambda^{\mathrm{even}}\CC^5$ has dimension sixteen. A two-plane $F$ and its complementary three-plane $C$ give the charges
\begin{equation}
 Y=\frac{\widehat n_F}{2}-\frac{\widehat n_C}{3},\qquad
 Q=\widehat n_k-\frac{\widehat n_C}{3}=T_3+Y.
\end{equation}
The spin representation of $U(5)$ acts as $z^{-1}\Lambda U$, $z^2=\det U$. Fixing the vacuum gives $SU(5)$; retaining the division gives
\begin{equation}
 S(U(2)\times U(3))\simeq
 \frac{SU(2)\times SU(3)\times U(1)}{\mathbb Z_6}.
\end{equation}
The quotient follows from the surjection $(g_3,g_2,z)\mapsto(z^{-2}g_3,z^3g_2)$, whose kernel has $z^6=1$. Here $\widehat n_F=\widehat n_h+\widehat n_k$ and $T_3=(\widehat n_k-\widehat n_h)/2$. An occupied exterior direction contributes its charge through $\widehat n_F,\widehat n_C,\widehat n_k$. Each of the sixteen basis states therefore has a definite charge. Written as left-handed fields, the states are
\begin{center}
\begin{tabular}{ccccccc}
\toprule
Class&$\nu^c$&$e^c$&$Q_L$&$u^c$&$L_L$&$d^c$\\\midrule
Multiplicity&1&1&6&3&2&3\\
$Y$&0&1&$1/6$&$-2/3$&$-1/2$&$1/3$\\\bottomrule
\end{tabular}
\end{center}
The rotations that keep the actual last direction $h$ fixed form $\mathfrak s(\mathfrak u(1)\oplus\mathfrak u(3))\simeq\mathfrak u(3)$. On $(h,k,c_1,c_2,c_3)$ their centre is
\begin{equation}
 \ii\diag(0,-3,1,1,1)=-3\ii Q.
 \label{eq:electromagneticcentre}
\end{equation}
This phase leaves $h$ undisplaced and moves no rest zero. It is the electromagnetic phase. The centre of the full $\mathfrak s(\mathfrak u(2)\oplus\mathfrak u(3))$ is generated by hypercharge.

The linear and cubic hypercharge traces vanish. The mixed colour trace is $2/6-2/3+1/3=0$; the mixed weak trace is $3/6-1/2=0$. There are four weak doublets. These are the charges of one chiral generation, including a neutral singlet.

Matter fields have Grassmann coefficients in this sixteen-dimensional chiral representation.

\subsection{A direction retained over repeated returns}
A particle can keep a preceding direction $h$ while many singlet returns pass. An observer following that persistent direction sees the same walk at a selected set of moments. The unobserved passages still carry amplitude between those moments.

Let ${\mathscr D}=\{(\emptyset,t):t=1,\ldots,5\}\cup\{(g_h,h)\}$ and let ${\mathscr C}$ be the 74 remaining remembered states. The only transitions internal to ${\mathscr D}$ go from $(\emptyset,t)$, $t\ne h$, to $(g_h,h)$, so $B_{{\mathscr D}{\mathscr D}}^2=0$. Eliminating these six states from $K_{80}=\Id-uB$ gives exactly
\begin{align}
 S_h(u)&=\Id-\mathsf W_{\mathscr C}(u),\notag\\
 \mathsf W_{\mathscr C}(u)&=uB_{{\mathscr C}{\mathscr C}}+u^2B_{{\mathscr C}{\mathscr D}}B_{{\mathscr D}{\mathscr C}}
                 +u^3B_{{\mathscr C}{\mathscr D}}B_{{\mathscr D}{\mathscr D}}B_{{\mathscr D}{\mathscr C}},\notag\\
 \det(\Id-uB)&=\det S_h(u).
 \label{eq:retainedschur}
\end{align}
The three terms retain paths of one, two and three original steps between successive observations in ${\mathscr C}$. The hidden block has determinant one, so this elimination preserves the complete determinant. At $u=1/4$ it gives the trace of the stationary Markov chain on ${\mathscr C}$:
\begin{equation}
 \mathsf Q_{\mathscr C}=\mathsf P_{{\mathscr C}{\mathscr C}}
 +\mathsf P_{{\mathscr C}{\mathscr D}}(\Id-\mathsf P_{{\mathscr D}{\mathscr D}})^{-1}\mathsf P_{{\mathscr D}{\mathscr C}}
 =\mathsf W_{\mathscr C}(1/4).
 \label{eq:tracechain}
\end{equation}
Its rows and columns each sum to one. Every retained state has stationary probability $1/74$, and each of its four original choices has probability $1/296$. A transition of $\mathsf Q_{\mathscr C}$ sums the probabilities of the paths it contains. The equal weights belong to the original choices made at the retained observations.

The original transition counts are 276 within ${\mathscr C}$, 20 from ${\mathscr C}$ to ${\mathscr D}$, 20 from ${\mathscr D}$ to ${\mathscr C}$, and four within ${\mathscr D}$. Thus there are 296 observed departures and 24 hidden steps in the full stationary count of 320. One observed departure spans $320/296=40/37$ original steps on average. Counting the directions of those observed departures gives
\begin{equation}
 M=296,\qquad C_h=60,\qquad C_j=59\quad(j\ne h),\qquad C_h+4C_j=M.
\end{equation}
Of the 74 retained states $(x,t)$, 34 leave $h$ undivided at their vertex: $h\notin x$. Hence
\begin{align}
 c_L&=\frac{C_h}{M}=\frac{15}{74},&
 \sigma&=\frac{5C_h}{M}=\frac{75}{74},\notag\\
 \eta_0&=\frac{5C_j}{M}=\frac{295}{296},&
 R_{ud}&=\frac{34}{74}=\frac{17}{37}.
 \label{eq:stationaryweights}
\end{align}
These weights obey $\sigma+4\eta_0=5$ and $\eta_0/\sigma=59/60$.

The mean interval between singlet returns remains sixteen original steps.\cite{Kac} A stationary interval traverses each original transition $1/20$ times on average. The full stationary count therefore covers twenty mean return intervals:
\begin{equation}
 \begin{array}{c|rr}
 &\text{stationary count}&\text{mean return interval}\\\hline
 \text{all original transitions}&320&16\\
 \text{departures observed in }\mathscr C&296&74/5\\
 \text{observed departures along }h&60&3
 \end{array}
 \qquad E_H=\frac{16\cdot60}{320}=3.
 \label{eq:EH}
\end{equation}
A stopped excursion conditioned on the preceding direction $h$ has a definite initial preparation: the four first states labelled by $s\ne h$ each have probability $1/4$. On its 75 transient states let $Q_*$ be the substochastic continuation matrix, with target index first, and $v$ the column of initial probabilities. Its occupation vector $(\Id-Q_*)^{-1}v$ has total 15. The expected number of all steps labelled $h$ is $368/121$, and that for each other label is $392/121$; the return label has probabilities $25/121$ and $24/121$, respectively. The pre-return fraction leaving $h$ undivided is $167/363$. The initial preparation and final return fix this excursion law. Indefinitely repeated comparisons retaining $h$ give the stationary law above.

The directional means also follow directly from the continuation rule. Let $s_0=h$, let $s_{N_{\rm ret}}$ be the return label, and let $n_j^{\rm dir}$ count every step labelled $j$, including the return. Put $p_j=\Pr(s_{N_{\rm ret}}=j)$. Each next label is uniform among the four different from its predecessor, so
\begin{equation}
 4\EE n_j^{\rm dir}=\EE N_{\rm ret}-\EE n_j^{\rm dir}-[j=h]+p_j.
 \label{eq:excursiondirection}
\end{equation}
The predecessor sequence contains the initial label and omits the terminal one; these are the last two terms. The return equations give $p_h=25/121$ and $p_j=24/121$ for $j\ne h$, hence
\begin{equation}
 \EE n_h^{\rm dir}=\frac{16-1+25/121}{5}=\frac{368}{121},\qquad
 \EE n_j^{\rm dir}=\frac{16+24/121}{5}=\frac{392}{121}\quad(j\ne h).
\end{equation}
Their sum is sixteen. Equation~\eqref{eq:EH} measures a different specified preparation: the stationary source over one mean-return interval.

\subsection{The response between internal states}
\label{sec:internalgreen}
A change introduced at one internal state can reach another through all the intervening states. Injecting equal and opposite sources compares their response while removing the common constant mode. The inverse of the disagreement on that subspace is
\begin{equation}
 {\mathcal G_0}=(5\Id-A)^+=\frac{7\Id+A}{32}-\frac{3J_{16}}{128}.
 \label{eq:green}
\end{equation}
For a symmetric state $F$, the effective resistance between the singlet vertex and $F$ is
\begin{equation}
 b={\mathcal G_0}_{00}+{\mathcal G_0}_{FF}-2{\mathcal G_0}_{0F}=\frac7{16};
\end{equation}
the singlet-to-vector resistance is $3/8$. To compare two vertices we inject the source $\lvert0\rangle-\lvert F\rangle$; its quadratic response is exactly this effective resistance.

\subsection{The clock overlap and the electric contact}
The phase of a charged state must survive comparison in both orientations. Closed histories carry a Clifford sign; the relative sign of alternative histories determines their interference. A two-plane comparison retains the product of two complementary amplitudes and averages over the ten symmetric states:
\begin{equation}
 f=\frac{r_0^2}{10}=\frac9{250},\qquad
 \nu=\frac f{\EE N_{\rm ret}}=\frac9{4000}.
 \label{eq:fnu}
\end{equation}
The clock source has four consecutive quarter marks for the four directions beside the retained direction. Its fifth mark is the return along that direction, $-1/[4\Phis(1)]=19/4$. The five clock marks are
\begin{equation}
 (q_0,q_1,q_2,q_3,q_h)=\frac14(0,1,2,3,-1/\Phis(1))
 =\frac14(0,1,2,3,19).
 \label{eq:clockmarks}
\end{equation}
Every pair of marks is a possible interval between two readings. Giving their ten positive gaps equal weight defines the Laplace transform
\begin{align}
 M_Q(s)&=\frac1{10}\sum_D e^{-sD},\notag\\
 10M_Q(s)&=3e^{-s/4}+2e^{-s/2}+e^{-3s/4}
 +e^{-4s}+e^{-17s/4}+e^{-9s/2}+e^{-19s/4}.
 \label{eq:MQ}
\end{align}
They have $\EE D=2$ and
\begin{equation}
 \Var D=\frac{61}{16}
 =\frac{(84-12)^2+4\cdot30^2}{(3\cdot16)^2}.
\end{equation}
Reversing the comparison reverses the clock's tilt of an interval. The two normalized interval distributions are $p_\pm(D)=e^{\mp\nu D}/[10M_Q(\pm\nu)]$. Their square-root amplitudes have inner product $\sum_D\sqrt{p_+(D)p_-(D)}$, the Bhattacharyya coefficient. An interval sharply favoured by one orientation and disfavoured by the other contributes little to their common reading. The three complementary directions each carry this overlap, giving
\begin{equation}
 \hover=\left(\sum_D\sqrt{p_+(D)p_-(D)}\right)^3
 =[M_Q(\nu)M_Q(-\nu)]^{-3/2}
 =0.9999710493296864.
 \label{eq:h}
\end{equation}
Its logarithm begins with $-\tfrac32\Var D\,\nu^2$. The variance of the physical intervals controls the loss of shared amplitude. A return repeated $j$ times carries $\hover^j$.

Put $z_j=4(\ii/4)^j$ and
\begin{equation}
 p_j=\begin{cases}-\Phis(z_j),&j\text{ odd},\\ \Fret(z_j),&j\text{ even},\end{cases}
 \qquad N_h=\sum_{j\geq1}\hover^jp_j=0.07708584056314083.
 \label{eq:Nh}
\end{equation}
All $p_j$ are real and positive; $p_1=1/13$ and $p_2=1/6096$. The $j=59$ truncation used by the programs has tail below $10^{-143}$. At $\hover=1$, $N_1=0.0770880771192248$.

An electric encounter compares the phase of a line while preserving its charge. In the plane $\{h,k\}$ the charge is $Q=\widehat n_k-\widehat n_C/3$. For the even-subset step $S\mapsto S\mathbin\triangle g_t$, the all-left-handed charges satisfy $Q(S\mathbin\triangle g_t)+Q(S)=q_t(2[t\in S]-1)$. Since $q_h=0$, a continuation along $h$ emits zero charge; reading the outgoing entry with conjugate orientation preserves the charged line. Of the $296$ stationary departures, exactly $60$ have this direction; their share is $c_L=60/296$. Each participant brings the plane reflection $r_0=3/5$, and the ten symmetric states are its possible encounters. Each equal alternative therefore carries the endpoint contribution $P_j(x)=r_0/10$. The two-source response is
\begin{equation}
 \sigma_{12}=\frac{\partial^2\log Z}{\partial J_1\partial J_2}
 =\sum_xP_1(x)P_2(x)=\frac{r_0^2}{10}=f.
 \label{eq:electrictwosource}
\end{equation}
The factor $c_L$ retains the charge-preserving departures of that encounter. The phase becomes a fact when the comparison closes; the amplitude returned inside the encounter must therefore be included with the direct one.

Normalize the retained endpoint to one. Its internal return contains the $M$ stationary departures, the endpoint itself and its repeated quarter returns $N_h$. The two orientations have the common overlap $\hover$ and couple with amplitude $\sqrt{\nu\hover}$. In these endpoint coordinates the electric inverse propagator is
\begin{equation}
 K_{\mathrm{el}}=
 \begin{pmatrix}1&\sqrt{\nu\hover}\\\sqrt{\nu\hover}&M+1+N_h\end{pmatrix}.
 \label{eq:electricblock}
\end{equation}
Eliminating the second coordinate leaves $1-x$, where
\begin{equation}
 x=\frac{\nu\hover}{297+N_h}=7.573572544734402\times10^{-6}.
\end{equation}
The returned amplitude changes the coefficient at the physical endpoint by $x$. Summing further returns gives $(1-x)^{-1}$, the inverse of the same Schur complement. The two-source strength and the internal return are thus one response. Including the charge-preserving share,
\begin{equation}
 \alpha=\frac{c_Lf}{1-x},\qquad
 \alpha^{-1}=137.035999177096.
 \label{eq:alpha}
\end{equation}
The current contact symmetric under the two orientations uses $-\Phis(z_j)$ at odd $j$ and $[\Fret(z_j)+\Fret(-z_j)]/2$ at even $j$, with unit clock overlap. It gives
\begin{equation}
 N_{\mathrm{con}}=0.07718341131261781.
 \label{eq:Ncon}
\end{equation}

\subsection{The forty squares and the complete incidence spectrum}\label{the-forty-squares-and-the-complete-incidence-spectrum}

The boundary of the five-cube, with antipodal faces identified, has \(16,40,40,20,5\) cells in dimensions zero through four. A \(k\)-cell has \(k\) free coordinates \(I\); antipodal identification reverses its orientation by \((-1)^k\). Its real cochains can consequently be represented by functions of the fixed coordinates obeying

\begin{equation}
 f_I(-x)=(-1)^k f_I(x).
 \label{eq:proof-SD-1}
\end{equation}

Expand in the characters \(\chi_S(x)=\prod_{s\in S}x_s\), where \(S\subset I^c\). The allowed characters have \(|S|\equiv k\pmod2\). Their union \(T=I\cup S\) is an even subset of the five directions. The coboundary moves one coordinate from \(S\) into \(I\), so it preserves \(T\).

Normalize every character to unit norm on the quotient cells. Differencing a fixed coordinate at its two ends gives a factor two, while the quotient cell count halves. Thus every nonzero coboundary coefficient is \(\sqrt2\). For a fixed even set \(T\), of size \(m\), the complete subcomplex is \(\Lambda^\bullet\mathbb R^T\), with

\begin{equation}
 d=\sqrt2\,b_T\wedge,
 \qquad b_T=\sum_{t\in T}\mathbf e_t,
 \qquad d^\dagger d+dd^\dagger=2mI.
 \label{eq:proof-SD-2}
\end{equation}

The last equation is the exterior identity \(\iota_b(b\wedge)+(b\wedge)\iota_b=\|b\|^2I\). In degree one, \(d^\dagger d\) vanishes on the single vector \(b_T\) and equals \(2m\) on its \(m-1\) orthogonal directions.

There are ten sets \(T\) of size two and five of size four. The ten former sets each contribute \(0,4\) to the degree-one spectrum; the five latter sets each contribute \(0,8,8,8\). Therefore

\begin{equation}
 \operatorname{spec}(C_\square^{\mathsf T}C_\square)
       =0^{15},4^{10},8^{15},\qquad \operatorname{rank}C_\square=25.
 \label{eq:proof-SD-3}
\end{equation}

The fifteen zero directions are the vertex gradients. The constant zero-cochain is the \(T=\emptyset\) component; all positive-\(m\) subcomplexes in \eqref{eq:proof-SD-2} are exact. The successive real coboundary ranks are \(15,25,15,5\). In particular, if \(D_\square\) maps square fluxes to cubes,

\begin{equation}
 C_\square\operatorname{grad}=0,\qquad
 D_\square C_\square=0,\qquad
 \ker D_\square=\operatorname{im}C_\square.
 \label{eq:proof-SD-4}
\end{equation}

A real closed flux is thus the flux of a link cochain. The homogeneous Maxwell equation is the second identity. For a quadratic coefficient $k_{\rm field}>0$, varying $k_{\rm field}\|C_\square\theta\|^2/2-j^{\mathsf T}\theta$ gives $k_{\rm field}C_\square^{\mathsf T}C_\square\theta=j$. Multiplication by the gradient transpose gives current conservation.

The global sign is also retained. A transport \(-1\) on every link has holonomy \(+1\) on a square and \(-1\) on the noncontractible five-step loop. This is the flat \(\mathbb Z_2\) holonomy of the antipodal quotient. The square curvature and the global loop record have their explicit places in the same cell complex.

For a non-Abelian transport the same completion keeps the ordering of the links. The connection on an oriented edge transforms as \(U_{yx}\mapsto g_yU_{yx}g_x^{-1}\). On a closed loop the intermediate frame matrices cancel, leaving conjugation at its base vertex. A trace or determinant of its holonomy is consequently invariant. This is precisely the invariance used in the closed-walk expansion Eq.\eqref{eq:eulerproduct}.

With the carried connection of Eq.~\eqref{eq:parallelcurvature}, transport a state by $U_{x+\epsilon a,x}=\exp[-\epsilon\mathcal A_a(x)]+O(\epsilon^2)$. For the positively oriented loop $a,b,-a,-b$, ordered composition gives $U_\square=\Id-\epsilon^2\mathcal F_{ab}+O(\epsilon^3)$, where $\mathcal F=\dd\mathcal A+\mathcal A\wedge\mathcal A=D^2$. The term \(\mathcal A\wedge\mathcal A\) retains the ordering information of the finite loop. Its commutator is the gauge field's self-interaction. Completion around the boundary of a three-cell gives \(D\mathcal F=0\), which reduces to \eqref{eq:proof-SD-4} at first order in a scalar link phase. The continuum curvature thus uses the same parallel composition and the same closed-boundary identity.

The complete incidence result can be collected as
\begin{equation}
 C_\square\,\mathrm{grad}=0,\quad\operatorname{rank}C_\square=25,\quad
 \operatorname{spec}(C_\square^{\mathsf T}C_\square)
 =0^{(15)}\oplus4^{(10)}\oplus8^{(15)}.
 \label{eq:curlsi}
\end{equation}
\subsection{The signed link Hessian, including the cancellation}\label{the-signed-link-hessian-including-the-cancellation}

Orient each of the forty links once, put \(\theta_{yx}=-\theta_{xy}\), and use the signed adjacency \(A_-\). The phased Hermitian adjacency has expansion

\begin{equation}
 {A_-}[\theta]={A_-}+A_{-,1}+A_{-,2}+O(\|\theta\|^3),\quad
 (A_{-,1})_{xy}=i{A_-}_{xy}\theta_{xy},\quad
 (A_{-,2})_{xy}=-\tfrac12{A_-}_{xy}\theta_{xy}^2.
 \label{eq:proof-SCa-1}
\end{equation}

Let \(a_0=1+4u^2\), \(d_-=a_0^2-5u^2\), and \(K_{-,0}=a_0I-u{A_-}\). Since \({A_-}^2=5I\), \(K_{-,0}^{-1}=(a_0I+u{A_-})/d_-\). The quadratic part of the logarithmic determinant is

\begin{equation}
 S_2=-u\operatorname{Tr}(K_{-,0}^{-1}A_{-,2})
       -\frac{u^2}{2}\operatorname{Tr}(K_{-,0}^{-1}A_{-,1}K_{-,0}^{-1}A_{-,1}).
 \label{eq:proof-SCa-2}
\end{equation}

Set \(s_2=\sum_e\theta_e^2\), summing unoriented links once. The elementary traces are

\begin{equation}
 \operatorname{Tr}A_{-,2}=0,\quad
 \operatorname{Tr}({A_-}A_{-,2})=-s_2,\quad
 \operatorname{Tr}A_{-,1}^2=2s_2,\quad
 \operatorname{Tr}({A_-}A_{-,1}^2)=0.
 \label{eq:proof-SCa-3}
\end{equation}

The last equality uses the absence of triangles. The remaining fourth-order walk trace is \(T_4=\operatorname{Tr}({A_-}A_{-,1}{A_-}A_{-,1})\). Compute it through the first variation \(C_1={A_-}A_{-,1}+A_{-,1}{A_-}\) of \({A_-}[\theta]^2\). Expanding its square and using \({A_-}^2=5I\) gives

\begin{equation}
 \operatorname{Tr}C_1^2=2T_4+20s_2.
 \label{eq:proof-SCa-4}
\end{equation}

There is a geometric evaluation of exactly the same trace. Two vertices at distance two have two intervening paths. Their signed products have opposite signs, because the square has Clifford holonomy \(-1\). The first variation of their cancellation is \(i\) times the oriented square flux, up to the sign of one path. A square has two diagonals and each diagonal has two orientations. Hence it contributes four times the square of its flux to \(\operatorname{Tr}C_1^2\). With \(C_\square\) the square incidence matrix,

\begin{equation}
 \operatorname{Tr}C_1^2=4\|C_\square\theta\|^2,
 \qquad T_4=2\|C_\square\theta\|^2-10s_2.
 \label{eq:proof-SCa-5}
\end{equation}

Substitute \eqref{eq:proof-SCa-3} and \eqref{eq:proof-SCa-5} into \eqref{eq:proof-SCa-2}:

\begin{equation}
 \begin{aligned}
 S_2
 &=\frac{u^2}{d_-}s_2
   -\frac{u^2}{2d_-^2}(2a_0^2s_2+u^2T_4)\\
 &=\frac{u^2}{d_-^2}
       (d_--a_0^2+5u^2)s_2
      -\frac{u^4}{d_-^2}\|C_\square\theta\|^2\\
 &=-\frac{u^4}{d_-^2}\|C_\square\theta\|^2.
 \end{aligned}
 \label{eq:proof-SCa-6}
\end{equation}

The individual link-square terms cancel exactly. The remaining response depends on flux. With \(w=u/(1+4u^2)\), the coefficient is

\begin{equation}
 \kappa(u)=\frac{2u^4}{(1+3u^2+16u^4)^2}
          =\frac{2w^4}{(1-5w^2)^2}.
 \label{eq:proof-SCa-7}
\end{equation}

The signed law with reciprocal link phases is
\begin{equation}
 D_-(u;\theta)=(1-u^2)^{24}
 \det[(1+4u^2)\Id-uA_-[\theta]].
 \label{eq:signedlinkdet}
\end{equation}
The Ihara--Bass factor \((1-u^2)^{24}\) carries no link phase, so the same Hessian belongs to the eighty-state signed determinant. At \(u=\pm1/4\), \(w=\pm1/5\) and \(\kappa=1/200\). At \(u=\pm i/4\), \(w=\pm i/3\) and \(\kappa=1/98\). These are substitutions at the count and amplitude phases of the one step.

There is also a uniform remainder. Transposition gives \({A_-}[-\theta]={A_-}[\theta]^{\mathsf T}\), so the determinant and its logarithm through zero are even in all the phases together. On a sufficiently small closed ball where the determinant is nonzero, its fourth derivative is bounded. Taylor's theorem then gives

\begin{equation}
 \log\det(I-w{A_-}[\theta])-\log\det(I-w{A_-})
   =-\frac\kappa2\|C_\square\theta\|^2+O(\|\theta\|^4)
 \label{eq:proof-SCa-8}
\end{equation}

uniformly on that ball. The positive quadratic return coefficient is therefore \(\kappa\|C_\square\theta\|^2/2\). The charge-normalized field coordinate in the main text converts this coefficient to the Maxwell form; its continuation supplies the displayed Lorentz action.

The conversion also fixes the factor from antisymmetric tensor components. In an orthonormal two-form basis, a flux has one component for each unordered pair of axes. The tensor contraction \(F_{\mu\nu}F^{\mu\nu}\) contains both orders of every pair. For the charge-normalized coordinate \(\mathcal A_{\rm em}=\sqrt{4\pi\alpha\kappa}\,\theta\), this gives

\begin{equation}
 \int F_{\mu\nu}F^{\mu\nu}
     =8\pi\alpha\kappa\,\|C_\square\theta\|^2,\qquad
 \frac{\kappa}{2}\|C_\square\theta\|^2
     =\frac1{16\pi\alpha}\int F_{\mu\nu}F^{\mu\nu}.
 \label{eq:proof-SCa-9}
\end{equation}

The first identity states the coordinate and cell metric used in this normalization. The value of \(\alpha\) has already been obtained from the unit-charge source and its Schur return. Thus the link Hessian and the electric current enter the same field action with their displayed normalizations. Its static inverse retains the two endpoint charges and supplies the Coulomb coefficient \(\alpha Q_1Q_2/r\). Replacing the scalar charge by the retained Lie algebra generators gives the non-Abelian current in the same trace convention.

\subsection{One normalization for the field and its source}
A field can be expressed as the rotation it carries or as a canonically normalized amplitude. The source measures the same work in either choice. Before the neutral rotation, use the independent $SU(3)$, $SU(2)$, and $U(1)$ fields. All component fields below are real. For the non-Abelian factors choose Hermitian fundamental generators $t^a$ with $\tr(t^at^b)=\delta^{ab}/2$. The transport connection and curvature are the anti-Hermitian matrices
\begin{equation}
 \mathcal A_X=-\ii\mathcal A_X^at^a,\qquad
 \mathcal F_X=-\ii\mathcal F_X^at^a,\qquad
 \mathcal F_{X\mu\nu}^a\mathcal F_X^{a\mu\nu}
 =-2\tr(\mathcal F_{X\mu\nu}\mathcal F_X^{\mu\nu}).
 \label{eq:gaugecomponents}
\end{equation}
Hypercharge uses its real component and the physical generator $Y$ in the matter representation. The kinetic terms are written in components for all three factors. Put $\mathcal A_X^a=g_XA_{X,{\rm can}}^a$ and $\mathcal F_X^a=g_XF_{X,{\rm can}}^a$, so that
\begin{equation}
 -\frac1{4g_X^2}\mathcal F_{X\mu\nu}^a\mathcal F_X^{a\mu\nu}
 =-\frac14F_{X,{\rm can}\,\mu\nu}^aF_{X,{\rm can}}^{a\mu\nu}.
 \label{eq:canonicalgauge}
\end{equation}
The canonical curvature contains the commutator with coefficient $g_X$, and the canonical covariant derivative contains the same coefficient. Writing the source work as $j_X\mathcal A_X=j_{X,{\rm can}}A_{X,{\rm can}}$ gives $j_{X,{\rm can}}=g_Xj_X$. The same generator coordinate appears in the transported field and in the current that produces it.

\subsection{The local loops on the resolved chart}
The folded cells have forty elementary squares. On the resolved two-sheet
simplex chart the word $stst$ has displacement $2\epsilon(v_s-v_t)$.
The shortest local closed curves there are ten unoriented hexagons. Their
area forms have isotropic sum $5\Id$ on two-forms and give the local
coefficient $\sqrt5u^6F_{\mu\nu}F^{\mu\nu}/2$. The two chart descriptions
retain the exact closed curves of their respective spatial resolutions.

Before pairing the end orientations, the two-form kinetic operator is
\begin{equation}
 \kappa\Id_V\left(\frac35+\frac{*}{\sqrt5}\right),\qquad
 \lambda_\pm=\kappa\left(\frac35\pm\frac1{\sqrt5}\right)>0.
 \label{eq:orientedtwoform}
\end{equation}
The two equally weighted ends have opposite orientation. Their
$F\wedge F$ terms cancel, leaving the common curvature square. The
identities $C_\square\operatorname{grad}=0$ and
$D_\square C_\square=0$ retain local charge conservation and the
homogeneous field equation in either description.

\section{The three directions outside the plane}
\label{sec:step-six}

\subsection{The normalized crossing}
An internal rotation compares how a state occupies one set of directions with how it occupies another. A weak change compares that rotation with the direction already retained by the particle. At fixed exterior grade $N$, the contrast associated with a set of directions $D$ is
\begin{equation}
 H_D=\frac{5|S\cap D|-|D|N}{4}.
\end{equation}
\paragraph*{Theorem: the normalized crossing response.}
Let $X_i=[i\in S]$. Uniformity among the $N$ occupied directions gives $\EE X_i=N/5$ and, for $i\ne j$, $\EE X_iX_j=N(N-1)/20$. Since
\[
 H_D=\frac14\sum_{i=1}^5(5[i\in D]-|D|)X_i,
 \qquad \sum_i(5[i\in D]-|D|)=0,
\]
the common occupation cancels and the covariance is
\begin{equation}
 \EE(H_DH_E\mid N)
 =\frac{N(5-N)}{64}\{5|D\cap E|-|D||E|\}.
\end{equation}
For general direction projectors the overlap is $\tr(P_DP_E)$. To obtain the same formula, diagonalize $P_D$; the off-diagonal part of $P_E$ has zero trace against the resulting occupation-diagonal contrast. Thus the coordinate-set calculation gives the invariant projector pairing.
The complete clock retains the reference state as well as the two nonzero grades. Its observable is $\widetilde N=N+8\Pi_0$. The uniform fifteen-state reference has grade probabilities $2/3,1/3$ and variance $8/9$; its unit clock response therefore normalizes the general variance by $9/8$. Define the two responses by
\begin{equation}
 \widetilde N=N+8\Pi_0,\qquad
 J_{DE}=\operatorname{Cov}(\widetilde N,H_DH_E),\qquad
 C_W=\frac98\operatorname{Var}\widetilde N.
\end{equation}
Every expectation and covariance here uses the same normalized central measure. The zero grade enters as $\widetilde N=8$ and the other grades as $\widetilde N=N$. Consequently $N(5-N)=8-\widetilde N$ on all three grades $0,2,4$, and
\begin{equation}
 \Cov(\widetilde N,N(5-N))=-\Var\widetilde N.
\end{equation}
Combining this identity with the conditional covariance gives, for every central grade distribution with nonzero clock variance,
\begin{equation}
 \frac{J_{DE}}{C_W}
 =\frac{|D||E|-5\tr(P_DP_E)}{64}\frac89
 =\frac{|D||E|-5\tr(P_DP_E)}{72}.
 \label{eq:crossing}
\end{equation}
The common spectral weights have cancelled because the same occupied states carry the response and its clock. For one direction on each side of a division the numerator is one.

The cancellation can also be followed through the spectral parameter. On the fifteen-state reduction at $\chi\geq0$, the grade weights are $(\chi+4)^{-1}$ and $(\chi+8)^{-1}$. The total grade probabilities are $2(\chi+8)/(3\chi+20)$ and $(\chi+4)/(3\chi+20)$. They give
\begin{equation}
 C_W(\chi)=\frac{9(\chi+4)(\chi+8)}{(3\chi+20)^2},\qquad
 J_{\rm pair}(\chi)=\frac{(\chi+4)(\chi+8)}{8(3\chi+20)^2}.
 \label{eq:crossingchi}
\end{equation}
At $\chi=0$ these are $18/25$ and $1/100$; their ratio is still $1/72$. It is unchanged throughout the regular branch.

\paragraph*{Definition: the colour current.}
The colour singlet gives equal amplitudes to the eight traceless rotations of its three-plane. Its retained line reflection is $r_{\rm v}=1/3$. The isotropic source with this total reflection is $T_C=(r_{\rm v}/8)\Id_8=\Id_8/24$. Its squared norm reads the same crossing:
\begin{equation}
 \Tr T_C^\dagger T_C=\frac1{72},\qquad
 \alpha_s=\frac1{\sqrt{72}}=0.117851130197758.
 \label{eq:alphas}
\end{equation}
The signed return gives the same ratio at either phase:
\begin{equation}
 \frac1{72}=-\frac{\Phis(1)}{4[1+\Phis(1)]}
 =-\frac{\Phis(\ii)}{6[1+\Phis(\ii)]}.
 \label{eq:strong_return_identity}
\end{equation}
The squared norm is $\alpha_s^2=1/72$. In the usual field convention $g_s^2=4\pi\alpha_s$. The classical comparison reads this current as $\alpha_s^{(5),\overline{\mathrm{MS}}}(M_Z)$; Sec.~\ref{sec:evolution} gives its subsequent evolution.

\subsection{The phase of the strong vacuum}

For a reciprocal unitary link field, its adjacency has operator norm at most
five. Indeed, summing the norm bound on each of the five neighbours gives
$|\langle v,A_Uv\rangle|\leq5\langle v,v\rangle$. Therefore
\begin{equation}
 (1+4u^2)\Id-uA_U\succeq(1-u)(1-4u)\Id,
 \qquad 0\leq u<1/4.
 \label{eq:positiveconnection}
\end{equation}
The determinant and its reciprocal are positive on this branch. Integration
over a sector with its positive invariant measure preserves positivity;
reversal of the connection gives the sector with opposite charge and equal
weight. The convergent sector sum is
\begin{equation}
 \mathcal Z_{\rm top}(\theta)=\sum_{q\in\mathbb Z}Z_qe^{\ii q\theta},\qquad
 Z_q=Z_{-q}\geq0,\qquad\sum_qZ_q<\infty.
\end{equation}
Its phase response is
\begin{equation}
 \mathcal Z_{\rm top}(0)-\mathcal Z_{\rm top}(\theta)
 =2\sum_{q>0}Z_q(1-\cos q\theta)\geq0.
 \label{eq:thetasectors}
\end{equation}
\subsection{The unit topological sector}
\label{sec:unittopologicalsector}

The unit directions of the five-dimensional distinction space form
$S^4$. At $n\in S^4$, the positive side of the distinction is the
rank-two projector
$P_+(n)=(\Id+\gamma(n))/2$.
A section $\psi$ of this side obeys $P_+\psi=\psi$; its induced
connection is $\nabla\psi=P_+\dd\psi$.
Differentiating $P_+^2=P_+$ gives
$P_+\dd P_+P_+=0$, so the curvature is
\begin{equation}
 F=P_+\,\dd P_+\wedge\dd P_+\,P_+.
 \label{eq:projectorcurvature}
\end{equation}
Choose an oriented orthonormal tangent frame
$v_1,\ldots,v_4$ at $n$, with the orientation of the five Clifford
matrices in Eq.~\eqref{eq:explicitclifford}. Set
$\gamma_a=\gamma(v_a)$ in the following local calculation.
Since $\dd P_+(v_a)=\gamma_a/2$ and tangent distinctions
anticommute with $\gamma(n)$,
\begin{equation}
 F_{ab}:=F(v_a,v_b)
 =\frac12P_+\gamma_a\gamma_bP_+\quad(a\ne b),
 \qquad F_{aa}=0.
 \label{eq:unitsectorcurvature}
\end{equation}
The product $\gamma_1\gamma_2\gamma_3\gamma_4=\gamma(n)$
acts as $\Id$ on the retained side. The Clifford relation then gives
$\tfrac12\epsilon_{abcd}F_{cd}=-F_{ab}$, so $*F=-F$.
For four distinct tangent directions,
$\tr_4(F_{ab}F_{cd})=\epsilon_{abcd}/2$.
Consequently
\begin{equation}
 \begin{gathered}
 \frac14\epsilon^{abcd}\tr_4(F_{ab}F_{cd})=3,
 \qquad \tr_4(F\wedge F)=3\,\mathrm{vol}_{S^4},\\
 \frac1{8\pi^2}\int_{S^4}\tr_4(F\wedge F)
 =\frac{3}{8\pi^2}\frac{8\pi^2}{3}=1.
 \end{gathered}
 \label{eq:unitsectornumber}
\end{equation}
Here the trace on the retained two-dimensional side equals the
ordinary four-spinor trace because $F=P_+FP_+$.
This connection supplies the unit sector. Its smooth neighbourhood keeps the integer charge and has positive invariant measure. The positive closed
determinant then gives this sector positive weight, so $Z_1>0$ in the sector sum.
The term $2Z_1(1-\cos\theta)$ then makes equality in
$\mathcal Z_{\rm top}(0)-\mathcal Z_{\rm top}(\theta)\geq0$
occur exactly at $\theta=0\pmod{2\pi}$.

The positive closed determinant has zero argument. The stationary free-energy minimum therefore has $\bar\theta=0$.

The stationary vacuum minimizes $-\log\mathcal Z_{\rm top}$ on this positive branch.

\section{The last direction}
\label{sec:step-seven}

\subsection{Rest blocks and physical source matrix elements}
What returns must reproduce the state that left. A charged fermion joins its left and right chiral components; a real scalar returns to its radial state; a neutral fermion can return through its charge conjugate. Their rest operators express these different states within Eq.~\eqref{eq:masszero}.

\paragraph*{Definition: the charged rest source.}
For energy-independent positive clock forms $C_L,C_R$ on equal-dimensional chiral spaces and square rest amplitude $A_X$, also independent of $E$, the Hermitian inverse propagator at rest is
\begin{equation}
 \mathcal K_X(E)=\begin{pmatrix}EC_L&-A_X\\-A_X^\dagger&EC_R\end{pmatrix}.
 \label{eq:restpencil}
\end{equation}
\paragraph*{Theorem: charged rest frequencies.}
Congruence by $\diag(C_L^{-1/2},C_R^{-1/2})$ gives equal unit clocks and off-diagonal amplitude $\widehat A_X=C_L^{-1/2}A_XC_R^{-1/2}$. Eliminating either chiral coordinate gives
\begin{equation}
 \det(E^2\Id-\widehat A_X\widehat A_X^\dagger)=0.
 \label{eq:chargednecessity}
\end{equation}
The positive zeros are the singular values of $\widehat A_X$; a scalar entry gives $E_X^2=|A_X|^2/(C_LC_R)$. With the clock fixed, multiplication of the continuing amplitude by $q>0$ multiplies its rest frequency by $q$. This is why the successive grade factors below become mass ratios. Their finite amplitudes retain the source's directions, invariant contractions and returns.

\paragraph*{Theorem: equal endpoint shares of the clock.}
At a nonzero real zero of Eq.~\eqref{eq:restpencil}, let the chiral state be $(x,y)^{\mathsf T}$. Its two equations are $EC_Lx=A_Xy$ and $EC_Ry=A_X^\dagger x$. Multiplication by $x^\dagger$ and $y^\dagger$ makes their right sides complex conjugates; their left sides are real. Consequently
\begin{equation}
 x^\dagger C_Lx=y^\dagger C_Ry,\qquad
 x^\dagger C_Lx+y^\dagger C_Ry=1
 \quad\Longrightarrow\quad
 x^\dagger C_Lx=y^\dagger C_Ry=\frac12.
 \label{eq:endpointclockhalf}
\end{equation}
Each end carries half the complete clock norm. For a scalar channel with $A_X\ne0$, eliminating $y$ gives $S(E)=EC_L-|A_X|^2/(EC_R)$; its derivative at the positive root is $S_E=2C_L$. The second $C_L$ is exactly the clock share of the eliminated end. The Majorana block below retains its heavy diagonal term and yields its own, unequal, component shares by the same full derivative.

\paragraph*{Definition and spectral consequence: bosonic and neutral sources.}
A vector or a real scalar has the second-order inverse propagator
\begin{equation}
 K_B(E)=E^2C_B-\mathcal M_B,\qquad
 Q_B=C_B^{-1/2}\mathcal M_BC_B^{-1/2}.
 \label{eq:bosonblock}
\end{equation}
Here $C_B$ is positive and $\mathcal M_B$ is the nonnegative displacement form of the selected state. The positive rest frequencies are $\sqrt{\lambda(Q_B)}$. If a first-order representation is wanted, the Hermitian matrix
\begin{equation}
 H_B=\begin{pmatrix}0&Q_B^{1/2}\\Q_B^{1/2}&0\end{pmatrix}
\end{equation}
has these positive eigenvalues and their negatives. The residue identity applies directly to $K_B(E)$, with derivative $2EC_B$, or to this first-order representation. At a massless direction the regular spectral variable is $\xi=E^2$, whose clock derivative is $C_B$. Neutral Majorana fields use their symmetric mass matrix and its Takagi factorization, given explicitly in Eq.~\eqref{eq:neutralblock}. These are the usual spin-dependent forms of one analytic inverse propagator.

The massless field directions preserve the selected state; the massive directions change it. In the finite geometry this is the statement that $[X,P_D]$ vanishes or does not vanish. A fermion's own clock generator is among these physical rotations. With perpendicular Clifford directions,
\begin{equation}
 H=\sum_{j=1}^3p_j\gamma_{c_j}+m\gamma_h+m_5\gamma_k,
 \qquad H^2=(|p|^2+m^2+m_5^2)\Id.
 \label{eq:fermiondispersion}
\end{equation}
Its rest mass is $\sqrt{m^2+m_5^2}$. The axial rotation $\exp(\theta\gamma_h\gamma_k/2)$ rotates $(m,m_5)$ and preserves that mass. The familiar relation $E^2=|p|^2+m_0^2$ expresses the null condition of the reference-relative quadratic form with a retained internal component. The numerical value of $m_0$ is the self-interaction amplitude evaluated for the persistent state.

\subsection{A marked source and its clock}
A clock compares the phase now with the phase at a previous occurrence. Give one of the 320 directed transitions the phase $e^{-\ii\vartheta}$ and follow its effect on the continuing state. Equal transition cofactors and Jacobi's identity give the exact polynomial
\begin{equation}
 D_\vartheta(u)=D(u)+(e^{-\ii\vartheta}-1)\frac{uD'(u)}{320}.
\end{equation}
At the continued critical zero,
\begin{equation}
 \log(4u_c(\vartheta))=
 \frac{\ii\vartheta}{320}
 +\frac{18889}{12288000}\vartheta^2+O(\vartheta^3).
 \label{eq:marked}
\end{equation}
The clock rate of any first-order phase source follows from the same stationary mode. Put
\begin{equation}
 (B_\vartheta)_{fe}=B_{fe}e^{-\ii\vartheta X_{fe}},\qquad
 r=l=\frac{\mathbf1}{\sqrt{80}},\qquad B_0r=4r.
\end{equation}
The integer matrix $B^5$ is positive entrywise, with entries between 10 and 24, so its Perron eigenvalue is simple and the continuation $\lambda_{\rm PF}(\vartheta)$ of 4 is analytic. Its derivative and $u_c=1/\lambda_{\rm PF}$ give
\begin{equation}
 \lambda_{\rm PF}'(0)=l^\dagger B'_0r
 =-\frac{\ii}{80}\sum_{\mathrm{arcs}}X_{fe},\qquad
 \left.\frac{\dd}{\dd\vartheta}\log(4u_c)\right|_0
 =\frac{\ii}{320}\sum_{\mathrm{arcs}}X_{fe}.
 \label{eq:markedsourceclock}
\end{equation}
The differential phase rate of the elementary mark is $1/320$ per step. Its finite continuation also changes the modulus. For a real mark $X$ on the transitions, the long weak-mark limit is
\begin{equation}
 (\mathsf P\odot e^{-\ii TX/n})^n
 \longrightarrow \Pi\exp\!\left(-\frac{\ii T}{320}\sum_{\mathrm{arcs}}X\right),
\end{equation}
where $\Pi$ projects onto the stationary state. This fixes the elementary clock used for rest-frequency ratios.

The dimensional calibration is
\begin{equation}
 m_ec^2=0.51099895069(16)\MeV,
 \label{eq:meunit}
\end{equation}

\subsection{Proton and neutron rest sources}
A proton preserves one common state of three colour constituents. The antisymmetric colour contraction fixes which state is returned. With the orientation $\gamma_h\gamma_k\gamma_{c_1}\gamma_{c_2}\gamma_{c_3}=\Id$,
\begin{equation}
 \frac1{3!}\epsilon_{abc}\gamma_{c_a}\gamma_{c_b}\gamma_{c_c}
 =\gamma_{c_1}\gamma_{c_2}\gamma_{c_3}
 =-\gamma_h\gamma_k.
 \label{eq:colourclosure}
\end{equation}
The product closes in the same two-plane that carries the common clock. Each exchanged colour index changes the epsilon sign and the Clifford sign together. The proton's internal return therefore compares one retained colour-singlet state with its own past.

\paragraph*{Definition and projection: the common nucleon state.}
The colour factor is antisymmetric. For three constituents in a symmetric spatial $S$-wave, the spin--isospin factor belongs to the symmetric subspace of
\begin{equation}
 \mathcal H_{\rm si}
 =(\CC^2_{\rm spin}\otimes\CC^2_{\rm isospin})^{\otimes3},
 \qquad
 P_{\rm sym}=\frac1{6}\sum_{\pi\in S_3}U_\pi,
 \label{eq:nucleonsymmetrizer}
\end{equation}
where $U_\pi$ permutes the three constituents. This subspace has dimension twenty and decomposes under spin and isospin as
\begin{equation}
 \operatorname{Sym}^3(\CC^2_{\rm spin}\otimes\CC^2_{\rm isospin})
 \simeq(4,4)\oplus(2,2).
 \label{eq:nucleondecomposition}
\end{equation}
The pair $(4,4)$ has total spin and isospin $3/2$; $(2,2)$ has both equal to $1/2$. With total spin $\mathbf S=\tfrac12\sum_{r=1}^3\boldsymbol\sigma_r$ and $I_3=\tfrac12\sum_{r=1}^3\tau_{3,r}$, define
\begin{equation}
 P_N=P_{\rm sym}\frac{15/4-\mathbf S^2}{3},\qquad
 P_p=P_N(1/2+I_3),\qquad P_n=P_N(1/2-I_3).
 \label{eq:nucleonprojectors}
\end{equation}
The eigenvalues of $\mathbf S^2$ are $15/4$ and $3/4$ on the two summands. Thus $P_N$ is the orthogonal projector onto the rank-four nucleon space $\mathcal H_N=\operatorname{ran}P_N$. Within that space $I_3=\pm1/2$, each with a spin doublet, so
\begin{equation}
 P_pP_n=0,\qquad P_p+P_n=P_N,\qquad
 \rank P_p=\rank P_n=2.
 \label{eq:nucleonprojectorranks}
\end{equation}
The electric charge on $\mathcal H_N$ is $1/2+I_3$. Both rest amplitudes and all nucleon currents below act on these same two doublets.

The colour contraction has returned to the plane $hk$, which contains $h$. It therefore retains the stationary history of $h$. All three colour indices have closed. The plane selects where the state returns; the comparison of each moment with itself is still $\Omega$. Its first contribution counts every retained departure once. The normalized direction clock and its insertion are
\begin{equation}
 C_{\rm dir}=\diag(\sigma,\eta_0,\eta_0,\eta_0,\eta_0),\qquad
 X_h=6\Id+P_h.
\end{equation}
One insertion at each observed departure gives
\begin{equation}
 \frac M5\Tr(C_{\rm dir}X_h)
 =\sum_s C_s(6+\delta_{sh})=6M+C_h=1836.
 \label{eq:protonsourcecount}
\end{equation}
The direction clock weights the contracted source numerator. This is the operation selected by the colour-singlet state. For an open directional pole the same algebra is $S(E,J)=cE-JA$, with zero $JA/c$ and residue $1/c$; it explains the inverse $1/\eta_0$ in an open quark mass. The fully contracted nucleon instead carries $C_{\rm dir}$ in Eq.~\eqref{eq:protonsourcecount}. Its current and its rest amplitude retain the same state.

The three invariant contractions of the retained ends are $\Id,-P,-E_\Omega$. With trace inner product divided by 16, the identities $P^2=\Id$, $E_\Omega^2=4E_\Omega$, $PE_\Omega=-E_\Omega$ and $\Tr P=\Tr E_\Omega=4$ make every diagonal Gram entry one and every off-diagonal entry $-1/4$:
\begin{equation}
 \begin{gathered}
 (G_3)_{ab}=\frac1{16}\Tr(X_a^\dagger X_b),\qquad
 (X_1,X_2,X_3)=(\Id,-P,-E_\Omega),\\
 G_3=\frac{5\Id-J_3}{4}.
 \end{gathered}
 \label{eq:protonpairgram}
\end{equation}
Its common direction has eigenvalue $1/2$ and its two contrasts have eigenvalue $5/4$. The normalized transmission form
\begin{equation}
 F_0=\frac45G_3=\Id-\frac15J_3
\end{equation}
has eigenvalues $1,1,2/5$. For each retained departure the comparison with the present relation contributes
\begin{equation}
 \langle\Omega,(L+\gamma_h\otimes\gamma_s)\Omega\rangle=6+[s=h].
\end{equation}
The rest amplitude on $P_p\mathcal H_N$ compares each stationary departure with the present singlet. The displayed matrix element gives six units for every departure and a seventh when it continues in the retained direction. There are 60 such departures. The three invariant contractions also carry their transmission determinant.

The stationary contribution $6M+C_h$ is the phase rate divided by the elementary rate, as shown in Eqs.~\eqref{eq:EH} and \eqref{eq:markedsourceclock}. The three invariant contractions give $\det F_0=t_0=2/5$. Multiplying every step by $t_0$ changes the continued zero exactly from $1/4$ to $1/(4t_0)$, so its logarithmic displacement is $\log(1/t_0)$. The comparison of a state with itself has eigenvalue $L(\Omega)=6$; its inverse normalizes this displacement:
\begin{equation}
 h_p=-\frac{\log\det F_0}{L(\Omega)}=\frac16\log\frac52.
 \label{eq:gramheight}
\end{equation}
The five vector states each have negative comparison sign, giving the oriented response volume $(-a)^5=-a^5$. The four directions other than $h$ have seven nonempty even subsets: six positive pair states and one negative four-direction state. Their oriented volume is $-a^7$; its first variation per mean return gives $-7a^6/16=-ba^6$. In operator form,
\begin{align}
 V_5(a)&=\det_{\mathcal H_5}(a\,\sgn L)=-a^5,\notag\\
 V_{\bar h}(a)&=\det_{\mathcal H_{\bar h}}(a\,\sgn L)=-a^7,
 \qquad \frac1{16}\frac{\dd V_{\bar h}}{\dd a}=-ba^6.
 \label{eq:protonvolumes}
\end{align}
Here $\mathcal H_5$ is the vector grade and $\mathcal H_{\bar h}$ is the seven-dimensional space just enumerated.

\paragraph*{The complete colour-singlet rest amplitude.}
One closed state carries the phase of its stationary departures, the transmission of its invariant contractions and their signed exterior volumes. The connected phase adds the logarithm of a multiplicative transmission; a specified set of occupied components contributes its determinant. The complete colour-singlet source retains each of these operations once. The transmission is measured per singlet comparison, and the opened volume per mean return. With these common normalizations its rest amplitude is
\begin{align}
 R_p
 &=6M+C_h-\frac16\log\det F_0-a^5-ba^6\notag\\
 &=1836+\frac16\log\frac52-a^5-ba^6\notag\\
 &=1836.152673412052.
 \label{eq:proton}
\end{align}
The complete source can now be written directly on the retained spaces:
\begin{equation}
 \mathcal R_p=
 \sum_{e\in\mathcal A_h}\langle\Omega,(L+\gamma_h\otimes\gamma_{s(e)})\Omega\rangle
 -\frac16\log\det F_0+V_5(a)+\frac1{16}V_{\bar h}'(a),
 \label{eq:completeprotonsource}
\end{equation}
where $\mathcal A_h$ consists of the $296$ retained original departures.
 The four terms in Eq.~\eqref{eq:proton} are respectively
\begin{equation}
 1836,\quad 0.152715121979026,\quad
 -0.0000394395579498235,\quad -0.00000227036928987471.
\end{equation}
They give the charged doublet's scalar matrix element in the common nucleon source. Its current acts on the same colour and spin--isospin state, with the complete normalization of Eq.~\eqref{eq:residue}.

The response of all ten symmetric states can be compared against the same singlet at once. Let $R$ have columns $\lvert F\rangle-\lvert0\rangle$, one for each symmetric vertex, and let $A_P$ be their Petersen adjacency: two two-element sets are adjacent when disjoint. Grounding the common reference gives
\begin{align}
 \Gamma_P=R^{\mathsf T}{\mathcal G_0}R&=\frac{7\Id+7J_{10}+A_P}{32},\notag\\
 \Gamma_P^{-1}&=\frac{24\Id-4A_P-J_{10}}5.
 \label{eq:groundedpetersen}
\end{align}
Indeed $A_P^2=2\Id-A_P+J_{10}$ and $A_PJ_{10}=3J_{10}$; multiplication proves the inverse. The positive inverse on these ten independent differences is the corresponding source form.


The neutron preserves the common colour state of the proton and changes its neutral internal comparison. The extra self-interaction resolves the relative comparisons among the ten symmetric states. Their adjacency is the Petersen graph, and their transmitted contrast is $a^{L_P}$, where $L_P=3\Id-A_P$. Let $P_0=J_{10}/10$ project onto its common mode. The source retains the departure and return on its orthogonal complement and the opened seven-state contact on all ten components:
\begin{equation}
 \mathcal C_n=\Tr\!\left[(\Id-P_0)(a^{L_P})^\dagger a^{L_P}\right]
 -\frac{10c_L}{16}\frac{\dd V_{\bar h}}{\dd a},\qquad
 R_n=\exp[-(\log t_0)\mathcal C_n].
 \label{eq:neutronsourcefunctional}
\end{equation}
The common mode already belongs to the nucleon amplitude, so the projector leaves nine contrasts. Since $\operatorname{spec}L_P=0,2^{(5)},5^{(4)}$,
\begin{equation}
 \Tr\bigl[(\Id-J_{10}/10)a^{2L_P}\bigr]
 =5a^4+4a^{10}.
 \label{eq:neutroncontrasts}
\end{equation}
The powers count the passage out and back through each contrast, and their coefficients count how many independent contrasts have that passage. The retained two-end contact adds the opened seven-state volume on each of the ten symmetric states with the selected-direction share: $10c_Lba^6$. Together they change the rest phase through the reciprocal singlet transmission. The neutron source is
\begin{equation}
 R_n
 =\left(\frac52\right)^{5a^4+4a^{10}+10c_Lba^6}
 =1.001378419564620.
 \label{eq:neutron}
\end{equation}
\paragraph*{Both nucleon masses from the same rest operator.}
On $\mathcal H_N$, the source with the common clock normalized is
\begin{equation}
 A_N=m_eR_p(P_p+R_nP_n),\qquad
 \mathcal K_N(E)=
 \begin{pmatrix}EP_N&-A_N\\-A_N&EP_N\end{pmatrix}
 \quad\hbox{on }\mathcal H_N\oplus\mathcal H_N.
 \label{eq:nucleonoperator}
\end{equation}
The orthogonal projectors in Eq.~\eqref{eq:nucleonprojectorranks} diagonalize the returning amplitude. The determinant is therefore
\begin{equation}
 \det\mathcal K_N(E)
 =(E^2-m_e^2R_p^2)^2(E^2-m_e^2R_p^2R_n^2)^2.
 \label{eq:nucleonrestzero}
\end{equation}
Its positive zeros are the proton and neutron rest energies, each with the spin multiplicity two:
\begin{equation}
 \frac{m_p}{m_e}=R_p,\qquad \frac{m_n}{m_p}=R_n.
 \label{eq:nucleonmassratios}
\end{equation}
The complete clock derivative is $\diag(P_N,P_N)$ on this space. Each chiral end carries half its normalized norm by Eq.~\eqref{eq:endpointclockhalf}. Equation~\eqref{eq:multipletresponse} then evaluates a perturbation on either doublet with the same clock. The two rest energies and the proton Compton length are
\begin{align}
 m_pc^2&=938.272089420197\MeV,&
 m_nc^2&=939.565422025190\MeV,\notag\\
 \lamP&=\frac{\hbar}{m_pc}=0.210308910053202\fm.
\end{align}

\subsection{The three complementary pairings}
The last direction leaves four transverse directions, say $a,b,c,d$.
Their disjoint complete pairings are
\begin{equation}
 (ab|cd),\quad(ac|bd),\quad(ad|bc),\qquad
 \frac{4!}{2^2\,2!}=3.
 \label{eq:threepairings}
\end{equation}
There are three ways to pair the four directions $a,b,c,d$ perpendicular to $h$. Embed them in the ten-dimensional symmetric sector by the isometry
\begin{equation}
 U_h=\frac1{\sqrt2}
 \bigl(\lvert ab\rangle+\lvert cd\rangle,\ \lvert ac\rangle+\lvert bd\rangle,\ \lvert ad\rangle+\lvert bc\rangle\bigr).
 \label{eq:pairingisometry}
\end{equation}
Its columns have disjoint supports. Inside each column the two symmetric states have disjoint direction sets and are adjacent; between different columns their direction sets intersect. Hence
\begin{equation}
 U_h^{\mathsf T}U_h=\Id_3,\quad
 U_h^{\mathsf T}A_PU_h=\Id_3,\quad
 U_h^{\mathsf T}J_{10}U_h=2J_3.
\end{equation}
The complete compressed source is therefore
\begin{equation}
 10U_h^{\mathsf T}\Gamma_P^{-1}U_h=40\Id_3-4J_3.
 \label{eq:pairingcompression}
\end{equation}
It has eigenvalue $28$ for the common pairing and $40$ for each of the two contrasts. The same comparison has supplied the three pairings, their common mode and their independent differences.

The common pairing and its two contrasts are independent. The cumulative
grades of the even exterior comparison are $0,2,4$, with dimensions
$1,10,5$. They give the three successive states of the charged source:
\begin{equation}
 \mathcal H_0,\qquad \mathcal H_0\oplus\mathcal H_2,\qquad
 \mathcal H_0\oplus\mathcal H_2\oplus\mathcal H_4,
 \qquad \Lambda^6\CC^5=0.
 \label{eq:threecumulativegrades}
\end{equation}
The isolated factor $3^{5/2}$ belongs to the last increment in this sequence.
The sequence's rest equation contains it after its earlier grades.

Matter fields have Grassmann coefficients in this representation. Their family space is $\CC^3$, indexed by the three pairings of the four directions beside $h$. The full internal matter space is therefore $\CC^3\otimes\Lambda^{\mathrm{even}}\CC^5$. The three family states are part of the physical matter specification. Within each family the exterior grades have dimensions $1,10,5$.

\subsection{Comparison grades and the charged masses}
The electron, muon and tau carry the same electric charge. What distinguishes their rest phases is how much of the comparison must return unchanged. The elementary charged line retains the reference; its two successive completions retain the ten symmetric states and then the five vector states. The three pair contractions index the family space; the cumulative grades zero, two and four belong to its internal pair representation, $1\oplus10\oplus5$. Each added grade must recover its whole amplitude, its signed endpoint comparison and its common clock.

Let $L_n=L_*\Id_n$ on $(n,L_*)=(10,2)$ or $(5,-2)$. Choose a unit common amplitude $s_n$ and its projector $P_{0,n}=|s_n\rangle\langle s_n|$. In a basis with equal common components, $P_{0,n}=J_n/n$. The even plane source fixes this amplitude and retains the relative components; the odd three-direction source carries the whole grade. Thus, for $p_S=(-1)^{|S|}$ and $\epsilon_S=(1+p_S)/2$, its response projector is
\begin{equation}
 Q_{n,p_S}=\Id_n-\epsilon_SP_{0,n},\qquad
 \Tr Q_{n,p_S}=n-\epsilon_S.
 \label{eq:gradesourceprojector}
\end{equation}
The selected plane and odd three-direction endpoints have $\Psi_+=1/28$ and $\Psi_-=1/2$, the two values of Eq.~\eqref{eq:support}. Their clock pieces are $c_{10}(p_S)=\log M_Q(p_S\nu)/6$ and $c_5(p_S)=(5-\epsilon_S)p_S\nu\,\EE D/4$: the first measures the logarithmic return per unit singlet comparison, and the second adds the mean interval over the retained components. The complete logarithmic source is the sum of these three contractions:
\begin{equation}
 \boxed{\log\mathcal L(n,L_*,p_S)=
 -\frac12\log\det\left|\frac{L_n}{L(\Omega)}\right|
 -\frac{\Psi_{p_S}}4\Tr(Q_{n,p_S}L_n)+c_n(p_S).}
 \label{eq:gradesourcefunctional}
\end{equation}
The first contraction retains the product of all normalized comparisons. The inverse square root converts that product into amplitude size. The second adds the signed responses of the components that remain independent at the two ends. The third compares their elapsed interval. These operations are invariant under a simultaneous unitary change of the grade basis and its common amplitude.

Their evaluation is immediate from the comparison, with $L(\Omega)=6$:
\begin{equation}
 \det|L_n/6|=3^{-n},\qquad
 \Tr(Q_{n,p_S}L_n)=L_*(n-\epsilon_S).
 \label{eq:gradecontractions}
\end{equation}
Since $L_n$ is scalar on its grade, the latter trace depends only on the rank of the retained contrast projector. With these endpoint and clock terms, Eq.~\eqref{eq:gradesourcefunctional} gives
\begin{align}
 \log\mathcal L(n,L_*,p_S)
 &=\frac n2\log3-\frac{\Psi_{p_S} L_*}{4}(n-\epsilon_S)+c_n(p_S),\notag\\
 \Psi_+&=\frac1{28},\qquad \Psi_- =\frac12,\notag\\
 c_{10}(p_S)&=\frac16\log M_Q(p_S\nu),\qquad
 c_5(p_S)=\frac{(5-\epsilon_S)p_S\nu\,\EE D}{4}.
 \label{eq:layers}
\end{align}
The signed return $\Psi_{p_S}$ compares the two endpoints through their intervening history, as in Eq.~\eqref{eq:support}. The endpoint contractions $-1/5$ and $\ii$ give same-orientation products $1/25$ and $-1$; the common rule $\Psi(w)=w/(1+3w)$ gives the displayed $1/28$ and $1/2$. A first-order contribution per step carries $|u|=1/4$. On the ten-dimensional grade the closed clock logarithm is measured per unit of the state's self-comparison, giving $c_{10}=\log M_Q(p_S\nu)/6$. The proton's Gram logarithm uses this same inverse singlet comparison. On the five-dimensional grade the mean interval accumulates over the retained components, giving $c_5$. Changing the parity changes the signed contraction and the interval's orientation within one source. Since the amplitude multiplies over retained components, each additional grade multiplies the rest amplitude in Eq.~\eqref{eq:restpencil}; the common clock makes that multiplier a mass ratio.

The single family expression also has a matrix form. Let
\begin{equation}
 F_{10}=\diag(0,1,1),\qquad F_5=\diag(0,0,1),
\end{equation}
so that a family retains the ten-dimensional grade once $F_{10}=1$ and the five-dimensional grade once $F_5=1$. Here $\kappa_u=(3/2)^{8/5}$ is the ratio of reflection to transmission, accumulated over $16/10$ steps per symmetric-state arrival. Define the two logarithmic increments of Eq.~\eqref{eq:layers},
\begin{align}
 \lambda_{10}^{\log}(p_S)&=5\log3-\frac{\Psi_{p_S}}{2}(10-\epsilon_S)
       +\frac16\log M_Q(p_S\nu),\notag\\
 \lambda_5^{\log}(p_S)&=\frac52\log3+\frac{\Psi_{p_S}}{2}(5-\epsilon_S)
       +\frac{(5-\epsilon_S)p_S\nu}{2}.
\end{align}
The charged-family source matrices are then
\begin{align}
 \log(M_\ell/m_e)&=\lambda_{10}^{\log}(+)F_{10}+\lambda_5^{\log}(+)F_5,\notag\\
 \log(M_d/m_e)&=\frac32\log3\,\Id+\lambda_{10}^{\log}(-)F_{10}
       +\lambda_5^{\log}(-)F_5-\log\eta_0\,(\Id-F_5),\notag\\
 M_u&=R_{ud}M_d\exp\!\left[\log\kappa_u\left(5F_{10}+\frac52F_5\right)\right].
 \label{eq:chargedsource}
\end{align}
The retained grade, direction and divided-set parity determine these diagonal mass forms. A change of preparation or detection basis rotates the source matrices, while the clock form in Eq.~\eqref{eq:restpencil} fixes their normalization.

The electron source retains the elementary mark. The muon source additionally retains the ten-dimensional grade, and the tau source also retains the five-dimensional grade. For the even divided set their rest-frequency ratios are
\begin{align}
 \frac{m_\mu}{m_e}
 &=\mathcal L(10,2,+1)=3^5e^{-9/56}M_Q(\nu)^{1/6}
 =206.768283945476,\notag\\
 \frac{m_\tau}{m_\mu}
 &=\mathcal L(5,-2,+1)=3^{5/2}e^{1/14+2\nu}
 =16.8181606410583.
 \label{eq:leptons}
\end{align}
Thus $m_\mu c^2=105.658376132110\MeV$ and $m_\tau c^2=1776.979542863196\MeV$. The three pairings of the four directions beside the retained direction index the three family source blocks. The factor $3^{5/2}m_ec^2=7.96568530688517\MeV$ by itself is the uncompleted five-dimensional contribution; the charged-lepton source includes its lower retained grades. Record L04 counts the three channels of Eq.~\eqref{eq:chargedsource}; L05 records the absence of that isolated five-dimensional scale from its spectrum. The Koide combination $\sum m/(\sum\sqrt m)^2$ is $0.666667259761091$.

The quark source adds the reference direction to the even plane. The resulting odd three-direction set changes the signed endpoint response from $\Psi_+=1/28$ to $\Psi_-=1/2$. The same ten and five internal components now accumulate their phases with these different return weights. The two factors become
\begin{align}
 {p_q}&=3^5e^{-5/2}M_Q(-\nu)^{1/6}=19.9616523937784,\notag\\
 {q_q}&=3^{5/2}e^{5/4-(5/2)\nu}=54.1038702261733,
\end{align}
and the symmetric-state visit conversion is
\begin{equation}
 \kappa_u=(r_0/t_0)^{\EE N_{\rm ret}/\EE n_F}=(3/2)^{8/5}=1.91313675093869.
\end{equation}
The down-type source includes three colour components and compares their retained direction with weight $\eta_0$. The up-type current preserves that direction in fraction $R_{ud}$ of the stationary states. Its further visits to the symmetric sector convert reflection to transmission by $\kappa_u$. The cumulative retained grade dimensions are $0,10,15$. Each amplitude takes the square root of its grade volume, giving powers $0,5,15/2$. In the native source coordinates, the quark mass sources are
\begin{align}
 (m_d^{\rm native},m_s^{\rm native},m_b^{\rm native})&=3^{3/2}m_e(\eta_0^{-1},{p_q}\eta_0^{-1},{p_q}{q_q}),\notag\\
 (m_u^{\rm native},m_c^{\rm native},m_t^{\rm native})&=R_{ud}(m_d^{\rm native},\kappa_u^5m_s^{\rm native},\kappa_u^{15/2}m_b^{\rm native}).
 \label{eq:quarks}
\end{align}
In MeV they are
\begin{align}
 (m_d^{\rm native},m_s^{\rm native},m_b^{\rm native})&=(2.66422920998645,53.1824173871003,2867.65374846158),\notag\\
 (m_u^{\rm native},m_c^{\rm native},m_t^{\rm native})&=(1.22410531269648,626.246546463788,170949.965605963).
\end{align}
Their invariant combinations include
\begin{equation}
 \frac{m_b^{\rm native}}{m_s^{\rm native}}={q_q}\eta_0=53.9210868808146,\quad
 \frac{m_c^{\rm native}}{m_s^{\rm native}}=R_{ud}\kappa_u^5=11.7754434121621,
\end{equation}
and
\begin{equation}
 \frac{(m_b^{\rm native})^2m_d^{\rm native}}{(m_s^{\rm native})^3}
 =\frac{(m_t^{\rm native})^2m_u^{\rm native}}{(m_c^{\rm native})^3}
 =\frac{{q_q}^2\eta_0^2}{{p_q}}=145.653453585.
 \label{eq:ladderinvariant}
\end{equation}
The up-type current preserves the invariant in Eq.~\eqref{eq:ladderinvariant}: its common $R_{ud}$ cancels, and the extra transmission exponent is $2(15/2)-3(5)=0$. This equality belongs to the common source-scale masses. To read them in an experiment we carry them to its stated energy and subtraction scheme. The running equations in Sec.~\ref{sec:evolution} specify that step.

All charged rest amplitudes in this source are positive. In the common orientation,
\begin{equation}
 \bar\theta=\theta_{\rm vac}+\arg\det(M_uM_d)=0.
 \label{eq:chargedvacuumorientation}
\end{equation}
The sector phase and the charged mass orientation therefore give the same strong vacuum.

\subsection{The weak contact}
A weak encounter changes the direction retained by the particle. Each endpoint contributes the same normalized stationary current $1/M$. Two endpoints changing that direction have nine distinct choices after their coincident choice is removed from the ten symmetric channels. Their factor is $9/10$. The state then returns through the same clock intervals as before. Four steps establish the first closed comparison; the further steps carry the overlap $\hover$. A joint endpoint occupation of probability $2b$ contributes one additional overlap. Averaging these passages gives
\begin{equation}
 W_W=\Fret(\hover)\hover^{-4}(1-2b+2b\hover)=0.9996273895675378.
\end{equation}
The two-end weak source is the projector away from the coincident symmetric channel,
\begin{equation}
 T_W=\frac{\sqrt{W_W}}{M\sqrt{10}}\left(\Id_{10}-\frac{J_{10}}{10}\right),
 \qquad G_Fm_p^2=\|T_W\|_{\rm F}^2.
 \label{eq:weakcontactmatrix}
\end{equation}
The projector has rank nine, so the normalized contact is
\begin{equation}
 G_Fm_p^2=\frac9{10}\frac{W_W}{M^2}
 =1.02682689305\times10^{-5}.
 \label{eq:GF}
\end{equation}
With the proton source already obtained, this gives $G_F=1.16637885870644\times10^{-5}\GeV^{-2}$ and $v=(\sqrt2G_F)^{-1/2}=246.219634042871\GeV$ in units $\hbar=c=1$.

\subsection{The neutral vector metric and the radial source}
The return function at the full and quarter turns gives
\begin{equation}
 \chi_A=\frac1{1+\Phis(\ii)}=\frac{13}{12},\quad
 w_h=\frac1{4\chi_A}=\frac3{13},\quad
 \chi_h=w_h(1-w_h)=\frac{30}{169},
\end{equation}

An internal rotation becomes massive when it changes the actual scalar state against which the particle compares itself. The selected two-plane leaves three complementary directions. Their preserving algebra is $\mathfrak u(2)\oplus\mathfrak u(3)$, of dimension 13. Fix an actual unit vector $h_0$ in the two-plane, including its phase. For anti-Hermitian generators orthonormal in $\operatorname{Re}\Tr(X^\dagger Y)$, its infinitesimal displacement is measured by
\begin{equation}
 \mathcal M(X,Y)=\operatorname{Re}\langle Xh_0,Yh_0\rangle.
 \label{eq:vectorgram}
\end{equation}
The projector onto the real line $\RR h_0\subset\RR^{10}$ measures the same displacement:
\begin{equation}
 \|[X,P_{\RR h_0}]\|_F^2=2\|Xh_0\|^2,
 \qquad \operatorname{Re}\langle h_0,Xh_0\rangle=0.
 \label{eq:realvector}
\end{equation}
A projector onto the complex line instead gives
$2(\|Xh_0\|^2-|\langle h_0,Xh_0\rangle|^2)$: a phase rotation preserves the complex line while changing the vector. The actual scalar state includes this phase direction and its neutral mass.

For $P_h=|h_0\rangle\langle h_0|$ and the real number
$\beta_X=-\ii\langle h_0,Xh_0\rangle$, polarization gives
\begin{equation}
 \Re\langle Xh_0,Yh_0\rangle
 =\frac12\Re\Tr([X,P_h]^\dagger[Y,P_h])+\beta_X\beta_Y.
 \label{eq:weakgrambridge}
\end{equation}
The projector term measures the two charged changes of direction; the last
term retains the actual vector's clock phase.

Use generators orthonormal in $\Re\Tr X^\dagger Y$ for this Gram form. On the thirteen-dimensional algebra it has spectrum $(1,1/2,1/2)$ and ten zeros. The moved fraction has two equivalent readings, $w_h=3/13=(1+\Phis(\ii))/4$: the count of moved rotations and the signed return at the weak phase. On the traceless twelve-dimensional subalgebra the spectrum is $(4/5,1/2,1/2)$ and nine zeros. The two charged generators each have squared displacement $1/2$. The normalized weak diagonal and hypercharge generators have squared displacements $1/2$ and $3/10$, which add to $4/5$ in their one moved neutral combination. With the field normalization $\tr(t^at^b)=\delta^{ab}/2$, the three nonzero Gram eigenvalues are instead $1/4,1/4,2/5$. Dividing by the kinetic form gives the same physical frequencies in either basis. The clock susceptibility is $\chi_h=w_h(1-w_h)=30/169$.

Three complementary clock comparisons give the neutral mixture $\widehat s^2=3N_h$. The weak and abelian fields have different kinetic norms because they compare different components of the same source. Their generalized mass equation uses the abelian norm and couplings
\begin{equation}
 Z_B=\frac{1+f-\widehat s^2}{1-\widehat s^2},\quad
 g_2^2=\frac{4\pi\alpha(1+f)}{(1-f)\widehat s^2},\quad
 g_{YJ}^2=\frac{4\pi\alpha(1+f)}{(1-f)(1-\widehat s^2)}.
 \label{eq:weak_coefficients}
\end{equation}
With $\xi=E^2$, the actual zeros are those of
\begin{equation}
 \det\left[
 \xi\begin{pmatrix}1&0\\0&Z_B\end{pmatrix}
 -\frac{v^2}{4}
 \begin{pmatrix}g_2^2&-g_2g_{YJ}\\-g_2g_{YJ}&g_{YJ}^2\end{pmatrix}
 \right]=0.
 \label{eq:weakpencil}
\end{equation}
Canonical normalization is $B_c=\sqrt{Z_B}B_J$ and $g_{Yc}=g_{YJ}/\sqrt{Z_B}$. It leaves a massless combination and gives
\begin{align}
 s_{\mathrm{OS}}^2&=\frac{\widehat s^2}{1+f}=0.223221546032261,\notag\\
 M_W^2&=\frac{\pi\alpha v^2}{(1-f)s_{\mathrm{OS}}^2},\qquad
 M_Z^2=\frac{M_W^2}{1-s_{\mathrm{OS}}^2},\notag\\
 M_W&=80.3662510403325\GeV,\qquad
 M_Z=91.1853615183457\GeV.
 \label{eq:weakmasses}
\end{align}
The three measured angle conventions use $\widehat s^2=0.231257521689422$, $s_{\mathrm{OS}}^2$ above, and $s_{\mathrm{eff}}^2=3N_{\mathrm{con}}=0.231550233937853$. The static electric response has $e_0^2=4\pi\alpha$; the vector-pole normalization has $e_{\mathrm{pole}}^2=4\pi\alpha/(1-f)$, giving the endpoint correction $\Delta r=f$.

For full $V-A$ currents the charged vertex is $g_2/(2\sqrt2)$ and $g_2^2/(8M_W^2)=G_F/\sqrt2$. With signature $(-+++)$, the term $-M_W^2W^+_\mu W^{-\mu}$ gives the eliminated contact $+(G_F/\sqrt2)j^+_\mu j^{-\mu}$. The normalization with left-only currents absorbs a factor four into their coefficient.

The Higgs excitation changes the length of the scalar vector. Its phase direction supplies a neutral vector mass; its radial direction supplies a scalar mass. The same fluctuation that counted the moved rotations fixes the inverse radial susceptibility. Explicitly, the Bernoulli cumulant function and its Legendre transform are
\begin{align}
 W_h(j)&=\log(1-w_h+w_he^j),\notag\\
 I_h(x)&=x\log\frac{x}{w_h}+(1-x)\log\frac{1-x}{1-w_h},\qquad
 I_h''(w_h)=\frac1{\chi_h}.
 \label{eq:radialsusceptibility}
\end{align}
Here $0\leq x\leq1$, with $0\log0$ defined by continuity; $W_h$ generates the Bernoulli cumulants. The physical scalar source uses the quadratic inverse susceptibility in $r_\varphi=\varphi^\dagger\varphi-v^2/2$. The specified local scalar interaction uses that susceptibility in the quartic potential:
\begin{equation}
 V_H(\varphi)=\lambda_H(\varphi^\dagger\varphi-v^2/2)^2,
 \qquad \lambda_H=\frac{\pi\alpha}{\chi_h}=0.129145908531315.
\end{equation}
The top source retains the same directional form $R_h=\diag(\sigma,\eta_0,\eta_0,\eta_0,\eta_0)$. Its radial current compares the complementary matrix element with the selected one:
\begin{equation}
 y_t^2=\frac{\langle c|R_h|c\rangle}{\langle h|R_h|h\rangle},\qquad
 A_t^{\rm rad}=\frac{v y_t}{\sqrt2}.
 \label{eq:radialfermionsource}
\end{equation}
The scalar Hessian and the positive zero of this fermion's rest pencil give
\begin{equation}
 m_H=v\sqrt{2\lambda_H}=125.134771593402\GeV,\qquad
 y_t^2=\frac{\eta_0}{\sigma}=\frac{59}{60},\qquad
 m_t^{\mathrm{rad}}=\frac{y_tv}{\sqrt2}=172.646613614442\GeV.
 \label{eq:radial}
\end{equation}
Writing $\varphi=(0,(v+H)/\sqrt2)^{\mathsf T}$ gives the displayed radial mass from the coefficient of $H^2/2$. The top Yukawa coefficient compares the stationary weight of a complementary direction with that of the retained direction: $\eta_0/\sigma=59/60$. The resulting radial top response and the native family mass of Eq.~\eqref{eq:quarks} refer to their stated source operators and retain their separate mass conventions in Sec.~\ref{sec:quarkconversions}.

\section{The whole}
\label{sec:step-eight}
\label{sec:whole}

\subsection{The nucleon determinant and the entire incoming state}
\label{sec:logtransfer}

An entire incoming system carries the history by which its three spatial directions and its clock have remained related. To compare that history with the receiving system, retain its transmission in each complementary direction. Let $E_h$ select remembered direction $h$ and $\Pi_{C,h}$ the 74 states retained in Eq.~\eqref{eq:retainedschur}. On the eighty-state space the directional weights and their logarithmic transfer are
\begin{equation}
 D_h=\sigma E_h+\eta_0(\Id-E_h),\qquad
 V_{\mathrm{in}}=(\log t_0)\mathsf P+E_h\mathsf P \Pi_{C,h}\log D_h.
 \label{eq:incoming}
\end{equation}
An incoming transition with last direction $h$ has a preceding direction different from $h$. Its stationary logarithmic rate is therefore
\begin{equation}
 \bar v=\frac1{80}{\bf1}^{\mathsf T}V_{\mathrm{in}}{\bf1}
 =\log t_0+\frac{60}{320}\log\eta_0.
\end{equation}
The receiving participant meets all three complementary directions of the incoming system. Each carries sixteen steps on average. To see what a persistent line transmits, let $a_n$ be its successive original departures and let $v(a_n)$ be $\log t_0$, with the extra $\log\eta_0$ on the $60$ marked retained departures along $h$. The product along $N$ steps and its rate obey
\begin{equation}
 T_N=\exp\!\left[\sum_{n=1}^Nv(a_n)\right],\qquad
 \frac1N\log T_N\longrightarrow
 \sum_a\pi(a)v(a)
 =\log t_0+\frac{60}{320}\log\eta_0=\bar v
 \quad\hbox{almost surely}.
 \label{eq:typicallogtransfer}
\end{equation}
Here $\pi(a)=1/320$. The finite walk is irreducible and aperiodic because $B^5$ is strictly positive; its ergodic theorem gives the stated limit. This proves the rate along a persisting stationary history. The incoming-frame source exponentiates that rate over one mean-return interval and takes its determinant over the three complementary directions:
\begin{equation}
 W=\det_{\CC^3}(e^{16\bar v}\Id)
 =t_0^{48}\eta_0^9=7.68514963665086\times10^{-20}.
 \label{eq:W}
\end{equation}
The proton determinant contribution and this entire incoming amplitude are directly related:
\begin{equation}
 h_p=\frac16\log\frac52,\qquad
 -\log W=288h_p+9\log\frac{296}{295}.
 \label{eq:protonbridge}
\end{equation}
The factor $288=L(\Omega)\cdot3\cdot16$ converts the logarithmic displacement per unit singlet comparison into transmission through three complete mean returns. The additional nine stationary weights are the three selected-direction exposures in each of the three complementary directions. The transfer in Eq.~\eqref{eq:typicallogtransfer} supplies the logarithmic rate retained by an entire incoming frame. The tilted-history mean instead has cumulant generator $\log\,r_{\rm spec}(B_s/4)$, with derivative $\bar v$ and higher derivatives given by fluctuations of the accumulated logarithm. Each quantity has its specified operation on the same history law. Varying the return variable gives the source's continuation in terms of the logarithmic return length:
\begin{equation}
 W(u)=\eta_0^{-3/5}\exp\!\left[
 3\left(\log t_0+\frac{\log\eta_0}{5}\right)
 \frac{4u\Fret'(4u)}{\Fret(4u)}\right].
\end{equation}
At $u=1/4$ it equals Eq.~\eqref{eq:W}; its logarithmic derivative with respect to $4u-1$ is $3\cdot156(\log t_0+\log\eta_0/5)$.

Gravity compares the incoming frame with the receiving frame at both ends of the interaction. Each endpoint carries the amplitude $W$. Their product is the dimensionless strength $W^2$, and its ratio to the proton mass fixes the unreduced Planck mass:
\begin{align}
 \alpha_G&=\frac{Gm_p^2}{\hbar c}=W^2=5.90615249377146\times10^{-39},\notag\\
 \MP&=\frac{m_p}{W}=1.22088981188542\times10^{19}\GeV,
 \qquad \ell_P=W\lamP.
 \label{eq:gravitystrength}
\end{align}
All dimensional quantities in Eq.~\eqref{eq:gravitystrength} use the same electron calibration as the finite masses.

\subsection{Transport of the complete frame}
The endpoint frame is transported on the sphere of distinctions. For unit
$r,s\in\RR^5$ with $s\ne-r$, the short-arc spin transport is
\begin{equation}
 U_{sr}=\frac{\Id+\gamma(s)\gamma(r)}{\sqrt{2(1+s\cdot r)}},\qquad
 U_{sr}\gamma(r)U_{sr}^\dagger=\gamma(s).
 \label{eq:sphere_transport}
\end{equation}
Anticommutation gives $U_{sr}^\dagger U_{sr}=\Id$ and
$U_{rs}=U_{sr}^{-1}$. Two successive arcs also compare antipodal endpoints.
For $s=r+\dd r$, its linear term is $\gamma(\dd r)\gamma(r)/2$;
the same carried comparison therefore supplies the spin connection of the frame.

\subsection{The four-spinor trace and the Einstein--Cartan variations}\label{the-four-spinor-trace-and-the-einsteincartan-variations}

The four tangent directions beside \(h\) and their six rotations give the connection

\begin{equation}
 \mathcal A=\frac14\omega^{ab}\gamma_a\gamma_b
      +\frac1{2\ell}\vartheta^a\gamma_h\gamma_a,
 \qquad R^{ab}=d\omega^{ab}+\omega^a{}_c\wedge\omega^{cb},
 \qquad T^a=d\vartheta^a+\omega^a{}_b\wedge \vartheta^b.
 \label{eq:proof-SE-1}
\end{equation}

Indices \(a,b\) run over the four tangent axes, with full ordered antisymmetric sums. Since \(\gamma_h\gamma_a\gamma_h\gamma_b=-\gamma_a\gamma_b\) for distinct tangent axes, the square of the coframe part lies in the six rotations. Computing \(d\mathcal A+\mathcal A\wedge\mathcal A\) yields

\begin{equation}
 \mathcal F=\frac14(R^{ab}-\ell^{-2}\vartheta^a\wedge \vartheta^b)\gamma_a\gamma_b
       +\frac1{2\ell}T^a\gamma_h\gamma_a.
 \label{eq:proof-SE-2}
\end{equation}

The coefficient of \(R-\vartheta\vartheta/\ell^2\) and the torsion coefficient are fixed by the same connection.

Choose the orientation with \(\operatorname{tr}_4(\gamma_h\gamma_a\gamma_b\gamma_c\gamma_d)=4\epsilon_{habcd}\). Products with the torsion component have zero trace after multiplication by \(\gamma_h\). With \(R_{\rm sp}=R^{ab}\gamma_a\gamma_b/4\), the remaining trace is

\begin{equation}
 \begin{aligned}
 \operatorname{tr}_4[\gamma_h(\mathcal F\wedge\mathcal F
                       -R_{\rm sp}\wedge R_{\rm sp})]
  &=-\frac1{2\ell^2}\epsilon_{abcd}R^{ab}\wedge \vartheta^c\wedge \vartheta^d\\
  &\quad+\frac1{4\ell^4}\epsilon_{abcd}\vartheta^a\wedge \vartheta^b\wedge \vartheta^c\wedge \vartheta^d\\
  &=-\ell^{-2}(R-6/\ell^2)\operatorname{vol}.
 \end{aligned}
 \label{eq:proof-SE-3}
\end{equation}

The last step uses \(\epsilon R\vartheta\vartheta=2R\operatorname{vol}\) and \(\epsilon\vartheta\vartheta\vartheta\vartheta=24\operatorname{vol}\). The ordinary four-spinor trace has supplied the factor \(1/4\) multiplying the epsilon contraction. The trace prefactor \(-\ell^2/(16\pi G)\) therefore gives

\begin{equation}
 I_g=\frac1{16\pi G}\int(R-2\Lambda)\operatorname{vol},
 \qquad \Lambda=\frac3{\ell^2}.
 \label{eq:proof-SE-4}
\end{equation}

The subtracted \(\epsilon R\wedge R\) is the Euler term. Its connection variation is \(2\epsilon D\delta\omega\wedge R\), a boundary term because \(DR=0\). It leaves the local coframe and connection equations unchanged.

For an explicit independent variation, use the first-order form

\begin{equation}
 I_g=\frac1{32\pi G}\int\epsilon_{abcd}
       \left(\vartheta^a\wedge \vartheta^b\wedge R^{cd}
             -\frac\Lambda6\vartheta^a\wedge \vartheta^b\wedge \vartheta^c\wedge \vartheta^d\right).
 \label{eq:proof-SE-5}
\end{equation}

Define the matter three-forms by

\begin{equation}
 \delta I_m=-\int\delta \vartheta^a\wedge\mathfrak t_a
             -\frac12\int\delta\omega^{ab}\wedge\mathfrak s_{ab}
\end{equation}

on the matter equations. The displayed signs define the stress three-form \(\mathfrak t_a\) and spin three-form \(\mathfrak s_{ab}\). There are two equal coframe variations in the curvature term and four in the volume term. The coframe equation is consequently

\begin{equation}
 \epsilon_{abcd}\vartheta^b\wedge
       (R^{cd}-\Lambda \vartheta^c\wedge \vartheta^d/3)=16\pi G\mathfrak t_a.
 \label{eq:proof-SE-6}
\end{equation}

For the connection, \(\delta R^{ab}=D\delta\omega^{ab}\). Integration by parts replaces \(D(\vartheta^a\wedge \vartheta^b)\) by \(T^a\wedge \vartheta^b-\vartheta^a\wedge T^b\). The two antisymmetric contributions agree. Reordering the resulting three-form past \(\delta\omega\) gives

\begin{equation}
 \epsilon_{abcd}T^c\wedge \vartheta^d=8\pi G\mathfrak s_{ab}.
 \label{eq:proof-SE-7}
\end{equation}

This equation is algebraic in the torsion. For an invertible coframe, zero spin current gives \(T^a=0\); a spin current fixes the corresponding contorsion through the same linear equation. Thus the vertical part of the curvature remains in the field variation.

With the metric stress definition \(T_{\mu\nu}=-2(\sqrt{-g})^{-1}\delta I_m/\delta g^{\mu\nu}\), the torsion-free coframe equation reads

\begin{equation}
 R_{\mu\nu}-\frac12Rg_{\mu\nu}+\Lambda g_{\mu\nu}
       =\frac{8\pi G}{c^4}T_{\mu\nu}.
 \label{eq:proof-SE-8}
\end{equation}

For a constant energy density, \(I_\Lambda=-\int\rho_\Lambda\sqrt{-g}\,d^4x\). Varying \(\sqrt{-g}\) gives \(T_{\mu\nu}^\Lambda=-\rho_\Lambda g_{\mu\nu}\). In the \((-+++)\) rest frame this has energy density \(\rho_\Lambda\) and pressure \(-\rho_\Lambda\), with \(\rho_\Lambda=\Lambda c^4/(8\pi G)\).

The frame variation with $\Lambda=3/\ell^2$ gives the weak-field equation for matter energy density $\rho_m$:
\begin{equation}
 \nabla^2\Phi_G=4\pi G\rho_m/c^2-\Lambda c^2,\qquad
 \Phi_G(r)=-\frac{Gm}{r}-\frac{\Lambda c^2r^2}{6}.
 \label{eq:weakfieldgravity}
\end{equation}

\subsection{The circulation moves the long-range zero}
\label{sec:circulation}
Orient every unsigned edge once and write a real link phase as
$\theta=\operatorname{grad}\varphi+\theta_{\rm cyc}$, with
$\operatorname{grad}^{\mathsf T}\theta_{\rm cyc}=0$. A change of vertex
phase removes the gradient part. The normalized constant state of the
sixteen-vertex adjacency is $v=\mathbf1/4$. For the divergence-free part,
the first adjacency variation has $A_1v=0$. The second variation gives
\begin{equation}
 v^\dagger A_2v=-\frac1{16}\sum_e\theta_{{\rm cyc},e}^2.
 \label{eq:circulationperturbation}
\end{equation}
The usual second-order eigenvector term contains $A_1v$ and vanishes.
Conjugating all phases changes $A_\theta$ to its transpose, so its leading
eigenvalue is even. Its separated Perron branch therefore has
\begin{equation}
 \lambda_{\max}(A_\theta)
 =5-\frac{\|\theta_{\rm cyc}\|^2}{16}+O(\|\theta_{\rm cyc}\|^4).
 \label{eq:circulationperron}
\end{equation}
The zero satisfies $u_c^{-1}+4u_c=\lambda_{\max}$. Differentiating this
equation in $\log u$ at $u=1/4$ gives $-3$, hence
\begin{equation}
 \log(4u_c)=\frac{\|\theta_{\rm cyc}\|^2}{48}
                 +O(\|\theta_{\rm cyc}\|^4),
 \qquad48=3\EE N_{\rm ret}=2(40-16).
 \label{eq:circulation48}
\end{equation}
The finite resolvent bound on this branch is
\begin{equation}
 \left|\log(4u_c)-\frac{\|\theta_{\rm cyc}\|^2}{48}\right|
 \leq\frac{\|\theta_{\rm cyc}\|^4}{20},
 \qquad\|\theta_{\rm cyc}\|^2\leq\frac1{100}.
 \label{eq:circulationbound}
\end{equation}
The unperturbed adjacency gap is four. Resolvent expansion on the
orthogonal complement of the constant state, together with
$|e^{i\theta}-1-i\theta+\theta^2/2|\leq|\theta|^3/6$, controls the
remainder uniformly on the displayed ball. The two-dimensional relation
between $u$ and $\lambda$ then gives the stated bound.

\subsection{The common clock of the complete state}
Let ${\mathcal N}$ be the local proper-time rate, $\dd\tau={\mathcal N}\dd t$. A stationary state of proper energy $E_*$ has phase $\exp(-\ii E_*{\mathcal N}t/\hbar)$ and coordinate energy $E={\mathcal N}E_*$. Referring every departure, internal return and arrival to that common clock gives
\begin{equation}
 K_X(E,{\mathcal N},J)=K_X(E/{\mathcal N},1,J),\qquad
 (K_X)_{\mathcal N}=-\frac E{\mathcal N}(K_X)_E.
 \label{eq:lapse}
\end{equation}
\paragraph*{Theorem: one proper-time change for the complete state.}
Exact Schur elimination preserves the relation because the energy derivative includes the full internal term in Eq.~\eqref{eq:selfenergyderivative}. In the Hermitian rest description, Eq.~\eqref{eq:fullstatenorm} identifies its denominator with the norm of the complete state. The lapse insertion has the same norm multiplied by $-E/{\mathcal N}$. Applying the simple-zero identity therefore yields
\begin{equation}
 \frac{\dd E_X}{\dd {\mathcal N}}
 =-\frac{l^\dagger(S_X)_{\mathcal N}r}{l^\dagger(S_X)_Er}
 =\frac{E_X}{{\mathcal N}},\qquad E_X({\mathcal N})={\mathcal N}E_X(1).
 \label{eq:fullclock}
\end{equation}
At ${\mathcal N}=1+\Phi_G/c^2$, the first energy shift is $E_X(1)\Phi_G/c^2=m_X\Phi_G$. The energy of the motion and the gravitational source are the same continuing phase. A composite's binding energy participates through its internal propagation in $S_E$; moving the common clock before or after elimination gives the same answer. The electron mark, the charged grades, the nucleon contractions, the vector displacement and the neutral return can all be compared with this one normalization.

\subsection{The retained nine-dimensional determinant}

The coherent comparison on the sixteen even exterior states is
\begin{equation}
 {\mathcal D}=\frac{|L|}{L(\Omega)}
  =\Pi_0+\frac13(\Id-\Pi_0).
 \label{eq:coherentcomparison}
\end{equation}
Order the five complex directions as $C\oplus\CC h\oplus\CC k$.
Preserving the colour--plane division, determinant one and the actual vector
$h$ gives
\begin{equation}
 \mathfrak k=\{X(A)=\diag(A,0,-\Tr A):A\in\mathfrak u(3)\},
 \qquad \dim_{\RR}\mathfrak k=3+2\binom32=9.
 \label{eq:neutralalgebra}
\end{equation}
On the exterior space, $c_a^\dagger$ adds coordinate direction $a$ and $c_a$ contracts it; $\{c_a,c_b^\dagger\}=\delta_{ab}$. The retained algebra acts by $\dd\Lambda(X)=\sum_{a,b}X_{ab}c_a^\dagger c_b$.
The action annihilates the vacuum and preserves degree. Therefore
\begin{equation}
 \dd\Lambda(X)\Omega=0,\qquad \Pi_0\dd\Lambda(X)=0,
 \qquad {\mathcal D}\dd\Lambda(X)=\frac13\dd\Lambda(X).
 \label{eq:retainedcomparison}
\end{equation}
The Hilbert--Schmidt inner product is positive on this faithful representation.
Its two Gram forms satisfy
\begin{align}
 g(X,Y)&=\Tr\dd\Lambda(X)^\dagger\dd\Lambda(Y),\notag\\
 g(X,\mathcal EY)&=\Tr\dd\Lambda(X)^\dagger {\mathcal D}\dd\Lambda(Y)
                =\frac13g(X,Y).
 \label{eq:neutralpullback}
\end{align}
Hence $\mathcal E=\Id_{\mathfrak k}/3$ on the central coordinate and the eight
traceless coordinates alike.

This evaluation has a direct finite matrix check. For colour indices
$c,d\in\{1,2,3\}$ and $k=5$, use the nine real basis matrices
\begin{equation}
 \ii(E_{cc}-E_{55}),\quad
 E_{cd}-E_{dc},\quad \ii(E_{cd}+E_{dc})\quad(c<d).
 \label{eq:neutralbasis}
\end{equation}
Their unweighted and weighted Gram matrices are
\begin{equation}
 {\mathsf G_r}=\begin{pmatrix}8&4&4\\4&8&4\\4&4&8\end{pmatrix}\oplus8\Id_6,
 \qquad {\mathsf H_r}=\frac {\mathsf G_r}3,
 \qquad \det {\mathsf G_r}=2^{26}.
 \label{eq:neutralgram}
\end{equation}
The complete alternating response is consequently
\begin{equation}
 \Lambda^9_{\RR}\mathcal E
 =\frac{\det {\mathsf H_r}}{\det {\mathsf G_r}}=3^{-9},
 \qquad M_R=M_{\rm Pl}\Lambda^9_{\RR}\mathcal E.
 \label{eq:neutraldeterminant}
\end{equation}
The real determinant line is one-dimensional and neutral under the connected
retained group. It supplies the heavy neutral entry of the following mass
equation.

\subsection{Neutral rest blocks and their complete currents}
\label{sec:neutralreturn}

A neutral state can meet its charge conjugate. Its return retains the amplitude $r_{\rm v}=1/3$ in each of the nine real generators of the algebra that keeps the colour space. Its mass is
\begin{equation}
 \MR=\MP r_{\rm v}^9=6.20276285060927\times10^{14}\GeV.
 \label{eq:neutralthresholds}
\end{equation} The neutral endpoint reads the amplitude square root of the vector-grade factor $k_4(u)=1+3u+4u^2$. At the odd quarter,
\begin{equation}
 k_4(\ii/4)=\frac34(1+\ii)
 =\frac{3\sqrt2}{4}e^{\ii\pi/4},\qquad
 \sqrt{k_4(\ii/4)}=\sqrt{\frac{3\sqrt2}{4}}\,e^{\ii\pi/8}.
 \label{eq:neutralhalfphase}
\end{equation}
Taking the square root halves the phase from $\pi/4$ to $\pi/8$. The ratio of its imaginary and real quadratures is $\tan(\pi/8)=\sqrt2-1$. At the even quarters, $k_4(1/4)=2$ and $k_4(-1/4)=1/2$; the amplitude square root is real and its imaginary quadrature is zero. The three neutral source states retain this zero quadrature, the odd-quarter quadrature ratio, and the full selected-direction amplitude. Their Dirac entries are
\begin{equation}
 (m_{D1},m_{D2},m_{D3})=\frac{\sigma v}{\sqrt2}(0,\sqrt{\rho_\nu},1),
 \qquad \rho_\nu=3-2\sqrt2=\tan^2(\pi/8).
\end{equation}
For each family the symmetric neutral block\cite{Minkowski} is
\begin{equation}
 \mathcal M_i=\begin{pmatrix}0&m_{Di}\\m_{Di}&\MR\end{pmatrix}.
 \label{eq:neutralblock}
\end{equation}
The full six-state source is the direct sum of these pairs, written in the three active and three singlet coordinates as
\begin{equation}
 \mathcal M_\nu^{(6)}=
 \begin{pmatrix}0&D_D\\D_D^{\mathsf T}&\MR\Id_3\end{pmatrix},
 \qquad D_D=\diag(m_{D1},m_{D2},m_{D3}).
 \label{eq:neutralfull}
\end{equation}
Each two-state block has eigenvalues $-m_i$ and $\MR+m_i$, with
\begin{equation}
 m_i(\MR+m_i)=m_{Di}^2,\qquad
 m_i=\frac{2m_{Di}^2}{\MR+\sqrt{\MR^2+4m_{Di}^2}}.
 \label{eq:neutralroots}
\end{equation}
The stable positive-root expression avoids subtraction of the nearly equal heavy quantities. It gives
\begin{align}
 (m_1,m_2,m_3)&=(0,.00861267197365250,.0501983309477775)\eV,\notag\\
 \Delta m_{21}^2&=7.41781185257386\times10^{-5}\,\mathrm{eV}^2,\notag\\
 \Delta m_{31}^2&=2.51987242994258\times10^{-3}\,\mathrm{eV}^2,\notag\\
 \Delta m_{32}^2&=2.44569431141684\times10^{-3}\,\mathrm{eV}^2,\notag\\
 \sum_i m_i&=.0588110029214298\eV.
 \label{eq:neutrinovalues}
\end{align}
The same quadratic equation relates the two nonzero light masses exactly:
\begin{equation}
 \frac{m_2}{m_3}=\rho_\nu\frac{\MR+m_3}{\MR+m_2}.
 \label{eq:neutralratioexact}
\end{equation}
The partner participating in the return shifts with the light state. The light mass and the heavy return must satisfy the same quadratic equation; consequently, in the light limit $m_2/m_3\to\rho_\nu$ and $\Delta m_{31}^2/\Delta m_{21}^2\to\rho_\nu^{-2}=17+12\sqrt2$. The exact squared-ratio correction is
\begin{equation}
 \frac{(m_3/m_2)^2}{\rho_\nu^{-2}}-1=-1.34087534783118\times10^{-25}.
\end{equation}

The heavy return makes the common clock particularly visible. Eliminate the heavy coordinate from $E\Id-\mathcal M_i$:
\begin{equation}
 S_i(E)=E-\frac{m_{Di}^2}{E-\MR},\qquad
 S_i'(E)=1+\frac{m_{Di}^2}{(E-\MR)^2}.
 \label{eq:neutrinoschur}
\end{equation}
At the light eigenvalue $E=-m_i$, use $m_{Di}^2=m_i(\MR+m_i)$. The pole residue is
\begin{equation}
 \frac1{S_i'(-m_i)}=\frac{\MR+m_i}{\MR+2m_i}.
 \label{eq:neutrinoresidue}
\end{equation}
The same internal return has fixed the light mass and the fraction of its state accessible to the active current. Put
\begin{equation}
 c_i^2=\frac{\MR+m_i}{\MR+2m_i},\qquad
 s_i^2=\frac{m_i}{\MR+2m_i}.
\end{equation}
Indeed, the full light state with unit retained coordinate is
\begin{equation}
 \psi_i=\begin{pmatrix}1\\-m_{Di}/(\MR+m_i)\end{pmatrix},\qquad
 \|\psi_i\|^2=1+\frac{m_i}{\MR+m_i}
 =S_i'(-m_i)=c_i^{-2}.
 \label{eq:neutralfullnorm}
\end{equation}
Its normalized form is $(c_i,-s_i)^{\mathsf T}$. The heavy coordinate's norm is exactly the extra time in the Schur derivative of the active state, as in Eq.~\eqref{eq:returnnorm}.
Let $U$ be the unitary change from charged-lepton flavour coordinates to neutral source coordinates. The Takagi phase of the negative real eigenvalue gives light and heavy charged-current matrices $U\diag(c_i)$ and $-\ii U\diag(s_i)$, up to a common convention. Their completeness is exact because $c_i^2+s_i^2=1$. The invisible-width factor of the light states is
\begin{equation}
 \sum_i c_i^4(1-4m_i^2/M_Z^2)^{3/2}
 \simeq3-2.06151458390405\times10^{-24}.
\end{equation}
The channel count is three, and the displayed width retains the exact current shares and phase space. The first heavy state has $m_{D1}=0$; it is uncoupled by the displayed Yukawa source and has mass $\MR$.

\subsection{Quark and lepton mixing}
A weak current produces a definite charged or neutral family combination. The rest operator propagates its mass eigenstates. Their relative frame is a unitary rotation. Let $R_{ij}(s,\varphi)$ be the identity outside the $ij$ plane, with plane block
\begin{equation}
 \begin{pmatrix}\sqrt{1-s^2}&se^{-\ii\varphi}\\-se^{\ii\varphi}&\sqrt{1-s^2}\end{pmatrix}.
\end{equation}
We write $R_{ij}(s)=R_{ij}(s,0)$. 

The five direction labels permute the ten symmetric states. Their Petersen contrasts split into the five-dimensional representation $[3,2]$ and the four-dimensional representation $[4,1]$. If $f(g)$ and $l(g)$ count fixed two-element sets and fixed directions, then $\chi_{[3,2]}(g)=f(g)-l(g)$. The central projector onto $[5]\oplus[1^5]\oplus[3,2]$ has convolution kernel
\begin{equation}
 k(g)=\frac{1+\sgn g+5\chi_{[3,2]}(g)}{120},\qquad
 k*k=k,\quad\sum_gk(g)=1.
 \label{eq:mixingcharacters}
\end{equation}
At the identity, an exchange, a five-cycle and a three-cycle, respectively,
\begin{equation}
 k(1)=\frac9{40},\quad k((ab))=\frac1{24},\quad
 k((abcde))=\frac1{60},\quad k((abc))=-\frac1{40}.
 \label{eq:mixingcharactervalues}
\end{equation}
The first three coefficients fix the family overlaps; the odd quarter gives the relative phase $\arg(b+\ii)$. Equivalently, the same coefficients are return moments:
For the quark family source, the remaining variance of the symmetric-state visits after conditioning on the number of vector-state visits, averaged over $n_L$, is $V_C=9$. Dividing by the mean return length and the singlet transmission factor fixes the first overlap. The second retains four changes against $\Tr L^2=96$. The third passes through their common two-plane. The quark current in the mass basis is
\begin{align}
 s_{12}&=\frac{V_C}{(\EE N_{\rm ret})(5/2)}=\frac9{40},\qquad
 s_{23}=\frac4{\Tr L^2}=\frac1{24},\qquad
 s_{13}=t_0s_{12}s_{23}=\frac3{800},\notag\\
 \delta_q&=\arg(7+16\ii)=1.15838588519751,\notag\\
 V&=R_{23}(s_{23})R_{13}(s_{13},\delta_q)R_{12}(s_{12}).
 \label{eq:V}
\end{align}
The lepton source first compares the two complementary directions equally and then compares their common frame with its charged endpoint. Its resistance weight is $b$ and its endpoint phase is the complementary quarter phase. With this order of physical comparisons its overlap is
\begin{equation}
 U=R_{12}\left(\sqrt{2b}\,s_{12},\frac\pi2-\delta_q\right)^\dagger
 R_{23}(1/\sqrt2)R_{12}(\sqrt b).
 \label{eq:U}
\end{equation}
These are exact unitary matrices; their moduli are
\begin{align}
 |V|&=\begin{pmatrix}
 .974351913&.224998418&.003750000\\
 .224865668&.973498508&.041666374\\
 .008589642&.040943464&.999124542
 \end{pmatrix},\notag\\
 |U|&=\begin{pmatrix}
 .824329552&.546198089&.148823511\\
 .318952646&.648396212&.691268083\\
 .467707173&.530330086&.707106781
 \end{pmatrix}.
\end{align}
The Chau--Keung angles and Jarlskog quartets\cite{ChauKeung,Jarlskog} give
\begin{align}
 J_q&=3.13550917578981\times10^{-5},\qquad
 J_l=-0.0144644347463,\notag\\
 (\alpha_q,\beta_q,\gamma_q)
 &=(90.713773064338,22.951143935389,66.335083000273)^\circ,\notag\\
 (s_{12,l}^2,s_{13,l}^2,s_{23,l}^2)
 &=(.305089610922636,.0221484375,.488674949067231),\notag\\
 \delta_l&=205.584225835201^\circ,
 \qquad \sin^2(2\theta_{13,l})=.0866315368652344.
\end{align}
In particular, $|U_{e3}|^2=bs_{12}^2=567/25600$ and
\begin{equation}
 (|U_{\tau1}|^2,|U_{\tau2}|^2,|U_{\tau3}|^2)
 =\left(\frac7{32},\frac9{32},\frac12\right).
 \label{eq:taurow}
\end{equation}

The rows of $U$ are charged-lepton flavour states and its columns are neutral-family states before the Majorana mass matrix is diagonalized. The light and heavy mass states and their complete current follow from the six-state source in Eq.~\eqref{eq:neutralfull}. With physical Dirac matrices $\Gamma^\mu$ and chirality $\Gamma_*$, the positive current is
\begin{equation}
 j_+^\mu=\bar u\Gamma^\mu(1-\Gamma_*)Vd
 +\bar\nu\Gamma^\mu(1-\Gamma_*)U^\dagger\ell.
 \label{eq:chargedcurrent}
\end{equation}
For example, the source mass matrices can be written $M_d=VD_d$ and $M_\nu=U^*D_\nu U^\dagger$ in the up-quark and charged-lepton diagonal bases. Diagonalizing these mass operators reproduces the same current overlaps.

A neutral wave created in charged-lepton flavour $\alpha$ propagates in the
mass basis. At energy $E\ll M_R$, its accessible current norm is
$n_\alpha=\sum_i|U_{\alpha i}|^2c_i^2$. The long-wave relative propagation
phase follows from $E_i=p c+m_i^2c^3/(2p)+\cdots$. For a baseline $L_{\rm base}$,
\begin{align}
 \mathcal A_{\alpha\to\beta}(L_{\rm base})
 &=\frac{\sum_i U_{\beta i}U_{\alpha i}^*c_i^2
 e^{-\ii m_i^2c^3L_{\rm base}/(2\hbar E)}}{\sqrt{n_\alpha n_\beta}},\notag\\
 P_{\alpha\to\beta}&=|\mathcal A_{\alpha\to\beta}|^2.
 \label{eq:oscillation}
\end{align}
In the light limit $c_i\to1$, unitarity gives $n_\alpha=1$ and the usual
three-family expression. The complete current norms in
Eq.~\eqref{eq:neutralfullnorm} give the finite correction.

The effective masses in the displayed Majorana basis are
\begin{align}
 m_\beta&=\sqrt{\sum_i|U_{ei}|^2c_i^2m_i^2}\simeq.00882841829799682\eV,\notag\\
 m_{\beta\beta}&=\left|\sum_iU_{ei}^2c_i^2m_i\right|\simeq.0033134118567138\eV.
 \label{eq:effectivemasses}
\end{align}
The factors $c_i^2$ are the exact light-current norms of the neutral blocks. They change these numbers below the displayed precision. The additional relative Majorana phases are zero in this source basis. Equations~\eqref{eq:taurow}, \eqref{eq:neutrinovalues}, and \eqref{eq:effectivemasses} restore the textual vector records L02, T05, and K04.

\subsection{Constant curvature and stationary occupations}
A constant curvature compares the whole clock and spatial frame after continuation around a closed region. The same comparison has a stationary population of particles and a phase that remains when those particles dilute. The proton period is $2\pi$ in its own phase units. The elementary rest frequency is smaller by $R_p$, and one mean return contains $16$ steps, so its phase per return slice is $2\pi/(16R_p)$. The other end has reflection $r_0$ shared equally among the ten symmetric states. These two endpoint factors are
\begin{equation}
 \theta_c=\frac{2\pi}{16R_p}=.000213870604217778,\qquad
 s_F=\frac{r_0}{10}=\frac3{50}.
\end{equation}
The constant-curvature source compares two entire frames. Its forward entry carries the elementary rest phase through three whole transfers; its reverse entry carries one reflected share in each complementary direction. The two-ended curvature source is
\begin{equation}
 \mathcal T_{\mathrm{whole}}=
 \begin{pmatrix}0&W^3\theta_c\\W^3s_F^3&0\end{pmatrix},\qquad
 \mathcal T_{\mathrm{whole}}^2=W^6\theta_cs_F^3\Id.
 \label{eq:whole_source}
\end{equation}
Its constant curvature is therefore
\begin{align}
 \Lambda\ell_P^2&=3W^6\theta_cs_F^3
 =2.85522324300807\times10^{-122},\notag\\
 \ell&=\frac{\lamP}{W^2\sqrt{\theta_cs_F^3}}
 =5369.08213305099\,\mathrm{Mpc},\qquad \Lambda\ell^2=3.
 \label{eq:curvaturewhole}
\end{align}
In coordinates measured in $\lamP$, the frame action is
\begin{equation}
 I_{\mathrm{frame}}=\frac1{16\pi W^2}
 \int\dd^4\widehat x\sqrt{|g|}
 (\widehat R-6W^4\theta_cs_F^3).
\end{equation}
A stationary population asks which continuing states are occupied together. In the population state used here, the sixteen internal basis states have independent Bernoulli occupations, each with probability $a$. For occupations $n_i\in\{0,1\}$ its full probability law is
\begin{equation}
 \Pr(n_1,\ldots,n_{16})=\prod_{i=1}^{16}a^{n_i}(1-a)^{1-n_i}.
 \label{eq:stationarypopulation}
\end{equation}
A specified set of $r$ occupied states has probability $a^r$ after summing over the others: every unrestricted state contributes $a+(1-a)=1$. Let $I_+=\prod_{i:L_i>0}n_i$ require the eleven states with positive comparison, the singlet and the ten symmetric states. Let $I_{16}=\prod_{i=1}^{16}n_i$ require all sixteen. Then $\EE I_+=a^{11}$ and $\EE I_{16}=a^{16}$.

Every first return begins and ends internally at a vector state. Complete
these two compulsory arrivals first and write $K=n_L-2$. The marginal of
Eq.~\eqref{eq:jointreturn} is
\begin{equation}
 \Fret(1;x,1)=\frac{x^2}{4-3x},\qquad
 G_K(x)=x^{-2}\Fret(1;x,1)=\frac1{4-3x}.
 \label{eq:connectedreturn}
\end{equation}
Thus $\Pr(K=k)=\tfrac14(\tfrac34)^k$, and logarithmic differentiation gives
$\EE K=3$ and $\Var K=12$.

The positive comparison space $\mathcal H_+$ has dimension eleven. Its full
exterior coefficient is
\begin{equation}
 q_+=\det_{\mathcal H_+}(a\,\sgn L)=a^{11}.
 \label{eq:positivevolume}
\end{equation}
Completing this volume by its connected internal return gives one source,
\begin{align}
 \mathcal C(j)
 &=\Tr_{\Lambda^{11}\mathcal H_+}
   [\Lambda^{11}(a\,\sgn L)\log G_K(e^j)]\notag\\
 &=-a^{11}\log(4-3e^j)
 =a^{11}\sum_{r\geq1}\frac{(3/4)^r}{r}(e^{rj}-1),
 \quad\Re j<\log(4/3).
 \label{eq:singlepopulation}
\end{align}
For real $j<\log(4/3)$ the exponential of Eq.~\eqref{eq:singlepopulation} is
\begin{equation}
 \begin{gathered}
 e^{\mathcal C(j)}
 =4^{-q_+}\sum_{k=0}^{\infty}
   \frac{(q_+)_k}{k!}\left(\frac34\right)^k e^{kj},\\
 (q_+)_0=1,\qquad
 (q_+)_k=q_+(q_++1)\cdots(q_++k-1)\quad(k\geq1).
 \end{gathered}
 \label{eq:positivestatisticalsource}
\end{equation}
To obtain the coefficients, write
$\mathcal Z_+(x)=(4-3x)^{-q_+}=\sum_{k\geq0}p_kx^k$.
The equation $(4-3x)\mathcal Z_+'(x)=3q_+\mathcal Z_+(x)$ gives
$p_0=4^{-q_+}$ and
$(k+1)p_{k+1}=\tfrac34(q_++k)p_k$.
It yields the coefficients displayed above. Every $p_k$ is positive
and $\sum_kp_k=\mathcal Z_+(1)=1$.

Thus $\mathcal C$ is the cumulant generating function of a normalized positive distribution.

The exterior volume weights the connected source once. The factor $1/r$
counts a closed repetition once, as in Eq.~\eqref{eq:eulerproduct}. Its first
and second insertions are
\begin{equation}
 \mathcal C'(0)=3a^{11}=6.14004762346375\times10^{-10},\qquad
 \mathcal C''(0)=12a^{11}=2.45601904938550\times10^{-9},
 \qquad \frac{\mathcal C''(0)}{\mathcal C'(0)}=4.
 \label{eq:threeandtwelve}
\end{equation}
They are the density and connected scalar-fluctuation coefficients. The
second coefficient is $q_+\Var K$. These derivatives refer to the source
$q_+\log G_K$, with its exterior weight outside the logarithm. The occupation
of all sixteen states has $\EE I_{16}=a^{16}$ and gives the heavy neutral
number relative to the baryon number.
\begin{equation}
 \eta_B=\mathcal C'(0),\qquad
 \frac{n_{N_1}}{n_b}=a^{16}.
 \label{eq:populationreadings}
\end{equation}

The three-cycle character of the symmetric-state comparison is the scalar tilt. Its return reading is exactly
\begin{equation}
 -k((abc))=\frac{t_0}{\EE N_{\rm ret}}=\frac1{40}.
 \label{eq:tiltidentity}
\end{equation}
Thus the occupied positive states carry the spectrum
\begin{equation}
 \mathcal P_{\mathcal R}(k_{\rm com})=\mathcal C^{\prime\prime}(0)(k_{\rm com}\ell)^{-1/40}
 =12a^{11}(k_{\rm com}\ell)^{-1/40},\qquad
 n_s=1+k((abc))=1-\frac{t_0}{\EE N_{\rm ret}}=\frac{39}{40},\qquad
 \frac{\dd n_s}{\dd\log k}=0.
 \label{eq:primordial}
\end{equation}
Here $k_{\rm com}$ is the comoving wave number in inverse-length units. At the reporting pivot $k_*=0.05\,\mathrm{Mpc}^{-1}$, $A_s=2.13555115835\times10^{-9}$ and $\log(10^{10}A_s)=3.06130986015343$.

Each light channel has occupation $\eta_0$. Its retained number per photon is $(3/11)\eta_0\sigma$. For a Fermi gas the same ratio is $(3/4)\eta_0(T_\nu/T_\gamma)^3$, hence
\begin{equation}
 \frac{T_\nu}{T_\gamma}=\left(\frac{4\sigma}{11}\right)^{1/3}.
 \label{eq:neutrinotemperatureratio}
\end{equation}
The radiation energy density carries the fourth power of this temperature ratio, so
\begin{equation}
 N_{\mathrm{eff}}=3\eta_0\sigma^{4/3}=3.04385727361762.
\end{equation}
The heavy neutral occupation carries the following energy per proton reference mass:
\begin{equation}
 R_c^{\mathrm{finite}}=a^{16}\frac{\MR}{m_p}=5.33628360801120.
 \label{eq:coldnative}
\end{equation}
The factor $\kappa_B$ converts baryon number counted in proton masses into the energy density of the resulting atoms, including nuclear binding and electrons: $\rho_b=\kappa_Bm_p n_b c^2$. Thus the cold-to-baryon ratio is $R_c^{\mathrm{finite}}/\kappa_B$. Part III evaluates this conversion for each declared composition input.

For a homogeneous flat metric
\begin{equation}
 ds^2=-c^2dt^2+a_{\rm exp}(t)^2d\bm x^2,\qquad
 H_{\rm exp}=\frac{\dot a_{\rm exp}}{a_{\rm exp}},
\end{equation}
the time--time component of the frame equation gives the expansion. The contracted Bianchi identity gives conservation of the same stress. For fluid energy density $\rho_{\rm E}$ and pressure $p_{\rm E}$ in the same units, these equations are
\begin{equation}
 H_{\rm exp}^2=\frac{8\pi G}{3c^2}(\rho_m+\rho_{\rm rad}+\rho_\Lambda),\qquad
 \dot\rho_{\rm E}+3H_{\rm exp}(\rho_{\rm E}+p_{\rm E})=0,\qquad
 \rho_\Lambda=\frac{\Lambda c^4}{8\pi G}.
 \label{eq:friedmann}
\end{equation}
The comparison value of $H_{\rm exp}$ is denoted by $H_{\rm comp}$.

The quadratic frame response per spinor component is $\Tr L^2/4=24$, and the positive comparison space has $1+\EE n_F=11$ states. Equal stationary response of these two coefficients gives
\begin{equation}
 \frac{\rho_\Lambda}{\rho_m}=\frac{\Tr L^2/4}{1+\EE n_F}=\frac{24}{11},\qquad\Omega_k=0.
 \label{eq:epochcount}
\end{equation} Matter dilutes as the volume grows while the constant-curvature density stays fixed, so this ratio locates a definite stage of that branch. At this epoch the remaining radiation fraction gives
\begin{equation}
 \Omega_m=\frac{11}{35}(1-\Omega_r),\quad
 \Omega_\Lambda=\frac{24}{35}(1-\Omega_r),\quad
 H_{\rm comp}=\frac{c}{\ell\sqrt{\Omega_\Lambda}}.
 \label{eq:epoch}
\end{equation}
Thermal and acoustic records from this epoch are specified in Sec.~\ref{sec:evolution}.

The de Sitter readings\cite{GibbonsHawking} are
\begin{align}
 k_BT_{\mathrm{dS}}&=\frac{\hbar c}{2\pi\ell}
 =1.89564063245871\times10^{-34}\eV,\notag\\
 \frac{S_{\mathrm{dS}}}{k_B}&=\frac{3\pi}{\Lambda\ell_P^2}
 =3.30089003858073\times10^{122},\notag\\
 r_{\mathrm{dS}}&=\frac{16\Lambda\ell_P^2}{3\pi\,12a^{11}}
 =1.97359117550\times10^{-113}.
\end{align}
The last ratio uses the de Sitter curvature in Eq.~\eqref{eq:curvaturewhole} and scalar amplitude at $k\ell=1$. It specifies the curvature and state of this tensor reading.

\subsection{Shared quantities and exact identities}
Several distant measurements depend on the same small set of finite quantities. Eliminating those quantities gives relations among the results. For example, the electric response and neutral current use the same repeated return $N_h$ and clock overlap $\hover$:
\begin{equation}
 \left(\frac{3700}{27}-\alpha^{-1}\right)
 (891+\widehat s^2)[M_Q(\nu)M_Q(-\nu)]^{3/2}
 =\frac{37}{40}.
\end{equation}
Here $1/(c_Lf)=3700/27$, $\widehat s^2=3N_h$, and $\nu/(c_Lf)=37/120$. The proton's Gram logarithm and the complete-frame amplitude give
\begin{equation}
 -\frac12\log\alpha_G-9\log\frac{296}{295}
 =288\left(R_p-1836+a^5+ba^6\right)
 =48\log\frac52.
\end{equation}
The logarithmic family amplitudes also imply
\begin{equation}
 \log\!\left[\eta_0^{-2}M_Q(-\nu)^{1/6}
 \frac{(m_b^{\rm native})^2m_d^{\rm native}}{(m_s^{\rm native})^3}\right]
 +3|V_{ub}|=5,
\end{equation}
because $q_q^2/p_q=\exp(5-5\nu)M_Q(-\nu)^{-1/6}$ and $3|V_{ub}|=5\nu$.
These equalities follow by substitution into the displayed sources; their physical content is that the same return and normalization appear in each reading.

\subsection{Ward identities and conservation}\label{ward-identities-and-conservation}

Put $\mathcal G_{\rm resp}=\mathcal K^{-1}$. Under a change of frame, \({\mathcal K}\mapsto U{\mathcal K}U^{-1}\). For its infinitesimal generator \(\lambda\), \(\delta {\mathcal K}=[\lambda,{\mathcal K}]\), and

\begin{equation}
 \delta\log\det {\mathcal K}=\operatorname{Tr}({\mathcal G_{\rm resp}}[\lambda,{\mathcal K}])=0.
 \label{eq:proof-SF-1}
\end{equation}

This is cyclicity of a finite trace. With matter sources \({\mathcal K}\psi=\mathfrak j\) and \(\bar\psi {\mathcal K}=\bar{\mathfrak j}\), the same identity resolves to

\begin{equation}
 [{\mathcal K},{\mathcal G_{\rm resp}}+\psi\bar\psi]=\mathfrak j\bar\psi-\psi\bar{\mathfrak j}.
 \label{eq:proof-SF-2}
\end{equation}

Indeed \({\mathcal K}{\mathcal G_{\rm resp}}={\mathcal G_{\rm resp}}{\mathcal K}=I\), and the remaining commutator is the displayed pair of matter equations. With the external insertions removed, the local diagonal components express current balance at every comparison.

On the links, a scalar frame change sends \(\theta\) to \(\theta+\operatorname{grad}\lambda\). Its variation is \(j^{\mathsf T}\operatorname{grad}\lambda\). The coefficient of arbitrary vertex values \(\lambda\) is \(\operatorname{grad}^{\mathsf T}j\), which vanishes on the field equations. In the local chart, the same variation of

\begin{equation}
 \mathcal L_A=-\frac{F^a_{\mu\nu}F^{a\mu\nu}}{16\pi\alpha}
               -j^{a\mu}A^a_\mu
\end{equation}

uses \(\delta F_{\mu\nu}=D_\mu\delta A_\nu-D_\nu\delta A_\mu\) and gives \(D_\mu F^{\mu\nu}=4\pi\alpha j^\nu\). Gauge variation \(\delta A=D\lambda\), together with the matter equations, gives \(D_\mu j^\mu=0\). The identity \(DF=0\) follows by expanding \(F=dA+A\wedge A\), or from the Jacobi identity for three covariant derivatives. These are the inhomogeneous and homogeneous field equations of the completed parallel comparison.

Let $E_a$ be dual to the coframe $\vartheta^a$. For the frame, take a vector field \(\xi=\xi^aE_a\). Combining a coordinate variation with its local Lorentz frame change gives the covariant variations

\begin{equation}
 \delta_\xi \vartheta^a=D\xi^a+\iota_\xi T^a,\qquad
 \delta_\xi\omega^{ab}=\iota_\xi R^{ab}.
 \label{eq:proof-SF-3}
\end{equation}

Insert these in the matter variation defining \(\mathfrak t_a,\mathfrak s_{ab}\) above. Integration by parts on \(D\xi^a\) yields the coefficient of each arbitrary \(\xi^a\):

\begin{equation}
 D\mathfrak t_a=(\iota_{E_a}T^b)\wedge\mathfrak t_b
             +\frac12(\iota_{E_a}R^{bc})\wedge\mathfrak s_{bc}.
 \label{eq:proof-SF-4}
\end{equation}

This is energy--momentum balance with the torsion and spin terms included. Local Lorentz variation \(\delta \vartheta^a=\lambda^a{}_b\vartheta^b\), \(\delta\omega^{ab}=-D\lambda^{ab}\) supplies the corresponding balance of spin and antisymmetric stress. With the torsion-free metric connection and the total symmetric stress tensor, coordinate invariance becomes \(\nabla_\mu T^{\mu\nu}=0\). It also follows from the contracted Bianchi identity applied to \eqref{eq:proof-SE-8}.

The spin-current identity can be written explicitly in the same variation convention. Substitution of the Lorentz variation into \(\delta I_m\), followed by integration by parts on \(D\lambda^{ab}\), gives

\begin{equation}
 0=-\frac12\int\lambda^{ab}
       \left(D\mathfrak s_{ab}+\vartheta_b\wedge\mathfrak t_a-\vartheta_a\wedge\mathfrak t_b\right).
 \label{eq:proof-SF-5}
\end{equation}

Every antisymmetric parameter is independent, so the three terms in parentheses sum to zero. This shows how the antisymmetric part of the coframe stress and the divergence of spin enter a single conservation equation.

In the torsion-free metric form, the proof is equally direct. A coordinate variation of the inverse metric is \(\delta g^{\mu\nu}=-\nabla^\mu\xi^\nu-\nabla^\nu\xi^\mu\). The stress definition gives \(\delta I_m=\int\sqrt{-g}\,T_{\mu\nu}\nabla^\mu\xi^\nu\). After integration by parts its coefficient is \(-\nabla^\mu T_{\mu\nu}\), proving conservation for arbitrary \(\xi\). For homogeneous isotropic energy and pressure this is \(\dot\rho_{\rm E}+3H_{\rm exp}(\rho_{\rm E}+p_{\rm E})=0\), the continuity equation used with the frame's constant volume term.

Finally, the mass and curvature parts of the matter equation remain in one operator. With the main text's Lorentz Clifford convention, the symmetric part of \(\Gamma^\mu\Gamma^\nu\) is \(-g^{\mu\nu}\), and its antisymmetric part contracts only the derivative commutator. Therefore

\begin{equation}
 (i\Gamma^\mu D_\mu+m)(i\Gamma^\nu D_\nu-m)
  =\Box_D-m^2-\frac14[\Gamma^\mu,\Gamma^\nu][D_\mu,D_\nu].
 \label{eq:proof-SF-6}
\end{equation}

The mass comes from the rest zero, and the commutator is the already completed parallel curvature. Their coefficients use the same normalized matter state from Eq.\eqref{eq:positiveclock}. This is the explicit return of the earlier pole norm in the complete local matter equation.

Finally, integration on the circle of a link phase gives a finite Schwinger--Dyson identity directly. For a differentiable observable \(O\), periodicity implies \(\int d(Oe^{-S})=0\), so \(\langle D_XO-O\,D_XS\rangle=0\). The determinant variation, the source equations, the conserved current, and this integration identity are four evaluations of the same stationary law. The rest zero and the parallel curvature have kept their source normalization through all of them.
\subsection{The common action and its variations}
\label{sec:commonaction}

The eight completed steps supply the coefficients and fields of the
following action. Use $c=\hbar=1$, $g=(-+++)$ and the real curvature
components in the charge convention of Eq.~\eqref{eq:gaugecomponents}.
The fields $F_C,F_W,F_Y$ in the action are the real charge-coordinate components $\mathcal F_X^a$; $F_{X,{\rm can}}^a=\mathcal F_X^a/g_X$ denotes their canonical form. For a non-Abelian factor the Hermitian generators satisfy
$\tr(t^at^b)=\delta^{ab}/2$.
The weak pole coefficients are
$\alpha_2=g_2^2/(4\pi)$ and $\alpha_Y=g_1^2/(4\pi)$, where
$g_1=g_{Yc}$ after the canonical normalization of
Eq.~\eqref{eq:weakpencil}. Thus
\begin{equation}
 \alpha_2=\frac{\alpha}{(1-f)s_{\rm OS}^2},\qquad
 \alpha_Y=\frac{\alpha_2s_{\rm OS}^2}{1-s_{\rm OS}^2},\qquad
 e_0^2=(1-f)e_{\rm pole}^2,
 \quad e_{\rm pole}^2=\frac{g_1^2g_2^2}{g_1^2+g_2^2}.
 \label{eq:actionpolecoefficients}
\end{equation}
The electric coefficient at zero transfer and the vector-pole coefficient
are connected by the endpoint norm already present in the weak mass
equation.

With $\widetilde\varphi=i\sigma_2\varphi^*$, the action is
\begin{equation}
\begin{split}
\Gamma=\int d^4x\sqrt{-g}\,\bigg[&
 \frac{R-2\Lambda}{16\pi G}
 -\frac{F_C^a{}_{\mu\nu}F_C^{a\mu\nu}}{16\pi\alpha_s}
 -\frac{F_W^i{}_{\mu\nu}F_W^{i\mu\nu}}{16\pi\alpha_2}
 -\frac{F_Y{}_{\mu\nu}F_Y^{\mu\nu}}{16\pi\alpha_Y}
 +\frac{\bar\theta}{32\pi^2}F_C^a{}_{\mu\nu}\widetilde F_C^{a\mu\nu}\\
 &+\bar\psi\,i\Gamma^\mu D_\mu\psi
 -\bigl(\bar q_LY_d\varphi d_R+\bar q_LY_u\widetilde\varphi u_R
       +\bar\ell_LY_\ell\varphi e_R+\bar\ell_LY_\nu\widetilde\varphi N_R
       +\mathrm{h.c.}\bigr)\\
 &-\frac12\bigl(M_R\overline{N_R^{\,c}}N_R+\mathrm{h.c.}\bigr)
 -(D_\mu\varphi)^\dagger D^\mu\varphi
 -\lambda_H(\varphi^\dagger\varphi-v^2/2)^2\bigg].
\end{split}
\label{eq:completeaction}
\end{equation}
This is the article's complete action, with each coefficient obtained
before it is used. The following table gives those dependencies.

\begingroup\small
\setlength{\tabcolsep}{4pt}
\begin{longtable}{@{}>{\raggedright\arraybackslash}p{.16\textwidth}
>{\raggedright\arraybackslash}p{.46\textwidth}
>{\raggedright\arraybackslash}p{.31\textwidth}@{}}
\caption{Coefficients of the common action.}\label{tab:actioncoefficients}\\
\toprule Term&Coefficient from the completed return&Earlier derivation\\\midrule
\endfirsthead
\toprule Term&Coefficient from the completed return&Earlier derivation\\\midrule
\endhead\bottomrule\endfoot
Frame curvature&$G=W^2/m_p^2$; $W=t_0^{48}\eta_0^9$.
&Eqs.~\eqref{eq:W}, \eqref{eq:gravitystrength}.\\[4pt]
Constant volume&$\Lambda\ell_P^2=3W^6\theta_cs_F^3$;
$\Lambda\ell^2=3$.
&Eqs.~\eqref{eq:curvaturewhole}, \eqref{eq:proof-SE-4}.\\[4pt]
Colour curvature&$\alpha_s=r_{\rm v}/\sqrt8=1/\sqrt{72}$;
$\bar\theta=0$.
&Eqs.~\eqref{eq:alphas}, \eqref{eq:thetasectors}.\\[4pt]
Weak and hypercharge curvature&$\alpha_2$, $\alpha_Y$ in
Eq.~\eqref{eq:actionpolecoefficients}; the zero-transfer input is
$\alpha=c_Lf/[1-\nu h_\nu/(297+N_h)]$.
&Eqs.~\eqref{eq:alpha}, \eqref{eq:weakpencil}, \eqref{eq:weakmasses}.\\[4pt]
Scalar scale&$v=(\sqrt2G_F)^{-1/2}$;
$G_Fm_p^2=9W_W/(10M^2)$.
&Eq.~\eqref{eq:GF}.\\[4pt]
Radial scalar&$\lambda_H=\pi\alpha/\chi_h$;
$\chi_h=(3/13)(10/13)$.
&Eqs.~\eqref{eq:radialsusceptibility}, \eqref{eq:radial}.\\[4pt]
Charged Yukawa sources&Singular values are $\sqrt2/v$ times the entries
of $M_d$, $\widehat M_u=\diag(m_u,m_c,m_t^{\rm rad})$ and $M_\ell$.
The left bases have overlap $V$.
&Eqs.~\eqref{eq:chargedsource}, \eqref{eq:radialfermionsource},
\eqref{eq:V}.\\[4pt]
Neutral Yukawa source&Singular values are $\sqrt2/v$ times $D_D$;
$D_D=(\sigma v/\sqrt2)\diag(0,\sqrt{\rho_\nu},1)$,
$\rho_\nu=\tan^2(\pi/8)$. Its left basis has overlap $U$ with the
charged-lepton basis.
&Eqs.~\eqref{eq:neutralhalfphase}, \eqref{eq:neutralfull},
\eqref{eq:U}.\\[4pt]
Neutral Majorana source&$M_R=M_{\rm Pl}\det_{\mathbb R}(\Id_9/3)$.
&Eq.~\eqref{eq:neutraldeterminant}.\\
\end{longtable}
\endgroup

The native top entry belongs to the colour family source. The scalar
action uses the radial entry $m_t^{\rm rad}=vy_t/\sqrt2$, with
$y_t^2=\eta_0/\sigma$. These are the two source coordinates used in the
article's two tables.

Variation keeps the field and its source in the same normalization.
For each gauge factor $X$, define its complete matter current by the
coefficient of $-\delta A_{X\nu}^a$ in the matter variation. Integrating
the curvature variation by parts gives
\begin{equation}
 (D_\mu F_X^{\mu\nu})^a=4\pi\alpha_Xj_X^{a\nu},\qquad
 D_{[\lambda}F_{X\mu\nu]}=0.
 \label{eq:actiongaugevariation}
\end{equation}
For the abelian factor $D_\mu$ is the ordinary derivative. The second
equation is the curvature identity, and the matter equations give
$D_\nu j_X^\nu=0$. The constant topological term has
$\bar\theta=0$ on the sector already determined in the colour step.

Varying $\bar\psi$ gives the covariant first-order equation with its
Yukawa rest operator. At $\varphi^\dagger\varphi=v^2/2$, its charged
blocks have the singular values in the coefficient table; its neutral
block is Eq.~\eqref{eq:neutralfull}. Varying $\varphi^\dagger$ gives
\begin{equation}
 D_\mu D^\mu\varphi
 -2\lambda_H(\varphi^\dagger\varphi-v^2/2)\varphi
 =J_\varphi,
 \qquad J_\varphi=-\frac{\partial\mathcal L_{\rm Yuk}}
                         {\partial\varphi^\dagger}.
 \label{eq:actionscalarvariation}
\end{equation}
Expanding the radial field about $v$ gives $m_H^2=2\lambda_Hv^2$.
The weak vector masses follow from the same scalar kinetic term with its
canonical gauge fields, as in Eq.~\eqref{eq:weakpencil}.

Independent coframe and spin-connection variations are
Eqs.~\eqref{eq:proof-SE-6} and \eqref{eq:proof-SE-7}. Their
torsion-free metric equation is Eq.~\eqref{eq:proof-SE-8}; its
homogeneous component and energy conservation give the Friedmann and
continuity equations. The constant volume has
$T_{\mu\nu}^{\Lambda}=-\rho_\Lambda g_{\mu\nu}$. The complete clock
identity of Eq.~\eqref{eq:fullclock} includes every rest contribution in
the gravitational source. Thus the coefficients of this action are the
same normalized responses that supplied the rest equations and the
interactions along the eight steps.
\subsection{The field coordinates of the common action}
With the canonical hypercharge coordinate of Eq.~\eqref{eq:weakpencil}, the charge-unit kinetic terms are
\begin{equation}
 \mathcal L_{\mathrm{gauge}}=
 -\frac{\mathcal F_C^a\mathcal F_C^a}{16\pi\alpha_s}
 -\frac{\mathcal F_2^i\mathcal F_2^i}{4g_2^2}
 -\frac{\mathcal F_Y\mathcal F_Y}{4g_{Yc}^2}.
 \label{eq:commongaugebasis}
\end{equation}
The neutral rotation then forms the photon and the $Z$ from the same weak and hypercharge fields. A local field functional uses one common matching scale. The electric response at zero transfer is $\alpha$; the vector response at its mass zero uses $\alpha/(1-f)$. Their matching relation is already contained in the generalized vector equation. Thus a change of field coordinates or response scale carries its source with it.

\newpage
\part{The one variable and its physical quantities}\label{part:maps}

\section{The map of the spectral variable}
\label{sec:maps}
\label{sec:classical}

\subsection{Classical quantities in the one spectral variable}
\label{sec:variablemap}

The phase of a step becomes an energy when compared with the time taken by that step. The phase along a spatial direction becomes momentum. A phase brought back by a closed path becomes a field strength. These quantities acquire their familiar names because of what a participant compares. The spectral weight is $u$ throughout. A source specifies the state and the applied change in $K_X$; its zero, residue or mixed derivative gives the physical reading. The source definitions and spectral theorems in Sec.~\ref{sec:common} make that operation explicit for each kind of matter. The map below follows the same variable through the exact finite identities, their long-wave equations and the convention of the final measurement.\cite{Repository}

Write $\Delta x=\lambda_p$ and $\Delta t=\lambda_p/c$, so one proper step is one radian of the proton clock. The two small phase angles of a moving state are $\mu=mc^2\Delta t/\hbar$ and $k=p\Delta x/\hbar$. The gravitational potential is denoted $\Phi_G$, distinct from the signed first-return function $\Phis(z)$.

\begingroup\small
\setlength{\tabcolsep}{4pt}
\renewcommand{\arraystretch}{1.12}
\begin{longtable}{@{}>{\raggedright\arraybackslash}p{.20\textwidth}>{\raggedright\arraybackslash}p{.74\textwidth}@{}}
\caption{The history variable and the quantities read from it.}\label{tab:variablemap}\\
\toprule Quantity&Reading and equation\\\midrule
\endfirsthead
\multicolumn{2}{l}{\textit{The spectral variable and classical quantities, continued}}\\
\toprule Quantity&Reading and equation\\\midrule
\endhead
\bottomrule\endfoot
Step weight&Four equal continuations give $|u|=1/4$; the divided set contributes $u=\ii^{|D|}/4$. Their determinant law is $D(u)=\det(\Id-uB)$, Eq.~\eqref{eq:bass}.\\[4pt]
Spectral gap&$\chi(u)=u^{-1}+4u-5$ is the scalar displacement of $\Delta=6\Id-L$. It is unchanged by $u\mapsto1/(4u)$; the four quarters and the exact ellipse are Eqs.~\eqref{eq:ellipse}--\eqref{eq:fourquarters}.\\[4pt]
Proper time&The retained direction has clock $\ii\gamma_h$, with $e^{\ii\pi\gamma_h/2}=\ii\gamma_h$. Its phase is counted along successive changes. In laboratory units $\Delta t=\hbar/(m_pc^2)\simeq7.01515\times10^{-25}\,\mathrm s$.\\[4pt]
Common time&A phase assignment $z_x=q_{xy}z_y$ exists precisely when every closed path has $\prod_\gamma q=1$. The central rotation $\gamma_h\gamma_k$ of the retained plane supplies the shared clock, as derived in Sec.~\ref{sec:step-two}.\\[4pt]
Metric&The same quarter-turn preserves the Hermitian norm and changes the bilinear square: $q_D^\dagger q_D=\Id$, $q_D^{\mathsf T}q_D=\theta_D$. Thus $g_D(v)=\|v_{D^\perp}\|^2-\|v_D\|^2$.\\[4pt]
Light cone&The balanced directions of the two-plus-four quadratic form give a local conformal Lorentz cone. On the selected two-spinor side $\det(t\Id+\bm x\cdot\bm\sigma)=t^2-|\bm x|^2$.\\[4pt]
Frequency and energy&A common step phase $e^{-\ii\omega}$ moves the critical zero to $u_c=e^{\ii\omega}/4$. The proper energy is $E=\hbar\omega/\Delta t$. An applied source moves the zero according to Eq.~\eqref{eq:residue}.\\[4pt]
Momentum and wavelength&A spatial step phase $k$ gives $p=\hbar k/\Delta x$ and wavelength $2\pi\Delta x/|k|$. Its exact ordered step is $e^{-\ii k\gamma_j}e^{-\ii\mu\gamma_h}$.\\[4pt]
Velocity&The ordered step gives $\cos\omega=\cos\mu\cos k$ and $v_g/c=\partial_k\omega=\cos\mu\sin k/\sin\omega$. At $\mu=0$ the principal massless branch has $\omega=|k|$ and speed $c=\Delta x/\Delta t$.\\[4pt]
Rest mass&$mc^2=\hbar\mu/\Delta t$ is the rest turn. More generally $\det K_X(E_X,0)=0$ and $m_X=E_X/c^2$. The reciprocal chart has $\epsilon^2M_\epsilon^2=2\chi+\chi^2/5$, Eq.~\eqref{eq:exactgap}.\\[4pt]
Mass ratios&The same clock divides two rest frequencies. The elementary marked-transition rate is $1/320$; the stationary proton contribution has rate $1836/320$. Its complete source also contains the Gram and signed-volume terms in Eq.~\eqref{eq:proton}.\\[4pt]
Moving clock&At small phase, $\omega^2=\mu^2+k^2-\mu^2k^2/3+O((|k|+|\mu|)^6)$. The long-wave limit gives $E^2=p^2c^2+m^2c^4$ and $\dd\tau/\dd t=mc^2/E$.\\[4pt]
Length and size&$\lambda_p=\hbar/(m_pc)=0.2103089101\fm$ converts the count of steps. The return slopes give $r_p^2=16\lambda_p^2$, $r_\pi^2=10\lambda_p^2$ and $r_A^2=9\lambda_p^2$.\\[4pt]
Planck scales&The entire-frame amplitude gives $\ell_P=W\lambda_p$ and $\MP=m_p/W$. The same $W$ appears in the gravitational source, Eq.~\eqref{eq:gravitystrength}.\\[4pt]
Static potential and range&The long-wave inverse is $\mathcal G_{M_{\rm inv}}(r)=e^{-M_{\rm inv}r}/(4\pi r)$. The common massless limit is taken at $u=1/4$. Sources normalized by $\sqrt{4\pi}$ give $V_{\rm el}=\alpha\hbar c/r$ and $V_G=-W^2\hbar c/r$ for proton reference sources. A nonzero gap gives range $1/M_{\rm inv}$.\\[4pt]
Action and $\hbar$&The closed history action is $\Gamma=\log\det(\Id-uB[U])$, with $\log Z=\sum_{r\ge1}u^r\Tr B[U]^r/r$. The dimensional action is measured in $\hbar$; $E\Delta t=\hbar\omega$. A Legendre transform changes source coordinates to field coordinates.\\[4pt]
Probability&The postulated positive pairing gives $\Tr(\rho F_r)$. The exact side overlap is $Q_nQ_mQ_n=\tfrac12(1+n\cdot m)Q_n$; the doubled source gives the nonnegative squared amplitude of Sec.~\ref{sec:common}.\\[4pt]
Spin&Anticommuting perpendicular changes return a minus sign: $\gamma_s\gamma_t\gamma_s\gamma_t=-\Id$. The corresponding spin rotation has $R_{ab}(2\pi)=-\Id$ and $R_{ab}(4\pi)=\Id$.\\[4pt]
Charge and isospin&Occupation counts give $Y=\widehat n_F/2-\widehat n_C/3$, $T_3=(\widehat n_k-\widehat n_h)/2$, and $Q=\widehat n_k-\widehat n_C/3=T_3+Y$. The step along $h$ has zero emitted charge. The full sixteen-state representation is given in Sec.~\ref{sec:step-five}.\\[4pt]
Coupling&A two-source response is a mixed derivative of $\log Z$. Its finite endpoint coefficients are $\alpha_s=72^{-1/2}$, $\alpha^{-1}=137.035999177$, $G_Fm_p^2=1.02682689\times10^{-5}$ and $\alpha_G=W^2$. The respective reading phases are $-\ii/4,-1/4,\ii/4,1/4$.\\[4pt]
Vector potential&A link carries $U_{xy}$, abelian phase $e^{\ii\theta_{xy}}$. Independent local frame changes give $U_{xy}\mapsto g_xU_{xy}g_y^{-1}$. A closed holonomy transforms by conjugation, leaving its trace fixed.\\[4pt]
Field strength&The difference between two routes is the phase around their square. Its quadratic action has $-\kappa(u)\|C_\square\theta\|^2/2$, where $\kappa(u)=2u^4/(1+3u^2+16u^4)^2$. It is $1/200$ at $u=\pm1/4$ and $1/98$ at $u=\pm\ii/4$.\\[4pt]
Field normalization&The transported square angle and the real charge-unit curvature obey $\int\mathcal F_X^2=8\pi\alpha_X\kappa\|C_\square\theta\|^2$. Consequently $-\int\mathcal F_X^2/(16\pi\alpha_X)=-\kappa\|C_\square\theta\|^2/2$; the source is transformed with its field.\\[4pt]
Spacetime curvature&The ten rotations split into the connection $\omega$ that preserves $h$ and the coframe $\vartheta/\ell$ that moves it. Their curvature is $\mathcal F^{ab}=R^{ab}-\vartheta^a\wedge \vartheta^b/\ell^2$, with torsion in the mixed components.\\[4pt]
Curvature response&A circulating link phase moves the common zero by $\log(4u_c)=\|\theta_{\rm cyc}\|^2/48+O(\theta^4)$; $48=2(40-16)=3\times16$.\\[4pt]
Newton constant&Each complete-frame endpoint transmits $W=t_0^{48}\eta_0^9$. Thus $G=\hbar cW^2/m_p^2\simeq6.6743035\times10^{-11}\,\mathrm{m^3\,kg^{-1}\,s^{-2}}$.\\[4pt]
Gravitational potential&A common lapse $\mathcal N$ replaces $E$ by $E/\mathcal N$ in the complete rest equation. Its zero obeys $\dd E_X/\dd\mathcal N=E_X/\mathcal N$; at $\mathcal N=1+\Phi_G/c^2$ the first shift is $m_X\Phi_G$. The whole-state identity is Eq.~\eqref{eq:fullclock}.\\[4pt]
Cosmological constant&The coframe action gives $\Lambda\ell^2=3$. The retained source fixes $\Lambda\ell_P^2=3W^6\theta_cs_F^3$, $s_F=3/50$, $\theta_c=2\pi/(16R_p)$, and $\ell=\lambda_p/[W^2\sqrt{\theta_cs_F^3}]$. The closure phase and reflected share are the two entries of the whole-frame source.\\[4pt]
Pressure&Varying the constant volume term gives $T_{\mu\nu}^{\Lambda}=-\rho_\Lambda g_{\mu\nu}$, hence $p_\Lambda=-\rho_\Lambda$ and $w_\Lambda=-1$.\\[4pt]
Temperature&The de Sitter horizon gives $k_BT_{\rm dS}=\hbar c/(2\pi\ell)$. The selected colour crossover is $k_BT_c=m_p/6$ in units $c=1$. The photon temperature is obtained by the thermal evolution at $\rho_\Lambda/\rho_m=24/11$.\\[4pt]
Entropy&The horizon entropy is $S_{\rm dS}/k_B=3\pi/(\Lambda\ell_P^2)$. For the normalized pairing weights $(\pi_1^{\rm pair},\pi_2^{\rm pair},\pi_3^{\rm pair})=(1/6,5/12,5/12)$, the proton Gram term is $h_p=\sum_i\pi_i^{\rm pair}\log[(5/12)/\pi_i^{\rm pair}]=\log(5/2)/6$.\\[4pt]
Vector masses&The scalar displacement has nonzero trace-normalized Gram eigenvalues $(1/2,1/2,4/5)$ and nine zeros. The physical frequencies solve $\det(E^2C_B-\mathcal M_B)=0$; their kinetic norms and exact $W,Z$ equations are Eq.~\eqref{eq:weakpencil}.\\[4pt]
Higgs field&The length of the retained scalar vector has $v=(\sqrt2G_F)^{-1/2}$. Its radial source uses $\lambda_H=\pi\alpha/\chi_h$, $\chi_h=30/169$; expanding its quartic gives $m_H^2=2\lambda_Hv^2$.\\[4pt]
Weak mixing&The repeated odd-quarter return gives $\widehat s^2=3N_h$, with $N_h=\Re\sum_{j\ge1}h_\nu^jp_j$. The return arguments are $4(\ii/4)^j$; the on-shell ratio uses $s_{\rm OS}^2=\widehat s^2/(1+f)$.\\[4pt]
Running and scale change&For $u\ne0$, $K_{16}(v)=(v/u)K_{16}(u)+[1+4v^2-(v/u)(1+4u^2)]\Id$. A response against a fixed clock varies; the completed normalized crossing is invariant. The conventional interacting running and threshold maps are given below.\\[4pt]
Density and scalar power&The independent occupied positive states give $\EE I_+=a^{11}$. The connected source $\mathcal C(j)=-a^{11}\log(4-3e^j)$ has first and second derivatives that give $\eta_B=3a^{11}$ and $\mathcal P_{\mathcal R}=12a^{11}(k_{\rm com}\ell)^{-1/40}$; $n_s=39/40$ and $N_{\rm eff}=3\eta_0\sigma^{4/3}$.\\[4pt]
Classical field equations&The long-wave terms are the Dirac kinetic term, gauge curvature, Einstein--Cartan curvature, the scalar displacement and its specified self-interaction. Their coefficients are the source responses listed above. Their equations follow by varying those long-wave terms; Part III carries their sources through the measurement conventions.\\
\end{longtable}
\endgroup

\subsection{Changing the comparison parameter}
The finite comparison already has an exact change of parameter. The relation is unchanged while the weight of its continuing steps changes. For $u\ne0$, substitution in $K_{16}(u)=(1+4u^2)\Id-uA$ gives
\begin{equation}
 K_{16}(v)=\frac vu K_{16}(u)
 +\left[1+4v^2-\frac vu(1+4u^2)\right]\Id.
 \label{eq:finiteparameter}
\end{equation}
The eigenspaces of the comparison remain the same, and their scalar weights change by the displayed affine transformation. A response measured against a fixed clock can therefore change with the parameter. When the clock is carried with the response, the completed clock $\widetilde N=N+8\Pi_0$ gives the common normalized crossing in Eq.~\eqref{eq:crossing} for every positive central grade distribution. Its cancellation of the grade weights expresses this invariance exactly. The interacting renormalization-group, threshold and cosmological calculations below are the evolution maps in their stated physical conventions.

\subsection{Units and common calibration}
Equation~\eqref{eq:meunit} fixes the dimensional scale. The exact SI defining values used for conversion are
\begin{align}
 c&=299792458\,\mathrm{m\,s^{-1}},&
 h_{\mathrm{Planck}}&=6.62607015\times10^{-34}\,\mathrm{J\,s},\notag\\
 k_B&=1.380649\times10^{-23}\,\mathrm{J\,K^{-1}},&
 1\,\mathrm{eV}&=1.602176634\times10^{-19}\,\mathrm J .
\end{align}
We use $\hbar=h_{\mathrm{Planck}}/(2\pi)$, a Julian year of 365.25 days, and the parsec conversion for the stated megaparsec units. These constants define units; the measured electron rest energy supplies the common mass calibration.\cite{CODATA,SI}

The relative standard uncertainty of the electron calibration is $u_e=3.1311\times10^{-10}$. A finite quantity $X\propto m_e^{p_X}$ has calibration covariance
\begin{equation}
 \frac{\operatorname{Cov}_{m_e}(X,Y)}{XY}=p_Xp_Yu_e^2.
\end{equation}
Masses have exponent $+1$, lengths and times $-1$, $G$ and $G_F$ exponent $-2$, and $\Lambda$ exponent $+2$. Dimensionless finite ratios have exponent zero. The scalar amplitude at a fixed dimensional pivot has exponent $1/40$. The last printed digits of an evaluated formula and a metrological uncertainty are different quantities; the register prints the evaluated expression and retains the uncertainty of each experimental comparison.

\subsection{Exact phase and range coordinates}
The rest phase and the inverse range are two coordinates of the same finite
kernel. Their conversion at a nonzero chart spacing is
Eq.~\eqref{eq:finite_mass_conversion}. The map's rest mass is the phase
frequency divided by the common proper time. Its static range is the
spatial root of the reciprocal chart. They have the same tangent value.
The exact relation retains their finite-spacing difference before that limit
is taken.

For a smooth change of spectral coordinate $u=u(v)$, a pole response obeys
\begin{equation}
 \frac{du_0}{dJ}
 =-\frac{l^\dagger K_Jr}{l^\dagger K_ur},\qquad
 \frac{dv_0}{dJ}
 =\frac{dv}{du}\frac{du_0}{dJ},\qquad
 K_v=K_u\frac{du}{dv}.
 \label{eq:coordinatecovariance}
\end{equation}
The residue and the response transform with the same clock derivative.
A common calibration fixes the energy represented by one phase and the
length represented by one step. It preserves ratios of two rest frequencies
and the dimensionless coefficients of the completed law.

The four-spinor trace uses the coefficient
$-\ell^{-2}(R-6/\ell^2)\operatorname{vol}$ derived in Part I.
If the trace is instead written as an epsilon contraction of tensor
curvatures, that contraction is four times larger. Its action prefactor is
therefore $-\ell^2/(64\pi G)$, while the ordinary spinor-trace prefactor
is $-\ell^2/(16\pi G)$. Both give the same $\Lambda=3/\ell^2$ and
the same Einstein--Cartan variation.

\part{Attempts}\label{part:attempts}

\section{Currents and numerical choices}
\label{sec:attemptchoices}

\subsection{The measured values in the construction of currents}
\label{sec:measurementfactors}
The finite comparison supplies the return moments, projectors and transmission
coefficients. The current expressions below also contain choices made while
their measured values were in view. The following table states those choices
at the point where these applications begin. Their later finite identities
and their evaluated comparisons are retained together.

\begingroup\small
\setlength{\tabcolsep}{4pt}\renewcommand{\arraystretch}{1.12}
\begin{longtable}{@{}>{\raggedright\arraybackslash}p{.23\textwidth}>{\raggedright\arraybackslash}p{.72\textwidth}@{}}
\caption{Factors chosen with measurements in view.}\label{tab:measuredchoices}\\
\toprule Quantity and factor&Origin of the choice and its present expression\\\midrule
\endfirsthead
\toprule Quantity and factor&Origin of the choice and its present expression\\\midrule
\endhead\bottomrule\endfoot
Neutron moment, $44/23$&The number first appeared as the continued-fraction convergent $[1;1,10,2]$ of the measured magnitude. Its finite rewriting is $2\Tr C/\Tr K$, with $\Tr C=22$ and $\Tr K=23$. The factor two retains this measured origin.\\[5pt]
Proton moment, $39/5$&The bare square was first written $8(39/40)$. The return moments give the identical $\Var N_{\rm ret}/(2N_2)=156/20$.\\[5pt]
Magnetic screening, $S_\mu$&The bare moments were compared with their measurements before the quadratic correction was selected. The current uses $\epsilon_2=\alpha^2/(1+\alpha/10)$ and $S_\mu=1-\epsilon_2/(16^2b)$. The later source contraction gives $16^2b=112$.\\[5pt]
Axial transmission, $2b/3$&The exponent in $g_A=(5/3)t_0^{2b/3}$ was selected with the measured axial residue visible. The coefficient $5/3$ follows independently from the rank-four nucleon projector.\\[5pt]
Axial radius, $3\lambda_p$&The factors four, two and three were compared with the measured size. The selected three is also the square root of the conditional return variance $\EE\Var(n_F\mid n_L)=9$.\\[5pt]
Pion current, $2\pi$, $\sqrt2$, $\eta_0$&The quark-level sigma-model relation $g_{\pi qq}=2\pi/\sqrt{k_C}$ supplies $2\pi$; the convention $f_\pi=\sqrt2F_\pi$ supplies $\sqrt2$. The stationary factor $\eta_0=295/296$ was chosen after the comparison of the unsuppressed current with the measured decay constant. The resulting expression is $\eta_0 m_p\sqrt{2k_C}/(2\pi\mu_p)$.\\[5pt]
Pion-pole residue, $1+a^2$&This multiplier of the pion-pole axial residue was selected with the measured pseudoscalar current in view. It is used with the same $g_A$, masses and radius as the other nucleon currents.\\[5pt]
Deuteron electric correction, $41/40$&The leading binding source was compared with the measured binding before the term $-(1+t_0/16)\epsilon_2$ was selected. Its contribution is retained in the displayed pole equation.\\[5pt]
Deuteron tensor direction, $7/18$&The assignment $g=7/18=\alpha_s^2/\Psi_+$ gives the normalized direction $(7,-\sqrt2)/\sqrt{51}$ and $P_D=2/51$. The value $7/18$ was taken with the measured deuteron tensor component in view.\\[5pt]
Light nuclei and bulk terms&The triangle and tetrahedron counts $13/16$ and $61/16$ are the clock variance minus three and the clock variance itself. The selected bulk and surface factors were compared through trial nuclear expressions. Their reference formulas follow below.\\[5pt]
Second pion construction, $5/36$&The target decay constant was known when the factor was selected. The resulting current is $5m_p/36$.\\[5pt]
Second scalar amplitude, $b^{24}$&The record found this expression among 513 forms compared with the scalar amplitude. The power is retained as that second construction.\\[5pt]
Second cold-matter count, $37/7$&The measured cold-to-baryon density ratio was available when $M/(M-4C_h)$ was selected.\\[5pt]
Second top expression, $m_p^2/(10m_e)$&The record searched 2,970 expressions with the mass comparison in view. The selected denominator is $\EE n_F\,m_e=10m_e$.\\
\end{longtable}\endgroup

\subsection{The nucleon projector and its currents}
The state whose phase gave the two nucleon masses also carries their electric, magnetic and weak currents. Its colour contraction is Eq.~\eqref{eq:colourclosure}; its rank-four spin--isospin projector and charge doublets are Eqs.~\eqref{eq:nucleonprojectors}--\eqref{eq:nucleonprojectorranks}. The symmetric spatial and spin--isospin factors together with the antisymmetric colour factor give the fermionic exchange sign of the full state. Let $A$ and $B$ be the normalized symmetric states made from $u_\uparrow u_\uparrow d_\downarrow$ and $u_\uparrow u_\downarrow d_\uparrow$. They span the subspace with $I_3=S_z=1/2$. On this subspace the total spin is
\begin{equation}
 \left.\mathbf S^2\right|_{A,B}=
 \begin{pmatrix}7/4&\sqrt2\\ \sqrt2&11/4\end{pmatrix},
 \qquad
 \psi_{p\uparrow}=\frac{\sqrt2A-B}{\sqrt3}.
 \label{eq:nucleonspinmatrix}
\end{equation}
The eigenvalues are $3/4$ and $15/4$; the displayed state is the normalized $3/4$ eigenvector selected by $P_N$. The axial operator $\sum_r\sigma_{z,r}\tau_{z,r}$ is diagonal on $A,B$, with entries $3,-1$. Its expectation on the same state is therefore $(2\cdot3-1)/3=5/3$. Throughout the nucleon subspace the projected axial operator is $(5/3)\sigma_N\tau_N$; the vector charge is one.

The quartet's invariant bilinear form is $\epsilon_{\mathrm{spin}}\otimes\epsilon_{\mathrm{isospin}}$. It is symmetric because the two antisymmetric factors supply two minus signs. This identity fixes the invariant pairing within the quartet. The antisymmetric colour wavefunction supplies the constituent exchange sign. Projecting the physical state with $P_N$ fixes the nucleon current before its source is evaluated.

The axial current asks how much of the spin--isospin change reaches the other end of the selected plane. Their comparison multiplies the projected residue $5/3$ by the plane transmission raised to $2b/3$. The electric correction uses the phase response $\alpha$, and the magnetic source retains the return-length fluctuation $156$. Their coefficients are
\begin{align}
 g_A&=\frac53t_0^{2b/3}=1.27580092811547,\notag\\
 \epsilon_2&=\frac{\alpha^2}{1+\alpha/10},\qquad
 S_\mu=1-\frac{\epsilon_2}{16^2b},\notag\\
 \mu_p&=\sqrt{\frac{156}{20}S_\mu}=2.79284734529616,\notag\\
 \mu_n&=-\frac{44}{23}\sqrt{S_\mu}=-1.91304302380604.
 \label{eq:nucleoncurrent}
\end{align}
Both numerical moments use the nuclear magneton $e_{\mathrm{charge}}\hbar/(2m_p)$. A Pauli vertex with denominator $2m_p$ has $F_2^p(0)=\mu_p-1$ and $F_2^n(0)=\mu_n$. If the denominator is instead $2m_N$, the conversion to nuclear magnetons is $(m_p/m_N)[F_1(0)+F_2(0)]$.

For momentum transfer $Q$ in inverse-length units, the electric source has form factor $G_E(Q^2)=\Fret(e^{-Q^2\lamP^2/6})$. Its slope and the axial and pion slopes are
\begin{equation}
 r_p^2=16\lamP^2,\qquad r_A^2=9\lamP^2,\qquad
 r_\pi^2=10\lamP^2.
\end{equation}
They give $r_p=0.841235640212810\fm$ and $r_A^2=0.398068538829894\,\mathrm{fm}^2$. The current at the pion endpoint is
\begin{equation}
 f_\pi=\frac{\eta_0 m_p\sqrt6}{2\pi\mu_p}=130.529195478157\MeV,
 \qquad f_\pi^{(2)}=\frac{5m_p}{36}=130.315567975027\MeV.
 \label{eq:fpi}
\end{equation}
The first expression evaluates the pion-current residue; its $|V_{ud}|$ product has records G08 and G09. The alternative $f_\pi^{(2)}$ uses the model coefficient $5/36$ and has records A002 and A003. Chiral formulas below use $F_\pi=f_\pi/\sqrt2$.

\subsection{The complete nucleon vertex}
\label{sec:nucleonvertex}
The physical projector acts on both ends of a current. With momentum transfer
$q=p'-p$ and the mass convention stated beside the Pauli term, its vertex is
\begin{equation}
 \Gamma_N^\mu=P_N\left[
 \gamma^\mu F_1(q^2)+\frac{i\sigma^{\mu\nu}q_\nu}{2m_N}F_2(q^2)
 +\gamma^\mu\gamma_5G_A(q^2)
 +\frac{q^\mu\gamma_5}{2m_N}G_P(q^2)\right]P_N.
 \label{eq:nucleonvertex}
\end{equation}
The zero-transfer vector charges are one and zero for proton and neutron.
The measured-origin factors in Table~\ref{tab:measuredchoices} give the axial
normalization and the two magnetic moments in Eq.~\eqref{eq:nucleoncurrent}.
The same state supplies their slopes:
\begin{equation}
 G_E(Q^2)=F_{\rm ret}(e^{-Q^2\lambda_p^2/6}),\qquad
 \frac{G_A(t_q)}{g_A}=\frac{w}{4-3w},\quad w=e^{r_A^2t_q/24},\quad t_q=-Q^2.
 \label{eq:nucleonslopes}
\end{equation}
Differentiating the first expression gives $r_p^2=\EE N_{\rm ret}\,\lambda_p^2$.
Conditioning the symmetric visits on the vector visits gives
$r_A^2=\EE\Var(n_F\mid n_L)\lambda_p^2=9\lambda_p^2$.
The assignment of this conditional length to the measured axial radius is
recorded in Table~\ref{tab:measuredchoices}.

The axial shape has slope $G_A'(0)/g_A=r_A^2/6$. Its pseudoscalar partner is
\begin{align}
 G_P(t_q)&=\frac{4m_N^2g_A(1+a^2)}{m_\pi^2-t_q}
 -\frac{4m_N^2}{t_q}[G_A(t_q)-g_A],\notag\\
 2m_NG_A(t_q)+\frac{t_q}{2m_N}G_P(t_q)
 &=2m_Ng_A\frac{m_\pi^2+a^2t_q}{m_\pi^2-t_q}.
 \label{eq:axialward}
\end{align}
The term at $t_q=0$ takes its continuous limit. This is the exact axial
Ward relation of the specified pion residue.

For the pion pole, use $F_\pi=f_\pi/\sqrt2$ and the source relation
\begin{equation}
 g_{\pi NN}=\frac{g_A m_N}{F_\pi}(1+a^2).
 \label{eq:pionnucleoncoupling}
\end{equation}
The multiplier $1+a^2$ was selected with the measured pseudoscalar current
visible. In the convention in which ordinary muon capture reads the induced
coupling at $q^2=-0.88m_\mu^2$, the first terms of the retained pole and axial-radius expansion are
\begin{equation}
 g_P(q^2)=\frac{2m_\mu F_\pi g_{\pi NN}}{m_{\pi^\pm}^2-q^2}
          -\frac{m_\mu m_Ng_A r_A^2}{3}.
 \label{eq:pseudoscalarcurrent}
\end{equation}
Here $r_A$ is expressed in inverse energy units. The normalization of
Eq.~\eqref{eq:nucleonvertex} and the conventional $g_P$ use their indicated
mass factors. Equation~\eqref{eq:pseudoscalarcurrent} is the pole-current
application of the chosen residue; its measurement belongs to the same
classical comparison register as the decay widths.


The full shape in Eq.~\eqref{eq:axialward}, evaluated through the
weak capture map, gives $g_P=8.36259039416$ at $Q^2=0.88m_\mu^2$ and a
singlet muonic-hydrogen capture rate $717.110240723\,\mathrm{s}^{-1}$.
These values are carried evaluations of that map.

\subsection{Four second constructions}
\label{sec:secondconstructions}
The record contains four further expressions built from quantities already
on the line. Their measured targets were available when they were selected.
They are evaluated here in the same electron mass unit.
\begin{center}\small
\begin{tabular}{@{}lll@{}}
\toprule Reading&Expression&Value\\\midrule
Pion current&$5m_p/36$&$130.315567975027\MeV$\\
Scalar power at $0.05\,\mathrm{Mpc}^{-1}$&$b^{24}(0.05\ell)^{-1/40}$&$2.10257562469885\times10^{-9}$\\
Cold-to-baryon count&$M/(M-4C_h)=37/7$&$5.28571428571429$\\
Top mass&$m_p^2/(\EE n_Fm_e)$&$172.281080537681\GeV$\\\bottomrule
\end{tabular}
\end{center}
The measured-origin factors are specified in Table~\ref{tab:measuredchoices}.
The second pion current gives its product with $|V_{ud}|$ and its lifetime
through the same leptonic decay formula. A002 and A003 retain the current
comparisons; Table~\ref{tab:furthercomparisons} gives the lifetime.

\section{Several bound world lines}
\label{sec:many}


A bound state has one continuing phase, although several constituents contribute to it. Its constituents can return independently or preserve a common comparison throughout the return. These are physical alternatives of the same joint inverse propagator. The selected examples below evaluate the shared attenuation, the light meson and baryon source matrix elements, the deuteron tensor state and a topological response. Their numerical entries are selected composite examples.

\subsection{The normal modes of a joint return}
A constituent that persists inside a common state reduces the amplitude of steps along its own direction by $e^{-S}$, with attenuation $S\geq0$. A step in that direction and a step in any of the other four directions are the two ways its continuing relation can evolve. The two continuation classes give the positive return eigenvalue
\begin{equation}
 4\lambda_1^2-3\lambda_1-e^{-S}=0,\qquad
 \lambda_1(S)=\frac{3+\sqrt{9+16e^{-S}}}{8},\qquad
 \delta_1(S)=-\log\lambda_1(S).
 \label{eq:rungroot}
\end{equation}
It has $\delta_1(0)=0$, $\delta_1'(0)=1/5$, and increases to $\log(4/3)$. The stable evaluation at small $S$ is
\begin{equation}
 \delta_1(S)=
 -\operatorname{log1p}\!\left[
 \frac{2\operatorname{expm1}(-S)}{\sqrt{9+16e^{-S}}+5}\right].
\end{equation}
The second root $(3-\sqrt{9+16e^{-S}})/8$ is the opposite branch of the same quadratic. The retained positive response is Eq.~\eqref{eq:rungroot}.

For $n$ such attenuations put $D=\diag(S_1,\ldots,S_n)$. Their common state has a positive Hermitian overlap matrix $G$, normalized by $G_{ii}=1$. Its retained response is
\begin{equation}
 R=D^{1/2}GD^{1/2},\qquad
 \delta(S_1,\ldots,S_n;G)=\Tr\delta_1(R).
 \label{eq:join}
\end{equation}
To see why the eigenvalues of $R$ carry the composition, let $J$ be a one-particle source insertion. The joint determinant has the factorization
\begin{equation}
 \det[\Id-(GD)\otimes K^{-1}J]
 =\prod_{\lambda\in\operatorname{spec}R}\det[\Id-\lambda K^{-1}J].
 \label{eq:joindet}
\end{equation}
The identity extends continuously to zero attenuations. Independent sources, $G=\Id$, give $\sum_i\delta_1(S_i)$; a fully locked source, $G=J_n$, gives $\delta_1(\sum_iS_i)$. Thus partial overlap is read in the spectrum of the actual joint matrix.

For two lines, 32 of the 60 stationary departures along the held direction move their clock. Their overlap coefficient is
\begin{equation}
 c_{\mathrm{pair}}=\frac{32}{60}=\frac8{15},\qquad
 R(S_1,S_2)=
 \begin{pmatrix}S_1&\sqrt{c_{\mathrm{pair}}S_1S_2}\\\sqrt{c_{\mathrm{pair}}S_1S_2}&S_2\end{pmatrix}.
 \label{eq:twoattenuations}
\end{equation}
Writing $\mathcal R(S_1,S_2)=\Tr\delta_1(R)$, the expansion begins
\begin{equation}
 \delta_1(S)=\frac S5-\frac{6S^2}{125}+\frac{2S^3}{3125}+O(S^4).
\end{equation}
The first joint contribution, $-12c_{\mathrm{pair}}S_1S_2/125$, lowers the squared mass relative to the sum of independent responses. At fourth order the spectral expression differs from an interpolation of the independent and locked endpoint responses by $-2c_{\mathrm{pair}}(1-c_{\mathrm{pair}})d_4S_1^2S_2^2$, where $d_4$ is the fourth Taylor coefficient of $\delta_1$. The exact matrix definition fixes all orders.

For three equal attenuations, relative phases also enter the closed overlap. If the off-diagonal entries have modulus $\sqrt{c_{\mathrm{pair}}}$, their product phase $\phi$ gives \begin{equation}
 \det G=1-3c_{\mathrm{pair}}+2c_{\mathrm{pair}}^{3/2}\cos\phi.
 \label{eq:threegram}
\end{equation} For a flat positive overlap and $n$ equal attenuations, $G=(1-\sqrt{c_{\mathrm{pair}}})\Id+\sqrt{c_{\mathrm{pair}}}J_n$ has one eigenvalue $1+(n-1)\sqrt{c_{\mathrm{pair}}}$ and $n-1$ eigenvalues $1-\sqrt{c_{\mathrm{pair}}}$. The phases around a triangle are part of the physical joint state. At $c_{\mathrm{pair}}=8/15$, positivity selects zero total phase from the four quarter-turn choices. More generally, normalized internal vectors $a_i$ give a positive overlap for any number of particles:
\begin{equation}
 G_{ij}=(1-\sqrt{c_{\mathrm{pair}}})\delta_{ij}+\sqrt{c_{\mathrm{pair}}}\,\langle a_i,a_j\rangle,
 \qquad G=(1-\sqrt{c_{\mathrm{pair}}})\Id+\sqrt{c_{\mathrm{pair}}}\,\mathsf A^\dagger\mathsf A.
 \label{eq:positiveoverlap}
\end{equation}
Here $\mathsf A$ has columns $a_i$. Their actual spin and flavour determine these vectors. For the flat equal overlap the common shift is explicitly
\begin{equation}
 \delta_n(S)=(n-1)\delta_1((1-\sqrt{c_{\mathrm{pair}}})S)
 +\delta_1((1+(n-1)\sqrt{c_{\mathrm{pair}}})S).
 \label{eq:equaloverlap}
\end{equation}
Each positive eigenmode saturates at $\log(4/3)$. A two-mode meson on this branch consequently has $m\leq m_p\sqrt{12\log(4/3)}$. The selected light states below lie on this branch. The full source inverse propagator retains its energy dependence as in Eq.~\eqref{eq:selfenergy}.

\subsection{The light mesons and the first baryon states}
A pseudoscalar pairs two light-quark directions with the sign of the even return. If its endpoint amplitudes are $x,y$, the alternating transfer $T=\left(\begin{smallmatrix}0&x\\y&0\end{smallmatrix}\right)$ obeys $T^2=xy\Id$. The rest frequency reads the product. Each endpoint carries the inverse signed return and both share the clock overlap, giving the completed source
\begin{equation}
 xy=\hover m_em_p\Tr_{\rm ends}\frac{\Id_2}{|\Phis(1)|},\qquad
 K_{\pi^0}(E)=E^2-xy.
 \label{eq:pionendpointsource}
\end{equation}
Here $m_e,m_p$ are the reference rest rates in which the source is read. The electric source adds the contact of the same two constituent endpoints. Evaluating the trace gives
\begin{equation}
 m_{\pi^0}^2=\frac{2\hover}{|\Phis(1)|}m_em_p=38\hover\,m_em_p,\qquad
 \Delta_{\mathrm{el}}=\frac{\pi\alpha m_p^2}{16[1-\alpha b\Phis(\ii)]}.
\end{equation}
The charged source has $m_{\pi^\pm}^2=m_{\pi^0}^2+\Delta_{\mathrm{el}}$, giving
\begin{equation}
 m_{\pi^0}=134.977044560\MeV,\qquad
 m_{\pi^\pm}=139.570386614\MeV.
 \label{eq:pions}
\end{equation}
The same neutral-pion rest frequency fixes how far the two light constituent sources attenuate their common return. The up/down current ratio and the strange/down grade ratio then determine the remaining light attenuations:
\begin{equation}
 6m_p^2\mathcal R(R_{ud}S_d,S_d)=m_{\pi^0}^2,\qquad
 S_u=R_{ud}S_d,\qquad S_s=p_qS_d.
 \label{eq:lightroot}
\end{equation}
Along the positive ray both eigenvalues of $R$ and both values of $\delta_1$ increase, so the selected root is unique. The program uses ninety bisections in $[0,1]$. Approximately,
$(S_u,S_d,S_s)=(.005447267,.011855816,.236661685)$. The calibration in Eq.~\eqref{eq:lightroot} is to Eq.~\eqref{eq:pions}, whose input is the same finite source coefficients.

The strange two-constituent mass zeros are
\begin{align}
 m_{K^0}^2&=6m_p^2\mathcal R(S_d,S_s),\notag\\
 m_{K^\pm}^2&=6m_p^2\mathcal R(S_u,S_s)+\frac{48}{25}\Delta_{\mathrm{el}}.\label{eq:kaons}
\end{align}
They give $m_{K^0}=497.613584451\MeV$ and $m_{K^\pm}=493.665032668308\MeV$. The factor $48/25$ is the electric source contraction in the symmetric sector.

The symmetric-state visit generating function for a first-return excursion is
\begin{equation}
 \EE x^{n_F}=\frac{x(1+x)}{16-11x-3x^2}.
\end{equation}
Its first positive singularity $x_c=(\sqrt{313}-11)/6$ gives the vector squared-pole reading. The other evaluated first composite states give
\begin{align}
 \Re s_\rho&=6m_p^2\log x_c=.576408555494\,\mathrm{GeV}^2,\notag\\
 m_\Lambda&=m_n(5/2)^{3/16}=1115.682900908\MeV,\notag\\
 m_{\Delta,\mathrm{count}}&=3bm_p=1231.482117364\MeV,\notag\\
 k_BT_c&=\frac{m_p}{6}=156.378681570\MeV.
\end{align}
The $\rho$ record reports $\Re s_\rho$. The square root of that real number is a different reading. The finite $\Delta$ value is compared here in the Breit--Wigner convention. Its separation from the nucleon also has the spin--spin form
\begin{equation}
 A_{\mathrm{spin}}=\frac{(3b-1)m_p}{6},\qquad
 U_{\mathrm{const}}=\frac{m_p}{3}+A_{\mathrm{spin}},\qquad
 m_p=3U_{\mathrm{const}}-3A_{\mathrm{spin}},\quad
 m_{\Delta,\mathrm{count}}=3U_{\mathrm{const}}+3A_{\mathrm{spin}}.
 \label{eq:spinmass}
\end{equation}
The few-MeV current masses in Eq.~\eqref{eq:quarks} describe the individual quark sources; the much larger constituent contribution is the energy retained in their common colour state. The neutral-pion identity makes the linear dependence on the light current masses explicit:
\begin{equation}
 m_{\pi^0}^2=C_\pi m_p(m_u^{\rm native}+m_d^{\rm native}),\qquad
 C_\pi=\frac{38\hover\eta_0}{3^{3/2}(1+R_{ud})}.
 \label{eq:chiralidentity}
\end{equation}
Here both quark masses use the common source scale; carrying the scalar density and the masses to another scale keeps their product in the same convention.

\subsection{The deuteron pole, tensor direction, and magnetic residue}
The deuteron source retains a common return of one proton and one neutron. Their three shared direction responses contribute $a^3$; the axial comparison contributes $\sqrt{\chi_A}$. The electric current changes the same binding source through $\epsilon_2$. Its binding-energy source is
\begin{equation}
 B_d=m_p\sqrt{\chi_A}\,a^3
 \left[1-\left(1+\frac{t_0}{16}\right)\epsilon_2\right]
 =2.22456646309919\MeV,
 \qquad m_{\rm deut}=m_p+m_n-B_d.
 \label{eq:deuteron}
\end{equation}
Its mass is $1875.61294498229\MeV$. The tensor interaction couples the $S$ and $D$ orbital states with $g=7/18$:
\begin{equation}
 T_{SD}=\begin{pmatrix}0&\sqrt8g\\\sqrt8g&6-2g\end{pmatrix}.
\end{equation}
Its lower eigenvalue is $-2/9$. The corresponding state has $D/S$ component ratio $-\sqrt2/7$ and internal $D$ fraction $p_{\rm D}=2/51$, so the tensor force determines how much orbital angular momentum is present in this retained state. Let $\mu_S=\mu_p+\mu_n$ and $g_l=m_n/(m_p+m_n)$. Projection of the magnetic operator gives
\begin{equation}
 (1-\tfrac32p_{\rm D})\mu_S+\tfrac32p_{\rm D}g_l=\frac{16\mu_S+g_l}{17}.
\end{equation}
The retained binding response adds its source derivative:
\begin{equation}
 \mu_d=\frac{16\mu_S+g_l}{17}
 -\frac{(B_d/m_p)(\mu_S-g_l)}{20+13B_d/m_p}
 =.857438233519.\label{eq:magneticdeuteron}
\end{equation}
The displayed $D/S$ component ratio is the internal tensor eigenvector. An asymptotic orbital ratio additionally includes spatial propagation outside the interaction region. The retained example here is the internal two-channel state and its magnetic residue.

\subsection{The selected bulk energy}
The extensive energy of a large symmetric population compares the occupied return of one nucleon with the same population at zero binding. The selected volume source retains two reversal factors and the one-return tilt. It gives
\begin{equation}
 \epsilon_{\rm bulk}=-m_pa^2\left(1-\frac{t_0}{16}\right)
 =-15.8382148058292\MeV.
 \label{eq:bulkenergy}
\end{equation}
The minus sign is the decrease of energy on binding; the coefficient is the energy per nucleon in the symmetric-matter limit. Row N22 compares this volume term with the extrapolation of nuclear masses. The light mesons, deuteron and this volume term constitute the selected composite calculations in the register.

\subsection{Topological susceptibility of the light source}
A topological phase can be shared among the three light flavours. The state distributes it so that its energy is least; the curvature of that energy measures the topological susceptibility. The relative light attenuations are $S=(R_{ud},1,p_q)$. On their two neutral contrasts define
\begin{equation}
 C_0=\begin{pmatrix}
 R_{ud}+1&(R_{ud}-1)/\sqrt3\\
 (R_{ud}-1)/\sqrt3&(R_{ud}+1+4p_q)/3
 \end{pmatrix},\qquad
 \lambda_{\mathrm{low}}=\min\operatorname{spec}C_0.
\end{equation}
Distribute the total phase $\theta$ between the three light attenuations by minimizing
$V=\sum_iS_i(1-\cos\phi_i)$ subject to $\sum_i\phi_i=\theta$. The multiplier satisfies $S_i\sin\phi_i=\lambda$. With $A_1=\sum_iS_i^{-1}$ and $A_3=\sum_iS_i^{-3}$,
\begin{equation}
 V(\theta)=\frac{\theta^2}{2A_1}
 -\frac{A_3\theta^4}{24A_1^4}+O(\theta^6).
\end{equation}
The physical source energy is $E_{\mathrm{top}}=(F_\pi^2m_{\pi^0}^2/\lambda_{\mathrm{low}})V$. Its susceptibility and fourth cumulant are therefore
\begin{equation}
 \frac{\chi_{\mathrm{top}}}{F_\pi^2m_{\pi^0}^2}
 =\frac1{\lambda_{\mathrm{low}}A_1}=.212911371927,\qquad
 b_2=-\frac{A_3}{12A_1^3}=-.028058473099,\label{eq:topology}
\end{equation}
where $E_{\mathrm{top}}(\theta)-E_{\mathrm{top}}(0)
=\frac12\chi_{\mathrm{top}}\theta^2(1+b_2\theta^2+\cdots)$ and $F_\pi=f_\pi/\sqrt2$.

\subsection{Both roots in the composition of several lines}
\label{sec:manyreference}
The connected current form on $r$-subsets is
\begin{equation}
 \mathcal C^{(r)}=\frac8{25}[r(5-r)\Id-A_{J(5,r)}].
 \label{eq:manycurrent}
\end{equation}
For the two roots define
\begin{equation}
 r_\pm(S)=\frac{3\pm\sqrt{9+16e^{-S}}}{8},\qquad
 \delta_\pm(S)=-\log|r_\pm(S)|.
\end{equation}
The complete two-root equation then has
\begin{equation}
 r_+r_-=-e^{-S}/4,\qquad \delta_++\delta_-=S+\log4.
 \label{eq:bothrestroots}
\end{equation}
Heavy-line continuation
retains both roots of Eq.~\eqref{eq:bothrestroots}. 

\subsubsection{The spatial two-nucleon kernel}\label{the-spatial-two-nucleon-kernel}

The spatial two-nucleon kernel in the retained nuclear calculation contains the pion pole of Eq.~\eqref{eq:pions}, the two axial form factors \eqref{eq:nucleonslopes}, the coupling \(g_{\pi NN}\), and a contact

\begin{equation}
C_{NN}=C_{NN}^{(0)}+C_\tau\boldsymbol\tau_1\cdot\boldsymbol\tau_2,
\qquad C_\tau=\frac1{8F_\pi^2}.
\label{eq:attempt-S10-5}
\end{equation}

The isoscalar contact \(C_{NN}^{(0)}\) is fixed by the already calculated pole \(B_d\).
The radial solution, its retained-state norm and the two-by-two angular block
give the pole residue, asymptotic D/S ratio and scattering observables.
The singlet pole also reads the isospin contact. The radial
outputs and their comparisons are collected in Table~\ref{tab:furthercomparisons}.

\subsubsection{Light nuclei and their exact mass-difference relations}\label{light-nuclei-and-their-exact-mass-difference-relations}

The ten clock gaps have variance

\begin{equation}
v_D=\frac{61}{16}=3+\chi_A\frac{\operatorname{Var}n_L}{\EE N_{\rm ret}}
=\frac{(\operatorname{Var}n_F-\operatorname{Var}n_L)^2+4\operatorname{Cov}(n_F,n_L)^2}{9(\EE N_{\rm ret})^2}.
\end{equation}

The recorded incidence polynomial for \(A\le4\) is

\begin{equation}
B_{\rm nuc}(A,Z)=B_d\left[\binom A2+(v_D-3)\binom A3+v_D\binom A4\right]
-\frac{\alpha m_p}{V_C}\binom Z2.
\label{eq:attempt-S11-1}
\end{equation}

Its pair, triple and quadruple coefficients are thus \(1\), \(13/16\) and \(61/16\). The Coulomb coefficient is \(\alpha m_p/9=0.760766915307\ \mathrm{MeV}\). The resulting nuclear masses are \(M(A,Z)=Zm_p+(A-Z)m_n-B_{\rm nuc}(A,Z)\).

\begin{longtable}[]{@{}
  >{\raggedright\arraybackslash}p{(\columnwidth - 6\tabcolsep) * \real{0.2000}}
  >{\raggedleft\arraybackslash}p{(\columnwidth - 6\tabcolsep) * \real{0.2667}}
  >{\raggedleft\arraybackslash}p{(\columnwidth - 6\tabcolsep) * \real{0.2667}}
  >{\raggedleft\arraybackslash}p{(\columnwidth - 6\tabcolsep) * \real{0.2667}}@{}}
\toprule\noalign{}
\begin{minipage}[b]{\linewidth}\raggedright
Nucleus
\end{minipage} & \begin{minipage}[b]{\linewidth}\raggedleft
Binding energy, MeV
\end{minipage} & \begin{minipage}[b]{\linewidth}\raggedleft
Nuclear mass, MeV
\end{minipage} & \begin{minipage}[b]{\linewidth}\raggedleft
Binding residual against the AME2020 values retained in the record
\end{minipage} \\
\midrule\noalign{}
\endhead
\bottomrule\noalign{}
\endlastfoot
Deuteron & 2.224566463099 & 1875.612944982 & Pole specified in \eqref{eq:deuteron} \\
Triton & 8.481159640566 & 2808.921773830 & \(-0.637\ \mathrm{keV}\) \\
Helion & 7.720392725258 & 2808.389208140 & \(+2.403\ \mathrm{keV}\) \\
Helium-4 & 28.297632508926 & 3727.377390382 & \(+2.022\ \mathrm{keV}\) \\
\end{longtable}

The mirror splitting in this polynomial is the single Coulomb coefficient
above. Binding energies and absolute masses have separate comparison
uncertainties. Table~\ref{tab:furthercomparisons} retains the mirror splitting
and the absolute-mass residuals of Proposition 186; row M3 gives the
deuteron mass comparison.

The six light-nucleus reaction energies can be defined explicitly as

\begin{equation}
\begin{array}{lll}
Q_1=B_d,&Q_2=B_h-B_d,&Q_3=B_h-2B_d,\\
Q_4=B_t-2B_d,&Q_5=B_\alpha-B_t-B_d,&Q_6=B_\alpha-B_h-B_d.
\end{array}
\end{equation}

These correspond respectively to \(np\to d\gamma\), \(dp\to h\gamma\), \(dd\to hn\), \(dd\to tp\), \(dt\to\alpha n\), and \(dh\to\alpha p\). Their energy releases satisfy

\begin{equation}
Q_2-Q_3-Q_1=0,\qquad Q_6-Q_5-Q_4+Q_3=0.
\label{eq:attempt-S11-2}
\end{equation}

These two closed threshold cycles follow from mass subtraction using the same four binding energies. Every further reaction threshold is obtained from the same masses and conserves charge and baryon number.

\subsubsection{Helium clusters, shells and bulk nuclear matter}\label{helium-clusters-shells-and-bulk-nuclear-matter}

For \(n\) helium-4 clusters, the retained bond formula is

\begin{equation}
B_{\rm nuc}(4n)=nB_\alpha+\beta_n\chi_AB_d,
\end{equation}

\begin{equation}
\beta_1=\beta_2=0,\quad\beta_3=3,\quad\beta_4=6,\quad\beta_5=8,
\qquad \beta_n=3n-6\quad(n\ge6).
\label{eq:attempt-S12-1}
\end{equation}

The coefficients are the stated contact and rigidity counts of the cluster centres on the \(A_3\) lattice. In particular

\begin{equation}
B_{\rm nuc}({}^{20}{\rm Ne})=5B_\alpha+8\chi_AB_d=160.767738558\ \mathrm{MeV}
\end{equation}

in the nuclear convention. For five clusters the specified contact count
is eight. Table~\ref{tab:furthercomparisons} retains the atomic comparison
and the silicon-28 and sulfur-32 bond comparisons.

The tetrahedral cone \(x_1+x_2+x_3\le N\) contains \(\binom{N+3}3\) lattice points. Including the two spin states gives closed counts

\begin{equation}
2\binom{N+3}3=2,8,20,40,70,112,168,\ldots .
\end{equation}

The stated aligned next-edge addition \(2N+4\) gives

\begin{equation}
2\binom{N+3}3+2N+4=28,50,82,126,184
\qquad(N=2,3,4,5,6).
\label{eq:attempt-S12-2}
\end{equation}

This supplies the shell counts used in the record. It specifies which boundary states are included; the full ordering of energy levels is a further property of the retained nuclear Hamiltonian.

The volume scale and saturation point are

\begin{equation}
r_{\rm sat}=\frac{3\lambda_p}{4a}=1.198760787303\ \mathrm{fm},\qquad
n_0=\frac3{4\pi r_{\rm sat}^3}=0.138584227666\ \mathrm{fm}^{-3},
\end{equation}

\begin{equation}
\epsilon_{\rm bulk}=-m_pa^2\frac{39}{40}=-15.8382148058\ \mathrm{MeV}.
\label{eq:attempt-S12-3}
\end{equation}

The historical bulk polynomial has the liquid-drop form

\begin{equation}
B_{\rm bulk}(A,Z)=a_vA-a_sA^{2/3}
-a_a\frac{(A-2Z)^2}{A}+a_p\frac{\delta_{AZ}}{\sqrt A}
-a_c\frac{Z(Z-1)}{A^{1/3}},
\label{eq:attempt-S12-4}
\end{equation}

where \(\delta_{AZ}=+1,0,-1\) for even-even, odd-\(A\), odd-odd nuclei. Its five specified coefficients are

\begingroup
\setlength{\LTpre}{3pt}
\begin{longtable}[]{@{}llr@{}}
\toprule\noalign{}
Coefficient & Formula & MeV \\
\midrule\noalign{}
\endhead
\bottomrule\noalign{}
\endlastfoot
Volume \(a_v\) & \(m_pa^2(39/40)\) & 15.838214806 \\
Surface \(a_s\) & \(m_pa^2\chi_A\) & 17.598016451 \\
Asymmetry \(a_a\) & \(m_pa^2 r_0/t_0\) & 24.366484317 \\
Pairing \(a_p\) & \((3/4)m_pa^2\) & 12.183242158 \\
Coulomb \(a_c\) & \((4/5)\alpha a m_p\) & 0.720726551 \\
\end{longtable}
\endgroup

The bulk polynomial and its specified shell correction are compared with
the same nuclear data in Table~\ref{tab:furthercomparisons}.

The folded charge-radius prescription is

\begin{equation}
r_c^2=\frac35r_{\rm sat}^2A^{2/3}+r_p^2+\frac{A-Z}{Z}r_n^2+\frac34\lambda_p^2.
\label{eq:attempt-S12-5}
\end{equation}

Its nuclear-radius comparisons, including the helium-4 value, are retained
in Table~\ref{tab:furthercomparisons}.

The later nuclear-matter calculation includes pion and two-pion exchange, the \(\Delta\), the \(\eta_8\) pole and the isospin contact \eqref{eq:attempt-S10-5}. Both Pauli blockings are carried on both sides of each meeting, giving the source factor \(R_\sigma(n)^2[1-f(n)]\). The retained evaluations are

\begin{equation}
K_0=9n_0^2e''(n_0)=219.89\ \mathrm{MeV},\quad
\frac{m^*}{m}=0.6261,\quad V-S=635.07\ \mathrm{MeV},
\end{equation}

\begin{equation}
J_{\rm sym}=30.67\ \mathrm{MeV},\qquad L_{\rm sym}=64.97\ \mathrm{MeV}.
\label{eq:attempt-S12-6}
\end{equation}

These replace the older one-strength values \(K_0\simeq550\), \(J_{\rm sym}=24.37\), \(L_{\rm sym}=84.30\ \mathrm{MeV}\). They are carried field-equation evaluations of the stated nuclear-matter model.

\paragraph*{The composite expressions as written.}
\label{sec:literalcompositesSI}

The two-mode positive return of Sec.~\ref{sec:many} can be evaluated for every displayed light-flavour pair. With $c_{\mathrm{pair}}=8/15$ define
\begin{align}
 \lambda_\pm(s_1,s_2)&=
 \frac{s_1+s_2\pm\sqrt{(s_1-s_2)^2+4c_{\mathrm{pair}}s_1s_2}}2,\notag\\
 \delta(s)&=-\log\frac{3+\sqrt{9+16e^{-s}}}{8},\notag\\
 P(s_1,s_2)&=m_p\sqrt{6[\delta(\lambda_+)+\delta(\lambda_-)]}.
 \label{eq:literalpairSI}
\end{align}
The attenuations obtained from the preceding light-pion and charged-line sources are
\begin{equation}
 (S_u,S_d,S_s)
 =(0.00544726697196406,\ 0.01185581635074530,\ 0.236661684838052).
 \label{eq:literalattenuationsSI}
\end{equation}
The neutral expressions in these two rows are exactly $P(S_d,S_d)$ and $P(S_s,S_s)$. They give $157.9250$ and $675.9592$ MeV. The earlier alternative prescription with neutral output $557.3912$ and $405.5233$ MeV uses different compositions. Each value belongs to its stated source.

For the vector expressions set
\begin{equation}
 s_R=\operatorname{Re}s_\rho
 =0.576408555494293\ {\rm GeV}^2
 =576408.555494293\ {\rm MeV}^2.
 \label{eq:realsquaredpoleSI}
\end{equation}
This convention takes the real part of the squared pole before taking its square root. A complex pole mass instead requires both the real and imaginary parts of the pole. The sums with a kaon or pair mass squared use this real squared-pole convention in $\mathrm{MeV}^2$.

\begingroup\small
\setlength{\tabcolsep}{4pt}
\renewcommand{\arraystretch}{1.1}
\begin{longtable}{@{}l p{.45\textwidth} rr@{}}
\caption{Completed evaluations of the stated composite expressions. The last column uses the experimental uncertainties retained in Part IV.}\label{tab:literalcompositesSI}\\
\toprule State&Expression&Mass (MeV)&Deviation ($\sigma$)\\\midrule
\endfirsthead
\toprule State&Expression&Mass (MeV)&Deviation ($\sigma$)\\\midrule
\endhead
\bottomrule\endfoot
$\pi^0$&$\sqrt{[P(S_u,S_u)^2+P(S_d,S_d)^2]/2}$&134.956486147&$-40.63$\\
$\eta$&$P(S_d,S_d)$&157.925020972&$-22937.47$\\
$\eta'$&$P(S_s,S_s)$&675.959220301&$-4697.01$\\
$\omega$&$\sqrt{s_R}$&759.215750294&$-180.34$\\
$K^{*\pm}$&$\sqrt{s_R+m_{K^\pm}^2}$&905.601302988&$+59.66$\\
$K^{*0}$&$\sqrt{s_R+m_{K^0}^2}$&907.759789220&$+61.00$\\
$\phi$&$\sqrt{s_R+P(S_s,S_s)^2}$&1016.528122092&$-183.24$\\
\pagebreak
$\Sigma^0$&$m_n(5/2)^{59/320}$&1112.492814405&$-3339.55$\\
$\Xi^0$&$m_n(5/2)^{118/320}$&1317.247562638&$+11.94$\\
$\Sigma^*$, mean&$3bm_p(5/2)^{20/320}$&1304.065364235&$-223.74$\\
$\Xi^*$, mean&$3bm_p(5/2)^{40/320}$&1380.926649457&$-448.45$\\
$\Omega^-$&$3bm_p(5/2)^{60/320}$&1462.318119536&$-724.59$\\
\end{longtable}
\endgroup

The displayed powers and the preceding proton, neutron, pion and kaon values give these discrepancies. The positive two-mode branch also has a strict upper limit. For $s\geq0$, $3/4<\exp[-\delta(s)]\leq1$, so
\begin{equation}
 0\leq\delta(s)<\log(4/3),\qquad
 P(s_1,s_2)<m_p\sqrt{12\log(4/3)}
 =1743.314811563\ {\rm MeV}.
 \label{eq:heavybranchboundSI}
\end{equation}
The comparison of this bound with the heavy mesons is retained in
Table~\ref{tab:furthercomparisons}. Heavy-line continuation uses both roots
of Eq.~\eqref{eq:bothrestroots}.

\section{Conventional evolution and measurement}
\label{sec:evolution}


A number read in an experiment belongs to the complete encounter. A quark examined at two GeV has contributed to a different range of quantum fluctuations from the same quark at the neutral-vector scale. A measured magnetic frequency includes the surrounding medium. An acoustic angle compares a ruler at emission with an angle at reception after the intervening expansion. The finite source must be followed through these circumstances. The equations in this section give that continuation, and the tables retain the scale and measurement convention of every result.

\subsection{Running of the strong coupling and quark masses}
A probe of shorter wavelength includes fluctuations that a longer-wavelength probe cannot resolve. The coupling and quark mass at its scale describe the resulting current. In a specified subtraction scheme their renormalization-group equations are
\begin{equation}
 \mu\frac{\dd\alpha_s}{\dd\mu}=\beta(\alpha_s),\qquad
 \mu\frac{\dd\log\bar m_q}{\dd\mu}=-\gamma_m(\alpha_s).
\end{equation}
At fixed flavour number and one loop,
\begin{align}
 \beta_0&=11-\frac23n_f,\qquad
 \alpha_s(\mu)=\frac{\alpha_s(\mu_0)}
 {1+\beta_0\alpha_s(\mu_0)\log(\mu/\mu_0)/(2\pi)},\notag\\
 \frac{\bar m_q(\mu)}{\bar m_q(\mu_0)}
 &=\left[\frac{\alpha_s(\mu)}{\alpha_s(\mu_0)}\right]^{4/\beta_0}.
\end{align}
The specified five-flavour input is $\alpha_s(M_Z)=1/\sqrt{72}$. The one-loop scale evaluated as S16 is
\begin{equation}
 \Lambda_5^{(1)}=M_Z\exp\!\left[-\frac{2\pi}{(23/3)\alpha_s(M_Z)}\right]
 =87.0577778265\MeV.
\end{equation}

For the four-loop integration\cite{Beta4}, put $a_s=\alpha_s/(4\pi)$ and
\begin{align}
 \frac{\dd a_s}{\dd\log\mu^2}&=-a_s^2P(a_s),\qquad
 P(x)=\beta_0+\beta_1x+\beta_2x^2+\beta_3x^3,\notag\\
 \beta_0&=\frac{23}{3},\quad \beta_1=\frac{116}{3},\quad
 \beta_2=\frac{9769}{54},\quad
 \beta_3=-\frac{598391}{1458}+\frac{352864}{81}\zeta(3).
\end{align}
The register's scale convention is fixed by the finite integral
\begin{align}
 \log\frac{\mu^2}{(\Lambda_5^{(4)})^2}
 &=\frac1{\beta_0a_s}+\frac{\beta_1}{\beta_0^2}\log(\beta_0a_s)
 +\int_0^{a_s}I(x)\,\dd x,\notag\\
 I(x)&=\frac{\beta_0\beta_2-\beta_1^2+
 (\beta_0\beta_3-\beta_1\beta_2)x-\beta_1\beta_3x^2}
 {\beta_0^2P(x)}.
 \label{eq:lambda4}
\end{align}
At $a_s=1/(4\pi\sqrt{72})$, the integral is $0.000528291238170220$ and the logarithm is $12.1769228349489$. Equation~\eqref{eq:lambda4} gives the three retained scale entries:
\begin{center}
\begin{tabular}{rrc}
\toprule
$\mu$ [GeV]&$\Lambda_5^{(4)}$ [MeV]&Record\\\midrule
$91.1853615183457$&$206.890203783025$&S14\\
$91.1876$&$206.895282667594$&S017\\
$91.188$&$206.896190226441$&S018\\\bottomrule
\end{tabular}
\end{center}
The corresponding truncated inverse-log expansion through order $\log^{-4}(\mu^2/\Lambda^2)$ gives $206.625917178\MeV$ at the constructed $M_Z$. The integral specifies which definition is used here.

The locally evaluated top sources are the native family mass
$m_t^{\rm native}=(17/37)\kappa_u^{15/2}m_b^{\rm native}$, the radial response
$m_t^{\rm rad}=v\sqrt{\eta_0/(2\sigma)}$, and the second construction S07,
$m_p^2/(10m_e)$. B07 records the up-family invariant $m_t^2m_u/m_c^3$,
equal to the down-family invariant of Eq.~\eqref{eq:ladderinvariant}.

I01--I05 carry the source's four-loop running and matched quark masses at $2\GeV$ or the self-mass scale. T08 records the vector $(\bar m_t(\bar m_t),y_t)=(162.984094960692\GeV,.936132971038294)$, and I07 records the matched top pole $172.76121248356\GeV$. Flavour thresholds, decoupling coefficients and the subtraction prescription form part of those imported calculations. The local programs evaluate the three source quantities above and retain T08 and I07 as records of the preceding calculations. The explicit convention and its reevaluated translation are given below.

% Insert in Part III after the four-loop running equations and before the remaining conventional records.
\subsection{The mass coordinates evaluated by the repository}
\label{sec:qcd-repository-protocol}\label{sec:quarkconversions}
The program \path{programs/article_results.py} evaluates the native quark roots
with CRunDec~3.1\cite{RunDec,CRunDec} through its \texttt{rundec}~0.7 Python binding. All masses in
the calls below are in GeV. The starting coupling and scale are calculated
before any comparison record is read:
\begin{equation}
 \alpha_s(M_Z)=0.11785113019775792,\qquad
 M_Z=91.18536151834573\ \mathrm{GeV}.
\end{equation}
The native roots are identified with $\overline{\mathrm{MS}}$ coordinates at
this scale. Five-flavour running first determines
$\overline m_b(\overline m_b)$. At that scale the coupling and the lighter
masses are matched to four flavours. The up, down and strange masses are
then evaluated at $2\ \mathrm{GeV}$, and the charm at its own running mass.
Running and matching use four loops. With \texttt{r=CRunDec()}, the calls are
\begin{verbatim}
mb = r.mMS2mSI(mb_native, alpha_s, MZ, 5, 4)
ab5 = r.AlphasExact(alpha_s, MZ, mb, 5, 4)
ab4 = r.DecAsDownSI(ab5, mb, mb, 4, 4)
a2 = r.AlphasExact(ab4, mb, 2.0, 4, 4)
m5 = r.mMS2mMS(m_native, alpha_s, ab5, 5, 4)
m4 = r.DecMqDownSI(m5, ab5, mb, mb, 4, 4)
m_at_2 = r.mMS2mMS(m4, ab4, a2, 4, 4)
mc_self = r.mMS2mSI(mc4, ab4, mb, 4, 4)
\end{verbatim}
Here \texttt{mc4} is the matched charm mass at the bottom scale. These calls
give
\begin{align}
 \overline m_b(\overline m_b)&=4.182482142009862\ \mathrm{GeV},\notag\\
 \alpha_s^{(5)}(\overline m_b)&=0.2240093689465201,\notag\\
 \alpha_s^{(4)}(\overline m_b)&=0.2242348160465256,\notag\\
 \alpha_s^{(4)}(2\ \mathrm{GeV})&=0.3004312312043321.
\end{align}
These four quantities appear at \texttt{AR QCD.thresholds}, line~80 of
\path{outputs/article_results.out}.

The top uses the explicit pole convention: five-flavour running to
$\overline m_t(\overline m_t)$, followed by the on-shell conversion with five
massless light flavours. For loop order $n$ the calls are
\begin{verbatim}
mt_self = r.mMS2mSI(mt_native, alpha_s, MZ, 5, n)
at = r.AlphasExact(alpha_s, MZ, mt_self, 5, n)
Mt = r.mMS2mOS(mt_self, None, at, mt_self, 6, n)
\end{verbatim}
The last argument $6$ counts the heavy top and its five light flavours.
The finite input is $m_t^{\rm native}=170.94996560596294\ \mathrm{GeV}$.
The order-by-order values are
\begin{center}\small
\begin{tabular}{rrrr}
\toprule
Loops&$\overline m_t(\overline m_t)$ [GeV]&$\alpha_s(\overline m_t)$&$M_t$ [GeV]\\
\midrule
1&163.881367381795&0.108688391276874&171.441015897706\\
2&163.084433275524&0.108379548012411&172.184593539940\\
3&163.024246903563&0.108369207513217&172.612763808905\\
4&163.021143078409&0.108365685769224&172.802744094651\\
\bottomrule
\end{tabular}
\end{center}
They appear at \texttt{AR QCD.top.1} through \texttt{AR QCD.top.4},
lines~76--79. The scalar radial response remains
$m_t^{\rm rad}=172.646613614442\ \mathrm{GeV}$; it is the distinct rest
quantity defined in the article.

Table~\ref{tab:qcd-repository-records} places the explicit prescription beside
the conventional values retained in the final article. The article records are inputs in \path{data/conventional_records.json}. The explicit prescription is the sequence of calls above. The differences compare these two recorded coordinates.
\begin{table}[htbp]\centering\small
\caption{Explicit QCD transport and the article's conventional records.
The difference is the explicit result minus the article record.}\label{tab:qcd-repository-records}
\begin{tabular}{lrrr}
\toprule
Coordinate&Explicit calculation&Article record&Difference\\
\midrule
$\overline m_u(2)$ [MeV]&2.11690723327643&2.11478&$+0.00212723327643$\\
$\overline m_d(2)$ [MeV]&4.60738633124870&4.60275&$+0.00463633124870$\\
$\overline m_s(2)$ [MeV]&91.9710443882324&91.8786&$+0.0924443882324$\\
$\overline m_c(\overline m_c)$ [GeV]&1.26312426800494&1.26243&$+0.000694268004941$\\
$\overline m_b(\overline m_b)$ [GeV]&4.18248214200986&4.18248214152943&$+4.80432\times10^{-10}$\\
$M_t$ [GeV]&172.802744094651&172.761212483560&$+0.0415316110907$\\
\bottomrule
\end{tabular}
\end{table}
The current values are printed at \texttt{AR QCD.u} through
\texttt{AR QCD.t}, lines~81--86. Each table coordinate also has its own address:
\texttt{T2.F1.u.carried}, \texttt{T2.F1.d.carried},
\texttt{T2.F2.c.carried}, \texttt{T2.F2.s.carried},
\texttt{T2.F3.t.carried}, and \texttt{T2.F3.b.carried}, at
lines~129, 133, 141, 145, 153 and 157, respectively. Its following two lines
print the article record and the difference. The records reside in
\path{data/conventional_records.json}; their measured comparisons remain
separate.

The four-loop scale integral in the preceding subsection is evaluated
directly with the same coupling. It gives
\begin{equation}
 \int_0^{a_s}I(x)\,\mathrm dx=0.0005282912381702425,
 \qquad \log\frac{\mu^2}{\Lambda_5^2}=12.176922834948924.
\end{equation}
At the calculated $M_Z$, at $91.1876\ \mathrm{GeV}$, and at
$91.188\ \mathrm{GeV}$, respectively,
\begin{equation}
 \Lambda_5=(206.890203783025,\ 206.895282667595,\ 206.896190226441)
 \ \mathrm{MeV}.
\end{equation}
The integral and three reporting coordinates are printed at
\texttt{AR QCD.Lambda.integral}, \texttt{AR QCD.S14},
\texttt{AR QCD.S017}, and \texttt{AR QCD.S018}, lines~87--90.

\subsection{Electric running, anomalies, and bound currents}
An electric probe also polarizes the charged states through which it propagates. Define the renormalized polarization by $\Pi_R(0)=0$, so the long-distance electric response is the reference. Its dispersion relation gives
\begin{equation}
 \alpha(q^2)=\frac{\alpha(0)}{1-\Delta\alpha(q^2)}.
\end{equation}
The vector-current spectrum enters this polarization and the hadronic terms in the lepton anomaly. The vertex map is
\begin{equation}
 a_\ell=F_2^\ell(0)=\sum_{n\geq1}C_n^{(\ell)}(\alpha/\pi)^n
 +a_\ell^{\mathrm{had}}+a_\ell^{\mathrm{weak}},
 \qquad C_1^{(\ell)}=\frac12.
\end{equation}
I10 is the five-loop electron-vertex record. Q04 retains the muon polarization in $10^{10}a_\mu$, and Q01 uses the same vector spectrum in the running electric coupling.

The magnetic-frequency comparison carries the proton shielding correction $\sigma'_p$ of the measurement medium:
\begin{equation}
 R'_\mu=\frac{a_\mu}
 {\mu_p(m_e/m_p)(1-\sigma'_p)(m_\mu/m_e)}.
\end{equation}
The reference is the shielded proton in spherical water at $25^\circ\mathrm C$.
The source's full muon anomaly, the calculated $\mu_p$, $m_e/m_p$ and $m_\mu/m_e$,
and the measured shielding $\sigma'_p$ give the recorded conversion
$R'_\mu=0.003707299080299$. Row Q12 gives the experimental value,
its width and the residual.
The complete QED, hadronic, weak and shielding inputs belong to this conventional
conversion.\cite{MuonG2Final}
 I11 is the hydrogen-binding series at $13.5984345996929\eV$. This carried record includes the atomic binding convention of the calculation.

The full axial source is written $w/(4-3w)$ with its axial Ward partner. A001 evaluates its induced pseudoscalar coupling at $Q^2=0.88m_\mu^2$ and S06 carries it through the singlet muon-capture map. Their output records include the momentum-dependent source, phase space, radiative factors, and bound-state conventions. The finite return is Eq.~\eqref{eq:returns}, the zero-momentum axial residue is Eq.~\eqref{eq:nucleoncurrent}, and $w=e^{r_A^2t_q/24}$ is the momentum-dependent current argument in Eq.~\eqref{eq:nucleonslopes}.

\subsection{Inclusive weak decay and heavy operator transport}
A decaying pseudoscalar transfers its axial current into the lepton pair. The helicity of that pair and the available phase space determine how much of the source can be emitted; the inclusive measurement also counts the specified photons. In the convention $f_\pi\simeq130\MeV$, the decay rate is
\begin{equation}
 \Gamma(P\to\ell\nu[\gamma])
 =\frac{G_F^2|V_{ij}|^2f_P^2m_Pm_\ell^2}{8\pi}
 \left(1-\frac{m_\ell^2}{m_P^2}\right)^2(1+\delta_{\mathrm{incl}}).
\end{equation}
The source supplies $f_P$, masses, $G_F$, and the family overlap; $\delta_{\mathrm{incl}}$ specifies the electromagnetic inclusion convention. The electron-to-muon ratio cancels the common residue and retains the helicity-suppression factor and the ratio of inclusive corrections. S15 uses the archived point-pseudoscalar correction through order $\alpha$.

The neutron lifetime S05 is evaluated from the source current with two classical factors:
\begin{equation}
 \tau_n=\frac{5099.33\,\mathrm s}
 {1.03983\,|V_{ud}|^2(1+3g_A^2)}
 =878.051411826579\,\mathrm s.
 \label{eq:neutronlife}
\end{equation}
The factor $5099.33\,\mathrm s$ contains the measured Fermi constant together with the retained phase space and unit conversion; $1.03983$ is the radiative factor. Thus this lifetime evaluates the finite axial current and family overlap using that measured weak normalization.

\paragraph*{The six complex directions.}
The mixed complex space is $\operatorname{Hom}_{\CC}(F,C)$ and has dimension
six. Its exterior action gives $G_{\rm off}=4\Id_6$ and
$H_{\rm off}=4\Id_6/3$. Its full complex amplitude is
\begin{equation}
 \det_{\CC}(G_{\rm off}^{-1}H_{\rm off})=3^{-6},\qquad
 M_{\rm off}=M_{\rm Pl}3^{-6}.
 \label{eq:offdiagonaldeterminant}
\end{equation}
The realification has determinant $3^{-12}$, the squared magnitude of this
complex amplitude. Thus $M_{\rm off}/M_R=27$ and
$M_{\rm off}^3=M_{\rm Pl}M_R^2$.

The resulting threshold is $M_{\rm off}=1.67474596966450\times10^{16}\GeV$. Its use as the exchange scale in the following proton-decay expression is the additional interpretation made in this attempt.

\begin{equation}
 \frac{\tau_p}{B(e^+\pi^0)}
 =\frac{\hbar}{\alpha_s^4m_p}\left(\frac{M_{\mathrm{off}}}{m_p}\right)^4
 =\frac{72^2\hbar}{3^{24}W^4m_p}
 =1.16971098945626\times10^{37}\,\mathrm{yr}.
 \label{eq:protonlife}
\end{equation}
Masses are rest energies in this expression. K06 and L12 are records for this same completion. Q11 separately carries the proton-decay matrix element through the archive's continuous four-loop operator evolution and computed bottom threshold.

\subsection{The Fermi reference in the neutron lifetime}

Resolving the condensed rate in Eq.~\eqref{eq:neutronlife} gives
\begin{equation}
 \tau_n^{-1}=\frac{G_F^2m_e^5}{2\pi^3\hbar}
 |V_{ud}|^2(1+3g_A^2)f_n(1+\Delta_R).
 \label{eq:neutronrate}
\end{equation}
The normalization $5099.33\,\mathrm s$ uses
$G_{F,\rm ref}=1.1663787\times10^{-5}\GeV^{-2}$. Carrying the calculated
contact into that rate gives
\begin{align}
 \tau_n^{(B,L)}
 &=\frac{5099.33\,\mathrm s}
 {1.03983|V_{ud}|^2(1+3g_A^2)}
 \left(\frac{G_{F,\rm ref}}{G_F}\right)^2\notag\\
 &=878.051172877794\,\mathrm s.
 \label{eq:neutronlifenative}
\end{align}
The squared ratio is $0.999999727864700$; the corresponding condensed
normalization is $5099.32861229230\,\mathrm s$. Holding its phase-space and
radiative factors fixed gives $f_n=1.68870216893486$. The comparison table
retains the recorded normalization as S05, and the displayed conversion
uses the Fermi coefficient of Eq.~\eqref{eq:GF}.

For the muon the leading inclusive rate is
\begin{equation}
 \tau_\mu^{-1}=\frac{G_F^2m_\mu^5}{192\pi^3\hbar}
 f(m_e^2/m_\mu^2)C_\mu,\qquad
 f(x)=1-8x+8x^3-x^4-12x^2\log x.
 \label{eq:muondecay}
\end{equation}
Here $C_\mu$ carries the electromagnetic and finite-vector corrections of
the inclusive measurement convention.

\subsection{Mass conventions for a resonance}
A real rest zero is the zero of the source operator in Eq.~\eqref{eq:restpencil}. A decaying channel has a complex continuation once its open-channel self-energy is supplied. A reported mass specifies the convention used to read that continuation.

For the running-width denominator
\begin{equation}
 D_{\mathrm{BW}}(s)=s-M_{\mathrm{BW}}^2
 +\ii s\,\Gamma_{\mathrm{BW}}/M_{\mathrm{BW}},
\end{equation}
the complex pole is
\begin{equation}
 s_p=\frac{M_{\mathrm{BW}}^2}{1+\ii\Gamma_{\mathrm{BW}}/M_{\mathrm{BW}}}
 =\bar M^2-\ii\bar M\bar\Gamma,\qquad
 \bar M=\frac{M_{\mathrm{BW}}}{\sqrt{1+\Gamma_{\mathrm{BW}}^2/M_{\mathrm{BW}}^2}},
\end{equation}
with the same factor relating $\bar\Gamma$ to $\Gamma_{\mathrm{BW}}$. Another convention writes $\sqrt{s_p}=\mu-\ii\gamma/2$. The width is an input to this conversion. The $W$ and $Z$ register rows identify their experimental mass convention, and the direct top row identifies the reconstructed quantity. Rows D11, D12, D16 and D17 give the adopted Particle Data Group values, their widths and the averaging conventions.\cite{PDG2026} The finite vector quantities in Eq.~\eqref{eq:weakpencil} are real generalized mass zeros. A comparison to the experimental Breit--Wigner parameters carries the open-channel self-energy and the width conversion above. The top direct measurement carries its reconstruction and Monte Carlo mass convention. The selected $\Delta$ example uses the Breit--Wigner comparison in the register.

\subsection{The thermal epoch and acoustic transport}
As the universe expands, the same particle population dilutes and cools. Radiation carries energy away more rapidly than matter, while the constant-curvature density remains fixed. The expansion and continuity equations are Eq.~\eqref{eq:friedmann}.
The radiation convention is
\begin{equation}
 \rho_{\mathrm{rad}}=
 \left[1+\frac78\left(\frac4{11}\right)^{4/3}N_{\mathrm{eff}}\right]\rho_\gamma.
\end{equation}
The source's nuclear-reaction, recombination, and transfer calculations carry the dimensionless occupations of Sec.~\ref{sec:whole} to the selected epoch. Its baryon-mass conversion is $\kappa_B$, and its helium reporting convention is $Y_p=4y/(1+4y)$.

The thermal record gives
\begin{align}
 T_\gamma&=2.72549728670312\,\mathrm K,\notag\\
 H_{\rm comp}&=67.4316192077386\,\mathrm{km\,s^{-1}\,Mpc^{-1}},\qquad
 t_{\rm age}=13.7944222597347\,\mathrm{Gyr},\notag\\
 \Omega_m&=.314264682713255,\qquad
 \Omega_\Lambda=.685668398647101,\qquad
 \Omega_r=6.69186396440\times10^{-5}.
\end{align}
The density fractions obey $\Omega_\Lambda/\Omega_m=24/11$. Substitution in Eq.~\eqref{eq:epoch} gives $H_{\rm comp}=67.4316192077390$ in the stated units. The zero-radiation value at the same curvature is $67.4293629538788$.

Q05--Q10 are the CAMB transfers. They retain the common-area acoustic angle, the optical-depth-one angle, the visibility-peak angle, the drag ruler, and the variance of the matter field. The derived matter statistic is $S_8=\sigma_8(\Omega_m/0.3)^{1/2}$. The numerical programs import the network, recombination, and transfer outputs with these conventions.

\subsection{The two baryonic composition evaluations}
\label{sec:composition-repository-protocol}\label{sec:composition}
The cold state gives
\begin{equation}
 R_c=a^{16}\frac{M_R}{m_p}=5.33628360801120.
\end{equation}
This counts baryons in proton-mass units. To express the result as the
ratio of cold and baryonic energy densities, the baryonic mass per count is
included. For the analytic hydrogen--helium composition,
\begin{align}
 X_n&=1-2b=\frac18,\notag\\
 M_H&=m_p+m_e-B_H,\notag\\
 \kappa_B^{\rm analytic}
 &=(1-2X_n)\frac{M_H}{m_p}+\frac{X_n}{2}\frac{M_{\rm He}}{m_p}.
\label{eq:atomicbaryons}
\end{align}
The numerical inputs are
$m_e=0.51099895069\ \mathrm{MeV}$,
$B_H=13.5984345996929\ \mathrm{eV}$ from atomic record I11,
and $M_{\rm He}=3.72839931079\ \mathrm{GeV}$ from the stated atomic
composition. The proton mass is calculated as
$m_p=938.2720894201967\ \mathrm{MeV}$. These operations give
\begin{equation}
 \kappa_B^{\rm analytic}=0.998763869852295,
 \qquad \frac{R_c}{\kappa_B^{\rm analytic}}=5.34288811308360.
\end{equation}
The network composition used in the article is the recorded input
\begin{equation}
 \kappa_B^{\rm network}=0.998785632760614,
 \qquad \frac{R_c}{\kappa_B^{\rm network}}=5.34277169492503.
\end{equation}
The network composition is an input in \path{data/inputs.json}. The program \path{programs/article_results.py} evaluates the density ratio from that input.
The five printed addresses \texttt{AR COMPOSITION.Xn},
\texttt{AR COMPOSITION.kappa.analytic},
\texttt{AR COMPOSITION.ratio.analytic},
\texttt{AR COMPOSITION.kappa.network}, and
\texttt{AR COMPOSITION.ratio.network} are lines~71--75 of
\path{outputs/article_results.out}. Table~I's cold ratio is also printed at
\texttt{AR T1.32}.

In units $k_B=\hbar=c=1$, the photon, baryon and cold densities are
\begin{equation}
 n_\gamma=\frac{2\zeta(3)}{\pi^2}T_\gamma^3,
 \quad n_b=\eta_Bn_\gamma,
 \quad \rho_b=\kappa_Bm_pn_b,
 \quad \rho_c=R_cm_pn_b.
 \label{eq:thermaldensities}
\end{equation}
The neutral temperature is $T_\nu/T_\gamma=(4\sigma/11)^{1/3}$.
With the two helicities of each occupied species, its exact energy is
\begin{equation}
 \rho_{\nu,\mathrm{th}}(T_\nu)=\frac{\eta_0}{\pi^2}\sum_i
 \int_0^\infty\frac{p^2\sqrt{p^2+m_i^2}}{e^{p/T_\nu}+1}\,dp.
 \label{eq:neutralthermal}
\end{equation}
The massless limit gives
$\rho_{\nu,\mathrm{th}}/\rho_\gamma=(7/8)3\eta_0(T_\nu/T_\gamma)^4$ and hence
$N_{\rm eff}=3\eta_0\sigma^{4/3}$. At small temperature each massive
term approaches $m_i n_i$. The matter energy increases with temperature,
so $\rho_\Lambda/\rho_m=24/11$ fixes one temperature for the stated
composition. The analytic composition gives the carried quadrature
$T_0=2.725500393881\,\mathrm K$; the network composition gives the
register's $2.72549728670312\,\mathrm K$.

The expansion and age then use
\begin{equation}
 H_{\rm comp}^2=\frac{8\pi G}{3}\rho_{\rm total},\qquad
 t_{\rm age}=\int_0^1\frac{dx}{xH_{\rm comp}(x)}.
 \label{eq:thermalage}
\end{equation}
The scale factor is $x$ here. Photon temperature varies as $x^{-1}$;
baryon and cold particle numbers vary as $x^{-3}$; the neutral integral
retains the transition from relativistic to massive energy. The network
composition and its thermal map yield
\begin{center}\small
\begin{tabular}{@{}ll@{}}
\toprule Quantity&Carried value\\\midrule
$H_{\rm comp}$&$67.4316192077386\,\mathrm{km\,s^{-1}\,Mpc^{-1}}$\\
$t_{\rm age}$&$13.7944222597347\,\mathrm{Gyr}$\\
$(\Omega_m,\Omega_\Lambda)$&$(0.314264682713255,0.685668398647101)$\\
$Y_p=4y/(1+4y)$, $y=n_{\rm He}/n_{\rm H}$&$0.24694744653623$\\
$(\omega_b,\omega_c)$&$(0.0224294823835265,0.119835603610525)$\\\bottomrule
\end{tabular}
\end{center}
Here $\omega_i=\Omega_i h_H^2$ and
$h_H=H_{\rm comp}/(100\,\mathrm{km\,s^{-1}\,Mpc^{-1}})$. The acoustic
angle also specifies its surface: the common-area reading is
$100\theta_*=1.04094999030065$, the optical-depth-one reading is
$1.04098061206$, and the visibility-peak reading is $1.041686318178$.
The carried transfer calculation gives $r_d=147.0826558961\,\mathrm{Mpc}$,
$\sigma_8=0.8213353983338$ and
$S_8=\sigma_8\sqrt{\Omega_m/0.3}=0.840635453047808$.
The thermal and acoustic records in Part IV retain these surfaces and their
measurement conventions.

\subsection{The 43 classical evolution records}
The following table identifies the physical evolution used by every classical record. The numeric outputs and comparison labels are printed in Sec.~\ref{sec:register}. Repeated records refer to repeated recorded observables.
\begingroup
\small
\setlength{\tabcolsep}{4pt}
\begin{longtable}{@{}>{\raggedright\arraybackslash}p{.20\textwidth}>{\raggedright\arraybackslash}p{.26\textwidth}>{\raggedright\arraybackslash}p{.48\textwidth}@{}}
\caption{The classical maps and their finite inputs.}\label{tab:classicalmaps}\\
\toprule
Records&Output&Source and map\\\midrule
\endfirsthead
\multicolumn{3}{l}{\textit{Transport table, continued}}\\
\toprule Records&Output&Source and map\\\midrule
\endhead
\bottomrule\endfoot
I10&Electron anomaly&Electric residue and charged rest sources through the five-loop vertex, with the hadronic and weak terms.\\[3pt]
I11&Hydrogen binding&Electric residue, proton/electron ratio, and bound-state series; also the declared atomic conversion anchor.\\[3pt]
Q01&$\alpha^{-1}(M_Z)$&Subtracted dispersion of the vector-current spectrum.\\[3pt]
Q04&Muon hadronic polarization&The vector-current spectrum in the muon anomaly.\\[3pt]
Q12&Muon magnetic-frequency ratio&The full anomaly and finite mass and proton-moment ratios, with the measured proton shielding in spherical water at $25^\circ\mathrm C$.\\[3pt]
I01, I02, I03&$m_{u,d,s}(2\GeV)$&Common $Z$-scale quark masses through four-loop running and threshold matching.\\[3pt]
I04, I05&$m_c(m_c),m_b(m_b)$&Same family masses through running and matching to the self-mass scale.\\[3pt]
S10--S12, T10--T12&Coloured/lepton statistics&Three earlier source mass combinations of the down, strange, and bottom entries, each with its repeated record.\\[3pt]
T08&Running top and Yukawa&Top source transported to $\bar m_t(\bar m_t)$ and $y_t=\sqrt2\bar m_t/v$.\\[3pt]
I07&Matched top pole&Four-loop pole matching of the top entry at the common source scale.\\[3pt]
A001, S06&$g_P$ and singlet capture&Full axial shape $w/(4-3w)$, Ward partner, specified momentum or capture map.\\[3pt]
S15&Pion leptonic ratio&Charged masses through point-pseudoscalar QED through order $\alpha$.\\[3pt]
L21&Neutral mass-sum endpoint&The neutral roots with the clock overlap set to $\hover=1$.\\[3pt]
I06, T09&Low-scale $\alpha_s$&Four-loop running to $m_p$, $m_b$, and $m_c$, with the recorded thresholds.\\[3pt]
S14, S017, S018&$\Lambda_5^{(4)}$&Integral convention of Eq.~\eqref{eq:lambda4} at the three printed scales.\\[3pt]
Q11&Proton-decay matrix element&Four-loop operator transport of the proton and pion source, with the computed bottom threshold.\\[3pt]
F02&Cold-to-baryon ratio&$a^{16}\MR/m_p$ through the network's baryon-mass factor $\kappa_B$.\\[3pt]
F05--F09, F11, F12&Temperature, expansion, age, and densities&Thermal history at $\rho_\Lambda/\rho_m=24/11$, with the same $\kappa_B$.\\[3pt]
F10&Primordial helium&The reaction network with its $Y_p=4y/(1+4y)$ convention.\\[3pt]
Q05, Q06, Q07&Three acoustic angles&CAMB transport with the network's helium: common area, optical depth one, visibility peak.\\[3pt]
Q08, Q09, Q10&$\sigma_8,S_8,r_d$&The same matter and radiation transfer, with $S_8=\sigma_8(\Omega_m/0.3)^{1/2}$.\\
\end{longtable}
\endgroup

\subsection{Inputs and measurement conventions}
The numerical evaluation fixes an encounter's inputs as well as its formula. The electron mass calibrates the dimensionful quantities. Unit conversions, classical evolution, and the reporting convention of an experiment enter at their stated places. The following record identifies them.

\begingroup\small
\setlength{\tabcolsep}{4pt}
\setlength{\LTpre}{3pt}
\begin{longtable}{@{}>{\raggedright\arraybackslash}p{.25\textwidth}>{\raggedright\arraybackslash}p{.69\textwidth}@{}}
\toprule Input or convention&Use in the evaluation\\\midrule
\endfirsthead
\multicolumn{2}{l}{\textit{Input record, continued}}\\
\toprule Input or convention&Use in the evaluation\\\midrule
\endhead
\bottomrule\endfoot
Electron mass&$m_ec^2=0.51099895069(16)\MeV$ is the measured common calibration. Its correlated uncertainty is given in Sec.~\ref{sec:classical}.\\[3pt]
Unit conversions&$\hbar$, $\hbar c$, the electronvolt, the megaparsec, the Julian year and $k_B$ convert the finite ratios into the displayed units. The program uses $1\,\mathrm{Mpc}=3.0856775814913673\times10^{37}\fm$.\\[3pt]
Scalar-spectrum pivot&$k_*=0.05\,\mathrm{Mpc}^{-1}$ specifies where $A_s$ is reported; the source is $12a^{11}(k_{\rm com}\ell)^{-1/40}$.\\[3pt]
Current and mass conventions&The quark family masses use their native source coordinates; running masses, pole masses and reconstructed masses retain their stated maps. The vector comparison uses the stated Breit--Wigner and complex-pole conversion.\\[3pt]
Classical evolution&Four-loop running, threshold matching, charged-current radiative terms, reaction-network, recombination and cosmological-transfer values enter through the 43 records in Sec.~\ref{sec:classicalregister}. The five-flavour one-loop coefficient is $\beta_0=23/3$.\\[3pt]
Neutron lifetime&The $5099.33\,\mathrm s$ normalization contains the measured $G_F$; $1.03983$ is the radiative factor in Eq.~\eqref{eq:neutronlife}.\\[3pt]
Alternative pion current&The coefficient $5/36$ is the model input of A002 and A003.\\[3pt]
Atomic conventions&The broader evaluation uses the Thomas--Fermi coefficient $1.588071$ and hydrogen binding $I_H=13.5984345996929\eV$ in its atomic-mass conversions. The hydrogen entry I11 retains that binding convention.\\[3pt]
Experimental comparisons&Centres, asymmetric widths, confidence intervals, data combinations and upper bounds are those identified beside each numerical row. Derived ratios and repeated rows remain correlated through their shared inputs.\\
\end{longtable}
\endgroup

\newpage
\part{The complete comparison register}\label{part:table}

\section{The seventy-three formulas and their derivations}
\label{sec:formulaindex}
Each formula has a permanent identifier F1--F73. The source calculations use $P=n_F$, $Q=n_L$, $h=h_\nu$, $e=r_0$, $e_\ell=r_{\rm v}$ and $t=t_0$; their program and proof identifiers retain those spellings. Equations here use the notation of Part I. The table records its
expression, physical role, detailed derivation and numerical rows. The
trace normalization, stationary departure count, native quark coordinates,
connected scalar source and mass-coordinate conversions use the exact forms
derived in Parts I--III. A row identifier links directly to its full comparison.
\begingroup\small
\setlength{\tabcolsep}{3pt}\renewcommand{\arraystretch}{1.08}
\begin{longtable}{@{}>{\raggedright\arraybackslash}p{.065\textwidth}>{\raggedright\arraybackslash}p{.9\textwidth}@{}}
\caption{All 73 formulas of the main line.}\label{tab:formulaindex}\\
\toprule ID&Expression, role and detailed location\\\midrule
\endfirsthead\toprule ID&Expression, role and detailed location\\\midrule
\endhead\bottomrule\endfoot
\hypertarget{formula-F1}{}\textbf{F1}\label{formula:F1} & $\gamma_a\gamma_b+\gamma_b\gamma_a=2\delta_{ab}\Id$, on $\CC^4$; $\gamma_1\gamma_2\gamma_3\gamma_4\gamma_5=\pm\Id$.\par
Four complex amplitudes and five real distinctions.\par
Derivation: \eqref{eq:explicitclifford}. \\[7pt]
\hypertarget{formula-F2}{}\textbf{F2}\label{formula:F2} & $L=\Id+\sum_t\gamma_t\otimes\gamma_t=2(P+E_\Omega)$; $\operatorname{spec}L=\{6^{(1)},2^{(10)},(-2)^{(5)}\}$.\par
The singlet, symmetric and vector comparisons.\par
Derivation: \eqref{eq:paircomparison}, \eqref{eq:onekernel}. Rows: \hyperlink{row-P01}{P01}.\\[7pt]
\hypertarget{formula-F3}{}\textbf{F3}\label{formula:F3} & $\langle\Omega,L\Omega\rangle=6$.\par
The paired agreement state.\par
Derivation: \eqref{eq:paircomparison}. Rows: \hyperlink{row-P01}{P01}.\\[7pt]
\hypertarget{formula-F4}{}\textbf{F4}\label{formula:F4} & $\binom d2+\binom{6-d}2=(15,10,7,6)$; $d(6-d)=(0,5,8,9)$ for $d=0,1,2,3$.\par
Preserved and moved rotations. The divided side has dimension $d\leq3$; only $2|4$ retains a central plane rotation.\par
Derivation: \eqref{eq:divisionroot}, Sec.~\ref{sec:step-two}. Rows: \hyperlink{row-P02}{P02}.\\[7pt]
\hypertarget{formula-F5}{}\textbf{F5}\label{formula:F5} & ${\mathsf T}_1=\ii\gamma_h$, ${\mathsf T}_2=\gamma_h\gamma_k$, ${\mathsf T}_3=\ii\gamma_k$; $H_k=(\Id+\ii\gamma_k)/\sqrt2$, $H_k{\mathsf T}_1H_k^\dagger={\mathsf T}_2$; ${\mathsf T}_a^2=-\Id$, ${\mathsf T}_1{\mathsf T}_2={\mathsf T}_3$.\par
The two clock planes are related by one quaternion rotation.\par
Derivation: \eqref{eq:commonphase}. Rows: \hyperlink{row-P03}{P03}.\\[7pt]
\hypertarget{formula-F6}{}\textbf{F6}\label{formula:F6} & $Y=\widehat n_F/2-\widehat n_C/3$, $Q=\widehat n_k-\widehat n_C/3=T_3+Y$; $\Tr Y=\Tr Y^3=0$.\par
The sixteen exterior states form one chiral family. The five direction charges are $(0,1,-1/3,-1/3,-1/3)$; the centre fixing $h$ generates $Q$.\par
Derivation: \eqref{eq:tenclifford}, \eqref{eq:paircomplexstructure}. Rows: \hyperlink{row-P04}{P04}.\\[7pt]
\hypertarget{formula-F7}{}\textbf{F7}\label{formula:F7} & $\operatorname{spec}\mathcal M=(1/2,1/2,4/5,0^{(9)})$ in $\Re\Tr X^\dagger Y$-orthonormal generators. In $\tr(t^at^b)=\delta^{ab}/2$ normalization the nonzero values are $(1/4,1/4,2/5)$.\par
Three of the twelve traceless rotations move the scalar. The source ratio is $M_W^2/M_Z^2=5/8$, with $\sin^2\theta=3/8$.\par
Derivation: \eqref{eq:vectorgram}, \eqref{eq:realvector}. Rows: \hyperlink{row-P05}{P05}.\\[7pt]
\hypertarget{formula-F8}{}\textbf{F8}\label{formula:F8} & $\mathcal N_{2,4}=(S^1\times S^3)/\mathbb Z_2$; $\operatorname{sign}T\mathcal N_{2,4}=(1,3)$.\par
The balanced directions of the held plane give one time and three spatial coordinates.\par
Derivation: \eqref{eq:nullquotient}, \eqref{eq:localcone}. Rows: \hyperlink{row-P06}{P06}.\\[7pt]
\hypertarget{formula-F9}{}\textbf{F9}\label{formula:F9} & $B_{(y,s),(x,t)}=[s\ne t][y=x\mathbin\triangle g_s]$, with $g_s=\{1,\ldots,5\}\setminus\{s\}$.\par
The remembered step has 80 states and 320 directed transitions.\par
Derivation: \eqref{eq:B}. Rows: \hyperlink{row-P07}{P07}.\\[7pt]
\hypertarget{formula-F10}{}\textbf{F10}\label{formula:F10} & $\det(\Id-uB)=(1-u^2)^{24}u^{16}\det[\Delta+\chi(u)\Id]$, $\Delta=6\Id-L$, $\chi=u^{-1}+4u-5$.\par
The history law is the comparison with one movable scalar displacement. The unitary character change $T=\mathscr H_0\mathscr F^{\mathsf T}$ gives $T^\dagger\Delta T=5\Id-A$.\par
Derivation: \eqref{eq:bass}, \eqref{eq:onekernel}. Rows: \hyperlink{row-P08}{P08}.\\[7pt]
\hypertarget{formula-F11}{}\textbf{F11}\label{formula:F11} & $k_N(u)=u(2N+\chi)=1+4u^2-u(5-2N)$; $D(u)=(1-u^2)^{24}k_0k_2^{10}k_4^5$.\par
The three comparison grades have disagreements $0,4,8$.\par
Derivation: \eqref{eq:onekernel}. Rows: \hyperlink{row-P08}{P08}.\\[7pt]
\hypertarget{formula-F12}{}\textbf{F12}\label{formula:F12} & $u=\ii^d/4$, $d=0,1,2,3$, gives $\chi=(0,-5-3\ii,-10,-5+3\ii)$.\par
The four divided sides give the four quarter values of one variable.\par
Derivation: \eqref{eq:fourquarters}. Rows: \hyperlink{row-P09}{P09}.\\[7pt]
\hypertarget{formula-F13}{}\textbf{F13}\label{formula:F13} & $k_N(-1/4)=(5-N)/2$, so $(k_0,k_2,k_4)=(5/2,3/2,1/2)$.\par $t_0=1/k_0=2/5$, $r_0=k_2/k_0=3/5$, $t_{\rm v}=1/k_2=2/3$, $r_{\rm v}=k_4/k_2=1/3$.\par
The plane and one-vector transmission and reflection amplitudes.\par
Derivation: \eqref{eq:scattering}, \eqref{eq:support}. Rows: \hyperlink{row-P10}{P10}.\\[7pt]
\hypertarget{formula-F14}{}\textbf{F14}\label{formula:F14} & $u=1/4-\epsilon^2m^2/24$ gives $\chi=\epsilon^2m^2/2+O(\epsilon^4)$ and $(1+4u^2)\Id-uA_\epsilon=(\epsilon^2/8)(m^2-\Delta_{\rm tan})+O(\epsilon^3)$.\par $\epsilon^2M_\epsilon^2=2\chi+\chi^2/5$, $M_\epsilon\to m$.\par
Here $\Delta_{\rm tan}=\sum_{\mu=1}^4\partial_\mu^2$ is the tangent Laplacian. The scalar long-wave equation is Klein--Gordon. Along $v_h$, the exact phase mass is the arsinh conversion of the inverse range.\par
Derivation: \eqref{eq:exactgap}, \eqref{eq:exactradialbranch}, \eqref{eq:finite_mass_conversion}. Rows: \hyperlink{row-P11}{P11}.\\[7pt]
\hypertarget{formula-F15}{}\textbf{F15}\label{formula:F15} & $U=e^{-\ii k\gamma_j}e^{-\ii\mu\gamma_h}$, $U+U^\dagger=2\cos k\cos\mu\,\Id$; $\cos\omega=\cos\mu\cos k$.\par
The first-order step gives the Dirac limit, with $c=\Delta x/\Delta t$ and energy $\hbar\omega/\Delta t$.\par
Derivation: \eqref{eq:finiteclockunitary}, \eqref{eq:finiteclockdispersion}. Rows: \hyperlink{row-P12}{P12}.\\[7pt]
\hypertarget{formula-F16}{}\textbf{F16}\label{formula:F16} & $\mathcal G_{\rm sp}(k)=[V_3(M_{\rm inv}^2+L_3(k)/5)]^{-1}$, $V_3=1/2$; $\mathcal G_{M_{\rm inv}}(r)=e^{-M_{\rm inv}r}/(4\pi r)$.\par
Two static sources have the inverse-distance response at the zero and Yukawa decay below it. The finite directional checks give coefficients from $0.99304$ to $0.99989$ relative to the continuum normalization.\par
Derivation: \eqref{eq:continuumgreen}. \\[7pt]
\hypertarget{formula-F17}{}\textbf{F17}\label{formula:F17} & $\log\det(\Id-wA_-[e^{\ii\theta}])=\log\det(\Id-wA_-)-\kappa\|C_\square\theta\|^2/2+O(\theta^4)$.\par $w=(\chi+5)^{-1}$, $\kappa=2/[(\chi+5)^2-5]^2$.\par
The signed return gives the curvature square: $\kappa=1/200$ at $u=\pm1/4$ and $1/98$ at $u=\pm\ii/4$.\par
Derivation: \eqref{eq:proof-SCa-5}, \eqref{eq:proof-SCa-8}. Rows: \hyperlink{row-P13}{P13}.\\[7pt]
\hypertarget{formula-F18}{}\textbf{F18}\label{formula:F18} & $\log(4u_c)=\|\theta_{\rm cyc}\|^2/48+O(\theta^4)$; $48=3\EE N_{\rm ret}=2(40-16)$.\par
A circulation of the link field moves the long-range zero.\par
Derivation: \eqref{eq:circulationperron}, \eqref{eq:circulation48}, \eqref{eq:circulationbound}. Rows: \hyperlink{row-P14}{P14}.\\[7pt]
\hypertarget{formula-F19}{}\textbf{F19}\label{formula:F19} & $\tr_4[\gamma_h(\mathcal F_{\rm sp}\wedge\mathcal F_{\rm sp}-R_{\rm sp}\wedge R_{\rm sp})]=-\ell^{-2}(R-6/\ell^2)\operatorname{vol}$.\par
The epsilon contraction is four times the ordinary four-spinor trace. Their consistently normalized Einstein--Cartan actions give $\Lambda\ell^2=3$.\par
Derivation: \eqref{eq:proof-SE-3}, \eqref{eq:proof-SE-1}, \eqref{eq:proof-SE-4}. Rows: \hyperlink{row-P15}{P15}.\\[7pt]
\hypertarget{formula-F20}{}\textbf{F20}\label{formula:F20} & $K_X(E,{\mathcal N})=K_X(E/{\mathcal N},1)$, $K_{\mathcal N}=-(E/{\mathcal N})K_E$, $\dd E_X/\dd {\mathcal N}=E_X({\mathcal N})/{\mathcal N}$ for ${\mathcal N}>0$.\par
A lapse couples to the complete rest energy, including internal propagation and binding.\par
Derivation: \eqref{eq:lapse}, \eqref{eq:fullclock}. Rows: \hyperlink{row-P16}{P16}.\\[7pt]
\hypertarget{formula-F21}{}\textbf{F21}\label{formula:F21} & $\Fret(z)=z^4(1+z)/(16-8z-3z^2-3z^3)$; $\Fret(1)=1$, $\EE N_{\rm ret}=16$, $\Var N_{\rm ret}=156$.\par
The first-return distribution closes with probability one.\par
Derivation: \eqref{eq:returns}, \eqref{eq:jointreturn}. Rows: \hyperlink{row-P17}{P17}.\\[7pt]
\hypertarget{formula-F22}{}\textbf{F22}\label{formula:F22} & $\Phis(z)=-z^4/(16+3z^2)$; $\Phis(\pm1)=-1/19$, $\Phis(\pm\ii)=-1/13$.\par
The same histories carry their Clifford loop sign.\par
Derivation: \eqref{eq:returns}, \eqref{eq:signedjointreturn}. Rows: \hyperlink{row-P18}{P18}.\\[7pt]
\hypertarget{formula-F23}{}\textbf{F23}\label{formula:F23} & $\EE n_F=10$, $\EE n_L=5$, $\Var n_F=84$, $\Var n_L=12$, $\Cov(n_F,n_L)=30$.\par $\EE x^{n_F}=x(1+x)/(16-11x-3x^2)$; $x_c=(\sqrt{313}-11)/6$.\par
The symmetric and vector arrivals in one stopped return, including their joint fluctuations.\par
Derivation: \eqref{eq:visitmoments}, \eqref{eq:conditional}, \eqref{eq:conditionalfactor}. Rows: \hyperlink{row-P19}{P19}.\\[7pt]
\hypertarget{formula-F24}{}\textbf{F24}\label{formula:F24} & 74 retained remembered states; $M=296$, $C_h=60$, $C_j=59$; $c_L=15/74$, $\sigma=75/74$, $\eta_0=295/296$, $R_{ud}=34/74=17/37$, $E_H=3$.\par
There are 34 retained remembered states with $h\notin x$. Stationary departures and excursion occupation retain their respective normalizations.\par
Derivation: \eqref{eq:retainedschur}, \eqref{eq:tracechain}, \eqref{eq:stationaryweights}, \eqref{eq:EH}. Rows: \hyperlink{row-P20}{P20}.\\[7pt]
\hypertarget{formula-F25}{}\textbf{F25}\label{formula:F25} & ${\mathcal G_0}=(5\Id-A)^+=(7\Id+A)/32-3J_{16}/128$; $b={\mathcal G_0}_{00}+{\mathcal G_0}_{FF}-2{\mathcal G_0}_{0F}=7/16$.\par
The response to equal and opposite sources at the singlet and a symmetric state.\par
Derivation: \eqref{eq:green}, \eqref{eq:groundedpetersen}. Rows: \hyperlink{row-P21}{P21}.\\[7pt]
\hypertarget{formula-F26}{}\textbf{F26}\label{formula:F26} & $a=[1-\Phis(1)]/8=5/38$, $\chi_A=[1+\Phis(\ii)]^{-1}=13/12$, $w_h=1/(4\chi_A)=3/13$, $\chi_h=w_h(1-w_h)=30/169$.\par
The even and odd signed returns give reversal, moved fraction and radial susceptibility.\par
Derivation: \eqref{eq:radialsusceptibility}, \eqref{eq:returns}. Rows: \hyperlink{row-P22}{P22}.\\[7pt]
\hypertarget{formula-F27}{}\textbf{F27}\label{formula:F27} & $J_{DE}/C_W=(|D||E|-5\tr P_DP_E)/72$.\par $1/72=-\Phis(1)/[4(1+\Phis(1))]=-\Phis(\ii)/[6(1+\Phis(\ii))]=r_{\rm v}^2/8$.\par
The crossing response divided by its complete clock is independent of the central grade weights. The three intersection values are $1/72$, $1/18$, $1/12$.\par
Derivation: \eqref{eq:crossing}, \eqref{eq:crossingchi}, \eqref{eq:alphas}. Rows: \hyperlink{row-P23}{P23}.\\[7pt]
\hypertarget{formula-F28}{}\textbf{F28}\label{formula:F28} & $\langle\Omega,(L+\gamma_h\otimes\gamma_s)\Omega\rangle=6+[s=h]$.\par
A retained departure compares its moment with itself, with one extra unit along its own time direction.\par
Derivation: \eqref{eq:protonsourcecount}. Rows: \hyperlink{row-P24}{P24}.\\[7pt]
\hypertarget{formula-F29}{}\textbf{F29}\label{formula:F29} & $D_\vartheta(u)=D(u)+(e^{-\ii\vartheta}-1)uD'(u)/320$; $\arg(4u_c)=\vartheta/320+O(\vartheta^3)$.\par
All 320 directed transitions have the same cofactor; the elementary source rate is $1/320$.\par
Derivation: \eqref{eq:marked}, \eqref{eq:markedsourceclock}. Rows: \hyperlink{row-P25}{P25}.\\[7pt]
\hypertarget{formula-F30}{}\textbf{F30}\label{formula:F30} & $f=r_0^2/10=9/250$, $\nu=f/\EE N_{\rm ret}=9/4000$.\par
Two endpoint reflections meet on one of ten symmetric states, then are counted per mean return.\par
Derivation: \eqref{eq:fnu}. Rows: \hyperlink{row-P26}{P26}.\\[7pt]
\hypertarget{formula-F31}{}\textbf{F31}\label{formula:F31} & $(q_0,q_1,q_2,q_3,q_h)=(0,1,2,3,19)/4$; $M_Q(s)=\frac1{10}\sum_De^{-sD}$, $\EE D=2$.\par $h_\nu=[M_Q(\nu)M_Q(-\nu)]^{-3/2}=0.99997104933$.\par
The ten pair intervals of the five clock marks carry the overlap of the two opposite tilts.\par
Derivation: \eqref{eq:clockmarks}, \eqref{eq:MQ}, \eqref{eq:h}. Rows: \hyperlink{row-P27}{P27}.\\[7pt]
\hypertarget{formula-F32}{}\textbf{F32}\label{formula:F32} & $N_h=\Re\sum_{j\geq1}h_\nu^jp_j$, where $p_j=-\Phis(4(\ii/4)^j)$ for odd $j$ and $p_j=\Fret(4(\ii/4)^j)$ for even $j$; $N_h=0.0770858406$.\par
Repeated traversals of the quarter return retain the sign at odd repetitions.\par
Derivation: \eqref{eq:Nh}. Rows: \hyperlink{row-P28}{P28}.\\[7pt]
\hypertarget{formula-F33}{}\textbf{F33}\label{formula:F33} & $x=\nu h_\nu/(M+1+N_h)=7.573573\times10^{-6}$.\par
The Schur complement closes the two-end electric contact.\par
Derivation: \eqref{eq:electricblock}. Rows: \hyperlink{row-P29}{P29}.\\[7pt]
\hypertarget{formula-F34}{}\textbf{F34}\label{formula:F34} & $W_W=\Fret(h_\nu)h_\nu^{-4}(1-2b+2bh_\nu)=0.9996273896$.\par
The weak contact retains its excursion beyond the four compulsory steps and its joint endpoint occupation.\par
Derivation: \eqref{eq:GF}. Rows: \hyperlink{row-P30}{P30}.\\[7pt]
\hypertarget{formula-F35}{}\textbf{F35}\label{formula:F35} & $\alpha_s=[r_{\rm v}^2/(3^2-1)]^{1/2}=1/\sqrt{72}=0.117851130$.\par
The one-vector reflection is shared over the eight traceless colour rotations.\par
Derivation: \eqref{eq:alphas}. Rows: \hyperlink{row-D01}{D01}.\\[7pt]
\hypertarget{formula-F36}{}\textbf{F36}\label{formula:F36} & $\alpha=c_Lf/(1-x)$; $\alpha^{-1}=137.035999177$. At $h_\nu=1$, $\alpha^{-1}=137.035999147$.\par
The electric response contracts the normalized endpoint current with the complete contact inverse.\par
Derivation: \eqref{eq:alpha}. Rows: \hyperlink{row-D02}{D02}, \hyperlink{row-L19}{L19}.\\[7pt]
\hypertarget{formula-F37}{}\textbf{F37}\label{formula:F37} & $G_Fm_p^2=(1-1/10)W_W/M^2$; $G_F=1.16637886\times10^{-5}\,\mathrm{GeV}^{-2}$.\par
Two stationary endpoints each contribute $1/M$; removing their coincident symmetric channel leaves nine of ten.\par
Derivation: \eqref{eq:GF}. Rows: \hyperlink{row-D03}{D03}.\\[7pt]
\hypertarget{formula-F38}{}\textbf{F38}\label{formula:F38} & $W=t_0^{3\EE N_{\rm ret}}\eta_0^{3E_H}=(2/5)^{48}(295/296)^9$.\par $\alpha_G=Gm_p^2/(\hbar c)=W^2=5.906152\times10^{-39}$; $M_{\rm Pl}=m_p/W=1.22088981\times10^{19}\GeV$.\par
Each of the three complementary directions carries sixteen steps and three retained departures.\par
Derivation: \eqref{eq:W}, \eqref{eq:gravitystrength}. Rows: \hyperlink{row-P31}{P31}, \hyperlink{row-D04}{D04}, \hyperlink{row-D05}{D05}.\\[7pt]
\hypertarget{formula-F39}{}\textbf{F39}\label{formula:F39} & $\left.\dd\arg(4u_c)/\dd\vartheta\right|_0=1/320$; $m_e c^2=0.51099895069\MeV$.\par
The one-transition rest source sets the electron mass unit.\par
Derivation: \eqref{eq:marked}, \eqref{eq:markedsourceclock}. Rows: \hyperlink{row-P25}{P25}.\\[7pt]
\hypertarget{formula-F40}{}\textbf{F40}\label{formula:F40} & $m_p/m_e=1836+\frac16\log(5/2)-a^5-ba^6=1836.152673412$; $m_pc^2=938.27208942\MeV$.\par
The integer is $6M+C_h$. The remaining terms are the pairing determinant and the signed five- and seven-state exterior responses.\par
Derivation: \eqref{eq:completeprotonsource}, \eqref{eq:protonsourcecount}, \eqref{eq:proton}. Rows: \hyperlink{row-J01}{J01}, \hyperlink{row-B12}{B12}, \hyperlink{row-M1}{M1}.\\[7pt]
\hypertarget{formula-F41}{}\textbf{F41}\label{formula:F41} & $m_n/m_p=(5/2)^{5a^4+4a^{10}+10c_Lba^6}=1.00137841956$; $m_nc^2=939.565422025\MeV$.\par
The neutral source carries nine Petersen contrasts and the opened seven-state contact.\par
Derivation: \eqref{eq:neutronsourcefunctional}, \eqref{eq:neutron}, \eqref{eq:nucleonmassratios}. Rows: \hyperlink{row-B13}{B13}, \hyperlink{row-J02}{J02}, \hyperlink{row-B14}{B14}, \hyperlink{row-M2}{M2}.\\[7pt]
\hypertarget{formula-F42}{}\textbf{F42}\label{formula:F42} & $\log\mathcal L(n,L_*,p)=\frac n2\log(1/r_{\rm v})-\frac{\Psi_pL_*}{4}(n-\epsilon_p)+c_n(p)$.\par $\Psi_+=1/28$, $\Psi_-=1/2$, $\epsilon_p=(1+p)/2$; $c_{10}(p)=\frac16\log M_Q(p\nu)$, $c_5(p)=(5-\epsilon_p)p\nu\EE D/4$.\par
A retained grade contributes its amplitude volume, signed endpoint association and clock term.\par
Derivation: \eqref{eq:gradesourcefunctional}, \eqref{eq:gradecontractions}, \eqref{eq:layers}, \eqref{eq:threecumulativegrades}. \\[7pt]
\hypertarget{formula-F43}{}\textbf{F43}\label{formula:F43} & $m_\mu/m_e=\mathcal L(10,2,+)=3^5e^{-9/56}M_Q(\nu)^{1/6}=206.768283945$.\par $m_\tau/m_\mu=\mathcal L(5,-2,+)=3^{5/2}e^{1/14+2\nu}=16.8181606411$.\par
The three pairings retain the reference, then the ten symmetric states, then the five vector states. The isolated scale $3^{5/2}m_e$ is absent from this spectrum. The Koide ratio is $0.66666726$.\par
Derivation: \eqref{eq:leptons}. Rows: \hyperlink{row-B01}{B01}, \hyperlink{row-B02}{B02}, \hyperlink{row-L01}{L01}, \hyperlink{row-B03}{B03}, \hyperlink{row-L17}{L17}, \hyperlink{row-L04}{L04}, \hyperlink{row-L05}{L05}.\\[7pt]
\hypertarget{formula-F44}{}\textbf{F44}\label{formula:F44} & $p_q=\mathcal L(10,2,-)=19.961652$, $q_q=\mathcal L(5,-2,-)$; $(m_d^{\rm native},m_s^{\rm native},m_b^{\rm native})=3^{3/2}m_e(\eta_0^{-1},p_q\eta_0^{-1},p_q q_q)$; $m_b^{\rm native}/m_s^{\rm native}=q_q\eta_0=53.921087$.\par
The odd lifted source gives native down-family masses; an open directional pole carries the reciprocal clock weight.\par
Derivation: \eqref{eq:quarks}, \eqref{eq:ladderinvariant}. Rows: \hyperlink{row-B04}{B04}, \hyperlink{row-B05}{B05}, \hyperlink{row-B06}{B06}, \hyperlink{row-B07}{B07}, \hyperlink{row-B08}{B08}.\\[7pt]
\hypertarget{formula-F45}{}\textbf{F45}\label{formula:F45} & $m_u^{\rm native}/m_d^{\rm native}=R_{ud}=17/37$, $\kappa_u=(r_0/t_0)^{\EE N_{\rm ret}/\EE n_F}=(3/2)^{8/5}=1.913137$.\par $m_c^{\rm native}/m_s^{\rm native}=R_{ud}\kappa_u^5$, $m_t^{\rm native}/m_b^{\rm native}=R_{ud}\kappa_u^{15/2}$.\par
The native up-family source preserves $h$ in 34 of 74 states; half the cumulative retained dimensions gives powers $0,5,15/2$.\par
Derivation: \eqref{eq:quarks}, \eqref{eq:ladderinvariant}. Rows: \hyperlink{row-B09}{B09}, \hyperlink{row-B10}{B10}, \hyperlink{row-B11}{B11}.\\[7pt]
\hypertarget{formula-F46}{}\textbf{F46}\label{formula:F46} & $v=(\sqrt2G_F)^{-1/2}=246.219634\GeV$.\par
The scalar length follows from the same finite Fermi contact.\par
Derivation: \eqref{eq:radial}, \eqref{eq:GF}. Rows: \hyperlink{row-D15}{D15}.\\[7pt]
\hypertarget{formula-F47}{}\textbf{F47}\label{formula:F47} & $\widehat s^2=3N_h=0.2312575$, $s_{\rm OS}^2=3N_h/(1+f)=0.2232215$, $s_{\rm eff}^2=3N_{\rm con}=0.2315502$. At $h_\nu=1$, $3N_h=0.2312642$.\par
The three complementary directions carry the repeated return; the pole adds the symmetric contact, and the current symmetrizes the even returns. The endpoint correction is $\Delta r=f$.\par
Derivation: \eqref{eq:Nh}, \eqref{eq:Ncon}, \eqref{eq:realvector}. Rows: \hyperlink{row-D08}{D08}, \hyperlink{row-D10}{D10}, \hyperlink{row-L20}{L20}, \hyperlink{row-D09}{D09}, \hyperlink{row-D14}{D14}.\\[7pt]
\hypertarget{formula-F48}{}\textbf{F48}\label{formula:F48} & $M_W^2=\pi\alpha v^2/[(1-f)s_{\rm OS}^2]$, $M_Z=M_W/\sqrt{1-s_{\rm OS}^2}$.\par $M_W=80.36625\GeV$, $M_Z=91.18536\GeV$.\par
The three moved directions of the scalar acquire their frequencies against the kinetic metric; the combination preserving $h_0$ is massless.\par
Derivation: \eqref{eq:weakpencil}, \eqref{eq:weakmasses}. Rows: \hyperlink{row-D11}{D11}, \hyperlink{row-L15}{L15}, \hyperlink{row-D12}{D12}, \hyperlink{row-D13}{D13}.\\[7pt]
\hypertarget{formula-F49}{}\textbf{F49}\label{formula:F49} & $\lambda_H=\pi\alpha/\chi_h=0.129146$, $m_H=v\sqrt{2\lambda_H}=125.1348\GeV$.\par
The inverse susceptibility of the moved fraction sets the radial scalar curvature.\par
Derivation: \eqref{eq:radialsusceptibility}, \eqref{eq:radial}. Rows: \hyperlink{row-D16}{D16}.\\[7pt]
\hypertarget{formula-F50}{}\textbf{F50}\label{formula:F50} & $y_t^2=\eta_0/\sigma=59/60$, $m_t^{\rm rad}=v\sqrt{\eta_0/(2\sigma)}=172.646613614442\GeV$.\par $m_t^{\rm native}=R_{ud}\kappa_u^{15/2}m_b^{\rm native}=170.949965605963\GeV$.\par
The radial current and native family mass retain their separate source meanings. Part III gives the stated QCD conversions; S07 retains the second mass $m_p^2/(10m_e)$ chosen with its measurement in view.\par
Derivation: \eqref{eq:radialfermionsource}, \eqref{eq:radial}, \eqref{eq:quarks}, Sec.~\ref{sec:secondconstructions}. Rows: \hyperlink{row-D17}{D17}, \hyperlink{row-S07}{S07}.\\[7pt]
\hypertarget{formula-F51}{}\textbf{F51}\label{formula:F51} & $s_{12}=V_C/[(\EE N_{\rm ret})(5/2)]=9/40$, $s_{23}=4/\Tr L^2=1/24$, $s_{13}=t_0s_{12}s_{23}=3/800$, $\delta_q=\arg(b+\ii)$.\par $V=R_{23}(s_{23})R_{13}(s_{13},\delta_q)R_{12}(s_{12})$.\par
The family overlaps are the identity, exchange and five-cycle coefficients of the Petersen character projector. Its three-cycle coefficient is $-1/40$.\par
Derivation: \eqref{eq:V}. Rows: \hyperlink{row-C01}{C01}, \hyperlink{row-C02}{C02}, \hyperlink{row-C03}{C03}, \hyperlink{row-C04}{C04}, \hyperlink{row-C05}{C05}, \hyperlink{row-C06}{C06}, \hyperlink{row-C07}{C07}, \hyperlink{row-C08}{C08}, \hyperlink{row-C09}{C09}, \hyperlink{row-K05}{K05}.\\[7pt]
\hypertarget{formula-F52}{}\textbf{F52}\label{formula:F52} & $U=R_{12}(\sqrt{2b}\,s_{12},\pi/2-\delta_q)^\dagger R_{23}(1/\sqrt2)R_{12}(\sqrt b)$.\par
The lepton frame compares the same endpoint from the opposite orientation. The squared angles are $0.30509,0.022148,0.48867$, with phase $205.6^\circ$.\par
Derivation: \eqref{eq:U}, \eqref{eq:taurow}. Rows: \hyperlink{row-C10}{C10}, \hyperlink{row-C11}{C11}, \hyperlink{row-C12}{C12}, \hyperlink{row-C13}{C13}, \hyperlink{row-C14}{C14}, \hyperlink{row-T06}{T06}.\\[7pt]
\hypertarget{formula-F53}{}\textbf{F53}\label{formula:F53} & $M_R=M_{\rm Pl}r_{\rm v}^9$, $m_{Di}=\sigma v(0,\sqrt{\rho_\nu},1)_i/\sqrt2$, $\rho_\nu=\tan^2(\pi/8)$; $m_i=2m_{Di}^2/[M_R+\sqrt{M_R^2+4m_{Di}^2}]$.\par $M_{\rm off}=M_{\rm Pl}r_{\rm v}^6$ in the complex-six construction of Part III.\par
The neutral determinant has nine real directions. The light masses are $(0,8.613,50.198)$ meV. The mixed threshold uses six complex directions, whose realification has determinant $3^{-12}$.\par
Derivation: \eqref{eq:neutralpullback}, \eqref{eq:neutraldeterminant}, \eqref{eq:offdiagonaldeterminant}, \eqref{eq:neutralthresholds}, \eqref{eq:neutralroots}, \eqref{eq:neutrinovalues}, \eqref{eq:effectivemasses}. Rows: \hyperlink{row-E07}{E07}, \hyperlink{row-E01}{E01}, \hyperlink{row-E02}{E02}, \hyperlink{row-E03}{E03}, \hyperlink{row-E04}{E04}, \hyperlink{row-E05}{E05}, \hyperlink{row-E06}{E06}, \hyperlink{row-K02}{K02}, \hyperlink{row-K03}{K03}, \hyperlink{row-L03}{L03}, \hyperlink{row-T01}{T01}, \hyperlink{row-T02}{T02}.\\[7pt]
\hypertarget{formula-F54}{}\textbf{F54}\label{formula:F54} & $P_N=P_{\rm sym}(15/4-\mathbf S^2)/3$; $\langle p\uparrow|\sum_r\sigma_{z,r}\tau_{z,r}|p\uparrow\rangle=5/3$; $g_A=(5/3)t_0^{2b/3}=1.275801$.\par
The rank-four nucleon projector gives $5/3$, and the conditional return variance gives $r_A^2=9\lambda_p^2$. Part III names the empirical axial exponent, records the measured origin of the radius assignment, and gives the measured weak normalization in the $878.05$ s lifetime.\par
Derivation: \eqref{eq:nucleonprojectors}, \eqref{eq:nucleonspinmatrix}, \eqref{eq:nucleoncurrent}, \eqref{eq:neutronlife}, \eqref{eq:neutronlifenative}. Rows: \hyperlink{row-P36}{P36}, \hyperlink{row-S01}{S01}, \hyperlink{row-S05}{S05}, \hyperlink{row-S02}{S02}, \hyperlink{row-S03}{S03}.\\[7pt]
\hypertarget{formula-F55}{}\textbf{F55}\label{formula:F55} & $\epsilon_2=\alpha^2/(1+\alpha/10)$, $S_\mu=1-\epsilon_2/[(\EE N_{\rm ret})^2b]$.\par $\mu_p=\sqrt{(\Var N_{\rm ret})S_\mu/20}=2.7928473453$, $\mu_n=-(44/23)\sqrt{S_\mu}=-1.9130430238$.\par
Part III names the measurement-guided factors $39/5$, $44/23$ and $S_\mu$, and the deuteron correction $41/40$.\par
Derivation: \eqref{eq:nucleoncurrent}, \eqref{eq:nucleonvertex}. Rows: \hyperlink{row-G03}{G03}, \hyperlink{row-G04}{G04}, \hyperlink{row-G05}{G05}, \hyperlink{row-G13}{G13}.\\[7pt]
\hypertarget{formula-F56}{}\textbf{F56}\label{formula:F56} & $\lambda_p=\hbar/(m_pc)=0.2103089101\fm$; $r_p=\sqrt{\EE N_{\rm ret}}\lambda_p=4\lambda_p=0.841236\fm$, $r_\pi=\sqrt{\EE n_F}\lambda_p=0.66506\fm$.\par
The returned form-factor slopes turn visit counts into squared radii.\par
Derivation: \eqref{eq:nucleoncurrent}, \eqref{eq:nucleonslopes}. Rows: \hyperlink{row-G01}{G01}, \hyperlink{row-S04}{S04}.\\[7pt]
\hypertarget{formula-F57}{}\textbf{F57}\label{formula:F57} & $f_\pi=\eta_0m_p\sqrt6/(2\pi\mu_p)=130.529\MeV$; $f_\pi^{(2)}=5m_p/36=130.316\MeV$.\par
The two current constructions retain their separate empirical factors in Part III.\par
Derivation: \eqref{eq:fpi}, Table~\ref{tab:measuredchoices}, Sec.~\ref{sec:secondconstructions}. Rows: \hyperlink{row-G08}{G08}, \hyperlink{row-G09}{G09}, \hyperlink{row-A002}{A002}, \hyperlink{row-A003}{A003}.\\[7pt]
\hypertarget{formula-F58}{}\textbf{F58}\label{formula:F58} & $(1+4u^2)\Id-uA_U\succeq(1-u)(1-4u)\Id>0$, $0\leq u<1/4$.\par $\mathcal Z_{\rm top}(\theta)=\sum_qZ_qe^{\ii q\theta}$, $Z_q=Z_{-q}\geq0$, $Z_1>0$; $\bar\theta=0$.\par
Positivity and reversal make the stationary free energy minimal at zero sector phase.\par
Derivation: \eqref{eq:positiveconnection}, \eqref{eq:thetasectors}. Rows: \hyperlink{row-P37}{P37}, \hyperlink{row-K01}{K01}, \hyperlink{row-L16}{L16}.\\[7pt]
\hypertarget{formula-F59}{}\textbf{F59}\label{formula:F59} & $\dim SU(3)=8$; $\Lambda_{\rm QCD}^{(1)}=M_Z\exp[-2\pi/((23/3)\alpha_s)]=87.1\MeV$.\par
The colour algebra is finite; the one-loop scale uses the conventional evolution of Part III.\par
Derivation: \eqref{eq:alphas}, \eqref{eq:lambda4}. Rows: \hyperlink{row-L14}{L14}, \hyperlink{row-S16}{S16}.\\[7pt]
\hypertarget{formula-F60}{}\textbf{F60}\label{formula:F60} & $\chi_{\rm top}/(F_\pi^2m_{\pi^0}^2)=[\lambda_{\rm low}\sum_jS_j^{-1}]^{-1}=0.2129$; $b_2=-\sum_jS_j^{-3}/[12(\sum_jS_j^{-1})^3]=-0.02806$, $S=(R_{ud},1,p_q)$.\par
The additional pseudoscalar source gives a vacuum curvature and fourth cumulant in Part III.\par
Derivation: \eqref{eq:topology}. Rows: \hyperlink{row-I08}{I08}, \hyperlink{row-I09}{I09}.\\[7pt]
\hypertarget{formula-F61}{}\textbf{F61}\label{formula:F61} & $\tau_p/B(e^+\pi^0)=72^2\hbar/(3^{24}W^4m_p)=1.170\times10^{37}\,\mathrm{yr}$.\par
The proton-decay attempt uses the six-complex-direction exchange scale and its stated operator convention.\par
Derivation: \eqref{eq:protonlife}. Rows: \hyperlink{row-K06}{K06}, \hyperlink{row-L12}{L12}.\\[7pt]
\hypertarget{formula-F62}{}\textbf{F62}\label{formula:F62} & $K=n_L-2$, $G_K(x)=(4-3x)^{-1}$, $\mathcal C(j)=-a^{11}\log(4-3e^j)$; $\eta_B=\mathcal C^\prime(0)=3a^{11}$.\par
The full positive comparison volume weights the connected return once. Its first derivative is the baryon source.\par
Derivation: \eqref{eq:connectedreturn}, \eqref{eq:positivevolume}, \eqref{eq:singlepopulation}, \eqref{eq:threeandtwelve}. Rows: \hyperlink{row-F01}{F01}.\\[7pt]
\hypertarget{formula-F63}{}\textbf{F63}\label{formula:F63} & $\mathcal C^{\prime\prime}(0)=12a^{11}$; $\mathcal P_{\mathcal R}(k_{\rm com})=12a^{11}(k_{\rm com}\ell)^{-1/40}$, $n_s=1+k((abc))=1-t_0/(\EE N_{\rm ret})=39/40$, $\dd n_s/\dd\log k_{\rm com}=0$.\par $A_s(0.05\,\mathrm{Mpc}^{-1})=2.13555115835\times10^{-9}$; $r=16\Lambda\ell_P^2/(3\pi\,12a^{11})$.\par
The second derivative of the same source gives the scalar fluctuation. Part III retains $b^{24}(k_{\rm com}\ell)^{-1/40}$ as its empirical alternative.\par
Derivation: \eqref{eq:singlepopulation}, \eqref{eq:threeandtwelve}, \eqref{eq:primordial}. Rows: \hyperlink{row-C15}{C15}, \hyperlink{row-C16}{C16}, \hyperlink{row-L06}{L06}, \hyperlink{row-F03}{F03}, \hyperlink{row-S08}{S08}, \hyperlink{row-L11}{L11}.\\[7pt]
\hypertarget{formula-F64}{}\textbf{F64}\label{formula:F64} & $\Lambda\ell_P^2=3W^6\theta_cs_F^3$, $\theta_c=2\pi/(16R_p)$, $s_F=r_0/10=3/50$.\par $\ell=\lambda_p/[W^2\sqrt{\theta_cs_F^3}]=5369.08\,\mathrm{Mpc}$, $k_BT_{\rm dS}=\hbar c/(2\pi\ell)=1.8956\times10^{-34}\eV$, $S_{\rm dS}/k_B=3\pi/(\Lambda\ell_P^2)$.\par
The two complete-frame entries give the constant curvature, horizon temperature and entropy; $w_\Lambda=-1$.\par
Derivation: \eqref{eq:curvaturewhole}, \eqref{eq:epoch}. Rows: \hyperlink{row-D06}{D06}, \hyperlink{row-T03}{T03}, \hyperlink{row-T04}{T04}, \hyperlink{row-L10}{L10}, \hyperlink{row-L08}{L08}.\\[7pt]
\hypertarget{formula-F65}{}\textbf{F65}\label{formula:F65} & $N_{\rm eff}=3\eta_0\sigma^{4/3}=3.0439$, $\Omega_k=0$, $R_c^{\rm finite}=a^{16}M_R/m_p$, $\Omega_c/\Omega_b=R_c^{\rm finite}/\kappa_B$.\par
The thermal number, whole-frame branch and neutral occupation. The independent count $M/(M-4C_h)=37/7$ remains an empirical alternative in Part III.\par
Derivation: \eqref{eq:epoch}, \eqref{eq:coldnative}, Sec.~\ref{sec:secondconstructions}. Rows: \hyperlink{row-F04}{F04}, \hyperlink{row-F13}{F13}, \hyperlink{row-L07}{L07}, \hyperlink{row-S09}{S09}.\\[7pt]
\hypertarget{formula-F66}{}\textbf{F66}\label{formula:F66} & $\dd\log\alpha/\dd(\Phi_G/c^2)=0$, $1+\gamma=2$, $\Omega_{\rm geo}/\Omega_{\rm Einstein}=1$; second-order bending coefficient $15\pi/4$.\par
The metric tests and coherent comparison of commuting or anticommuting orders, whose ideal success probability is one.\par
Derivation: \eqref{eq:lapse}, \eqref{eq:fullclock}, \eqref{eq:orderstate}, \eqref{eq:orderprobability}, Table~\ref{tab:variablemap}. Rows: \hyperlink{row-D07}{D07}, \hyperlink{row-L09}{L09}, \hyperlink{row-H01}{H01}, \hyperlink{row-H02}{H02}, \hyperlink{row-L18}{L18}, \hyperlink{row-H03}{H03}.\\[7pt]
\hypertarget{formula-F67}{}\textbf{F67}\label{formula:F67} & $z_S=\ii^{|S|}$, $p_S=z_S^2$, $r_{\rm occ}=|S\setminus\{n\}|$, $L_S=5-r_{\rm occ}+r_{\rm occ}p_S$, $w_S=z_S/L_S$.\par $\Psi_S=-\Phis(4w_S)/(16w_S^2)=p_S/(L_S^2+3p_S)$; $r_S=1-2|w_S|$, $t_S=2|w_S|$, $N(S)=|S|(6-|S|)+[|S|\text{ odd}]$.\par
All 64 divided sets: the four classes have sizes $32,2,10,20$, with associations $1/28,-1/22,-1/6,1/2$.\par
Derivation: \eqref{eq:support}. Rows: \hyperlink{row-P32}{P32}.\\[7pt]
\hypertarget{formula-F68}{}\textbf{F68}\label{formula:F68} & $\alpha_s^2=r_{\rm v}^2/8=1/72$, $f=r_0^2/10=9/250$, $\nu=r_0^2/[(\EE n_F)(\EE N_{\rm ret})]=9/4000$.\par
Each endpoint brings its support reflection, shared among the available encounter channels.\par
Derivation: \eqref{eq:support}, \eqref{eq:fnu}, \eqref{eq:alphas}. \\[7pt]
\hypertarget{formula-F69}{}\textbf{F69}\label{formula:F69} & $\Psi_+=1/28$, $\Psi_-=1/2$, $r_{\rm v}^{-1/2}=\sqrt3$.\par
The layer association and amplitude per retained component are values of the divided-support response.\par
Derivation: \eqref{eq:support}, \eqref{eq:layers}, \eqref{eq:pairingcompression}. \\[7pt]
\hypertarget{formula-F70}{}\textbf{F70}\label{formula:F70} & $-\log h_\nu/\nu^2\to\frac32\Var D$ as $\nu\to0$; $\Var D=[(\Var n_F-\Var n_L)^2+4\Cov(n_F,n_L)^2]/(3\EE N_{\rm ret})^2=61/16$.\par $\EE D=(\Var n_F+\Var n_L)/(3\EE N_{\rm ret})=2$; clock marks $(0,1,2,3,19)/4$.\par
The clock interval mean and spread are the first-return covariance invariants.\par
Derivation: \eqref{eq:MQ}, \eqref{eq:h}, \eqref{eq:associationfrommoments}. Rows: \hyperlink{row-P33}{P33}.\\[7pt]
\hypertarget{formula-F71}{}\textbf{F71}\label{formula:F71} & $R_p=\sum_{e\in\mathcal A_h}\langle\Omega,(L+\gamma_h\otimes\gamma_{s(e)})\Omega\rangle-\log\det F_0/6+V_5(a)+V_{\bar h}^\prime(a)/16$.\par $F_0=(4/5)G_3$, $\det F_0=t_0=2/5$, $V_5=-a^5$, $V_{\bar h}=-a^7$.\par
The complete proton functional: multiplying each step by $t_0$ moves its zero by $\log(1/t_0)$; the opened seven-volume gives $-ba^6$.\par
Derivation: \eqref{eq:completeprotonsource}, \eqref{eq:protonpairgram}, \eqref{eq:gramheight}, \eqref{eq:protonvolumes}. Rows: \hyperlink{row-P34}{P34}.\\[7pt]
\hypertarget{formula-F72}{}\textbf{F72}\label{formula:F72} & $R_n=t_0^{-[\Tr((\Id-J_{10}/10)a^{2L_P})+10c_Lba^6]}$; $\Tr((\Id-J_{10}/10)a^{2L_P})=5a^4+4a^{10}$.\par
The Petersen Laplacian has spectrum $0,2^{(5)},5^{(4)}$. The two-ended neutral response retains its nine contrasts.\par
Derivation: \eqref{eq:neutronsourcefunctional}, \eqref{eq:neutroncontrasts}. Rows: \hyperlink{row-P35}{P35}.\\[7pt]
\hypertarget{formula-F73}{}\textbf{F73}\label{formula:F73} & $W=\det_{\CC^3}(t_0^{\EE N_{\rm ret}}\eta_0^{E_H}\Id_3)=(2/5)^{48}(295/296)^9$.\par
The typical whole-frame amplitude is the determinant over the three complementary directions.\par
Derivation: \eqref{eq:incoming}, \eqref{eq:typicallogtransfer}, \eqref{eq:W}. Rows: \hyperlink{row-P31}{P31}.\\[7pt]
\end{longtable}
\endgroup

\subsection{The states in the rest equations}
\paragraph*{Definitions: the physical states in the rest equations.}
A source tells the equation what must return unchanged. Its physical data are the retained state, its positive clock form and its returning amplitude. The common clock counts the complete return; the numerator retains the charge, internal directions and contractions of that state. Table~\ref{tab:masssources} keeps those definitions beside the factors they produce. The spectral theorems determine the frequencies and currents of these sources. Their finite ingredients have been derived in Parts I and III.
\begingroup\small
\setlength{\tabcolsep}{4pt}
\begin{longtable}{@{}>{\raggedright\arraybackslash}p{.16\textwidth}>{\raggedright\arraybackslash}p{.49\textwidth}>{\raggedright\arraybackslash}p{.29\textwidth}@{}}
\caption{One return equation, with the state that fixes its source. Energies in the nucleon row are in electron units. Every response uses the complete derivative in Eq.~\eqref{eq:residue}.}
\label{tab:masssources}\\
\toprule Retained state&Origin of the source factors&Rest equation\\\midrule
\endfirsthead
\multicolumn{3}{l}{\textit{The retained mass sources, continued}}\\
\toprule Retained state&Origin of the source factors&Rest equation\\\midrule
\endhead
\bottomrule\endfoot
Elementary charged state&One of the $320$ equal departures carries the elementary phase. Its simple zero turns at $1/320$; one common change of clock fixes $m_e$ as the reference, Eq.~\eqref{eq:markedsourceclock}.&One marked zero; after the common normalization, $E_e/m_e=1$.\\[6pt]
Charged families&The successive source retains the reference, the ten symmetric states, then the five vector states. Each added component supplies $\sqrt3$; its signed endpoint return and common clock give the single increment in Eq.~\eqref{eq:layers}.&$E^2\in\operatorname{spec}(\widehat A_X\widehat A_X^\dagger)$, with $A_X=M_\ell,M_d,M_u$ in Eq.~\eqref{eq:chargedsource}.\\[6pt]
Common nucleon state&The colour-antisymmetric contraction and the rank-four $P_N$ retain the two spin doublets. The stationary count, pairing determinant and signed volumes form $R_p$; the neutron retains the nine Petersen contrasts. Both enter $A_N$ in Eq.~\eqref{eq:nucleonoperator}.&$E^2\in\operatorname{spec}[(A_N/m_e)^2]$: $R_p^2$ and $(R_pR_n)^2$, each twice.\\[6pt]
Weak vector fields&The actual scalar vector $h_0$, with its phase, is displaced by three rotations. Its Gram form has eigenvalues $1/2,1/2,4/5$ on the traceless algebra; the other nine rotations preserve it. The kinetic form fixes the clock.&$\det(E^2C_B-\mathcal M_B)=0$, with the neutral block in Eq.~\eqref{eq:weakpencil}.\\[6pt]
Radial scalar and top current&The length of the same scalar fluctuates about $v$. The source quartic uses the inverse susceptibility $1/\chi_h$ of the moved fraction $3/13$. The associated top current compares $\eta_0$ with $\sigma$, Eq.~\eqref{eq:radial}.&$m_H^2=2\lambda_Hv^2$; $m_t^{\rm rad}=v\sqrt{\eta_0/(2\sigma)}$.\\[6pt]
Neutral fermions&The present amplitude returns through a charge-conjugate partner. Nine reflected directions give $M_R$; the real and half-phase quadratures give $(0,\sqrt{\rho_\nu},1)$, Eq.~\eqref{eq:neutralfull}.&$m_i(M_R+m_i)=m_{Di}^2$; the same return has active weight in Eq.~\eqref{eq:neutrinoresidue}.\\[6pt]
Light pseudoscalars&Two constituent attenuations meet with overlap $G$. The eigenmodes of $R=D^{1/2}GD^{1/2}$ carry their joint return. The pion also fixes the light attenuation scale, Eqs.~\eqref{eq:pions}--\eqref{eq:kaons}.&$m^2=6m_p^2\Tr\delta_1(R)$ on the selected light branch, with its electric contact.\\
\end{longtable}
\endgroup

\section{Complete comparison table}
\label{sec:register}
The complete table contains 252 distinct records: 37 structural identities,
113 main-line scalar readings, eighteen additional scalar readings, three
evaluated lists, 43 classical conversions, twelve further evaluated composite
expressions and 26 analytic comparison records.
Each entry retains its formula, evaluated value,
measurement, source, comparison width and operation. Named Lean declarations identify the proved finite statements; the calculation and input records identify the evaluated comparisons.
The independent classical evolutions follow in a separate continuation of the
table. Quantities with a common source retain their correlations.

For a value $x$ and comparison $x_0$ with the stated one-sided widths,
\begin{equation}
 r=\frac{x-x_0}{w_+}\quad(x\geq x_0),\qquad
 r=\frac{x-x_0}{w_-}\quad(x<x_0).
 \label{eq:tableresidual}
\end{equation}
The comparison column states the confidence convention. Bounds, identities and
predictions have their own labels. A dash denotes an entry without a scalar
two-sided residual. The names of Lean declarations identify the proof
record; the analytic derivations above give their role in the present article.
\newpage
\begingroup\small
\setlength{\tabcolsep}{3pt}
\setlength{\LTcapwidth}{\textwidth}
\renewcommand{\arraystretch}{1.07}
\begin{longtable}{@{}>{\raggedright\arraybackslash}p{.065\textwidth}>{\raggedright\arraybackslash}p{.9\textwidth}@{}}
\caption{The complete finite comparison register.}\label{tab:numberscomplete}\\
\toprule ID&Quantity, calculation, measurement and proof\\\midrule
\endfirsthead
\multicolumn{2}{l}{\textit{Complete finite comparison register, continued}}\\
\toprule ID&Quantity, calculation, measurement and proof\\\midrule
\endhead
\bottomrule\endfoot
\multicolumn{2}{@{}l}{\textbf{The numbers the geometry and the law give directly}}\\*
\hypertarget{row-P01}{}\textbf{P01}\label{row:P01} & \textbf{the comparison of two systems}\par
\textit{Formula:} $L=\Id+\sum_t\gamma_t\otimes\gamma_t=2(P+E_\Omega)$; $\operatorname{spec}L=\{6^{(1)},2^{(10)},(-2)^{(5)}\}$.\par
\textit{Value:} 6, 2, -2; L(\ensuremath{\Omega}) = 6\quad\textit{Comparison:} Exact finite identity\quad\textit{Residual:} --\par
\textit{Operation:} Computed finite identity.\par
\textit{Calculation:} \path{programs/article_results.py}; \path{outputs/article_results.out}, line 299, identifier \path{SI.P01}.\par
{\footnotesize\textit{Lean:} \path{Comparison.System.prop_10} (\path{lean/Comparison/System/Geometry/TwoThings.lean}:337); \path{Comparison.System.prop_11} (\path{lean/Comparison/System/Geometry/TwoThings.lean}:359); \path{Comparison.Formulas.LOmega_inner} (\path{lean/Comparison/Formulas/MainLine.lean}:22); \path{Comparison.System.Completion.Num.LOmega_eq} (\path{lean/Comparison/System/Completion/Numbers.lean}:226); \path{Comparison.System.S3_1} (\path{lean/Comparison/System/Finite/Geometry.lean}:74)}\\[7pt]
\hypertarget{row-P02}{}\textbf{P02}\label{row:P02} & \textbf{the types of distinction}\par
\textit{Formula:} $\binom d2+\binom{6-d}2=(15,10,7,6)$; $d(6-d)=(0,5,8,9)$ for $d=0,1,2,3$.\par
\textit{Value:} 15, 10, 7, 6; 0, 5, 8, 9\quad\textit{Comparison:} Exact finite identity\quad\textit{Residual:} --\par
\textit{Operation:} Computed finite identity.\par
\textit{Calculation:} \path{programs/article_results.py}; \path{outputs/article_results.out}, line 300, identifier \path{SI.P02}.\par
{\footnotesize\textit{Lean:} \path{Comparison.Formulas.central_kept_iff_card} (\path{lean/Comparison/Formulas/Clocks.lean}:292); \path{Comparison.Formulas.no_clock_one_five} (\path{lean/Comparison/Formulas/Clocks.lean}:313); \path{Comparison.System.S3_2} (\path{lean/Comparison/System/Finite/Geometry.lean}:111)}\\[7pt]
\hypertarget{row-P03}{}\textbf{P03}\label{row:P03} & \textbf{the two clocks}\par
\textit{Formula:} ${\mathsf T}_1=\ii\gamma_h$, ${\mathsf T}_2=\gamma_h\gamma_k$, ${\mathsf T}_3=\ii\gamma_k$; $H_k=(\Id+\ii\gamma_k)/\sqrt2$, $H_k{\mathsf T}_1H_k^\dagger={\mathsf T}_2$; ${\mathsf T}_a^2=-\Id$, ${\mathsf T}_1{\mathsf T}_2={\mathsf T}_3$.\par
\textit{Value:} T\_\allowbreak{}a\ensuremath{{}^{2}} = -1\quad\textit{Comparison:} Exact finite identity\quad\textit{Residual:} --\par
\textit{Operation:} Computed finite identity.\par
\textit{Calculation:} \path{programs/article_results.py}; \path{outputs/article_results.out}, line 301, identifier \path{SI.P03}.\par
{\footnotesize\textit{Lean:} \path{Comparison.Formulas.quaternion_triple} (\path{lean/Comparison/Formulas/Clocks.lean}:92); \path{Comparison.Formulas.halfStep_carries_clock} (\path{lean/Comparison/Formulas/Clocks.lean}:81); \path{Comparison.System.prop_91} (\path{lean/Comparison/System/Existence/Numbered.lean}:200)}\\[7pt]
\hypertarget{row-P04}{}\textbf{P04}\label{row:P04} & \textbf{one generation}\par
\textit{Formula:} $Y=\widehat n_F/2-\widehat n_C/3$, $Q=\widehat n_k-\widehat n_C/3=T_3+Y$; $\Tr Y=\Tr Y^3=0$.\par
\textit{Value:} 0\quad\textit{Comparison:} Exact finite identity\quad\textit{Residual:} --\par
\textit{Operation:} Computed finite identity.\par
\textit{Calculation:} \path{programs/article_results.py}; \path{outputs/article_results.out}, line 302, identifier \path{SI.P04}.\par
{\footnotesize\textit{Lean:} \path{Comparison.System.prop_45} (\path{lean/Comparison/System/Existence/Numbered.lean}:298); \path{Comparison.System.S3_3} (\path{lean/Comparison/System/Finite/Geometry.lean}:140)}\\[7pt]
\hypertarget{row-P05}{}\textbf{P05}\label{row:P05} & \textbf{the Gram of h at the source}\par
\textit{Formula:} $\operatorname{spec}\mathcal M=(1/2,1/2,4/5,0^{(9)})$ in $\Re\Tr X^\dagger Y$-orthonormal generators. In $\tr(t^at^b)=\delta^{ab}/2$ normalization the nonzero values are $(1/4,1/4,2/5)$.\par
\textit{Value:} 5/8, 3/8\quad\textit{Comparison:} Exact finite identity\quad\textit{Residual:} --\par
\textit{Operation:} Computed finite identity.\par
\textit{Calculation:} \path{programs/article_results.py}; \path{outputs/article_results.out}, line 303, identifier \path{SI.P05}.\par
{\footnotesize\textit{Lean:} \path{Comparison.System.prop_47} (\path{lean/Comparison/System/Existence/Numbered.lean}:332); \path{Comparison.System.S3_4} (\path{lean/Comparison/System/Finite/Geometry.lean}:174)}\\[7pt]
\hypertarget{row-P06}{}\textbf{P06}\label{row:P06} & \textbf{one time and three space}\par
\textit{Formula:} $\mathcal N_{2,4}=(S^1\times S^3)/\mathbb Z_2$; $\operatorname{sign}T\mathcal N_{2,4}=(1,3)$.\par
\textit{Value:} $(1,3)$, negative directions first\quad\textit{Comparison:} Exact finite identity\quad\textit{Residual:} --\par
\textit{Operation:} Computed finite identity.\par
\textit{Calculation:} \path{programs/article_results.py}; \path{outputs/article_results.out}, line 304, identifier \path{SI.P06}.\par
{\footnotesize\textit{Lean:} \path{Comparison.System.prop_21} (\path{lean/Comparison/System/Geometry/RealFormSU4.lean}:357)}\\[7pt]
\hypertarget{row-P07}{}\textbf{P07}\label{row:P07} & \textbf{the step B}\par
\textit{Formula:} $B_{(y,s),(x,t)}=[s\ne t][y=x\mathbin\triangle g_s]$, with $g_s=\{1,\ldots,5\}\setminus\{s\}$.\par
\textit{Value:} 80, 320\quad\textit{Comparison:} Exact finite identity\quad\textit{Residual:} --\par
\textit{Operation:} Computed finite identity.\par
\textit{Calculation:} \path{programs/article_results.py}; \path{outputs/article_results.out}, line 305, identifier \path{SI.P07}.\par
{\footnotesize\textit{Lean:} \path{Comparison.System.prop_26} (\path{lean/Comparison/System/Existence/Numbered.lean}:56)}\\[7pt]
\hypertarget{row-P08}{}\textbf{P08}\label{row:P08} & \textbf{the law}\par
\textit{Formula:} $\det(\Id-uB)=(1-u^2)^{24}u^{16}\det[\Delta+\chi(u)\Id]$, $\Delta=6\Id-L$, $\chi=u^{-1}+4u-5$.\par
\textit{Value:} one leading zero, k\ensuremath{{}_{0}}(1/4) = 0\quad\textit{Comparison:} Exact finite identity\quad\textit{Residual:} --\par
\textit{Operation:} Computed finite identity.\par
\textit{Calculation:} \path{programs/article_results.py}; \path{outputs/article_results.out}, line 306, identifier \path{SI.P08}.\par
{\footnotesize\textit{Lean:} \path{Comparison.Formulas.law_gap} (\path{lean/Comparison/Formulas/Kernel.lean}:186); \path{Comparison.Formulas.law_bass_form} (\path{lean/Comparison/Formulas/Kernel.lean}:178); \path{Comparison.Formulas.kN_gap} (\path{lean/Comparison/Formulas/Kernel.lean}:215); \path{Comparison.Formulas.lawPoly_eq} (\path{lean/Comparison/Formulas/MainLine.lean}:70); \path{Comparison.System.prop_54} (\path{lean/Comparison/System/Law/Determinant.lean}:246); \path{Comparison.System.S3_5} (\path{lean/Comparison/System/Finite/Law.lean}:24)}\\[7pt]
\hypertarget{row-P09}{}\textbf{P09}\label{row:P09} & \textbf{the four points u = i\textasciicircum{}d/4}\par
\textit{Formula:} $u=\ii^d/4$, $d=0,1,2,3$, gives $\chi=(0,-5-3\ii,-10,-5+3\ii)$.\par
\textit{Value:} 0, -5 - 3i, -10, -5 + 3i\quad\textit{Comparison:} Exact finite identity\quad\textit{Residual:} --\par
\textit{Operation:} Computed finite identity.\par
\textit{Calculation:} \path{programs/article_results.py}; \path{outputs/article_results.out}, line 307, identifier \path{SI.P09}.\par
{\footnotesize\textit{Lean:} \path{Comparison.Formulas.chi_quarters} (\path{lean/Comparison/Formulas/Kernel.lean}:221); \path{Comparison.System.table_VII} (\path{lean/Comparison/System/Law/Quarters.lean}:104)}\\[7pt]
\hypertarget{row-P10}{}\textbf{P10}\label{row:P10} & \textbf{grade factors, transmission and reflection at u = -1/4}\par
\textit{Formula:} $k_N(-1/4)=(5-N)/2$, so $(k_0,k_2,k_4)=(5/2,3/2,1/2)$.\par $t_0=1/k_0=2/5$, $r_0=k_2/k_0=3/5$, $t_{\rm v}=1/k_2=2/3$, $r_{\rm v}=k_4/k_2=1/3$.\par
\textit{Value:} 5/2, 3/2, 1/2; 2/5, 3/5, 2/3, 1/3\quad\textit{Comparison:} Exact finite identity\quad\textit{Residual:} --\par
\textit{Operation:} Computed finite identity.\par
\textit{Calculation:} \path{programs/article_results.py}; \path{outputs/article_results.out}, line 308, identifier \path{SI.P10}.\par
{\footnotesize\textit{Lean:} \path{Comparison.Formulas.kN_plane_gap} (\path{lean/Comparison/Formulas/Kernel.lean}:240); \path{Comparison.Formulas.transmission_echo} (\path{lean/Comparison/Formulas/Kernel.lean}:249); \path{Comparison.System.Completion.Num.k_zero} (\path{lean/Comparison/System/Completion/Numbers.lean}:46); \path{Comparison.System.Completion.Num.k_two} (\path{lean/Comparison/System/Completion/Numbers.lean}:47); \path{Comparison.System.Completion.Num.k_four} (\path{lean/Comparison/System/Completion/Numbers.lean}:48); \path{Comparison.System.Completion.Num.t_eq} (\path{lean/Comparison/System/Completion/Numbers.lean}:61); \path{Comparison.System.Completion.Num.e_eq} (\path{lean/Comparison/System/Completion/Numbers.lean}:62); \path{Comparison.System.Completion.Num.tLine_eq} (\path{lean/Comparison/System/Completion/Numbers.lean}:63); \path{Comparison.System.Completion.Num.eLine_eq} (\path{lean/Comparison/System/Completion/Numbers.lean}:64); \path{Comparison.System.prop_44} (\path{lean/Comparison/System/Existence/Numbered.lean}:272)}\\[7pt]
\hypertarget{row-P11}{}\textbf{P11}\label{row:P11} & \textbf{the long-wave limit}\par
\textit{Formula:} $u=1/4-\epsilon^2m^2/24$ gives $\chi=\epsilon^2m^2/2+O(\epsilon^4)$ and $(1+4u^2)\Id-uA_\epsilon=(\epsilon^2/8)(m^2-\Delta_{\rm tan})+O(\epsilon^3)$.\par
$\epsilon^2M_\epsilon^2=2\chi+\chi^2/5$, $M_\epsilon\to m$.\par
\textit{Value:} $\chi/(\epsilon^2m^2/2)\to1$ for $m>0$\quad\textit{Comparison:} Exact finite identity\quad\textit{Residual:} --\par
\textit{Operation:} Computed finite identity.\par
\textit{Calculation:} \path{programs/article_results.py}; \path{outputs/article_results.out}, line 309, identifier \path{SI.P11}.\par
{\footnotesize\textit{Lean:} \path{Comparison.Formulas.chi_mass_shell} (\path{lean/Comparison/Formulas/MainLine.lean}:40); \path{Comparison.Formulas.chart_symbol_second_order} (\path{lean/Comparison/Formulas/OwnChart.lean}:210); \path{Comparison.System.prop_59} (\path{lean/Comparison/System/Law/Critical.lean}:202)}\\[7pt]
\hypertarget{row-P12}{}\textbf{P12}\label{row:P12} & \textbf{the first-order step}\par
\textit{Formula:} $U=e^{-\ii k\gamma_j}e^{-\ii\mu\gamma_h}$, $U+U^\dagger=2\cos k\cos\mu\,\Id$; $\cos\omega=\cos\mu\cos k$.\par
\textit{Value:} cos \ensuremath{\omega} = cos \ensuremath{\mu} cos k\quad\textit{Comparison:} Exact finite identity\quad\textit{Residual:} --\par
\textit{Operation:} Computed finite identity.\par
\textit{Calculation:} \path{programs/article_results.py}; \path{outputs/article_results.out}, line 310, identifier \path{SI.P12}.\par
{\footnotesize\textit{Lean:} \path{Comparison.System.prop_60} (\path{lean/Comparison/System/Law/Motion.lean}:289)}\\[7pt]
\hypertarget{row-P13}{}\textbf{P13}\label{row:P13} & \textbf{the gauge field's stiffness}\par
\textit{Formula:} $\log\det(\Id-wA_-[e^{\ii\theta}])=\log\det(\Id-wA_-)-\kappa\|C_\square\theta\|^2/2+O(\theta^4)$.\par $w=(\chi+5)^{-1}$, $\kappa=2/[(\chi+5)^2-5]^2$.\par
\textit{Value:} 1/200 at \ensuremath{\chi} = 0, -10; 1/98 at \ensuremath{\chi} = -5 \ensuremath{\mp} 3i\quad\textit{Comparison:} Exact finite identity\quad\textit{Residual:} --\par
\textit{Operation:} Computed finite identity.\par
\textit{Calculation:} \path{programs/article_results.py}; \path{outputs/article_results.out}, line 311, identifier \path{SI.P13}.\par
{\footnotesize\textit{Lean:} \path{Comparison.Formulas.kappa_gap} (\path{lean/Comparison/Formulas/Kernel.lean}:260); \path{Comparison.Formulas.kappa_gap_values} (\path{lean/Comparison/Formulas/Kernel.lean}:288); \path{Comparison.System.prop_62} (\path{lean/Comparison/System/Law/PhaseLaw.lean}:138); \path{Comparison.System.S3_15} (\path{lean/Comparison/System/Finite/Completion.lean}:392)}\\[7pt]
\hypertarget{row-P14}{}\textbf{P14}\label{row:P14} & \textbf{a circulation moves the leading zero}\par
\textit{Formula:} $\log(4u_c)=\|\theta_{\rm cyc}\|^2/48+O(\theta^4)$; $48=3\EE N_{\rm ret}=2(40-16)$.\par
\textit{Value:} 48\quad\textit{Comparison:} Exact finite identity\quad\textit{Residual:} --\par
\textit{Operation:} Computed finite identity.\par
\textit{Calculation:} \path{programs/article_results.py}; \path{outputs/article_results.out}, line 312, identifier \path{SI.P14}.\par
{\footnotesize\textit{Lean:} \path{Comparison.System.S3_17} (\path{lean/Comparison/System/Finite/Completion.lean}:538)}\\[7pt]
\hypertarget{row-P15}{}\textbf{P15}\label{row:P15} & \textbf{the constant term}\par
\textit{Formula:} $\tr_4[\gamma_h(\mathcal F_{\rm sp}\wedge\mathcal F_{\rm sp}-R_{\rm sp}\wedge R_{\rm sp})]=-\ell^{-2}(R-6/\ell^2)\operatorname{vol}$.\par
\textit{Value:} \ensuremath{\Lambda}\ensuremath{\ell}\ensuremath{{}^{2}} = 3\quad\textit{Comparison:} Exact finite identity\quad\textit{Residual:} --\par
\textit{Operation:} Computed finite identity.\par
\textit{Calculation:} \path{programs/article_results.py}; \path{outputs/article_results.out}, line 313, identifier \path{SI.P15}.\par
{\footnotesize\textit{Lean:} \path{Comparison.System.prop_61} (\path{lean/Comparison/System/Law/LastDirection.lean}:32); \path{Comparison.System.S3_15} (\path{lean/Comparison/System/Finite/Completion.lean}:392)}\\[7pt]
\hypertarget{row-P16}{}\textbf{P16}\label{row:P16} & \textbf{the lapse couples to the rest energy}\par
\textit{Formula:} $K_X(E,{\mathcal N})=K_X(E/{\mathcal N},1)$, $K_{\mathcal N}=-(E/{\mathcal N})K_E$, $\dd E_X/\dd {\mathcal N}=E_X({\mathcal N})/{\mathcal N}$ for ${\mathcal N}>0$.\par
\textit{Value:} 1\quad\textit{Comparison:} Exact finite identity\quad\textit{Residual:} --\par
\textit{Operation:} Computed finite identity.\par
\textit{Calculation:} \path{programs/article_results.py}; \path{outputs/article_results.out}, line 314, identifier \path{SI.P16}.\par
{\footnotesize\textit{Lean:} \path{Comparison.Formulas.lapse_displacement} (\path{lean/Comparison/Formulas/Lapse.lean}:66); \path{Comparison.Formulas.lapse_rate} (\path{lean/Comparison/Formulas/Lapse.lean}:61)}\\[7pt]
\hypertarget{row-P17}{}\textbf{P17}\label{row:P17} & \textbf{the first return and the mean return time}\par
\textit{Formula:} $\Fret(z)=z^4(1+z)/(16-8z-3z^2-3z^3)$; $\Fret(1)=1$, $\EE N_{\rm ret}=16$, $\Var N_{\rm ret}=156$.\par
\textit{Value:} 16, 156\quad\textit{Comparison:} Exact finite identity\quad\textit{Residual:} --\par
\textit{Operation:} Computed finite identity.\par
\textit{Calculation:} \path{programs/article_results.py}; \path{outputs/article_results.out}, line 315, identifier \path{SI.P17}.\par
{\footnotesize\textit{Lean:} \path{Comparison.System.prop_29} (\path{lean/Comparison/System/Existence/Numbered.lean}:87); \path{Comparison.System.Completion.Num.El_eq} (\path{lean/Comparison/System/Completion/Numbers.lean}:83); \path{Comparison.System.Completion.Num.VarEll_eq} (\path{lean/Comparison/System/Completion/Numbers.lean}:84); \path{Comparison.System.S3_9} (\path{lean/Comparison/System/Finite/Excursion.lean}:60)}\\[7pt]
\hypertarget{row-P18}{}\textbf{P18}\label{row:P18} & \textbf{the signed first return}\par
\textit{Formula:} $\Phis(z)=-z^4/(16+3z^2)$; $\Phis(\pm1)=-1/19$, $\Phis(\pm\ii)=-1/13$.\par
\textit{Value:} -1/19, -1/13\quad\textit{Comparison:} Exact finite identity\quad\textit{Residual:} --\par
\textit{Operation:} Computed finite identity.\par
\textit{Calculation:} \path{programs/article_results.py}; \path{outputs/article_results.out}, line 316, identifier \path{SI.P18}.\par
{\footnotesize\textit{Lean:} \path{Comparison.System.prop_30} (\path{lean/Comparison/System/Existence/Numbered.lean}:98); \path{Comparison.System.Completion.Num.Phi_closed} (\path{lean/Comparison/System/Completion/Numbers.lean}:239); \path{Comparison.System.Completion.Num.Phi_one} (\path{lean/Comparison/System/Completion/Numbers.lean}:246); \path{Comparison.System.Completion.Num.Phi_I} (\path{lean/Comparison/System/Completion/Numbers.lean}:249)}\\[7pt]
\hypertarget{row-P19}{}\textbf{P19}\label{row:P19} & \textbf{the visits of one excursion}\par
\textit{Formula:} $\EE n_F=10$, $\EE n_L=5$, $\Var n_F=84$, $\Var n_L=12$, $\Cov(n_F,n_L)=30$.\par $\EE x^{n_F}=x(1+x)/(16-11x-3x^2)$; $x_c=(\sqrt{313}-11)/6$.\par
\textit{Value:} 10, 5, 84, 12, 30\quad\textit{Comparison:} Exact finite identity\quad\textit{Residual:} --\par
\textit{Operation:} Computed finite identity.\par
\textit{Calculation:} \path{programs/article_results.py}; \path{outputs/article_results.out}, line 317, identifier \path{SI.P19}.\par
{\footnotesize\textit{Lean:} \path{Comparison.System.Completion.Num.EP_eq} (\path{lean/Comparison/System/Completion/Numbers.lean}:100); \path{Comparison.System.Completion.Num.EQ_eq} (\path{lean/Comparison/System/Completion/Numbers.lean}:101); \path{Comparison.System.Completion.Num.VarP_eq} (\path{lean/Comparison/System/Completion/Numbers.lean}:102); \path{Comparison.System.Completion.Num.VarQ_eq} (\path{lean/Comparison/System/Completion/Numbers.lean}:103); \path{Comparison.System.Completion.Num.Cov_eq} (\path{lean/Comparison/System/Completion/Numbers.lean}:104); \path{Comparison.System.Completion.Num.xc_root} (\path{lean/Comparison/System/Completion/Numbers.lean}:387); \path{Comparison.System.S3_10} (\path{lean/Comparison/System/Finite/Excursion.lean}:95)}\\[7pt]
\hypertarget{row-P20}{}\textbf{P20}\label{row:P20} & \textbf{the available transitions}\par
\textit{Formula:} 74 retained remembered states; $M=296$, $C_h=60$, $C_j=59$; $c_L=15/74$, $\sigma=75/74$, $\eta_0=295/296$, $R_{ud}=34/74=17/37$, $E_H=3$.\par
\textit{Value:} 296, 60, 59\quad\textit{Comparison:} Exact finite identity\quad\textit{Residual:} --\par
\textit{Operation:} Computed finite identity.\par
\textit{Calculation:} \path{programs/article_results.py}; \path{outputs/article_results.out}, line 318, identifier \path{SI.P20}.\par
{\footnotesize\textit{Lean:} \path{Comparison.System.prop_31} (\path{lean/Comparison/System/Existence/Numbered.lean}:109); \path{Comparison.System.Completion.Num.Nlife_eq} (\path{lean/Comparison/System/Completion/Numbers.lean}:195); \path{Comparison.System.Completion.Num.M_eq} (\path{lean/Comparison/System/Completion/Numbers.lean}:191); \path{Comparison.System.Completion.Num.Ch_eq} (\path{lean/Comparison/System/Completion/Numbers.lean}:192); \path{Comparison.System.Completion.Num.Cj_eq} (\path{lean/Comparison/System/Completion/Numbers.lean}:193); \path{Comparison.System.Completion.Num.Nup_eq} (\path{lean/Comparison/System/Completion/Numbers.lean}:194); \path{Comparison.System.Completion.Num.cL_eq} (\path{lean/Comparison/System/Completion/Numbers.lean}:203); \path{Comparison.System.Completion.Num.eta_eq} (\path{lean/Comparison/System/Completion/Numbers.lean}:205); \path{Comparison.System.Completion.Num.sigma_eq} (\path{lean/Comparison/System/Completion/Numbers.lean}:204); \path{Comparison.System.Completion.Num.Rud_eq} (\path{lean/Comparison/System/Completion/Numbers.lean}:206); \path{Comparison.System.Completion.Num.EH_eq} (\path{lean/Comparison/System/Completion/Numbers.lean}:212)}\\[7pt]
\hypertarget{row-P21}{}\textbf{P21}\label{row:P21} & \textbf{the effective resistance b}\par
\textit{Formula:} ${\mathcal G_0}=(5\Id-A)^+=(7\Id+A)/32-3J_{16}/128$; $b={\mathcal G_0}_{00}+{\mathcal G_0}_{FF}-2{\mathcal G_0}_{0F}=7/16$.\par
\textit{Value:} 7/16\quad\textit{Comparison:} Exact finite identity\quad\textit{Residual:} --\par
\textit{Operation:} Computed finite identity.\par
\textit{Calculation:} \path{programs/article_results.py}; \path{outputs/article_results.out}, line 319, identifier \path{SI.P21}.\par
{\footnotesize\textit{Lean:} \path{Comparison.System.Completion.Num.b_eq} (\path{lean/Comparison/System/Completion/Numbers.lean}:280); \path{Comparison.System.S3_11} (\path{lean/Comparison/System/Finite/Excursion.lean}:145)}\\[7pt]
\hypertarget{row-P22}{}\textbf{P22}\label{row:P22} & \textbf{the signed returns at the points}\par
\textit{Formula:} $a=[1-\Phis(1)]/8=5/38$, $\chi_A=[1+\Phis(\ii)]^{-1}=13/12$, $w_h=1/(4\chi_A)=3/13$, $\chi_h=w_h(1-w_h)=30/169$.\par
\textit{Value:} 5/38, 13/12, 3/13, 30/169\quad\textit{Comparison:} Exact finite identity\quad\textit{Residual:} --\par
\textit{Operation:} Computed finite identity.\par
\textit{Calculation:} \path{programs/article_results.py}; \path{outputs/article_results.out}, line 320, identifier \path{SI.P22}.\par
{\footnotesize\textit{Lean:} \path{Comparison.System.Completion.Num.a_eq} (\path{lean/Comparison/System/Completion/Numbers.lean}:275); \path{Comparison.System.Completion.Num.chiA_eq} (\path{lean/Comparison/System/Completion/Numbers.lean}:287); \path{Comparison.System.Completion.Num.wh_eq} (\path{lean/Comparison/System/Completion/Numbers.lean}:288); \path{Comparison.System.Completion.Num.chih_eq} (\path{lean/Comparison/System/Completion/Numbers.lean}:289); \path{Comparison.System.S3_14} (\path{lean/Comparison/System/Finite/Completion.lean}:341)}\\[7pt]
\hypertarget{row-P23}{}\textbf{P23}\label{row:P23} & \textbf{the response against its own clock}\par
\textit{Formula:} $J_{DE}/C_W=(|D||E|-5\tr P_DP_E)/72$.\par $1/72=-\Phis(1)/[4(1+\Phis(1))]=-\Phis(\ii)/[6(1+\Phis(\ii))]=r_{\rm v}^2/8$.\par
\textit{Value:} 1/72, 1/18, 1/12\quad\textit{Comparison:} Exact finite identity\quad\textit{Residual:} --\par
\textit{Operation:} Computed finite identity.\par
\textit{Calculation:} \path{programs/article_results.py}; \path{outputs/article_results.out}, line 321, identifier \path{SI.P23}.\par
{\footnotesize\textit{Lean:} \path{Comparison.Formulas.sweep} (\path{lean/Comparison/Formulas/Sweep.lean}:273); \path{Comparison.Formulas.sweep_ratios} (\path{lean/Comparison/Formulas/Sweep.lean}:322); \path{Comparison.Formulas.crossing_g_sq} (\path{lean/Comparison/Formulas/Sweep.lean}:363); \path{Comparison.System.S3_12} (\path{lean/Comparison/System/Finite/Completion.lean}:260)}\\[7pt]
\hypertarget{row-P24}{}\textbf{P24}\label{row:P24} & \textbf{the weight of one transition}\par
\textit{Formula:} $\langle\Omega,(L+\gamma_h\otimes\gamma_s)\Omega\rangle=6+[s=h]$.\par
\textit{Value:} 1836\quad\textit{Comparison:} Exact finite identity\quad\textit{Residual:} --\par
\textit{Operation:} Computed finite identity.\par
\textit{Calculation:} \path{programs/article_results.py}; \path{outputs/article_results.out}, line 322, identifier \path{SI.P24}.\par
{\footnotesize\textit{Lean:} \path{Comparison.System.Completion.Num.transitionSum_eq} (\path{lean/Comparison/System/Completion/Numbers.lean}:219); \path{Comparison.System.prop_78} (\path{lean/Comparison/System/Completion/Colour.lean}:108)}\\[7pt]
\hypertarget{row-P25}{}\textbf{P25}\label{row:P25} & \textbf{the unit: one marked transition}\par
\textit{Formula:} $D_\vartheta(u)=D(u)+(e^{-\ii\vartheta}-1)uD'(u)/320$; $\arg(4u_c)=\vartheta/320+O(\vartheta^3)$.\par
\textit{Value:} 1/320\quad\textit{Comparison:} Exact finite identity\quad\textit{Residual:} --\par
\textit{Operation:} Computed finite identity.\par
\textit{Calculation:} \path{programs/article_results.py}; \path{outputs/article_results.out}, line 323, identifier \path{SI.P25}.\par
{\footnotesize\textit{Lean:} \path{Comparison.Formulas.one_arc_explicit} (\path{lean/Comparison/Formulas/MainLine.lean}:78); \path{Comparison.System.prop_203} (\path{lean/Comparison/System/Completion/DividedSet.lean}:143); \path{Comparison.System.S3_13} (\path{lean/Comparison/System/Finite/Completion.lean}:294)}\\[7pt]
\hypertarget{row-P26}{}\textbf{P26}\label{row:P26} & \textbf{the two-source coefficients at the plane hk}\par
\textit{Formula:} $f=r_0^2/10=9/250$, $\nu=f/\EE N_{\rm ret}=9/4000$.\par
\textit{Value:} 9/250, 9/4000\quad\textit{Comparison:} Exact finite identity\quad\textit{Residual:} --\par
\textit{Operation:} Computed finite identity.\par
\textit{Calculation:} \path{programs/article_results.py}; \path{outputs/article_results.out}, line 324, identifier \path{SI.P26}.\par
{\footnotesize\textit{Lean:} \path{Comparison.System.Completion.Num.f_eq} (\path{lean/Comparison/System/Completion/Numbers.lean}:298); \path{Comparison.System.Completion.Num.nu_eq} (\path{lean/Comparison/System/Completion/Numbers.lean}:299); \path{Comparison.System.prop_64} (\path{lean/Comparison/System/Completion/Constants.lean}:31)}\\[7pt]
\hypertarget{row-P27}{}\textbf{P27}\label{row:P27} & \textbf{the clock's overlap}\par
\textit{Formula:} $(q_0,q_1,q_2,q_3,q_h)=(0,1,2,3,19)/4$; $M_Q(s)=\frac1{10}\sum_De^{-sD}$, $\EE D=2$.\par $h_\nu=[M_Q(\nu)M_Q(-\nu)]^{-3/2}=0.99997104933$.\par
\textit{Value:} 0.99997104933\quad\textit{Comparison:} Exact finite identity\quad\textit{Residual:} --\par
\textit{Operation:} Computed finite identity.\par
\textit{Calculation:} \path{programs/article_results.py}; \path{outputs/article_results.out}, line 325, identifier \path{SI.P27}.\par
{\footnotesize\textit{Lean:} \path{Comparison.System.Completion.Num.h_closed} (\path{lean/Comparison/System/Completion/Numbers.lean}:307); \path{Comparison.Formulas.clock_sweep_limit} (\path{lean/Comparison/Formulas/Rules.lean}:242); \path{Comparison.System.prop_206} (\path{lean/Comparison/System/Completion/Whole.lean}:169)}\\[7pt]
\hypertarget{row-P28}{}\textbf{P28}\label{row:P28} & \textbf{the repeated return along h}\par
\textit{Formula:} $N_h=\Re\sum_{j\geq1}h_\nu^jp_j$, where $p_j=-\Phis(4(\ii/4)^j)$ for odd $j$ and $p_j=\Fret(4(\ii/4)^j)$ for even $j$; $N_h=0.0770858406$.\par
\textit{Value:} 0.0770858406\quad\textit{Comparison:} Exact finite identity\quad\textit{Residual:} --\par
\textit{Operation:} Computed finite identity.\par
\textit{Calculation:} \path{programs/article_results.py}; \path{outputs/article_results.out}, line 326, identifier \path{SI.P28}.\par
{\footnotesize\textit{Lean:} \path{Comparison.System.Completion.Num.Nh_eq_q} (\path{lean/Comparison/System/Completion/Numbers.lean}:319); \path{Comparison.System.prop_72} (\path{lean/Comparison/System/Completion/Constants.lean}:244)}\\[7pt]
\hypertarget{row-P29}{}\textbf{P29}\label{row:P29} & \textbf{the contact interaction}\par
\textit{Formula:} $x=\nu h_\nu/(M+1+N_h)=7.573573\times10^{-6}$.\par
\textit{Value:} 7.573573 \ensuremath{\times} 10\ensuremath{{}^{-6}}\quad\textit{Comparison:} Exact finite identity\quad\textit{Residual:} --\par
\textit{Operation:} Computed finite identity.\par
\textit{Calculation:} \path{programs/article_results.py}; \path{outputs/article_results.out}, line 327, identifier \path{SI.P29}.\par
{\footnotesize\textit{Lean:} \path{Comparison.System.Completion.Num.x_eq_q} (\path{lean/Comparison/System/Completion/Numbers.lean}:330)}\\[7pt]
\hypertarget{row-P30}{}\textbf{P30}\label{row:P30} & \textbf{the weak contact's retained excursion}\par
\textit{Formula:} $W_W=\Fret(h_\nu)h_\nu^{-4}(1-2b+2bh_\nu)=0.9996273896$.\par
\textit{Value:} 0.9996273896\quad\textit{Comparison:} Exact finite identity\quad\textit{Residual:} --\par
\textit{Operation:} Computed finite identity.\par
\textit{Calculation:} \path{programs/article_results.py}; \path{outputs/article_results.out}, line 328, identifier \path{SI.P30}.\par
{\footnotesize\textit{Lean:} \path{Comparison.System.Completion.Num.WW_eq_q} (\path{lean/Comparison/System/Completion/Numbers.lean}:350)}\\[7pt]
\hypertarget{row-P31}{}\textbf{P31}\label{row:P31} & \textbf{the transmission factor W}\par
\textit{Formula:} $W=t_0^{3\EE N_{\rm ret}}\eta_0^{3E_H}=(2/5)^{48}(295/296)^9$.\par $\alpha_G=Gm_p^2/(\hbar c)=W^2=5.906152\times10^{-39}$; $M_{\rm Pl}=m_p/W=1.22088981\times10^{19}\GeV$.\par
\textit{Value:} W\ensuremath{{}^{2}} = 5.906152 \ensuremath{\times} 10\ensuremath{{}^{-39}}\quad\textit{Comparison:} Exact finite identity\quad\textit{Residual:} --\par
\textit{Operation:} Computed finite identity.\par
\textit{Calculation:} \path{programs/article_results.py}; \path{outputs/article_results.out}, line 329, identifier \path{SI.P31}.\par
{\footnotesize\textit{Lean:} \path{Comparison.System.Completion.Num.W_closed} (\path{lean/Comparison/System/Completion/Numbers.lean}:358); \path{Comparison.Formulas.W_minor} (\path{lean/Comparison/Formulas/Rules.lean}:334); \path{Comparison.System.prop_66} (\path{lean/Comparison/System/Completion/Constants.lean}:103); \path{Comparison.System.prop_83} (\path{lean/Comparison/System/Completion/Colour.lean}:224)}\\[7pt]
\hypertarget{row-P32}{}\textbf{P32}\label{row:P32} & \textbf{the rule at every divided set}\par
\textit{Formula:} $z_S=\ii^{|S|}$, $p_S=z_S^2$, $r_{\rm occ}=|S\setminus\{n\}|$, $L_S=5-r_{\rm occ}+r_{\rm occ}p_S$, $w_S=z_S/L_S$.\par
$\Psi_S=-\Phis(4w_S)/(16w_S^2)=p_S/(L_S^2+3p_S)$; $r_S=1-2|w_S|$, $t_S=2|w_S|$, $N(S)=|S|(6-|S|)+[|S|\text{ odd}]$.\par
\textit{Value:} 1/28, -1/22, -1/6, 1/2\quad\textit{Comparison:} Exact finite identity\quad\textit{Residual:} --\par
\textit{Operation:} Computed finite identity.\par
\textit{Calculation:} \path{programs/article_results.py}; \path{outputs/article_results.out}, line 330, identifier \path{SI.P32}.\par
{\footnotesize\textit{Lean:} \path{Comparison.Formulas.rule_signed_return} (\path{lean/Comparison/Formulas/Rules.lean}:35); \path{Comparison.Formulas.sqrt_three_rule} (\path{lean/Comparison/Formulas/Rules.lean}:65); \path{Comparison.System.prop_200} (\path{lean/Comparison/System/Completion/DividedSet.lean}:85); \path{Comparison.System.S3_14} (\path{lean/Comparison/System/Finite/Completion.lean}:341)}\\[7pt]
\hypertarget{row-P33}{}\textbf{P33}\label{row:P33} & \textbf{the clock's spread is the excursion's}\par
\textit{Formula:} $-\log h_\nu/\nu^2\to\frac32\Var D$ as $\nu\to0$; $\Var D=[(\Var n_F-\Var n_L)^2+4\Cov(n_F,n_L)^2]/(3\EE N_{\rm ret})^2=61/16$.\par $\EE D=(\Var n_F+\Var n_L)/(3\EE N_{\rm ret})=2$; clock marks $(0,1,2,3,19)/4$.\par
\textit{Value:} 61/16, 2\quad\textit{Comparison:} Exact finite identity\quad\textit{Residual:} --\par
\textit{Operation:} Computed finite identity.\par
\textit{Calculation:} \path{programs/article_results.py}; \path{outputs/article_results.out}, line 331, identifier \path{SI.P33}.\par
{\footnotesize\textit{Lean:} \path{Comparison.Formulas.clock_spread_is_excursion} (\path{lean/Comparison/Formulas/Rules.lean}:92); \path{Comparison.Formulas.clock_sweep_limit} (\path{lean/Comparison/Formulas/Rules.lean}:242); \path{Comparison.System.prop_206} (\path{lean/Comparison/System/Completion/Whole.lean}:169)}\\[7pt]
\hypertarget{row-P34}{}\textbf{P34}\label{row:P34} & \textbf{the Gram-determinant term}\par
\textit{Formula:} $R_p=\sum_{e\in\mathcal A_h}\langle\Omega,(L+\gamma_h\otimes\gamma_{s(e)})\Omega\rangle-\log\det F_0/6+V_5(a)+V_{\bar h}^\prime(a)/16$.\par $F_0=(4/5)G_3$, $\det F_0=t_0=2/5$, $V_5=-a^5$, $V_{\bar h}=-a^7$.\par
\textit{Value:} 2/5\quad\textit{Comparison:} Exact finite identity\quad\textit{Residual:} --\par
\textit{Operation:} Computed finite identity.\par
\textit{Calculation:} \path{programs/article_results.py}; \path{outputs/article_results.out}, line 332, identifier \path{SI.P34}.\par
{\footnotesize\textit{Lean:} \path{Comparison.Formulas.detF0_is_transmission} (\path{lean/Comparison/Formulas/Rules.lean}:316); \path{Comparison.System.Completion.Num.protonPole_closed} (\path{lean/Comparison/System/Completion/Numbers.lean}:406); \path{Comparison.System.prop_80} (\path{lean/Comparison/System/Completion/Colour.lean}:150); \path{Comparison.System.prop_81} (\path{lean/Comparison/System/Completion/Colour.lean}:171); \path{Comparison.System.prop_82} (\path{lean/Comparison/System/Completion/Colour.lean}:193)}\\[7pt]
\hypertarget{row-P35}{}\textbf{P35}\label{row:P35} & \textbf{the neutral contrasts}\par
\textit{Formula:} $R_n=t_0^{-[\Tr((\Id-J_{10}/10)a^{2L_P})+10c_Lba^6]}$; $\Tr((\Id-J_{10}/10)a^{2L_P})=5a^4+4a^{10}$.\par
\textit{Value:} 5a\ensuremath{{}^{4}} + 4a\ensuremath{{}^{10}}\quad\textit{Comparison:} Exact finite identity\quad\textit{Residual:} --\par
\textit{Operation:} Computed finite identity.\par
\textit{Calculation:} \path{programs/article_results.py}; \path{outputs/article_results.out}, line 333, identifier \path{SI.P35}.\par
{\footnotesize\textit{Lean:} \path{Comparison.System.Completion.Num.neutronExponent_closed} (\path{lean/Comparison/System/Completion/Numbers.lean}:425); \path{Comparison.System.prop_73} (\path{lean/Comparison/System/Completion/Constants.lean}:303)}\\[7pt]
\hypertarget{row-P36}{}\textbf{P36}\label{row:P36} & \textbf{the projected nucleon}\par
\textit{Formula:} $P_N=P_{\rm sym}(15/4-\mathbf S^2)/3$; $\langle p\uparrow|\sum_r\sigma_{z,r}\tau_{z,r}|p\uparrow\rangle=5/3$; $g_A=(5/3)t_0^{2b/3}=1.275801$.\par
\textit{Value:} 5/3\quad\textit{Comparison:} Exact finite identity\quad\textit{Residual:} --\par
\textit{Operation:} Computed finite identity.\par
\textit{Calculation:} \path{programs/article_results.py}; \path{outputs/article_results.out}, line 334, identifier \path{SI.P36}.\par
{\footnotesize\textit{Lean:} \path{Comparison.Formulas.Nucleon.rank_PN} (\path{lean/Comparison/Formulas/Nucleon.lean}:412); \path{Comparison.Formulas.Nucleon.axial_Pp} (\path{lean/Comparison/Formulas/Nucleon.lean}:506); \path{Comparison.System.prop_120} (\path{lean/Comparison/System/Completion/Spectrum.lean}:167)}\\[7pt]
\hypertarget{row-P37}{}\textbf{P37}\label{row:P37} & \textbf{no strong phase}\par
\textit{Formula:} $(1+4u^2)\Id-uA_U\succeq(1-u)(1-4u)\Id>0$, $0\leq u<1/4$.\par $\mathcal Z_{\rm top}(\theta)=\sum_qZ_qe^{\ii q\theta}$, $Z_q=Z_{-q}\geq0$, $Z_1>0$; $\bar\theta=0$.\par
\textit{Value:} \ensuremath{\bar{\theta}} = 0\quad\textit{Comparison:} Exact finite identity\quad\textit{Residual:} --\par
\textit{Operation:} Computed finite identity.\par
\textit{Calculation:} \path{programs/article_results.py}; \path{outputs/article_results.out}, line 335, identifier \path{SI.P37}.\par
{\footnotesize\textit{Lean:} \path{Comparison.System.prop_115} (\path{lean/Comparison/System/Law/Sectors.lean}:84); \path{Comparison.System.S3_8} (\path{lean/Comparison/System/Finite/Law.lean}:122)}\\[7pt]
\multicolumn{2}{@{}l}{\textbf{The 113 rows of the main line}}\\*
\hypertarget{row-D01}{}\textbf{D01}\label{row:D01} & \textbf{strong coupling alpha\_\allowbreak{}s(M\_\allowbreak{}Z)}\par
\textit{Formula:} \ensuremath{\alpha_{s}} = [\ensuremath{r_{\rm v}}\ensuremath{{}^{2}}/(\ensuremath{k_{C}}\ensuremath{{}^{2}} - 1)]\ensuremath{{}^{1/2}} = 72\ensuremath{{}^{-1/2}}\par
\textit{Value:} 0.117851130197758\quad\textit{Comparison:} 0.118 +0.0009/-0.0009 [PDG 2024]\quad\textit{Residual:} -0.165\par
\textit{Operation:} computed.\par
\textit{Calculation:} \path{programs/article_results.py}; \path{outputs/article_results.out}, line 336, identifier \path{SI.D01}.\par
{\footnotesize\textit{Lean:} \path{Comparison.System.row_D01} (\path{lean/Comparison/System/Table/Couplings.lean}:93)}\\[7pt]
\hypertarget{row-D02}{}\textbf{D02}\label{row:D02} & \textbf{inverse electric coupling 1/alpha}\par
\textit{Formula:} $\alpha^{-1}=(1-x)/(c_Lf)$, $x=\nu\hover/(M+1+N_h)$.\par
\textit{Value:} 137.035999177096\quad\textit{Comparison:} 137.035999177 +2.1e-08/-2.1e-08 [CODATA 2022]\quad\textit{Residual:} +0.005\par
\textit{Operation:} computed.\par
\textit{Calculation:} \path{programs/article_results.py}; \path{outputs/article_results.out}, line 337, identifier \path{SI.D02}.\par
{\footnotesize\textit{Lean:} \path{Comparison.System.row_D02} (\path{lean/Comparison/System/Table/Couplings.lean}:100)}\\[7pt]
\hypertarget{row-L19}{}\textbf{L19}\label{row:L19} & \textbf{clock-independent lower bound on inverse alpha}\par
\textit{Formula:} $\alpha^{-1}$ evaluated at $\hover=1$.\par
\textit{Value:} 137.035999147056\quad\textit{Comparison:} bound for every clock\quad\textit{Residual:} -\par
\textit{Operation:} computed.\par
\textit{Calculation:} \path{programs/article_results.py}; \path{outputs/article_results.out}, line 338, identifier \path{SI.L19}.\par
{\footnotesize\textit{Lean:} \path{Comparison.System.row_L19} (\path{lean/Comparison/System/Table/Couplings.lean}:214)}\\[7pt]
\hypertarget{row-D03}{}\textbf{D03}\label{row:D03} & \textbf{Fermi coefficient G\_\allowbreak{}F}\par
\textit{Formula:} $G_Fm_p^2=(9/10)W_W/M^2$, $W_W=\Fret(\hover)\hover^{-4}(1-2b+2b\hover)$.\par
\textit{Value:} 1.16637885870644\quad\textit{Unit:} $10^{-5}\,\mathrm{GeV}^{-2}$\quad\textit{Comparison:} 1.1663787 +6e-07/-6e-07 [CODATA 2022]\quad\textit{Residual:} +0.265\par
\textit{Operation:} computed.\par
\textit{Calculation:} \path{programs/article_results.py}; \path{outputs/article_results.out}, line 339, identifier \path{SI.D03}.\par
{\footnotesize\textit{Lean:} \path{Comparison.System.row_D03} (\path{lean/Comparison/System/Table/Couplings.lean}:113)}\\[7pt]
\hypertarget{row-D04}{}\textbf{D04}\label{row:D04} & \textbf{gravitational coupling G m\_\allowbreak{}p\textasciicircum{}2/(hbar c)}\par
\textit{Formula:} $Gm_p^2/(\hbar c)=W^2$, $W=t_0^{3\EE N_{\rm ret}}\eta_0^{3E_H}=t_0^{48}\eta_0^9$.\par
\textit{Value:} 5.90615249377149\quad\textit{Unit:} $10^{-39}$\quad\textit{Comparison:} 5.90614943 +0.000133/-0.000133 [CODATA 2022 G]\quad\textit{Residual:} +0.023\par
\textit{Operation:} computed.\par
\textit{Calculation:} \path{programs/article_results.py}; \path{outputs/article_results.out}, line 340, identifier \path{SI.D04}.\par
{\footnotesize\textit{Lean:} \path{Comparison.System.row_D04} (\path{lean/Comparison/System/Table/Couplings.lean}:119)}\\[7pt]
\hypertarget{row-D05}{}\textbf{D05}\label{row:D05} & \textbf{Planck mass, m\_\allowbreak{}p over the square root of the gravitational coupling}\par
\textit{Formula:} \ensuremath{M_{\mathrm{Pl}}} = \ensuremath{m_{p}}/W [10\ensuremath{{}^{19}} GeV]\par
\textit{Value:} 1.22088981188542\quad\textit{Unit:} $10^{19}\,\mathrm{GeV}$\quad\textit{Comparison:} 1.22089 +1.4e-05/-1.4e-05 [CODATA 2022]\quad\textit{Residual:} -0.013\par
\textit{Operation:} computed.\par
\textit{Calculation:} \path{programs/article_results.py}; \path{outputs/article_results.out}, line 341, identifier \path{SI.D05}.\par
{\footnotesize\textit{Lean:} \path{Comparison.System.row_D05} (\path{lean/Comparison/System/Table/Couplings.lean}:128)}\\[7pt]
\hypertarget{row-J01}{}\textbf{J01}\label{row:J01} & \textbf{m\_\allowbreak{}p/m\_\allowbreak{}e}\par
\textit{Formula:} \ensuremath{m_{p}}/\ensuremath{m_{e}} = 6M + \ensuremath{C_{\ensuremath{h_\nu}}} + (1/6) log k\ensuremath{{}_{0}}(-1/4) - a\ensuremath{{}^{5}} - ba\ensuremath{{}^{6}}\par
\textit{Value:} 1836.15267341205\quad\textit{Comparison:} 1836.152673426 +3.2e-08 / -3.2e-08 [CODATA 2022]\quad\textit{Residual:} -0.435882\par
\textit{Operation:} computed.\par
\textit{Calculation:} \path{programs/article_results.py}; \path{outputs/article_results.out}, line 342, identifier \path{SI.J01}.\par
{\footnotesize\textit{Lean:} \path{Comparison.System.row_J01} (\path{lean/Comparison/System/Table/Leptons.lean}:240)}\\[7pt]
\hypertarget{row-B12}{}\textbf{B12}\label{row:B12} & \textbf{proton/electron m\_\allowbreak{}p/m\_\allowbreak{}e}\par
\textit{Formula:} \ensuremath{m_{p}}/\ensuremath{m_{e}}, as J01\par
\textit{Value:} 1836.15267341205\quad\textit{Comparison:} 1836.152673426 +3.2e-08/-3.2e-08 [CODATA 2022]\quad\textit{Residual:} -0.436\par
\textit{Operation:} computed.\par
\textit{Calculation:} \path{programs/article_results.py}; \path{outputs/article_results.out}, line 343, identifier \path{SI.B12}.\par
{\footnotesize\textit{Lean:} \path{Comparison.System.row_B12} (\path{lean/Comparison/System/Table/Leptons.lean}:221)}\\[7pt]
\hypertarget{row-M1}{}\textbf{M1}\label{row:M1} & \textbf{proton mass MeV}\par
\textit{Formula:} \ensuremath{m_{p}} = \ensuremath{m_{e}}(\ensuremath{m_{p}}/\ensuremath{m_{e}}) [MeV]\par
\textit{Value:} 938.272089420197\quad\textit{Unit:} $\mathrm{MeV}$\quad\textit{Comparison:} 938.27208943 \ensuremath{\pm} 2.9e-07 [CODATA 2022]\quad\textit{Residual:} -0.0338046\par
\textit{Operation:} computed.\par
\textit{Calculation:} \path{programs/article_results.py}; \path{outputs/article_results.out}, line 344, identifier \path{SI.M1}.\par
{\footnotesize\textit{Lean:} \path{Comparison.System.row_M1} (\path{lean/Comparison/System/Table/Nuclei.lean}:158)}\\[7pt]
\hypertarget{row-B13}{}\textbf{B13}\label{row:B13} & \textbf{neutron/proton m\_\allowbreak{}n/m\_\allowbreak{}p}\par
\textit{Formula:} \ensuremath{m_{n}}/\ensuremath{m_{p}} = k\ensuremath{{}_{0}}(-1/4)\textasciicircum{}(5a\ensuremath{{}^{4}} + 4a\ensuremath{{}^{10}} + 10\ensuremath{c_{L}} b a\ensuremath{{}^{6}})\par
\textit{Value:} 1.00137841956462\quad\textit{Comparison:} 1.00137841946 +4e-10/-4e-10 [CODATA 2022]\quad\textit{Residual:} +0.262\par
\textit{Operation:} computed.\par
\textit{Calculation:} \path{programs/article_results.py}; \path{outputs/article_results.out}, line 345, identifier \path{SI.B13}.\par
{\footnotesize\textit{Lean:} \path{Comparison.System.row_B13} (\path{lean/Comparison/System/Table/Leptons.lean}:227)}\\[7pt]
\hypertarget{row-J02}{}\textbf{J02}\label{row:J02} & \textbf{m\_\allowbreak{}n/m\_\allowbreak{}p}\par
\textit{Formula:} \ensuremath{m_{n}}/\ensuremath{m_{p}}, as B13\par
\textit{Value:} 1.00137841956462\quad\textit{Comparison:} 1.00137841946 +4e-10 / -4e-10 [CODATA 2022]\quad\textit{Residual:} +0.261549\par
\textit{Operation:} computed.\par
\textit{Calculation:} \path{programs/article_results.py}; \path{outputs/article_results.out}, line 346, identifier \path{SI.J02}.\par
{\footnotesize\textit{Lean:} \path{Comparison.System.row_J02} (\path{lean/Comparison/System/Table/Leptons.lean}:244)}\\[7pt]
\hypertarget{row-B14}{}\textbf{B14}\label{row:B14} & \textbf{neutron/electron m\_\allowbreak{}n/m\_\allowbreak{}e}\par
\textit{Formula:} \ensuremath{m_{n}}/\ensuremath{m_{e}} = (\ensuremath{m_{p}}/\ensuremath{m_{e}})(\ensuremath{m_{n}}/\ensuremath{m_{p}})\par
\textit{Value:} 1838.68366218071\quad\textit{Comparison:} 1838.683662 +7.4e-07/-7.4e-07 [CODATA 2022]\quad\textit{Residual:} +0.244\par
\textit{Operation:} computed.\par
\textit{Calculation:} \path{programs/article_results.py}; \path{outputs/article_results.out}, line 347, identifier \path{SI.B14}.\par
{\footnotesize\textit{Lean:} \path{Comparison.System.row_B14} (\path{lean/Comparison/System/Table/Leptons.lean}:234)}\\[7pt]
\hypertarget{row-M2}{}\textbf{M2}\label{row:M2} & \textbf{neutron mass MeV}\par
\textit{Formula:} \ensuremath{m_{n}} = \ensuremath{m_{p}}(\ensuremath{m_{n}}/\ensuremath{m_{p}}) [MeV]\par
\textit{Value:} 939.56542202519\quad\textit{Unit:} $\mathrm{MeV}$\quad\textit{Comparison:} 939.56542194 \ensuremath{\pm} 4.8e-07 [CODATA 2022]\quad\textit{Residual:} +0.177479\par
\textit{Operation:} computed.\par
\textit{Calculation:} \path{programs/article_results.py}; \path{outputs/article_results.out}, line 348, identifier \path{SI.M2}.\par
{\footnotesize\textit{Lean:} \path{Comparison.System.row_M2} (\path{lean/Comparison/System/Table/Nuclei.lean}:162)}\\[7pt]
\hypertarget{row-B01}{}\textbf{B01}\label{row:B01} & \textbf{muon/electron}\par
\textit{Formula:} \ensuremath{m_{\mu}}/\ensuremath{m_{e}} = \ensuremath{\mathcal L}(10, 2, +) = 3\ensuremath{{}^{5}}e\ensuremath{{}^{-9/56}}\ensuremath{M_{Q}}(\ensuremath{\nu})\ensuremath{{}^{1/6}}\par
\textit{Value:} 206.768283945476\quad\textit{Comparison:} 206.7682827 +4.6e-06/-4.6e-06 [CODATA 2022]\quad\textit{Residual:} +0.271\par
\textit{Operation:} computed.\par
\textit{Calculation:} \path{programs/article_results.py}; \path{outputs/article_results.out}, line 349, identifier \path{SI.B01}.\par
{\footnotesize\textit{Lean:} \path{Comparison.System.row_B01} (\path{lean/Comparison/System/Table/Leptons.lean}:146)}\\[7pt]
\hypertarget{row-B02}{}\textbf{B02}\label{row:B02} & \textbf{tau/muon}\par
\textit{Formula:} \ensuremath{m_{\tau}}/\ensuremath{m_{\mu}} = \ensuremath{\mathcal L}(5, -2, +) = 3\ensuremath{{}^{5/2}}e\textasciicircum{}(1/14 + 2\ensuremath{\nu})\par
\textit{Value:} 16.8181606410583\quad\textit{Comparison:} 16.817691845 +0.000852/-0.000852 [PDG 2024 m\_\allowbreak{}tau]\quad\textit{Residual:} +0.550\par
\textit{Operation:} computed.\par
\textit{Calculation:} \path{programs/article_results.py}; \path{outputs/article_results.out}, line 350, identifier \path{SI.B02}.\par
{\footnotesize\textit{Lean:} \path{Comparison.System.row_B02} (\path{lean/Comparison/System/Table/Leptons.lean}:152)}\\[7pt]
\hypertarget{row-L01}{}\textbf{L01}\label{row:L01} & \textbf{tau pole mass}\par
\textit{Formula:} \ensuremath{m_{\tau}} = \ensuremath{m_{e}}(\ensuremath{m_{\mu}}/\ensuremath{m_{e}})(\ensuremath{m_{\tau}}/\ensuremath{m_{\mu}}) [MeV]\par
\textit{Value:} 1776.9795428632\quad\textit{Unit:} $\mathrm{MeV}$\quad\textit{Comparison:} from B01 and B02\quad\textit{Residual:} -\par
\textit{Operation:} computed.\par
\textit{Calculation:} \path{programs/article_results.py}; \path{outputs/article_results.out}, line 351, identifier \path{SI.L01}.\par
{\footnotesize\textit{Lean:} \path{Comparison.System.row_L01} (\path{lean/Comparison/System/Table/Leptons.lean}:248)}\\[7pt]
\hypertarget{row-B03}{}\textbf{B03}\label{row:B03} & \textbf{Koide statistic}\par
\textit{Formula:} (1 + \ensuremath{r_{\mu}} + \ensuremath{r_{\mu}}\ensuremath{r_{\tau}})/(1 + \ensuremath{r_{\mu}}\ensuremath{{}^{1/2}} + (\ensuremath{r_{\mu}}\ensuremath{r_{\tau}})\ensuremath{{}^{1/2}})\ensuremath{{}^{2}}, \ensuremath{r_{\mu}} = \ensuremath{m_{\mu}}/\ensuremath{m_{e}}, \ensuremath{r_{\tau}} = \ensuremath{m_{\tau}}/\ensuremath{m_{\mu}}\par
\textit{Value:} 0.666667259761091\quad\textit{Comparison:} 0.6666645 +6.8e-06/-6.8e-06 [PDG masses]\quad\textit{Residual:} +0.406\par
\textit{Operation:} computed.\par
\textit{Calculation:} \path{programs/article_results.py}; \path{outputs/article_results.out}, line 352, identifier \path{SI.B03}.\par
{\footnotesize\textit{Lean:} \path{Comparison.System.row_B03} (\path{lean/Comparison/System/Table/Leptons.lean}:159)}\\[7pt]
\hypertarget{row-L17}{}\textbf{L17}\label{row:L17} & \textbf{Koide discrimination}\par
\textit{Formula:} the same Koide statistic\par
\textit{Value:} 0.666667259761091\quad\textit{Comparison:} from B03\quad\textit{Residual:} -\par
\textit{Operation:} computed.\par
\textit{Calculation:} \path{programs/article_results.py}; \path{outputs/article_results.out}, line 353, identifier \path{SI.L17}.\par
{\footnotesize\textit{Lean:} \path{Comparison.System.row_L17} (\path{lean/Comparison/System/Table/Leptons.lean}:267)}\\[7pt]
\hypertarget{row-L04}{}\textbf{L04}\label{row:L04} & \textbf{number of charged families}\par
\textit{Formula:} three charged families: the three pairings of the four directions other than \ensuremath{h_\nu}\par
\textit{Value:} 3\quad\textit{Comparison:} prediction: no fourth family\quad\textit{Residual:} -\par
\textit{Operation:} computed.\par
\textit{Calculation:} \path{programs/article_results.py}; \path{outputs/article_results.out}, line 354, identifier \path{SI.L04}.\par
{\footnotesize\textit{Lean:} \path{Comparison.System.row_L04} (\path{lean/Comparison/System/Table/Neutrinos.lean}:111); \path{Comparison.System.row_L04_pairings} (\path{lean/Comparison/System/Completion/Families.lean}:183)}\\[7pt]
\hypertarget{row-L05}{}\textbf{L05}\label{row:L05} & \textbf{forbidden charged-lepton threshold}\par
\textit{Formula:} 3\ensuremath{{}^{5/2}}\ensuremath{m_{e}}, with no rest pole [MeV]\par
\textit{Value:} 7.96568530688517\quad\textit{Unit:} $\mathrm{MeV}$\quad\textit{Comparison:} prediction: no charged lepton at this threshold\quad\textit{Residual:} -\par
\textit{Operation:} computed.\par
\textit{Calculation:} \path{programs/article_results.py}; \path{outputs/article_results.out}, line 355, identifier \path{SI.L05}.\par
{\footnotesize\textit{Lean:} \path{Comparison.System.row_L05} (\path{lean/Comparison/System/Table/Leptons.lean}:256)}\\[7pt]
\hypertarget{row-B04}{}\textbf{B04}\label{row:B04} & \textbf{strange/down}\par
\textit{Formula:} $m_s^{\rm native}/m_d^{\rm native}=p_q=\mathcal L(10,2,-)$.\par
\textit{Value:} 19.9616523937784\quad\textit{Comparison:} 19.946 +0.335/-0.335 [FLAG 2021 N\_\allowbreak{}f = 2+1+1, m\_\allowbreak{}s/m\_\allowbreak{}ud = 27.23(10) times (1 + m\_\allowbreak{}u/m\_\allowbreak{}d)/2 with m\_\allowbreak{}u/m\_\allowbreak{}d = 0.465(24)]\quad\textit{Residual:} +0.047\par
\textit{Operation:} computed.\par
\textit{Calculation:} \path{programs/article_results.py}; \path{outputs/article_results.out}, line 356, identifier \path{SI.B04}.\par
{\footnotesize\textit{Lean:} \path{Comparison.System.row_B04} (\path{lean/Comparison/System/Table/Leptons.lean}:166)}\\[7pt]
\hypertarget{row-B05}{}\textbf{B05}\label{row:B05} & \textbf{bottom/strange}\par
\textit{Formula:} $m_b^{\rm native}/m_s^{\rm native}=q_q\eta_0$, $q_q=\mathcal L(5,-2,-)$.\par
\textit{Value:} 53.9210868808146\quad\textit{Comparison:} 53.88 +0.12/-0.12 [PDG 2026]\quad\textit{Residual:} +0.34\par
\textit{Operation:} computed.\par
\textit{Calculation:} \path{programs/article_results.py}; \path{outputs/article_results.out}, line 357, identifier \path{SI.B05}.\par
{\footnotesize\textit{Lean:} \path{Comparison.System.row_B05} (\path{lean/Comparison/System/Table/Leptons.lean}:172)}\\[7pt]
\hypertarget{row-B06}{}\textbf{B06}\label{row:B06} & \textbf{down statistic m\_\allowbreak{}b\textasciicircum{}2 m\_\allowbreak{}d/m\_\allowbreak{}s\textasciicircum{}3}\par
\textit{Formula:} $(m_b^{\rm native})^2m_d^{\rm native}/(m_s^{\rm native})^3=q_q^2\eta_0^2/p_q$.\par
\textit{Value:} 145.653453584562\quad\textit{Comparison:} 145.45 +2.63/-2.63 [(m\_\allowbreak{}b/m\_\allowbreak{}s)\textasciicircum{}2 (m\_\allowbreak{}d/m\_\allowbreak{}s) of the two FLAG ratios above]\quad\textit{Residual:} +0.08\par
\textit{Operation:} computed.\par
\textit{Calculation:} \path{programs/article_results.py}; \path{outputs/article_results.out}, line 358, identifier \path{SI.B06}.\par
{\footnotesize\textit{Lean:} \path{Comparison.System.row_B06} (\path{lean/Comparison/System/Table/Leptons.lean}:178)}\\[7pt]
\hypertarget{row-B07}{}\textbf{B07}\label{row:B07} & \textbf{up statistic m\_\allowbreak{}t\textasciicircum{}2 m\_\allowbreak{}u/m\_\allowbreak{}c\textasciicircum{}3}\par
\textit{Formula:} $(m_t^{\rm native})^2m_u^{\rm native}/(m_c^{\rm native})^3=q_q^2\eta_0^2/p_q$.\par
\textit{Value:} 145.653453584562\quad\textit{Comparison:} 144.156 +6.06/-5.72 [PDG 2026 masses, four-loop running to m\_\allowbreak{}t(m\_\allowbreak{}t)]\quad\textit{Residual:} +0.25\par
\textit{Operation:} computed.\par
\textit{Calculation:} \path{programs/article_results.py}; \path{outputs/article_results.out}, line 359, identifier \path{SI.B07}.\par
{\footnotesize\textit{Lean:} \path{Comparison.System.row_B07} (\path{lean/Comparison/System/Table/Leptons.lean}:188)}\\[7pt]
\hypertarget{row-B08}{}\textbf{B08}\label{row:B08} & \textbf{lepton/quark symmetric-state quotient}\par
\textit{Formula:} (\ensuremath{m_{\mu}}/\ensuremath{m_{e}})/(\ensuremath{m_{s}}/\ensuremath{m_{d}})\par
\textit{Value:} 10.3582749497192\quad\textit{Comparison:} 10.3664 +0.174/-0.174 [CODATA m\_\allowbreak{}mu/m\_\allowbreak{}e over the FLAG strange/down ratio]\quad\textit{Residual:} -0.047\par
\textit{Operation:} computed.\par
\textit{Calculation:} \path{programs/article_results.py}; \path{outputs/article_results.out}, line 360, identifier \path{SI.B08}.\par
{\footnotesize\textit{Lean:} \path{Comparison.System.row_B08} (\path{lean/Comparison/System/Table/Leptons.lean}:194)}\\[7pt]
\hypertarget{row-B09}{}\textbf{B09}\label{row:B09} & \textbf{up/down m\_\allowbreak{}u/m\_\allowbreak{}d}\par
\textit{Formula:} \ensuremath{m_{u}}/\ensuremath{m_{d}} = \ensuremath{R_{\mathrm{ud}}} = 34/74\par
\textit{Value:} 0.459459459459459\quad\textit{Comparison:} 0.4595745 +0.01/-0.01 [PDG 2.16/4.70]\quad\textit{Residual:} -0.012\par
\textit{Operation:} computed.\par
\textit{Calculation:} \path{programs/article_results.py}; \path{outputs/article_results.out}, line 361, identifier \path{SI.B09}.\par
{\footnotesize\textit{Lean:} \path{Comparison.System.row_B09} (\path{lean/Comparison/System/Table/Leptons.lean}:200)}\\[7pt]
\hypertarget{row-B10}{}\textbf{B10}\label{row:B10} & \textbf{up factor per visit to a symmetric state, kappa}\par
\textit{Formula:} $\kappa_u=(r_0/t_0)^{\EE N_{\rm ret}/\EE n_F}=(3/2)^{8/5}$\par
\textit{Value:} 1.91313675093869\quad\textit{Comparison:} 1.9131 +0.0062/-0.0062 [up-sector masses]\quad\textit{Residual:} +0.006\par
\textit{Operation:} computed.\par
\textit{Calculation:} \path{programs/article_results.py}; \path{outputs/article_results.out}, line 362, identifier \path{SI.B10}.\par
{\footnotesize\textit{Lean:} \path{Comparison.System.row_B10} (\path{lean/Comparison/System/Table/Leptons.lean}:207)}\\[7pt]
\hypertarget{row-B11}{}\textbf{B11}\label{row:B11} & \textbf{charm/strange}\par
\textit{Formula:} \ensuremath{m_{c}}/\ensuremath{m_{s}} = \ensuremath{R_{\mathrm{ud}}} \ensuremath{\kappa_u}\ensuremath{{}^{5}}\par
\textit{Value:} 11.7754434121621\quad\textit{Comparison:} 11.761 +0.028/-0.028 [PDG 2026 m\_\allowbreak{}c/m\_\allowbreak{}s]\quad\textit{Residual:} +0.516\par
\textit{Operation:} computed.\par
\textit{Calculation:} \path{programs/article_results.py}; \path{outputs/article_results.out}, line 363, identifier \path{SI.B11}.\par
{\footnotesize\textit{Lean:} \path{Comparison.System.row_B11} (\path{lean/Comparison/System/Table/Leptons.lean}:213)}\\[7pt]
\hypertarget{row-D15}{}\textbf{D15}\label{row:D15} & \textbf{electroweak scale v = (sqrt2 G\_\allowbreak{}F)\textasciicircum{}-1/2}\par
\textit{Formula:} v = (\ensuremath{\surd}2 \ensuremath{G_{F}})\ensuremath{{}^{-1/2}} [GeV]\par
\textit{Value:} 246.219634042871\quad\textit{Unit:} $\mathrm{GeV}$\quad\textit{Comparison:} 246.219650794 +6.33293e-05/-6.33293e-05 [$(\sqrt2G_F)^{-1/2}$ from CODATA 2022 $G_F=1.1663787(6)\times10^{-5}\,\mathrm{GeV}^{-2}$, the same record as D03]\quad\textit{Residual:} -0.264508\par
\textit{Operation:} computed.\par
\textit{Calculation:} \path{programs/article_results.py}; \path{outputs/article_results.out}, line 364, identifier \path{SI.D15}.\par
{\footnotesize\textit{Lean:} \path{Comparison.System.row_D15} (\path{lean/Comparison/System/Table/Couplings.lean}:185)}\\[7pt]
\hypertarget{row-D08}{}\textbf{D08}\label{row:D08} & \textbf{MSbar weak angle 3 N\_\allowbreak{}h}\par
\textit{Formula:} $\widehat s^{\,2}=3N_h$, with $N_h=\Re\sum_{j\ge1}\hover^j p_j$ in Eq.~\eqref{eq:Nh}.\par
\textit{Value:} 0.231257521689422\quad\textit{Comparison:} 0.23122 +6e-05/-6e-05 [PDG 2026, top decoupled]\quad\textit{Residual:} +0.625\par
\textit{Operation:} computed.\par
\textit{Calculation:} \path{programs/article_results.py}; \path{outputs/article_results.out}, line 365, identifier \path{SI.D08}.\par
{\footnotesize\textit{Lean:} \path{Comparison.System.row_D08} (\path{lean/Comparison/System/Table/Couplings.lean}:137)}\\[7pt]
\hypertarget{row-D10}{}\textbf{D10}\label{row:D10} & \textbf{on-shell weak angle 3 N\_\allowbreak{}h/(1 + f)}\par
\textit{Formula:} s\ensuremath{{}^{2}}\_\allowbreak{}OS = 3\ensuremath{N_{\ensuremath{h_\nu}}}/(1 + f) (on shell)\par
\textit{Value:} 0.223221546032261\quad\textit{Comparison:} 0.2233373 +0.000153/-0.000153 [direct M\_\allowbreak{}W/M\_\allowbreak{}Z]\quad\textit{Residual:} -0.758\par
\textit{Operation:} computed.\par
\textit{Calculation:} \path{programs/article_results.py}; \path{outputs/article_results.out}, line 366, identifier \path{SI.D10}.\par
{\footnotesize\textit{Lean:} \path{Comparison.System.row_D10} (\path{lean/Comparison/System/Table/Couplings.lean}:149)}\\[7pt]
\hypertarget{row-L20}{}\textbf{L20}\label{row:L20} & \textbf{clock-independent upper bound on the MS-bar weak angle}\par
\textit{Formula:} $3N_h$ at $\hover=1$.\par
\textit{Value:} 0.231264231357674\quad\textit{Comparison:} bound for every clock\quad\textit{Residual:} -\par
\textit{Operation:} computed.\par
\textit{Calculation:} \path{programs/article_results.py}; \path{outputs/article_results.out}, line 367, identifier \path{SI.L20}.\par
{\footnotesize\textit{Lean:} \path{Comparison.System.row_L20} (\path{lean/Comparison/System/Table/Couplings.lean}:225)}\\[7pt]
\hypertarget{row-D09}{}\textbf{D09}\label{row:D09} & \textbf{effective weak angle 3 N\_\allowbreak{}con}\par
\textit{Formula:} 3\ensuremath{N_{\mathrm{con}}} (effective)\par
\textit{Value:} 0.231550233937853\quad\textit{Comparison:} 0.23153 +0.00016/-0.00016 [LEP/SLD, Phys. Rep. 427, 257 (2006)]\quad\textit{Residual:} +0.126\par
\textit{Operation:} computed.\par
\textit{Calculation:} \path{programs/article_results.py}; \path{outputs/article_results.out}, line 368, identifier \path{SI.D09}.\par
{\footnotesize\textit{Lean:} \path{Comparison.System.row_D09} (\path{lean/Comparison/System/Table/Couplings.lean}:143)}\\[7pt]
\hypertarget{row-D14}{}\textbf{D14}\label{row:D14} & \textbf{weak correction delta\_\allowbreak{}inv = f}\par
\textit{Formula:} f = \ensuremath{r_0}\ensuremath{{}^{2}}/N\ensuremath{{}_{2}} = 9/250\par
\textit{Value:} 0.036\quad\textit{Comparison:} 0.03641 +0.00046/-0.00046 [1 - pi alpha/(sqrt2 G\_\allowbreak{}F M\_\allowbreak{}W\textasciicircum{}2 s2\_\allowbreak{}OS) of the records of 1/alpha, G\_\allowbreak{}F, M\_\allowbreak{}W and M\_\allowbreak{}Z]\quad\textit{Residual:} -0.891\par
\textit{Operation:} computed.\par
\textit{Calculation:} \path{programs/article_results.py}; \path{outputs/article_results.out}, line 369, identifier \path{SI.D14}.\par
{\footnotesize\textit{Lean:} \path{Comparison.System.row_D14} (\path{lean/Comparison/System/Table/Couplings.lean}:182)}\\[7pt]
\hypertarget{row-D11}{}\textbf{D11}\label{row:D11} & \textbf{W mass}\par
\textit{Formula:} \ensuremath{M_{W}} = [\ensuremath{\pi}\ensuremath{\alpha}v\ensuremath{{}^{2}}/((1 - f)s\ensuremath{{}^{2}}\_\allowbreak{}OS)]\ensuremath{{}^{1/2}} [GeV]\par
\textit{Value:} 80.3662510403325\quad\textit{Unit:} $\mathrm{GeV}$\quad\textit{Comparison:} 80.3625 +0.0077/-0.0077 [PDG 2026 evaluation without the CDF II measurement]\quad\textit{Residual:} +0.487\par
\textit{Operation:} computed.\par
\textit{Calculation:} \path{programs/article_results.py}; \path{outputs/article_results.out}, line 370, identifier \path{SI.D11}.\par
{\footnotesize\textit{Lean:} \path{Comparison.System.row_D11} (\path{lean/Comparison/System/Table/Couplings.lean}:155)}\\[7pt]
\hypertarget{row-L15}{}\textbf{L15}\label{row:L15} & \textbf{W pole mass test}\par
\textit{Formula:} \ensuremath{M_{W}}, as D11 [GeV]\par
\textit{Value:} 80.3662510403325\quad\textit{Unit:} $\mathrm{GeV}$\quad\textit{Comparison:} 80.3570(60) GeV [PDG 2026 SM indirect determination]\quad\textit{Residual:} -\par
\textit{Operation:} computed.\par
\textit{Calculation:} \path{programs/article_results.py}; \path{outputs/article_results.out}, line 371, identifier \path{SI.L15}.\par
{\footnotesize\textit{Lean:} \path{Comparison.System.row_L15} (\path{lean/Comparison/System/Table/Couplings.lean}:209)}\\[7pt]
\hypertarget{row-D12}{}\textbf{D12}\label{row:D12} & \textbf{Z mass}\par
\textit{Formula:} \ensuremath{M_{Z}} = \ensuremath{M_{W}}/(1 - s\ensuremath{{}^{2}}\_\allowbreak{}OS)\ensuremath{{}^{1/2}} [GeV]\par
\textit{Value:} 91.1853615183457\quad\textit{Unit:} $\mathrm{GeV}$\quad\textit{Comparison:} 91.1879 +0.002/-0.002 [PDG 2026 world average (LEP 1, CDF, LHCb), running-width convention]\quad\textit{Residual:} -1.269\par
\textit{Operation:} computed.\par
\textit{Calculation:} \path{programs/article_results.py}; \path{outputs/article_results.out}, line 372, identifier \path{SI.D12}.\par
{\footnotesize\textit{Lean:} \path{Comparison.System.row_D12} (\path{lean/Comparison/System/Table/Couplings.lean}:161)}\\[7pt]
\hypertarget{row-D13}{}\textbf{D13}\label{row:D13} & \textbf{W/Z mass ratio}\par
\textit{Formula:} \ensuremath{M_{W}}/\ensuremath{M_{Z}} = (1 - s\ensuremath{{}^{2}}\_\allowbreak{}OS)\ensuremath{{}^{1/2}}\par
\textit{Value:} 0.881350358238844\quad\textit{Comparison:} 0.88128469 +8.66e-05/-8.66e-05 [direct masses]\quad\textit{Residual:} +0.758\par
\textit{Operation:} computed.\par
\textit{Calculation:} \path{programs/article_results.py}; \path{outputs/article_results.out}, line 373, identifier \path{SI.D13}.\par
{\footnotesize\textit{Lean:} \path{Comparison.System.row_D13} (\path{lean/Comparison/System/Table/Couplings.lean}:167)}\\[7pt]
\hypertarget{row-D16}{}\textbf{D16}\label{row:D16} & \textbf{scalar mass m\_\allowbreak{}H}\par
\textit{Formula:} \ensuremath{m_{H}} = v(2\ensuremath{\pi}\ensuremath{\alpha}/\ensuremath{\chi_{\ensuremath{h_\nu}}})\ensuremath{{}^{1/2}}, \ensuremath{\chi_{\ensuremath{h_\nu}}} = 30/169 [GeV]\par
\textit{Value:} 125.134771593402\quad\textit{Unit:} $\mathrm{GeV}$\quad\textit{Comparison:} 125.13 +0.11/-0.11 [PDG 2026 average, error with scale factor S = 1.5]\quad\textit{Residual:} +0.043\par
\textit{Operation:} computed.\par
\textit{Calculation:} \path{programs/article_results.py}; \path{outputs/article_results.out}, line 374, identifier \path{SI.D16}.\par
{\footnotesize\textit{Lean:} \path{Comparison.System.row_D16} (\path{lean/Comparison/System/Table/Couplings.lean}:191)}\\[7pt]
\hypertarget{row-D17}{}\textbf{D17}\label{row:D17} & \textbf{Top radial response}\par
\textit{Formula:} $m_t^{\rm rad}=v\sqrt{\eta_0/(2\sigma)}$ [GeV]\par
\textit{Value:} 172.646613614442\quad\textit{Unit:} $\mathrm{GeV}$\quad\textit{Comparison:} 172.6 +0.27/-0.27 [PDG 2026 average of direct measurements, error with scale factor S = 1.5]\quad\textit{Residual:} +0.173\par
\textit{Operation:} computed.\par
\textit{Calculation:} \path{programs/article_results.py}; \path{outputs/article_results.out}, line 375, identifier \path{SI.D17}.\par
{\footnotesize\textit{Lean:} \path{Comparison.System.row_D17} (\path{lean/Comparison/System/Table/Couplings.lean}:199)}\\[7pt]
\hypertarget{row-S07}{}\textbf{S07}\label{row:S07} & \textbf{top mass, second construction}\par
\textit{Formula:} \ensuremath{m_{p}}\ensuremath{{}^{2}}/(E \ensuremath{n_{F}} \ensuremath{m_{e}}) [GeV]\par
\textit{Value:} 172.281080537681\quad\textit{Unit:} $\mathrm{GeV}$\quad\textit{Comparison:} 172.57 +0.29/-0.29 [PDG 2024 direct]\quad\textit{Residual:} -0.996\par
\textit{Operation:} Second construction; its factor was chosen with the measured value available.\par
\textit{Calculation:} \path{programs/article_results.py}; \path{outputs/article_results.out}, line 376, identifier \path{SI.S07}.\par
{\footnotesize\textit{Lean:} \path{Comparison.System.row_S07} (\path{lean/Comparison/System/Table/Couplings.lean}:231)}\\[7pt]
\hypertarget{row-C01}{}\textbf{C01}\label{row:C01} & \textbf{first quark mixing s12}\par
\textit{Formula:} s\ensuremath{{}_{12}} = (Var \ensuremath{n_{F}} - Cov\ensuremath{{}^{2}}/Var \ensuremath{n_{L}})/(E\ensuremath{N_{\rm ret}}\ensuremath{\cdot}5/2) = 9/40\par
\textit{Value:} 0.225\quad\textit{Comparison:} 0.22517 +0.00068/-0.00068 [PDG 2026 CKM global fit]\quad\textit{Residual:} -0.250\par
\textit{Operation:} computed.\par
\textit{Calculation:} \path{programs/article_results.py}; \path{outputs/article_results.out}, line 377, identifier \path{SI.C01}.\par
{\footnotesize\textit{Lean:} \path{Comparison.System.row_C01} (\path{lean/Comparison/System/Table/Mixings.lean}:101)}\\[7pt]
\hypertarget{row-C02}{}\textbf{C02}\label{row:C02} & \textbf{second quark mixing s23}\par
\textit{Formula:} s\ensuremath{{}_{23}} = 4/Tr L\ensuremath{{}^{2}} = 1/24\par
\textit{Value:} 0.0416666666666667\quad\textit{Comparison:} 0.04189 +0.00081/-0.00069 [PDG 2026 CKM fit]\quad\textit{Residual:} -0.324\par
\textit{Operation:} computed.\par
\textit{Calculation:} \path{programs/article_results.py}; \path{outputs/article_results.out}, line 378, identifier \path{SI.C02}.\par
{\footnotesize\textit{Lean:} \path{Comparison.System.row_C02} (\path{lean/Comparison/System/Table/Mixings.lean}:104)}\\[7pt]
\hypertarget{row-C03}{}\textbf{C03}\label{row:C03} & \textbf{far quark mixing s13}\par
\textit{Formula:} s\ensuremath{{}_{13}} = s\ensuremath{{}_{12}} \ensuremath{t_0} s\ensuremath{{}_{23}} = 3/800\par
\textit{Value:} 0.00375\quad\textit{Comparison:} 0.003763 +8.8e-05/-8.3e-05 [PDG 2026 CKM fit]\quad\textit{Residual:} -0.157\par
\textit{Operation:} computed.\par
\textit{Calculation:} \path{programs/article_results.py}; \path{outputs/article_results.out}, line 379, identifier \path{SI.C03}.\par
{\footnotesize\textit{Lean:} \path{Comparison.System.row_C03} (\path{lean/Comparison/System/Table/Mixings.lean}:107)}\\[7pt]
\hypertarget{row-C04}{}\textbf{C04}\label{row:C04} & \textbf{quark phase delta\_\allowbreak{}q}\par
\textit{Formula:} \ensuremath{\delta_{q}} = arg(b + i) = arg(7 + 16i) [rad]\par
\textit{Value:} 1.15838588519751\quad\textit{Unit:} $\mathrm{rad}$\quad\textit{Comparison:} 1.154 +0.025/-0.025 [PDG 2026 CKM fit]\quad\textit{Residual:} +0.175\par
\textit{Operation:} computed.\par
\textit{Calculation:} \path{programs/article_results.py}; \path{outputs/article_results.out}, line 380, identifier \path{SI.C04}.\par
{\footnotesize\textit{Lean:} \path{Comparison.System.row_C04} (\path{lean/Comparison/System/Table/Mixings.lean}:111)}\\[7pt]
\hypertarget{row-C05}{}\textbf{C05}\label{row:C05} & \textbf{quark quartet J\_\allowbreak{}q}\par
\textit{Formula:} \ensuremath{J_{q}} of V = R\ensuremath{{}_{23}}(s\ensuremath{{}_{23}})R\ensuremath{{}_{13}}(s\ensuremath{{}_{13}}, \ensuremath{\delta_{q}})R\ensuremath{{}_{12}}(s\ensuremath{{}_{12}}) [\ensuremath{\times}10\ensuremath{{}^{-5}}]\par
\textit{Value:} 3.13550917578981\quad\textit{Unit:} $10^{-5}$\quad\textit{Comparison:} 3.16 +0.13/-0.11 [PDG]\quad\textit{Residual:} -0.223\par
\textit{Operation:} computed.\par
\textit{Calculation:} \path{programs/article_results.py}; \path{outputs/article_results.out}, line 381, identifier \path{SI.C05}.\par
{\footnotesize\textit{Lean:} \path{Comparison.System.row_C05} (\path{lean/Comparison/System/Table/Mixings.lean}:117)}\\[7pt]
\hypertarget{row-C06}{}\textbf{C06}\label{row:C06} & \textbf{triangle angle alpha}\par
\textit{Formula:} the angle \ensuremath{\alpha} of the unitarity triangle of V [deg]\par
\textit{Value:} 90.713773064338\quad\textit{Unit:} $\mathrm{deg}$\quad\textit{Comparison:} 91.04 +1.61/-1.6 [PDG-derived]\quad\textit{Residual:} -0.204\par
\textit{Operation:} computed.\par
\textit{Calculation:} \path{programs/article_results.py}; \path{outputs/article_results.out}, line 382, identifier \path{SI.C06}.\par
{\footnotesize\textit{Lean:} \path{Comparison.System.row_C06} (\path{lean/Comparison/System/Table/Mixings.lean}:121)}\\[7pt]
\hypertarget{row-C07}{}\textbf{C07}\label{row:C07} & \textbf{triangle angle beta}\par
\textit{Formula:} the angle \ensuremath{\beta} of the unitarity triangle of V [deg]\par
\textit{Value:} 22.9511439353888\quad\textit{Unit:} $\mathrm{deg}$\quad\textit{Comparison:} 22.86 +0.71/-0.7 [PDG-derived]\quad\textit{Residual:} +0.128\par
\textit{Operation:} computed.\par
\textit{Calculation:} \path{programs/article_results.py}; \path{outputs/article_results.out}, line 383, identifier \path{SI.C07}.\par
{\footnotesize\textit{Lean:} \path{Comparison.System.row_C07} (\path{lean/Comparison/System/Table/Mixings.lean}:125)}\\[7pt]
\hypertarget{row-C08}{}\textbf{C08}\label{row:C08} & \textbf{triangle angle gamma}\par
\textit{Formula:} the angle \ensuremath{\gamma} of the unitarity triangle of V [deg]\par
\textit{Value:} 66.3350830002732\quad\textit{Unit:} $\mathrm{deg}$\quad\textit{Comparison:} 66.09 +1.43/-1.43 [PDG-derived]\quad\textit{Residual:} +0.171\par
\textit{Operation:} computed.\par
\textit{Calculation:} \path{programs/article_results.py}; \path{outputs/article_results.out}, line 384, identifier \path{SI.C08}.\par
{\footnotesize\textit{Lean:} \path{Comparison.System.row_C08} (\path{lean/Comparison/System/Table/Mixings.lean}:129)}\\[7pt]
\hypertarget{row-C09}{}\textbf{C09}\label{row:C09} & \textbf{two-row W current sum}\par
\textit{Formula:} \ensuremath{\Sigma}\ensuremath{|}V\ensuremath{|}\ensuremath{{}^{2}} over the first two rows\par
\textit{Value:} 2\quad\textit{Comparison:} 1.984 +0.021/-0.021 [CMS 2022, Phys. Rev. D 105, 072008]\quad\textit{Residual:} +0.762\par
\textit{Operation:} computed.\par
\textit{Calculation:} \path{programs/article_results.py}; \path{outputs/article_results.out}, line 385, identifier \path{SI.C09}.\par
{\footnotesize\textit{Lean:} \path{Comparison.System.row_C09} (\path{lean/Comparison/System/Table/Mixings.lean}:133)}\\[7pt]
\hypertarget{row-K05}{}\textbf{K05}\label{row:K05} & \textbf{diagonal quark current}\par
\textit{Formula:} \ensuremath{\Sigma_{i}} \ensuremath{|}\ensuremath{V_{\mathrm{ii}}}\ensuremath{|}\ensuremath{{}^{2}}\par
\textit{Value:} 2.89531084487576\quad\textit{Comparison:} >= 2.84295 [direct current rows]\quad\textit{Residual:} -\par
\textit{Operation:} computed.\par
\textit{Calculation:} \path{programs/article_results.py}; \path{outputs/article_results.out}, line 386, identifier \path{SI.K05}.\par
{\footnotesize\textit{Lean:} \path{Comparison.System.row_K05} (\path{lean/Comparison/System/Table/Mixings.lean}:170)}\\[7pt]
\hypertarget{row-C10}{}\textbf{C10}\label{row:C10} & \textbf{solar lepton angle s\textasciicircum{}2\_\allowbreak{}12}\par
\textit{Formula:} s\ensuremath{{}^{2}}\ensuremath{{}_{12}} of U = R\ensuremath{{}_{12}}(\ensuremath{\surd}(2b)s\ensuremath{{}_{12}}, \ensuremath{\pi}/2 - \ensuremath{\delta_{q}})\ensuremath{\dagger}R\ensuremath{{}_{23}}(1/\ensuremath{\surd}2)R\ensuremath{{}_{12}}(\ensuremath{\surd}b)\par
\textit{Value:} 0.305089610922635\quad\textit{Comparison:} 0.3036 +0.0064/-0.0064 [JUNO 2026, Ref.~\cite{JUNO2026}]\quad\textit{Residual:} +0.233\par
\textit{Operation:} computed.\par
\textit{Calculation:} \path{programs/article_results.py}; \path{outputs/article_results.out}, line 387, identifier \path{SI.C10}.\par
{\footnotesize\textit{Lean:} \path{Comparison.System.row_C10} (\path{lean/Comparison/System/Table/Mixings.lean}:137)}\\[7pt]
\hypertarget{row-C11}{}\textbf{C11}\label{row:C11} & \textbf{reactor angle s\textasciicircum{}2\_\allowbreak{}13}\par
\textit{Formula:} s\ensuremath{{}^{2}}\ensuremath{{}_{13}} = \ensuremath{|}\ensuremath{U_{e3}}\ensuremath{|}\ensuremath{{}^{2}} = b s\ensuremath{{}_{12}}\ensuremath{{}^{2}}\par
\textit{Value:} 0.0221484375\quad\textit{Comparison:} 0.02215 +0.00056/-0.00058 [NuFIT 6.0, normal ordering, IC24 with SK atmospheric data, Ref.~\cite{NuFIT6}]\quad\textit{Residual:} -0.003\par
\textit{Operation:} computed.\par
\textit{Calculation:} \path{programs/article_results.py}; \path{outputs/article_results.out}, line 388, identifier \path{SI.C11}.\par
{\footnotesize\textit{Lean:} \path{Comparison.System.row_C11} (\path{lean/Comparison/System/Table/Mixings.lean}:141)}\\[7pt]
\hypertarget{row-C12}{}\textbf{C12}\label{row:C12} & \textbf{reactor angle sin\textasciicircum{}2 2theta\_\allowbreak{}13}\par
\textit{Formula:} sin\ensuremath{{}^{2}}2\ensuremath{\theta}\ensuremath{{}_{13}} = 4s\ensuremath{{}^{2}}\ensuremath{{}_{13}}(1 - s\ensuremath{{}^{2}}\ensuremath{{}_{13}})\par
\textit{Value:} 0.0866315368652343\quad\textit{Comparison:} 0.0833 +0.0022/-0.0022 [Daya Bay, Phys. Rev. Lett. 133, 151801 (2024)]\quad\textit{Residual:} +1.514\par
\textit{Operation:} computed.\par
\textit{Calculation:} \path{programs/article_results.py}; \path{outputs/article_results.out}, line 389, identifier \path{SI.C12}.\par
{\footnotesize\textit{Lean:} \path{Comparison.System.row_C12} (\path{lean/Comparison/System/Table/Mixings.lean}:145)}\\[7pt]
\hypertarget{row-C13}{}\textbf{C13}\label{row:C13} & \textbf{atmospheric angle s\textasciicircum{}2\_\allowbreak{}23}\par
\textit{Formula:} s\ensuremath{{}^{2}}\ensuremath{{}_{23}} of U\par
\textit{Value:} 0.488674949067231\quad\textit{Comparison:} 0.47 +0.017/-0.013 [NuFIT 6.0, normal ordering, IC24 with SK atmospheric data, Ref.~\cite{NuFIT6}]\quad\textit{Residual:} +1.099\par
\textit{Operation:} computed.\par
\textit{Calculation:} \path{programs/article_results.py}; \path{outputs/article_results.out}, line 390, identifier \path{SI.C13}.\par
{\footnotesize\textit{Lean:} \path{Comparison.System.row_C13} (\path{lean/Comparison/System/Table/Mixings.lean}:151)}\\[7pt]
\hypertarget{row-C14}{}\textbf{C14}\label{row:C14} & \textbf{lepton phase delta\_\allowbreak{}l}\par
\textit{Formula:} \ensuremath{\delta_{l}}, the phase of U [deg]\par
\textit{Value:} 205.584225835201\quad\textit{Unit:} $\mathrm{deg}$\quad\textit{Comparison:} 212 +26/-41 [NuFIT 6.0, normal ordering, IC24 with SK atmospheric data, Ref.~\cite{NuFIT6}]\quad\textit{Residual:} -0.156\par
\textit{Operation:} computed.\par
\textit{Calculation:} \path{programs/article_results.py}; \path{outputs/article_results.out}, line 391, identifier \path{SI.C14}.\par
{\footnotesize\textit{Lean:} \path{Comparison.System.row_C14} (\path{lean/Comparison/System/Table/Mixings.lean}:159)}\\[7pt]
\hypertarget{row-T06}{}\textbf{T06}\label{row:T06} & \textbf{lepton quartet J\_\allowbreak{}l}\par
\textit{Formula:} \ensuremath{J_{l}} of U\par
\textit{Value:} -0.0144644347463475\quad\textit{Comparison:} prediction\quad\textit{Residual:} -\par
\textit{Operation:} computed.\par
\textit{Calculation:} \path{programs/article_results.py}; \path{outputs/article_results.out}, line 392, identifier \path{SI.T06}.\par
{\footnotesize\textit{Lean:} \path{Comparison.System.row_T06} (\path{lean/Comparison/System/Table/Mixings.lean}:155)}\\[7pt]
\hypertarget{row-E07}{}\textbf{E07}\label{row:E07} & \textbf{light neutrino channels}\par
\textit{Formula:} the three light neutrino channels\par
\textit{Value:} 3\quad\textit{Comparison:} 2.9963 +0.0074/-0.0074 [Janot and Jadach, Phys. Lett. B 803, 135319 (2020)]\quad\textit{Residual:} +0.500\par
\textit{Operation:} computed.\par
\textit{Calculation:} \path{programs/article_results.py}; \path{outputs/article_results.out}, line 393, identifier \path{SI.E07}.\par
{\footnotesize\textit{Lean:} \path{Comparison.System.row_E07} (\path{lean/Comparison/System/Table/Neutrinos.lean}:92)}\\[7pt]
\hypertarget{row-E01}{}\textbf{E01}\label{row:E01} & \textbf{solar splitting Delta m\textasciicircum{}2\_\allowbreak{}21}\par
\textit{Formula:} $\Delta m_{21}^2=m_2^2-m_1^2$, with the light roots of Eq.~\eqref{eq:neutralroots}.\par
\textit{Value:} 7.41781185257392\quad\textit{Unit:} $10^{-5}\,\mathrm{eV}^{2}$\quad\textit{Comparison:} 7.388 +0.078/-0.078 [JUNO 2026, Ref.~\cite{JUNO2026}]\quad\textit{Residual:} +0.382\par
\textit{Operation:} computed.\par
\textit{Calculation:} \path{programs/article_results.py}; \path{outputs/article_results.out}, line 394, identifier \path{SI.E01}.\par
{\footnotesize\textit{Lean:} \path{Comparison.System.row_E01} (\path{lean/Comparison/System/Table/Neutrinos.lean}:67)}\\[7pt]
\hypertarget{row-E02}{}\textbf{E02}\label{row:E02} & \textbf{splitting Delta m\textasciicircum{}2\_\allowbreak{}31}\par
\textit{Formula:} \ensuremath{\Delta}m\ensuremath{{}^{2}}\ensuremath{{}_{31}} = m\ensuremath{{}_{3}}\ensuremath{{}^{2}} [\ensuremath{\times}10\ensuremath{{}^{-3}} eV\ensuremath{{}^{2}}]\par
\textit{Value:} 2.5198724299426\quad\textit{Unit:} $10^{-3}\,\mathrm{eV}^{2}$\quad\textit{Comparison:} 2.509 +0.027/-0.025 [JUNO 2026 with Daya Bay constraint, Ref.~\cite{JUNO2026}]\quad\textit{Residual:} +0.403\par
\textit{Operation:} computed.\par
\textit{Calculation:} \path{programs/article_results.py}; \path{outputs/article_results.out}, line 395, identifier \path{SI.E02}.\par
{\footnotesize\textit{Lean:} \path{Comparison.System.row_E02} (\path{lean/Comparison/System/Table/Neutrinos.lean}:71)}\\[7pt]
\hypertarget{row-E03}{}\textbf{E03}\label{row:E03} & \textbf{splitting Delta m\textasciicircum{}2\_\allowbreak{}32}\par
\textit{Formula:} \ensuremath{\Delta}m\ensuremath{{}^{2}}\ensuremath{{}_{32}} = m\ensuremath{{}_{3}}\ensuremath{{}^{2}} - m\ensuremath{{}_{2}}\ensuremath{{}^{2}} [\ensuremath{\times}10\ensuremath{{}^{-3}} eV\ensuremath{{}^{2}}]\par
\textit{Value:} 2.44569431141686\quad\textit{Unit:} $10^{-3}\,\mathrm{eV}^{2}$\quad\textit{Comparison:} 2.438 +0.021/-0.019 [NuFIT 6.0]\quad\textit{Residual:} +0.366\par
\textit{Operation:} computed.\par
\textit{Calculation:} \path{programs/article_results.py}; \path{outputs/article_results.out}, line 396, identifier \path{SI.E03}.\par
{\footnotesize\textit{Lean:} \path{Comparison.System.row_E03} (\path{lean/Comparison/System/Table/Neutrinos.lean}:75)}\\[7pt]
\hypertarget{row-E04}{}\textbf{E04}\label{row:E04} & \textbf{neutrino mass ratio m\_\allowbreak{}2/m\_\allowbreak{}3}\par
\textit{Formula:} m\ensuremath{{}_{2}}/m\ensuremath{{}_{3}}\par
\textit{Value:} 0.17157287525381\quad\textit{Comparison:} 0.172645 +0.00229/-0.00229 [NuFIT 6.0-derived, m\_\allowbreak{}1 = 0]\quad\textit{Residual:} -0.469\par
\textit{Operation:} computed.\par
\textit{Calculation:} \path{programs/article_results.py}; \path{outputs/article_results.out}, line 397, identifier \path{SI.E04}.\par
{\footnotesize\textit{Lean:} \path{Comparison.System.row_E04} (\path{lean/Comparison/System/Table/Neutrinos.lean}:79)}\\[7pt]
\hypertarget{row-E05}{}\textbf{E05}\label{row:E05} & \textbf{splitting ratio Delta m\textasciicircum{}2\_\allowbreak{}31/Delta m\textasciicircum{}2\_\allowbreak{}21}\par
\textit{Formula:} \ensuremath{\Delta}m\ensuremath{{}^{2}}\ensuremath{{}_{31}}/\ensuremath{\Delta}m\ensuremath{{}^{2}}\ensuremath{{}_{21}}\par
\textit{Value:} 33.9705627484772\quad\textit{Comparison:} 33.55 +0.94/-0.94 [PDG-derived]\quad\textit{Residual:} +0.447\par
\textit{Operation:} computed.\par
\textit{Calculation:} \path{programs/article_results.py}; \path{outputs/article_results.out}, line 398, identifier \path{SI.E05}.\par
{\footnotesize\textit{Lean:} \path{Comparison.System.row_E05} (\path{lean/Comparison/System/Table/Neutrinos.lean}:83)}\\[7pt]
\hypertarget{row-E06}{}\textbf{E06}\label{row:E06} & \textbf{heaviest light neutrino m\_\allowbreak{}3}\par
\textit{Formula:} $m_3=2y_3^2m_p/[X_R+\sqrt{X_R^2+4y_3^2}]$, $y_i=\sigma(v/m_p)(0,\sqrt{\rho_\nu},1)_i/\sqrt2$, $X_R=r_{\rm v}^9/W$.\par
\textit{Value:} 50.1983309477775\quad\textit{Unit:} $\mathrm{meV}$\quad\textit{Comparison:} 50.09 +0.269/-0.25 [square root of the JUNO 2026 atmospheric splitting with Daya Bay constraint, Ref.~\cite{JUNO2026}]\quad\textit{Residual:} +0.403\par
\textit{Operation:} computed.\par
\textit{Calculation:} \path{programs/article_results.py}; \path{outputs/article_results.out}, line 399, identifier \path{SI.E06}.\par
{\footnotesize\textit{Lean:} \path{Comparison.System.row_E06} (\path{lean/Comparison/System/Table/Neutrinos.lean}:87)}\\[7pt]
\hypertarget{row-K02}{}\textbf{K02}\label{row:K02} & \textbf{neutrino mass sum}\par
\textit{Formula:} m\ensuremath{{}_{1}} + m\ensuremath{{}_{2}} + m\ensuremath{{}_{3}} [eV]\par
\textit{Value:} 0.05881100292143 (58.8 meV)\quad\textit{Unit:} $\mathrm{eV}$\quad\textit{Comparison:} $<0.081\eV$ [SPT-3G D1 + DESI, $95\%$ confidence, $\Lambda$CDM+$\sum m_\nu$]\quad\textit{Residual:} below this 81 meV bound; the stronger 55 meV combined-CMB bound in K03 lies below the calculated sum\par
\textit{Operation:} calculated sum; observational bound from Ref.~\cite{Camphuis}.\par
\textit{Calculation:} \path{programs/article_results.py}; \path{outputs/article_results.out}, line 400, identifier \path{SI.K02}.\par
{\footnotesize\textit{Lean:} \path{Comparison.System.row_K02} (\path{lean/Comparison/System/Table/Neutrinos.lean}:95)}\\[7pt]
\hypertarget{row-K03}{}\textbf{K03}\label{row:K03} & \textbf{neutrino mass sum}\par
\textit{Formula:} m\ensuremath{{}_{1}} + m\ensuremath{{}_{2}} + m\ensuremath{{}_{3}} [eV]\par
\textit{Value:} 0.05881100292143 (58.8 meV)\quad\textit{Unit:} $\mathrm{eV}$\quad\textit{Comparison:} $<0.055\eV$ [CMB-SPA + DESI, $95\%$ confidence, $\Lambda$CDM+$\sum m_\nu$]\quad\textit{Residual:} 3.811 meV above the upper bound\par
\textit{Operation:} calculated sum; observational bound from Ref.~\cite{Camphuis}.\par
\textit{Calculation:} \path{programs/article_results.py}; \path{outputs/article_results.out}, line 401, identifier \path{SI.K03}.\par
{\footnotesize\textit{Lean:} \path{Comparison.System.row_K03} (\path{lean/Comparison/System/Table/Neutrinos.lean}:100)}\\[7pt]
\hypertarget{row-L03}{}\textbf{L03}\label{row:L03} & \textbf{light-limit neutrino splitting ratio}\par
\textit{Formula:} \ensuremath{\rho_{\nu}}\ensuremath{{}^{-2}}, \ensuremath{\rho_{\nu}} = tan\ensuremath{{}^{2}}(\ensuremath{\pi}/8)\par
\textit{Value:} 33.9705627484771\quad\textit{Comparison:} prediction\quad\textit{Residual:} -\par
\textit{Operation:} computed.\par
\textit{Calculation:} \path{programs/article_results.py}; \path{outputs/article_results.out}, line 402, identifier \path{SI.L03}.\par
{\footnotesize\textit{Lean:} \path{Comparison.System.row_L03} (\path{lean/Comparison/System/Table/Neutrinos.lean}:103)}\\[7pt]
\hypertarget{row-T01}{}\textbf{T01}\label{row:T01} & \textbf{heavy neutral threshold M\_\allowbreak{}R}\par
\textit{Formula:} \ensuremath{M_{R}} = \ensuremath{M_{\mathrm{Pl}}} \ensuremath{r_{\rm v}}\ensuremath{{}^{9}} [GeV]\par
\textit{Value:} 620276285060925\quad\textit{Unit:} $\mathrm{GeV}$\quad\textit{Comparison:} threshold\quad\textit{Residual:} -\par
\textit{Operation:} computed.\par
\textit{Calculation:} \path{programs/article_results.py}; \path{outputs/article_results.out}, line 403, identifier \path{SI.T01}.\par
{\footnotesize\textit{Lean:} \path{Comparison.System.row_T01} (\path{lean/Comparison/System/Table/Couplings.lean}:246)}\\[7pt]
\hypertarget{row-T02}{}\textbf{T02}\label{row:T02} & \textbf{heavy threshold M\_\allowbreak{}off of the rotations mixing the three-plus-two split}\par
\textit{Formula:} \ensuremath{M_{\mathrm{off}}} = \ensuremath{M_{\mathrm{Pl}}} \ensuremath{r_{\rm v}}\ensuremath{{}^{6}} [GeV]\par
\textit{Value:} 1.6747459696645e+16\quad\textit{Unit:} $\mathrm{GeV}$\quad\textit{Comparison:} threshold\quad\textit{Residual:} -\par
\textit{Operation:} computed.\par
\textit{Calculation:} \path{programs/article_results.py}; \path{outputs/article_results.out}, line 404, identifier \path{SI.T02}.\par
{\footnotesize\textit{Lean:} \path{Comparison.System.row_T02} (\path{lean/Comparison/System/Table/Couplings.lean}:255)}\\[7pt]
\hypertarget{row-S01}{}\textbf{S01}\label{row:S01} & \textbf{axial coupling g\_\allowbreak{}A}\par
\textit{Formula:} \ensuremath{g_{A}} = (5/3)\ensuremath{t_0}\ensuremath{{}^{2b/3}}, 5/3 the axial current of the projected nucleon\par
\textit{Value:} 1.27580092811547\quad\textit{Comparison:} 1.2753 +0.0013/-0.0013 [PDG 2026]\quad\textit{Residual:} +0.385\par
\textit{Operation:} computed.\par
\textit{Calculation:} \path{programs/article_results.py}; \path{outputs/article_results.out}, line 405, identifier \path{SI.S01}.\par
{\footnotesize\textit{Lean:} \path{Comparison.System.row_S01} (\path{lean/Comparison/System/Table/Hadrons.lean}:307)}\\[7pt]
\hypertarget{row-S05}{}\textbf{S05}\label{row:S05} & \textbf{neutron lifetime}\par
\textit{Formula:} \ensuremath{\tau_{n}} = 5099.33 s/[1.03983\ensuremath{|}\ensuremath{V_{\mathrm{ud}}}\ensuremath{|}\ensuremath{{}^{2}}(1 + 3\ensuremath{g_{A}}\ensuremath{{}^{2}})] [s]\par
\textit{Value:} 878.051411826579\quad\textit{Unit:} $\mathrm{s}$\quad\textit{Comparison:} 878.3 +0.4/-0.4 [PDG 2026 neutron listing]\quad\textit{Residual:} -0.621\par
\textit{Operation:} Finite axial and mixing coefficients evaluated with the declared neutron-rate normalization of Part III.\par
\textit{Calculation:} \path{programs/article_results.py}; \path{outputs/article_results.out}, line 406, identifier \path{SI.S05}.\par
{\footnotesize\textit{Lean:} \path{Comparison.System.row_S05} (\path{lean/Comparison/System/Table/Hadrons.lean}:338)}\\[7pt]
\hypertarget{row-S02}{}\textbf{S02}\label{row:S02} & \textbf{axial radius squared}\par
\textit{Formula:} (Var \ensuremath{n_{F}} - Cov\ensuremath{{}^{2}}/Var \ensuremath{n_{L}})\ensuremath{\lambda_{p}}\ensuremath{{}^{2}} [fm\ensuremath{{}^{2}}]\par
\textit{Value:} 0.398068538829894\quad\textit{Unit:} $\mathrm{fm}^{2}$\quad\textit{Comparison:} 0.339 +0.0484/-0.0484 [ETMC 2024 lattice; errors 0.048 and 0.006]\quad\textit{Residual:} +1.221\par
\textit{Operation:} computed.\par
\textit{Calculation:} \path{programs/article_results.py}; \path{outputs/article_results.out}, line 407, identifier \path{SI.S02}.\par
{\footnotesize\textit{Lean:} \path{Comparison.System.row_S02} (\path{lean/Comparison/System/Table/Hadrons.lean}:313)}\\[7pt]
\hypertarget{row-S03}{}\textbf{S03}\label{row:S03} & \textbf{axial radius}\par
\textit{Formula:} (Var \ensuremath{n_{F}} - Cov\ensuremath{{}^{2}}/Var \ensuremath{n_{L}})\ensuremath{{}^{1/2}}\ensuremath{\lambda_{p}} [fm]\par
\textit{Value:} 0.630926730159607\quad\textit{Unit:} $\mathrm{fm}$\quad\textit{Comparison:} construction (S02 is its square)\quad\textit{Residual:} -\par
\textit{Operation:} computed.\par
\textit{Calculation:} \path{programs/article_results.py}; \path{outputs/article_results.out}, line 408, identifier \path{SI.S03}.\par
{\footnotesize\textit{Lean:} \path{Comparison.System.row_S03} (\path{lean/Comparison/System/Table/Hadrons.lean}:321)}\\[7pt]
\hypertarget{row-G03}{}\textbf{G03}\label{row:G03} & \textbf{proton moment mu\_\allowbreak{}p}\par
\textit{Formula:} $\mu_p=\sqrt{(\Var N_{\rm ret})S_\mu/20}$, $S_\mu=1-\epsilon_2/[(\EE N_{\rm ret})^2b]$, $\epsilon_2=\alpha^2/(1+\alpha/10)$.\par
\textit{Value:} 2.79284734529616\quad\textit{Unit:} $\text{nuclear magneton}$\quad\textit{Comparison:} 2.79284734463 +8.2e-10/-8.2e-10 [CODATA 2022]\quad\textit{Residual:} +0.812\par
\textit{Operation:} computed.\par
\textit{Calculation:} \path{programs/article_results.py}; \path{outputs/article_results.out}, line 409, identifier \path{SI.G03}.\par
{\footnotesize\textit{Lean:} \path{Comparison.System.row_G03} (\path{lean/Comparison/System/Table/Hadrons.lean}:189)}\\[7pt]
\hypertarget{row-G04}{}\textbf{G04}\label{row:G04} & \textbf{neutron moment mu\_\allowbreak{}n}\par
\textit{Formula:} $\mu_n=-(44/23)\sqrt{S_\mu}$, with $S_\mu$ as in G03.\par
\textit{Value:} -1.91304302380603\quad\textit{Unit:} $\text{nuclear magneton}$\quad\textit{Comparison:} -1.91304276 +4.5e-07/-4.5e-07 [CODATA 2022]\quad\textit{Residual:} -0.586\par
\textit{Operation:} computed.\par
\textit{Calculation:} \path{programs/article_results.py}; \path{outputs/article_results.out}, line 410, identifier \path{SI.G04}.\par
{\footnotesize\textit{Lean:} \path{Comparison.System.row_G04} (\path{lean/Comparison/System/Table/Hadrons.lean}:196)}\\[7pt]
\hypertarget{row-G05}{}\textbf{G05}\label{row:G05} & \textbf{moment ratio mu\_\allowbreak{}n/mu\_\allowbreak{}p}\par
\textit{Formula:} \ensuremath{\mu_{n}}/\ensuremath{\mu_{p}}\par
\textit{Value:} -0.684979444733371\quad\textit{Comparison:} -0.68497935 +1.6e-07/-1.6e-07 [CODATA 2022]\quad\textit{Residual:} -0.592\par
\textit{Operation:} computed.\par
\textit{Calculation:} \path{programs/article_results.py}; \path{outputs/article_results.out}, line 411, identifier \path{SI.G05}.\par
{\footnotesize\textit{Lean:} \path{Comparison.System.row_G05} (\path{lean/Comparison/System/Table/Hadrons.lean}:202)}\\[7pt]
\hypertarget{row-G13}{}\textbf{G13}\label{row:G13} & \textbf{deuteron moment mu\_\allowbreak{}d}\par
\textit{Formula:} \ensuremath{\mu_{d}} = (16\ensuremath{\mu_{S}} + \ensuremath{g_{l}})/17 - (\ensuremath{B_{d}}/\ensuremath{m_{p}})(\ensuremath{\mu_{S}} - \ensuremath{g_{l}})/(20 + 13\ensuremath{B_{d}}/\ensuremath{m_{p}}), \ensuremath{\mu_{S}} = \ensuremath{\mu_{p}} + \ensuremath{\mu_{n}}, \ensuremath{g_{l}} = \ensuremath{m_{n}}/(\ensuremath{m_{p}} + \ensuremath{m_{n}}) [\ensuremath{\mu_{N}}]\par
\textit{Value:} 0.857438233518713\quad\textit{Unit:} $\text{nuclear magneton}$\quad\textit{Comparison:} 0.8574382335 +2.2e-09/-2.2e-09 [CODATA 2022]\quad\textit{Residual:} +0.009\par
\textit{Operation:} computed.\par
\textit{Calculation:} \path{programs/article_results.py}; \path{outputs/article_results.out}, line 412, identifier \path{SI.G13}.\par
{\footnotesize\textit{Lean:} \path{Comparison.System.row_G13} (\path{lean/Comparison/System/Table/Hadrons.lean}:264)}\\[7pt]
\hypertarget{row-G01}{}\textbf{G01}\label{row:G01} & \textbf{proton charge radius}\par
\textit{Formula:} $r_p=\sqrt{\EE N_{\rm ret}}\lambda_p=4\lambda_p$, $\lambda_p=\hbar/(m_pc)$.\par
\textit{Value:} 0.84123564021281\quad\textit{Unit:} $\mathrm{fm}$\quad\textit{Comparison:} 0.84075 +0.00064/-0.00064 [CODATA 2022]\quad\textit{Residual:} +0.759\par
\textit{Operation:} computed.\par
\textit{Calculation:} \path{programs/article_results.py}; \path{outputs/article_results.out}, line 413, identifier \path{SI.G01}.\par
{\footnotesize\textit{Lean:} \path{Comparison.System.row_G01} (\path{lean/Comparison/System/Table/Hadrons.lean}:164)}\\[7pt]
\hypertarget{row-S04}{}\textbf{S04}\label{row:S04} & \textbf{pion charge radius}\par
\textit{Formula:} \ensuremath{r_{\pi}} = (E \ensuremath{n_{F}})\ensuremath{{}^{1/2}}\ensuremath{\lambda_{p}} [fm]\par
\textit{Value:} 0.665055167995603\quad\textit{Unit:} $\mathrm{fm}$\quad\textit{Comparison:} 0.659 +0.004/-0.004 [PDG]\quad\textit{Residual:} +1.514\par
\textit{Operation:} computed.\par
\textit{Calculation:} \path{programs/article_results.py}; \path{outputs/article_results.out}, line 414, identifier \path{SI.S04}.\par
{\footnotesize\textit{Lean:} \path{Comparison.System.row_S04} (\path{lean/Comparison/System/Table/Hadrons.lean}:330)}\\[7pt]
\hypertarget{row-G08}{}\textbf{G08}\label{row:G08} & \textbf{pion decay constant f\_\allowbreak{}pi}\par
\textit{Formula:} \ensuremath{f_{\pi}} = \ensuremath{\eta_0}\ensuremath{m_{p}}(2\ensuremath{k_{C}})\ensuremath{{}^{1/2}}/(2\ensuremath{\pi}\ensuremath{\mu_{p}}) [MeV]\par
\textit{Value:} 130.529195478157\quad\textit{Unit:} $\mathrm{MeV}$\quad\textit{Comparison:} 130.56 +0.138/-0.138 [PDG, (2)(4)(13) combined]\quad\textit{Residual:} -0.224\par
\textit{Operation:} computed.\par
\textit{Calculation:} \path{programs/article_results.py}; \path{outputs/article_results.out}, line 415, identifier \path{SI.G08}.\par
{\footnotesize\textit{Lean:} \path{Comparison.System.row_G08} (\path{lean/Comparison/System/Table/Hadrons.lean}:230)}\\[7pt]
\hypertarget{row-G09}{}\textbf{G09}\label{row:G09} & \textbf{pion constant times V\_\allowbreak{}ud}\par
\textit{Formula:} \ensuremath{f_{\pi}}\ensuremath{|}\ensuremath{V_{\mathrm{ud}}}\ensuremath{|} [MeV]\par
\textit{Value:} 127.181371252679\quad\textit{Unit:} $\mathrm{MeV}$\quad\textit{Comparison:} 127.13 +0.13/-0.13 [PDG, (1)(13) combined]\quad\textit{Residual:} +0.394\par
\textit{Operation:} computed.\par
\textit{Calculation:} \path{programs/article_results.py}; \path{outputs/article_results.out}, line 416, identifier \path{SI.G09}.\par
{\footnotesize\textit{Lean:} \path{Comparison.System.row_G09} (\path{lean/Comparison/System/Table/Hadrons.lean}:236)}\\[7pt]
\hypertarget{row-A002}{}\textbf{A002}\label{row:A002} & \textbf{pion decay constant, second construction 5m\_\allowbreak{}p/36}\par
\textit{Formula:} 5\ensuremath{m_{p}}/36 [MeV]\par
\textit{Value:} 130.315567975027\quad\textit{Unit:} $\mathrm{MeV}$\quad\textit{Comparison:} 130.2 +0.8/-0.8 [FLAG 2+1]\quad\textit{Residual:} +0.144\par
\textit{Operation:} Finite mass evaluated with the $5/36$ pion ansatz of Part III.\par
\textit{Calculation:} \path{programs/article_results.py}; \path{outputs/article_results.out}, line 417, identifier \path{SI.A002}.\par
{\footnotesize\textit{Lean:} \path{Comparison.System.row_A002} (\path{lean/Comparison/System/Table/Hadrons.lean}:342)}\\[7pt]
\hypertarget{row-A003}{}\textbf{A003}\label{row:A003} & \textbf{pion current product, second construction 5m\_\allowbreak{}p/36}\par
\textit{Formula:} (5\ensuremath{m_{p}}/36)\ensuremath{|}\ensuremath{V_{\mathrm{ud}}}\ensuremath{|} [MeV]\par
\textit{Value:} 126.97322288644\quad\textit{Unit:} $\mathrm{MeV}$\quad\textit{Comparison:} 127.13 +0.13/-0.13 [PDG inclusive pion decay; correlated with lifetime]\quad\textit{Residual:} -1.202\par
\textit{Operation:} Finite mass evaluated with the $5/36$ pion ansatz of Part III.\par
\textit{Calculation:} \path{programs/article_results.py}; \path{outputs/article_results.out}, line 418, identifier \path{SI.A003}.\par
{\footnotesize\textit{Lean:} \path{Comparison.System.row_A003} (\path{lean/Comparison/System/Table/Hadrons.lean}:349)}\\[7pt]
\hypertarget{row-K01}{}\textbf{K01}\label{row:K01} & \textbf{strong phase theta}\par
\textit{Formula:} \ensuremath{\bar{\theta}} = 0\par
\textit{Value:} 0\quad\textit{Comparison:} \ensuremath{|}\ensuremath{\theta}\ensuremath{|} < 1e-10 [neutron EDM]\quad\textit{Residual:} -\par
\textit{Operation:} computed.\par
\textit{Calculation:} \path{programs/article_results.py}; \path{outputs/article_results.out}, line 419, identifier \path{SI.K01}.\par
{\footnotesize\textit{Lean:} \path{Comparison.System.row_K01} (\path{lean/Comparison/System/Table/Neutrinos.lean}:121)}\\[7pt]
\hypertarget{row-L16}{}\textbf{L16}\label{row:L16} & \textbf{strong phase}\par
\textit{Formula:} \ensuremath{\bar{\theta}} = 0\par
\textit{Value:} 0\quad\textit{Comparison:} \ensuremath{|}\ensuremath{\bar{\theta}}\ensuremath{|} < 1e-10 [neutron EDM]\quad\textit{Residual:} -\par
\textit{Operation:} computed.\par
\textit{Calculation:} \path{programs/article_results.py}; \path{outputs/article_results.out}, line 420, identifier \path{SI.L16}.\par
{\footnotesize\textit{Lean:} \path{Comparison.System.row_L16} (\path{lean/Comparison/System/Table/Neutrinos.lean}:127)}\\[7pt]
\hypertarget{row-L14}{}\textbf{L14}\label{row:L14} & \textbf{colour component count N}\par
\textit{Formula:} \ensuremath{k_{C}}\ensuremath{{}^{2}} - 1 = 8 colour components\par
\textit{Value:} 8\quad\textit{Comparison:} identity\quad\textit{Residual:} -\par
\textit{Operation:} computed.\par
\textit{Calculation:} \path{programs/article_results.py}; \path{outputs/article_results.out}, line 421, identifier \path{SI.L14}.\par
{\footnotesize\textit{Lean:} \path{Comparison.System.row_L14} (\path{lean/Comparison/System/Table/Neutrinos.lean}:116)}\\[7pt]
\hypertarget{row-S16}{}\textbf{S16}\label{row:S16} & \textbf{five-flavour one-loop scale marker}\par
\textit{Formula:} \ensuremath{M_{Z}} e\textasciicircum{}(-2\ensuremath{\pi}/((23/3)\ensuremath{\alpha_{s}})) [MeV]\par
\textit{Value:} 87.0577778264853\quad\textit{Unit:} $\mathrm{MeV}$\quad\textit{Comparison:} one-loop convention\quad\textit{Residual:} -\par
\textit{Operation:} computed, with the one-loop coefficient \ensuremath{\beta}\ensuremath{{}_{0}} = 23/3.\par
\textit{Calculation:} \path{programs/article_results.py}; \path{outputs/article_results.out}, line 422, identifier \path{SI.S16}.\par
{\footnotesize\textit{Lean:} \path{Comparison.System.row_S16} (\path{lean/Comparison/System/Table/Couplings.lean}:239)}\\[7pt]
\hypertarget{row-I08}{}\textbf{I08}\label{row:I08} & \textbf{topological curvature chi/(F\_\allowbreak{}pi\textasciicircum{}2 m\_\allowbreak{}pi0\textasciicircum{}2)}\par
\textit{Formula:} $\chi_{\rm top}/(F_\pi^2m_{\pi^0}^2)=[\lambda_{\rm low}\sum_jS_j^{-1}]^{-1}$, $S=(R_{ud},1,p_q)$.\par
\textit{Value:} 0.212911371927165\quad\textit{Comparison:} 0.21054 +0.0224/-0.0224 [lattice, Borsanyi et al., Nature 539, 69 (2016)]\quad\textit{Residual:} +0.106\par
\textit{Operation:} computed.\par
\textit{Calculation:} \path{programs/article_results.py}; \path{outputs/article_results.out}, line 423, identifier \path{SI.I08}.\par
{\footnotesize\textit{Lean:} \path{Comparison.System.row_I08} (\path{lean/Comparison/System/Table/Hadrons.lean}:360)}\\[7pt]
\hypertarget{row-I09}{}\textbf{I09}\label{row:I09} & \textbf{fourth strong cumulant b\_\allowbreak{}2}\par
\textit{Formula:} b\ensuremath{{}_{2}} = -\ensuremath{\Sigma_{j}} \ensuremath{S_{j}}\ensuremath{{}^{-3}}/[12(\ensuremath{\Sigma_{j}} \ensuremath{S_{j}}\ensuremath{{}^{-1}})\ensuremath{{}^{3}}]\par
\textit{Value:} -0.0280584730987574\quad\textit{Comparison:} -0.029 +0.002/-0.002 [NLO chiral, Grilli di Cortona et al., JHEP 01 (2016) 034]\quad\textit{Residual:} +0.471\par
\textit{Operation:} computed.\par
\textit{Calculation:} \path{programs/article_results.py}; \path{outputs/article_results.out}, line 424, identifier \path{SI.I09}.\par
{\footnotesize\textit{Lean:} \path{Comparison.System.row_I09} (\path{lean/Comparison/System/Table/Hadrons.lean}:368)}\\[7pt]
\hypertarget{row-K06}{}\textbf{K06}\label{row:K06} & \textbf{proton partial lifetime tau/B(e pi)}\par
\textit{Formula:} \ensuremath{\tau}/B(e\ensuremath{{}^{+}}\ensuremath{\pi}\ensuremath{{}^{0}}) = 72\ensuremath{{}^{2}}\ensuremath{\hbar}/(3\ensuremath{{}^{24}}W\ensuremath{{}^{4}}\ensuremath{m_{p}}) [yr]\par
\textit{Value:} 1.16971098945626e+37\quad\textit{Unit:} $\mathrm{year}$\quad\textit{Comparison:} > 2.4e34 yr [Super-Kamiokande]\quad\textit{Residual:} -\par
\textit{Operation:} computed.\par
\textit{Calculation:} \path{programs/article_results.py}; \path{outputs/article_results.out}, line 425, identifier \path{SI.K06}.\par
{\footnotesize\textit{Lean:} \path{Comparison.System.row_K06} (\path{lean/Comparison/System/Table/Cosmos.lean}:213)}\\[7pt]
\hypertarget{row-L12}{}\textbf{L12}\label{row:L12} & \textbf{proton partial lifetime}\par
\textit{Formula:} \ensuremath{\tau}/B(e\ensuremath{{}^{+}}\ensuremath{\pi}\ensuremath{{}^{0}}), as K06 [yr]\par
\textit{Value:} 1.16971098945626e+37\quad\textit{Unit:} $\mathrm{year}$\quad\textit{Comparison:} > 2.4e34 yr [Super-Kamiokande]\quad\textit{Residual:} -\par
\textit{Operation:} computed.\par
\textit{Calculation:} \path{programs/article_results.py}; \path{outputs/article_results.out}, line 426, identifier \path{SI.L12}.\par
{\footnotesize\textit{Lean:} \path{Comparison.System.row_L12} (\path{lean/Comparison/System/Table/Cosmos.lean}:230)}\\[7pt]
\hypertarget{row-F01}{}\textbf{F01}\label{row:F01} & \textbf{baryons per photon eta\_\allowbreak{}B}\par
\textit{Formula:} \ensuremath{\eta_{B}} = 3a\ensuremath{{}^{11}} [\ensuremath{\times}10\ensuremath{{}^{-10}}]\par
\textit{Value:} 6.14004762346375\quad\textit{Unit:} $10^{-10}$\quad\textit{Comparison:} 6.12 +0.04/-0.04 [PDG 2024 BBN review, Planck omega\_\allowbreak{}b = 0.02237(15)]\quad\textit{Residual:} +0.501\par
\textit{Operation:} computed.\par
\textit{Calculation:} \path{programs/article_results.py}; \path{outputs/article_results.out}, line 427, identifier \path{SI.F01}.\par
{\footnotesize\textit{Lean:} \path{Comparison.System.row_F01} (\path{lean/Comparison/System/Table/Cosmos.lean}:112)}\\[7pt]
\hypertarget{row-C15}{}\textbf{C15}\label{row:C15} & \textbf{scalar tilt n\_\allowbreak{}s}\par
\textit{Formula:} \ensuremath{n_{s}} = 1 - \ensuremath{t_0}/E\ensuremath{N_{\rm ret}} = 39/40\par
\textit{Value:} 0.975\quad\textit{Comparison:} 0.9752 +0.003/-0.003 [P-ACT-LB2, Ref.~\cite{ACT2025}, Table 5]\quad\textit{Residual:} -0.067\par
\textit{Operation:} computed.\par
\textit{Calculation:} \path{programs/article_results.py}; \path{outputs/article_results.out}, line 428, identifier \path{SI.C15}.\par
{\footnotesize\textit{Lean:} \path{Comparison.System.row_C15} (\path{lean/Comparison/System/Table/Cosmos.lean}:100)}\\[7pt]
\hypertarget{row-C16}{}\textbf{C16}\label{row:C16} & \textbf{running of the tilt dn\_\allowbreak{}s/dlnk}\par
\textit{Formula:} d\ensuremath{n_{s}}/d ln k = 0\par
\textit{Value:} 0\quad\textit{Comparison:} -0.0045 +0.0067/-0.0067 [Planck 2018]\quad\textit{Residual:} +0.672\par
\textit{Operation:} The stated continuum equation evaluated with its finite coefficients.\par
\textit{Calculation:} \path{programs/article_results.py}; \path{outputs/article_results.out}, line 429, identifier \path{SI.C16}.\par
{\footnotesize\textit{Lean:} \path{Comparison.System.row_C16} (\path{lean/Comparison/System/Table/Cosmos.lean}:105)}\\[7pt]
\hypertarget{row-L06}{}\textbf{L06}\label{row:L06} & \textbf{scalar-index running}\par
\textit{Formula:} d\ensuremath{n_{s}}/d ln k = 0\par
\textit{Value:} 0\quad\textit{Comparison:} prediction\quad\textit{Residual:} -\par
\textit{Operation:} The stated continuum equation evaluated with its finite coefficients.\par
\textit{Calculation:} \path{programs/article_results.py}; \path{outputs/article_results.out}, line 430, identifier \path{SI.L06}.\par
{\footnotesize\textit{Lean:} \path{Comparison.System.row_L06} (\path{lean/Comparison/System/Table/Cosmos.lean}:108)}\\[7pt]
\hypertarget{row-F03}{}\textbf{F03}\label{row:F03} & \textbf{primordial amplitude ln(1e10 A\_\allowbreak{}s)}\par
\textit{Formula:} ln(10\ensuremath{{}^{10}}\ensuremath{A_{s}}), \ensuremath{A_{s}} = 12a\ensuremath{{}^{11}}(0.05\ensuremath{\ell})\ensuremath{{}^{-1/40}}\par
\textit{Value:} 3.06130986015343\quad\textit{Comparison:} 3.056 +0.013/-0.013 [P-ACT, Ref.~\cite{ACT2025}, Table 5]\quad\textit{Residual:} +0.408\par
\textit{Operation:} computed, at the pivot k\_\allowbreak{}* = 0.05 Mpc\ensuremath{{}^{-1}} of the measurement.\par
\textit{Calculation:} \path{programs/article_results.py}; \path{outputs/article_results.out}, line 431, identifier \path{SI.F03}.\par
{\footnotesize\textit{Lean:} \path{Comparison.System.row_F03} (\path{lean/Comparison/System/Table/Cosmos.lean}:120)}\\[7pt]
\hypertarget{row-S08}{}\textbf{S08}\label{row:S08} & \textbf{primordial amplitude, second construction}\par
\textit{Formula:} b\ensuremath{{}^{24}}(0.05\ensuremath{\ell})\ensuremath{{}^{-1/40}} [\ensuremath{\times}10\ensuremath{{}^{-9}}]\par
\textit{Value:} 2.10257562469885\quad\textit{Unit:} $10^{-9}$\quad\textit{Comparison:} 2.099 +0.029/-0.029 [Planck 2018]\quad\textit{Residual:} +0.123\par
\textit{Operation:} Second construction; its factor was chosen with the measured value available.\par
\textit{Calculation:} \path{programs/article_results.py}; \path{outputs/article_results.out}, line 432, identifier \path{SI.S08}.\par
{\footnotesize\textit{Lean:} \path{Comparison.System.row_S08} (\path{lean/Comparison/System/Table/Cosmos.lean}:184)}\\[7pt]
\hypertarget{row-L11}{}\textbf{L11}\label{row:L11} & \textbf{tensor-to-scalar ratio in the de Sitter convention}\par
\textit{Formula:} r = 16\ensuremath{\Lambda}\ensuremath{\ell_{P}}\ensuremath{{}^{2}}/(3\ensuremath{\pi}\ensuremath{\cdot}12a\ensuremath{{}^{11}})\par
\textit{Value:} 1.97359117549973e-113\quad\textit{Comparison:} prediction\quad\textit{Residual:} -\par
\textit{Operation:} computed.\par
\textit{Calculation:} \path{programs/article_results.py}; \path{outputs/article_results.out}, line 433, identifier \path{SI.L11}.\par
{\footnotesize\textit{Lean:} \path{Comparison.System.row_L11} (\path{lean/Comparison/System/Table/Cosmos.lean}:174)}\\[7pt]
\hypertarget{row-D06}{}\textbf{D06}\label{row:D06} & \textbf{cosmological constant Lambda l\_\allowbreak{}P\textasciicircum{}2}\par
\textit{Formula:} $\Lambda\ell_P^2=3W^6\theta_c s_F^3$, $\theta_c=2\pi/(16R_p)$, $s_F=r_0/10=3/50$.\par
\textit{Value:} 2.85522324300807\quad\textit{Unit:} $10^{-122}$\quad\textit{Comparison:} 2.8451 +0.076/-0.076 [Planck 2018 base model, 3 Omega\_\allowbreak{}Lambda H\_\allowbreak{}0\textasciicircum{}2 l\_\allowbreak{}P\textasciicircum{}2/c\textasciicircum{}2 with Omega\_\allowbreak{}Lambda = 0.6847 and H\_\allowbreak{}0 = 67.36]\quad\textit{Residual:} +0.133\par
\textit{Operation:} computed.\par
\textit{Calculation:} \path{programs/article_results.py}; \path{outputs/article_results.out}, line 434, identifier \path{SI.D06}.\par
{\footnotesize\textit{Lean:} \path{Comparison.System.row_D06} (\path{lean/Comparison/System/Table/Cosmos.lean}:151)}\\[7pt]
\hypertarget{row-T03}{}\textbf{T03}\label{row:T03} & \textbf{de Sitter radius}\par
\textit{Formula:} $\ell=\lambda_p/[W^2\sqrt{\theta_cs_F^3}]$.\par
\textit{Value:} 5369.08213305099\quad\textit{Unit:} $\mathrm{Mpc}$\quad\textit{Comparison:} de Sitter convention\quad\textit{Residual:} -\par
\textit{Operation:} computed.\par
\textit{Calculation:} \path{programs/article_results.py}; \path{outputs/article_results.out}, line 435, identifier \path{SI.T03}.\par
{\footnotesize\textit{Lean:} \path{Comparison.System.row_T03} (\path{lean/Comparison/System/Table/Cosmos.lean}:197)}\\[7pt]
\hypertarget{row-T04}{}\textbf{T04}\label{row:T04} & \textbf{de Sitter temperature}\par
\textit{Formula:} $k_B T_{\rm dS}=\hbar c/(2\pi\ell)$.\par
\textit{Value:} 1.89564063245871e-34\quad\textit{Unit:} $\mathrm{eV}$\quad\textit{Comparison:} de Sitter convention\quad\textit{Residual:} -\par
\textit{Operation:} computed.\par
\textit{Calculation:} \path{programs/article_results.py}; \path{outputs/article_results.out}, line 436, identifier \path{SI.T04}.\par
{\footnotesize\textit{Lean:} \path{Comparison.System.row_T04} (\path{lean/Comparison/System/Table/Cosmos.lean}:203)}\\[7pt]
\hypertarget{row-L10}{}\textbf{L10}\label{row:L10} & \textbf{de Sitter horizon entropy}\par
\textit{Formula:} $S_{\rm dS}/k_B=3\pi/(\Lambda\ell_P^2)$.\par
\textit{Value:} 3.30089003858068e+122\quad\textit{Comparison:} identity with D06\quad\textit{Residual:} -\par
\textit{Operation:} computed.\par
\textit{Calculation:} \path{programs/article_results.py}; \path{outputs/article_results.out}, line 437, identifier \path{SI.L10}.\par
{\footnotesize\textit{Lean:} \path{Comparison.System.row_L10} (\path{lean/Comparison/System/Table/Cosmos.lean}:167)}\\[7pt]
\hypertarget{row-L08}{}\textbf{L08}\label{row:L08} & \textbf{the equation of state of the constant term}\par
\textit{Formula:} $w_\Lambda=-1$.\par
\textit{Value:} -1\quad\textit{Comparison:} prediction\quad\textit{Residual:} -\par
\textit{Operation:} The stated continuum equation evaluated with its finite coefficients.\par
\textit{Calculation:} \path{programs/article_results.py}; \path{outputs/article_results.out}, line 438, identifier \path{SI.L08}.\par
{\footnotesize\textit{Lean:} \path{Comparison.System.row_L08} (\path{lean/Comparison/System/Table/Cosmos.lean}:148)}\\[7pt]
\hypertarget{row-F04}{}\textbf{F04}\label{row:F04} & \textbf{effective neutrinos N\_\allowbreak{}eff}\par
\textit{Formula:} \ensuremath{N_{\mathrm{eff}}} = 3\ensuremath{\eta_0}\ensuremath{\sigma}\ensuremath{{}^{4/3}}\par
\textit{Value:} 3.04385727361762\quad\textit{Comparison:} 2.99 +0.17/-0.17 [Planck + BAO]\quad\textit{Residual:} +0.317\par
\textit{Operation:} computed.\par
\textit{Calculation:} \path{programs/article_results.py}; \path{outputs/article_results.out}, line 439, identifier \path{SI.F04}.\par
{\footnotesize\textit{Lean:} \path{Comparison.System.row_F04} (\path{lean/Comparison/System/Table/Cosmos.lean}:129)}\\[7pt]
\hypertarget{row-F13}{}\textbf{F13}\label{row:F13} & \textbf{spatial curvature Omega\_\allowbreak{}k}\par
\textit{Formula:} \ensuremath{\Omega_{k}} = 0\par
\textit{Value:} 0\quad\textit{Comparison:} 0.0007 +0.0019/-0.0019 [Planck + BAO]\quad\textit{Residual:} -0.368\par
\textit{Operation:} The stated continuum equation evaluated with its finite coefficients.\par
\textit{Calculation:} \path{programs/article_results.py}; \path{outputs/article_results.out}, line 440, identifier \path{SI.F13}.\par
{\footnotesize\textit{Lean:} \path{Comparison.System.row_F13} (\path{lean/Comparison/System/Table/Cosmos.lean}:141)}\\[7pt]
\hypertarget{row-L07}{}\textbf{L07}\label{row:L07} & \textbf{spatial curvature fraction}\par
\textit{Formula:} \ensuremath{\Omega_{k}} = 0\par
\textit{Value:} 0\quad\textit{Comparison:} prediction\quad\textit{Residual:} -\par
\textit{Operation:} The stated continuum equation evaluated with its finite coefficients.\par
\textit{Calculation:} \path{programs/article_results.py}; \path{outputs/article_results.out}, line 441, identifier \path{SI.L07}.\par
{\footnotesize\textit{Lean:} \path{Comparison.System.row_L07} (\path{lean/Comparison/System/Table/Cosmos.lean}:144)}\\[7pt]
\hypertarget{row-S09}{}\textbf{S09}\label{row:S09} & \textbf{cold/baryon count, second construction}\par
\textit{Formula:} M/(M - 4\ensuremath{C_{\ensuremath{h_\nu}}}) = 37/7\par
\textit{Value:} 5.28571428571429\quad\textit{Comparison:} 5.3643 +0.0646/-0.0646 [Planck 2018 ratio]\quad\textit{Residual:} -1.216\par
\textit{Operation:} Second construction; its factor was chosen with the measured value available.\par
\textit{Calculation:} \path{programs/article_results.py}; \path{outputs/article_results.out}, line 442, identifier \path{SI.S09}.\par
{\footnotesize\textit{Lean:} \path{Comparison.System.row_S09} (\path{lean/Comparison/System/Table/Cosmos.lean}:193)}\\[7pt]
\hypertarget{row-D07}{}\textbf{D07}\label{row:D07} & \textbf{d log alpha / d(Phi/c\textasciicircum{}2)}\par
\textit{Formula:} d log \ensuremath{\alpha}/d(\ensuremath{\Phi_G}/c\ensuremath{{}^{2}}) = 0\par
\textit{Value:} 0\quad\textit{Comparison:} 14 +11/-11 [atomic clocks, Lange et al., Phys. Rev. Lett. 126, 011102 (2021)]\quad\textit{Residual:} -1.273\par
\textit{Operation:} The stated continuum equation evaluated with its finite coefficients.\par
\textit{Calculation:} \path{programs/article_results.py}; \path{outputs/article_results.out}, line 443, identifier \path{SI.D07}.\par
{\footnotesize\textit{Lean:} \path{Comparison.System.row_D07} (\path{lean/Comparison/System/Table/Cosmos.lean}:161)}\\[7pt]
\hypertarget{row-L09}{}\textbf{L09}\label{row:L09} & \textbf{a recorded electric fact does not drift}\par
\textit{Formula:} the drift of \ensuremath{\alpha}, 0\par
\textit{Value:} 0\quad\textit{Comparison:} prediction\quad\textit{Residual:} -\par
\textit{Operation:} The stated continuum equation evaluated with its finite coefficients.\par
\textit{Calculation:} \path{programs/article_results.py}; \path{outputs/article_results.out}, line 444, identifier \path{SI.L09}.\par
{\footnotesize\textit{Lean:} \path{Comparison.System.row_L09} (\path{lean/Comparison/System/Table/Cosmos.lean}:164)}\\[7pt]
\hypertarget{row-H01}{}\textbf{H01}\label{row:H01} & \textbf{light bending 1 + gamma}\par
\textit{Formula:} 1 + \ensuremath{\gamma} = 2\par
\textit{Value:} 2\quad\textit{Comparison:} 2.000021 +2.3e-05/-2.3e-05 [Cassini, Bertotti et al., Nature 425, 374 (2003)]\quad\textit{Residual:} -0.913\par
\textit{Operation:} The stated continuum equation evaluated with its finite coefficients.\par
\textit{Calculation:} \path{programs/article_results.py}; \path{outputs/article_results.out}, line 445, identifier \path{SI.H01}.\par
{\footnotesize\textit{Lean:} \path{Comparison.System.row_H01} (\path{lean/Comparison/System/Table/Neutrinos.lean}:134)}\\[7pt]
\hypertarget{row-H02}{}\textbf{H02}\label{row:H02} & \textbf{geodetic drift / its Einstein value}\par
\textit{Formula:} geodetic drift over its Einstein value = 1\par
\textit{Value:} 1\quad\textit{Comparison:} 0.99934908645 +0.00277/-0.00277 [Gravity Probe B, Everitt et al., Phys. Rev. Lett. 106, 221101 (2011)]\quad\textit{Residual:} +0.235\par
\textit{Operation:} The stated continuum equation evaluated with its finite coefficients.\par
\textit{Calculation:} \path{programs/article_results.py}; \path{outputs/article_results.out}, line 446, identifier \path{SI.H02}.\par
{\footnotesize\textit{Lean:} \path{Comparison.System.row_H02} (\path{lean/Comparison/System/Table/Neutrinos.lean}:137)}\\[7pt]
\hypertarget{row-L18}{}\textbf{L18}\label{row:L18} & \textbf{second-order light-bending coefficient}\par
\textit{Formula:} second-order light bending 15\ensuremath{\pi}/4\par
\textit{Value:} 11.7809724509617\quad\textit{Comparison:} prediction\quad\textit{Residual:} -\par
\textit{Operation:} The stated continuum equation evaluated with its finite coefficients.\par
\textit{Calculation:} \path{programs/article_results.py}; \path{outputs/article_results.out}, line 447, identifier \path{SI.L18}.\par
{\footnotesize\textit{Lean:} \path{Comparison.System.row_L18} (\path{lean/Comparison/System/Table/Neutrinos.lean}:141)}\\[7pt]
\hypertarget{row-H03}{}\textbf{H03}\label{row:H03} & \textbf{two orders in coherent superposition told apart (quantum switch)}\par
\textit{Formula:} the two orders in coherent superposition tell commuting and anticommuting pairs apart with certainty: 1\par
\textit{Value:} 1\quad\textit{Comparison:} 0.976 +0.015/-0.015 [quantum switch, Procopio et al., Nat. Commun. 6, 7913 (2015)]\quad\textit{Residual:} +1.600\par
\textit{Operation:} computed.\par
\textit{Calculation:} \path{programs/article_results.py}; \path{outputs/article_results.out}, line 448, identifier \path{SI.H03}.\par
{\footnotesize\textit{Lean:} \path{Comparison.System.row_H03} (\path{lean/Comparison/System/Table/Neutrinos.lean}:146)}\\[7pt]
\pagebreak
\multicolumn{2}{@{}l}{\textbf{Beyond the main line: several world lines bound together}}\\*
\hypertarget{row-G02}{}\textbf{G02}\label{row:G02} & \textbf{Delta mass}\par
\textit{Formula:} pole: the spin-aligned pole of the first return: \ensuremath{m_{\Delta}} = 3b \ensuremath{m_{p}}\par
\textit{Value:} 1231.48211736401\quad\textit{Unit:} $\mathrm{MeV}$\quad\textit{Comparison:} 1232 +2/-2 [PDG 2026]\quad\textit{Residual:} -0.259\par
\textit{Operation:} computed.\par
\textit{Calculation:} \path{programs/article_results.py}; \path{outputs/article_results.out}, line 449, identifier \path{SI.G02}.\par
{\footnotesize\textit{Lean:} \path{Comparison.System.row_G02} (\path{lean/Comparison/System/Table/Hadrons.lean}:172)}\\[7pt]
\hypertarget{row-G06}{}\textbf{G06}\label{row:G06} & \textbf{neutral pion mass}\par
\textit{Formula:} $m_{\pi^0}=\sqrt{38\hover m_em_p}$.\par
\textit{Value:} 134.977044559769\quad\textit{Unit:} $\mathrm{MeV}$\quad\textit{Comparison:} 134.9768 +0.0005/-0.0005 [PDG 2026]\quad\textit{Residual:} +0.489\par
\textit{Operation:} computed.\par
\textit{Calculation:} \path{programs/article_results.py}; \path{outputs/article_results.out}, line 450, identifier \path{SI.G06}.\par
{\footnotesize\textit{Lean:} \path{Comparison.System.row_G06} (\path{lean/Comparison/System/Table/Hadrons.lean}:209)}\\[7pt]
\hypertarget{row-G07}{}\textbf{G07}\label{row:G07} & \textbf{charged pion mass}\par
\textit{Formula:} $m_{\pi^\pm}=\sqrt{m_{\pi^0}^2+\pi\alpha m_p^2/[16(1-\alpha b\Phis(\ii))]}$.\par
\textit{Value:} 139.570386613809\quad\textit{Unit:} $\mathrm{MeV}$\quad\textit{Comparison:} 139.57039 +0.00018/-0.00018 [PDG 2024]\quad\textit{Residual:} -0.019\par
\textit{Operation:} computed.\par
\textit{Calculation:} \path{programs/article_results.py}; \path{outputs/article_results.out}, line 451, identifier \path{SI.G07}.\par
{\footnotesize\textit{Lean:} \path{Comparison.System.row_G07} (\path{lean/Comparison/System/Table/Hadrons.lean}:220)}\\[7pt]
\hypertarget{row-G10}{}\textbf{G10}\label{row:G10} & \textbf{deuteron binding m\_\allowbreak{}p + m\_\allowbreak{}n - m\_\allowbreak{}d}\par
\textit{Formula:} $B_d=m_p+m_n-m_{\rm deut}=m_p\sqrt{\chi_A}a^3[1-(41/40)\epsilon_2]$.\par
\textit{Value:} 2.22456646309919\quad\textit{Unit:} $\mathrm{MeV}$\quad\textit{Comparison:} 2.22456636409 +3.8e-07/-3.8e-07 [CODATA 2022 mass ratios]\quad\textit{Residual:} +0.261\par
\textit{Operation:} computed.\par
\textit{Calculation:} \path{programs/article_results.py}; \path{outputs/article_results.out}, line 452, identifier \path{SI.G10}.\par
{\footnotesize\textit{Lean:} \path{Comparison.System.row_G10} (\path{lean/Comparison/System/Table/Hadrons.lean}:242)}\\[7pt]
\hypertarget{row-G11}{}\textbf{G11}\label{row:G11} & \textbf{deuteron mass}\par
\textit{Formula:} $m_{\rm deut}=m_p+m_n-B_d$.\par
\textit{Value:} 1875.61294498229\quad\textit{Unit:} $\mathrm{MeV}$\quad\textit{Comparison:} 1875.612945 +5.8e-07/-5.8e-07 [CODATA 2022]\quad\textit{Residual:} -0.031\par
\textit{Operation:} computed.\par
\textit{Calculation:} \path{programs/article_results.py}; \path{outputs/article_results.out}, line 453, identifier \path{SI.G11}.\par
{\footnotesize\textit{Lean:} \path{Comparison.System.row_G11} (\path{lean/Comparison/System/Table/Hadrons.lean}:252)}\\[7pt]
\hypertarget{row-G12}{}\textbf{G12}\label{row:G12} & \textbf{deuteron/proton m\_\allowbreak{}d/m\_\allowbreak{}p}\par
\textit{Formula:} $m_{\rm deut}/m_p$.\par
\textit{Value:} 1.99900750126897\quad\textit{Comparison:} 1.99900750127 +8.4e-12/-8.4e-12 [CODATA 2022]\quad\textit{Residual:} -0.122019\par
\textit{Operation:} computed.\par
\textit{Calculation:} \path{programs/article_results.py}; \path{outputs/article_results.out}, line 454, identifier \path{SI.G12}.\par
{\footnotesize\textit{Lean:} \path{Comparison.System.row_G12} (\path{lean/Comparison/System/Table/Hadrons.lean}:258)}\\[7pt]
\hypertarget{row-J03}{}\textbf{J03}\label{row:J03} & \textbf{\ensuremath{\Lambda} (MeV)}\par
\textit{Formula:} \ensuremath{m_{n}}(5/2)\ensuremath{{}^{60/320}}\par
\textit{Value:} 1115.68290090769\quad\textit{Unit:} $\mathrm{MeV}$\quad\textit{Comparison:} 1115.683 +0.006 / -0.006 [source not recorded]\quad\textit{Residual:} -0.0165154\par
\textit{Operation:} computed.\par
\textit{Calculation:} \path{programs/article_results.py}; \path{outputs/article_results.out}, line 455, identifier \path{SI.J03}.\par
{\footnotesize\textit{Lean:} \path{Comparison.System.row_J03} (\path{lean/Comparison/System/Table/Hadrons.lean}:271)}\\[7pt]
\hypertarget{row-J04}{}\textbf{J04}\label{row:J04} & \textbf{\ensuremath{\Delta}, Breit--Wigner (MeV)}\par
\textit{Formula:} 3b\ensuremath{m_{p}}\par
\textit{Value:} 1231.48211736401\quad\textit{Unit:} $\mathrm{MeV}$\quad\textit{Comparison:} 1232 +2 / -2 [PDG 2026]\quad\textit{Residual:} -0.258941\par
\textit{Operation:} computed.\par
\textit{Calculation:} \path{programs/article_results.py}; \path{outputs/article_results.out}, line 456, identifier \path{SI.J04}.\par
{\footnotesize\textit{Lean:} \path{Comparison.System.row_J04} (\path{lean/Comparison/System/Table/Hadrons.lean}:180)}\\[7pt]
\hypertarget{row-J05}{}\textbf{J05}\label{row:J05} & \textbf{\ensuremath{\pi}\ensuremath{{}^{0}} (MeV)}\par
\textit{Formula:} $m_{\pi^0}=\sqrt{38\hover m_e m_p}$.\par
\textit{Value:} 134.977044559769\quad\textit{Unit:} $\mathrm{MeV}$\quad\textit{Comparison:} 134.9768 +0.0005 / -0.0005 [PDG 2026]\quad\textit{Residual:} +0.48912\par
\textit{Operation:} computed.\par
\textit{Calculation:} \path{programs/article_results.py}; \path{outputs/article_results.out}, line 457, identifier \path{SI.J05}.\par
{\footnotesize\textit{Lean:} \path{Comparison.System.row_J05} (\path{lean/Comparison/System/Table/Hadrons.lean}:215)}\\[7pt]
\hypertarget{row-J06}{}\textbf{J06}\label{row:J06} & \textbf{\ensuremath{\pi}\ensuremath{\pm} (MeV)}\par
\textit{Formula:} \ensuremath{\surd}(\ensuremath{m_{\pi}}0\ensuremath{{}^{2}} + \ensuremath{\Delta_{\mathrm{el}}})\par
\textit{Value:} 139.570386613809\quad\textit{Unit:} $\mathrm{MeV}$\quad\textit{Comparison:} 139.57039 +0.00018 / -0.00018 [PDG 2024]\quad\textit{Residual:} -0.0188122\par
\textit{Operation:} computed.\par
\textit{Calculation:} \path{programs/article_results.py}; \path{outputs/article_results.out}, line 458, identifier \path{SI.J06}.\par
{\footnotesize\textit{Lean:} \path{Comparison.System.row_J06} (\path{lean/Comparison/System/Table/Hadrons.lean}:226)}\\[7pt]
\hypertarget{row-J07}{}\textbf{J07}\label{row:J07} & \textbf{\ensuremath{\rho} pole, Re s [GeV\ensuremath{{}^{2}}] (GeV\ensuremath{{}^{2}})}\par
\textit{Formula:} 6\ensuremath{m_{p}}\ensuremath{{}^{2}} ln[(\ensuremath{\surd}313 - 11)/6]\par
\textit{Value:} 0.576408555494293\quad\textit{Unit:} $\mathrm{GeV}^{2}$\quad\textit{Comparison:} 0.577879 +0.002449 / -0.002449 [source not recorded]\quad\textit{Residual:} -0.600427\par
\textit{Operation:} computed.\par
\textit{Calculation:} \path{programs/article_results.py}; \path{outputs/article_results.out}, line 459, identifier \path{SI.J07}.\par
{\footnotesize\textit{Lean:} \path{Comparison.System.row_J07} (\path{lean/Comparison/System/Table/Hadrons.lean}:280)}\\[7pt]
\hypertarget{row-J08}{}\textbf{J08}\label{row:J08} & \textbf{K\ensuremath{{}^{0}} (MeV)}\par
\textit{Formula:} P(\ensuremath{S_{d}},\ensuremath{S_{s}}) = [2\ensuremath{k_{C}} \ensuremath{m_{p}}\ensuremath{{}^{2}} tr \ensuremath{\delta}\ensuremath{{}_{1}}(\ensuremath{R_{\mathrm{ds}}})]\ensuremath{{}^{1/2}}\par
\textit{Value:} 497.613584450934\quad\textit{Unit:} $\mathrm{MeV}$\quad\textit{Comparison:} 497.611 +0.013 / -0.013 [source not recorded]\quad\textit{Residual:} +0.198804\par
\textit{Operation:} computed.\par
\textit{Calculation:} \path{programs/article_results.py}; \path{outputs/article_results.out}, line 460, identifier \path{SI.J08}.\par
{\footnotesize\textit{Lean:} \path{Comparison.System.row_J08} (\path{lean/Comparison/System/Table/Hadrons.lean}:289)}\\[7pt]
\hypertarget{row-J09}{}\textbf{J09}\label{row:J09} & \textbf{K\ensuremath{\pm} (MeV)}\par
\textit{Formula:} [2\ensuremath{k_{C}} \ensuremath{m_{p}}\ensuremath{{}^{2}} tr \ensuremath{\delta}\ensuremath{{}_{1}}(\ensuremath{R_{\mathrm{us}}}) + (48/25)\ensuremath{\Delta_{\mathrm{el}}}]\ensuremath{{}^{1/2}}, the marked current at the plane hk\par
\textit{Value:} 493.665032668308\quad\textit{Unit:} $\mathrm{MeV}$\quad\textit{Comparison:} 493.677 +0.015 / -0.015 [source not recorded]\quad\textit{Residual:} -0.798\par
\textit{Operation:} computed.\par
\textit{Calculation:} \path{programs/article_results.py}; \path{outputs/article_results.out}, line 461, identifier \path{SI.J09}.\par
{\footnotesize\textit{Lean:} \path{Comparison.System.row_J09} (\path{lean/Comparison/System/Table/Hadrons.lean}:294)}\\[7pt]
\hypertarget{row-J17}{}\textbf{J17}\label{row:J17} & \textbf{kT\_\allowbreak{}c, the colour crossover temperature (MeV)}\par
\textit{Formula:} \ensuremath{m_{p}}/L(\ensuremath{\Omega}) = \ensuremath{m_{p}}/6\par
\textit{Value:} 156.378681570033\quad\textit{Unit:} $\mathrm{MeV}$\quad\textit{Comparison:} 156.5 +1.5 / -1.5 [source not recorded]\quad\textit{Residual:} -0.080879\par
\textit{Operation:} computed.\par
\textit{Calculation:} \path{programs/article_results.py}; \path{outputs/article_results.out}, line 462, identifier \path{SI.J17}.\par
{\footnotesize\textit{Lean:} \path{Comparison.System.row_J17} (\path{lean/Comparison/System/Table/Hadrons.lean}:299)}\\[7pt]
\hypertarget{row-L13}{}\textbf{L13}\label{row:L13} & \textbf{completed deuteron binding}\par
\textit{Formula:} pole: pole difference \ensuremath{B_{d}}\par
\textit{Value:} 2.22456646309919\quad\textit{Unit:} $\mathrm{MeV}$\quad\textit{Comparison:} from G10\quad\textit{Residual:} -\par
\textit{Operation:} computed.\par
\textit{Calculation:} \path{programs/article_results.py}; \path{outputs/article_results.out}, line 463, identifier \path{SI.L13}.\par
{\footnotesize\textit{Lean:} \path{Comparison.System.row_L13} (\path{lean/Comparison/System/Table/Hadrons.lean}:248)}\\[7pt]
\hypertarget{row-M3}{}\textbf{M3}\label{row:M3} & \textbf{deuteron nuclear mass MeV}\par
\textit{Formula:} \ensuremath{m_{p}} + \ensuremath{m_{n}} - \ensuremath{B_{d}}\par
\textit{Value:} 1875.61294498229\quad\textit{Unit:} $\mathrm{MeV}$\quad\textit{Comparison:} 1875.612945 \ensuremath{\pm} 5.8e-07 [CODATA 2022]\quad\textit{Residual:} -0.0305386\par
\textit{Operation:} computed.\par
\textit{Calculation:} \path{programs/article_results.py}; \path{outputs/article_results.out}, line 464, identifier \path{SI.M3}.\par
{\footnotesize\textit{Lean:} \path{Comparison.System.row_M3} (\path{lean/Comparison/System/Table/Nuclei.lean}:166)}\\[7pt]
\hypertarget{row-N22}{}\textbf{N22}\label{row:N22} & \textbf{saturation energy e\_\allowbreak{}0}\par
\textit{Formula:} pole: the bulk volume term -\ensuremath{a_{V}} = -\ensuremath{m_{p}} a\ensuremath{{}^{2}}(1 - \ensuremath{t_0}/E\ensuremath{N_{\rm ret}})\par
\textit{Value:} -15.8382148058291\quad\textit{Unit:} $\mathrm{MeV}$\quad\textit{Comparison:} -16 +0.5/-0.5 [extrapolation of masses to symmetric matter]\quad\textit{Residual:} +0.324\par
\textit{Operation:} computed.\par
\textit{Calculation:} \path{programs/article_results.py}; \path{outputs/article_results.out}, line 465, identifier \path{SI.N22}.\par
{\footnotesize\textit{Lean:} \path{Comparison.System.row_N22} (\path{lean/Comparison/System/Table/Nuclei.lean}:370)}\\[7pt]
\hypertarget{row-T07}{}\textbf{T07}\label{row:T07} & \textbf{deuteron D-state share}\par
\textit{Formula:} eigenvector: D-state component of the deuteron's eigenvector at its pole: \ensuremath{P_{D}} = 2/51\par
\textit{Value:} 0.0392156862745098\quad\textit{Comparison:} construction\quad\textit{Residual:} -\par
\textit{Operation:} computed.\par
\textit{Calculation:} \path{programs/article_results.py}; \path{outputs/article_results.out}, line 466, identifier \path{SI.T07}.\par
{\footnotesize\textit{Lean:} \path{Comparison.System.row_T07} (\path{lean/Comparison/System/Table/Hadrons.lean}:356)}\\[7pt]
\multicolumn{2}{@{}l}{\textbf{Three rows that are lists of values}}\\*
\hypertarget{row-K04}{}\textbf{K04}\label{row:K04} & \textbf{m\_\allowbreak{}beta and m\_\allowbreak{}betabeta}\par
\textit{Formula:} $m_\beta=\sqrt{\sum_i|U_{ei}|^2c_i^2m_i^2}$, $m_{\beta\beta}=|\sum_iU_{ei}^2c_i^2m_i|$.\par
\textit{Value:} [0.00882841829799682, 0.0033134118567138]\quad\textit{Unit:} $\mathrm{eV}$\quad\textit{Comparison:} < 0.45 eV and < 0.028 eV [KATRIN; KamLAND-Zen]\quad\textit{Residual:} --\par
\textit{Operation:} computed.\par
\textit{Calculation:} \path{programs/article_results.py}; \path{outputs/article_results.out}, line 467, identifier \path{SI.K04}.\par
{\footnotesize\textit{Derivation:} Eq.~\eqref{eq:effectivemasses}.}\\[7pt]
\hypertarget{row-T05}{}\textbf{T05}\label{row:T05} & \textbf{light neutrinos m1,m2}\par
\textit{Formula:} $(m_1,m_2)$, the two lightest roots of Eq.~\eqref{eq:neutralroots}.\par
\textit{Value:} [0, 0.00861267197365246]\quad\textit{Unit:} $\mathrm{eV}$\quad\textit{Comparison:} prediction\quad\textit{Residual:} --\par
\textit{Operation:} computed.\par
\textit{Calculation:} \path{programs/article_results.py}; \path{outputs/article_results.out}, line 468, identifier \path{SI.T05}.\par
{\footnotesize\textit{Derivation:} Eqs.~\eqref{eq:neutralroots} and \eqref{eq:neutrinovalues}.}\\[7pt]
\hypertarget{row-L02}{}\textbf{L02}\label{row:L02} & \textbf{tau row of the lepton current}\par
\textit{Formula:} $(|U_{\tau1}|^2,|U_{\tau2}|^2,|U_{\tau3}|^2)=(b/2,(1-b)/2,1/2)$.\par
\textit{Value:} [0.21875, 0.28125, 0.5]\quad\textit{Comparison:} prediction\quad\textit{Residual:} --\par
\textit{Operation:} computed.\par
\textit{Calculation:} \path{programs/article_results.py}; \path{outputs/article_results.out}, line 469, identifier \path{SI.L02}.\par
{\footnotesize\textit{Derivation:} Eq.~\eqref{eq:taurow}.}\\[7pt]
\end{longtable}
\endgroup

\subsection{The forty-three classical evolution records}\label{sec:classicalregister}
The finite inputs are carried through the conventional maps of Part III. The evaluated values and measurement conventions below are those of the calculation record.
\begingroup\small
\setlength{\tabcolsep}{3pt}\renewcommand{\arraystretch}{1.07}
\begin{longtable}{@{}>{\raggedright\arraybackslash}p{.065\textwidth}>{\raggedright\arraybackslash}p{.90\textwidth}@{}}
\toprule ID&Quantity, calculation and comparison\\\midrule
\endfirsthead\toprule ID&Quantity, calculation and comparison\\\midrule
\endhead\bottomrule\endfoot
\hypertarget{row-I01}{}\textbf{I01}\label{row:I01} & \textbf{up mass m\_\allowbreak{}u(2 GeV)}\par
\textit{Reading:} this closed world line's clock turning the leading zero \ensuremath{\div} the reference clock, the electron's phase on one directed edge or the world line the term names\par
\textit{Formula:} pole: C[four-loop running and matching to 2 GeV] of the up pole at \ensuremath{M_{Z}} (the open levels carry 1/\ensuremath{\eta_0})\par
\textit{Value:} 2.11478\quad\textit{Unit:} $\mathrm{MeV}$\quad\textit{Comparison:} 2.16 +0.04/-0.04 [PDG 2026]\quad\textit{Residual:} -1.13\par
\textit{Operation:} Declared conventional calculation record, with the inputs and reporting convention in Part III.\par
\textit{Calculation:} \path{programs/article_results.py}; \path{outputs/article_results.out}, line 470, identifier \path{SI.I01}.\par
{\footnotesize\textit{Derivation:} the finite inputs and the stated conventional calculation of Part III.}\\[7pt]
\hypertarget{row-I02}{}\textbf{I02}\label{row:I02} & \textbf{down mass m\_\allowbreak{}d(2 GeV)}\par
\textit{Reading:} this closed world line's clock turning the leading zero \ensuremath{\div} the reference clock, the electron's phase on one directed edge or the world line the term names\par
\textit{Formula:} pole: C[four-loop running and matching to 2 GeV] of the down pole at \ensuremath{M_{Z}} (the open levels carry 1/\ensuremath{\eta_0})\par
\textit{Value:} 4.60275\quad\textit{Unit:} $\mathrm{MeV}$\quad\textit{Comparison:} 4.7 +0.04/-0.04 [PDG 2026]\quad\textit{Residual:} -2.43\par
\textit{Operation:} Declared conventional calculation record, with the inputs and reporting convention in Part III.\par
\textit{Calculation:} \path{programs/article_results.py}; \path{outputs/article_results.out}, line 471, identifier \path{SI.I02}.\par
{\footnotesize\textit{Derivation:} the finite inputs and the stated conventional calculation of Part III.}\\[7pt]
\hypertarget{row-I03}{}\textbf{I03}\label{row:I03} & \textbf{strange mass m\_\allowbreak{}s(2 GeV)}\par
\textit{Reading:} this closed world line's clock turning the leading zero \ensuremath{\div} the reference clock, the electron's phase on one directed edge or the world line the term names\par
\textit{Formula:} pole: C[four-loop running and matching to 2 GeV] of the strange pole at \ensuremath{M_{Z}} (the open levels carry 1/\ensuremath{\eta_0})\par
\textit{Value:} 91.8786\quad\textit{Unit:} $\mathrm{MeV}$\quad\textit{Comparison:} 92.9 +0.4/-0.4 [PDG 2026]\quad\textit{Residual:} -2.55\par
\textit{Operation:} Declared conventional calculation record, with the inputs and reporting convention in Part III.\par
\textit{Calculation:} \path{programs/article_results.py}; \path{outputs/article_results.out}, line 472, identifier \path{SI.I03}.\par
{\footnotesize\textit{Derivation:} the finite inputs and the stated conventional calculation of Part III.}\\[7pt]
\hypertarget{row-I04}{}\textbf{I04}\label{row:I04} & \textbf{charm mass m\_\allowbreak{}c(m\_\allowbreak{}c)}\par
\textit{Reading:} this closed world line's clock turning the leading zero \ensuremath{\div} the reference clock, the electron's phase on one directed edge or the world line the term names\par
\textit{Formula:} pole: C[four-loop running and matching to \ensuremath{m_{c}}] of the charm pole at \ensuremath{M_{Z}} (the open levels carry 1/\ensuremath{\eta_0})\par
\textit{Value:} 1.26243\quad\textit{Unit:} $\mathrm{GeV}$\quad\textit{Comparison:} 1.2729 +0.0027/-0.0027 [PDG 2026]\quad\textit{Residual:} -3.88\par
\textit{Operation:} Declared conventional calculation record, with the inputs and reporting convention in Part III.\par
\textit{Calculation:} \path{programs/article_results.py}; \path{outputs/article_results.out}, line 473, identifier \path{SI.I04}.\par
{\footnotesize\textit{Derivation:} the finite inputs and the stated conventional calculation of Part III.}\\[7pt]
\hypertarget{row-I05}{}\textbf{I05}\label{row:I05} & \textbf{bottom mass m\_\allowbreak{}b(m\_\allowbreak{}b)}\par
\textit{Reading:} this closed world line's clock turning the leading zero \ensuremath{\div} the reference clock, the electron's phase on one directed edge or the world line the term names\par
\textit{Formula:} pole: C[four-loop running and matching to \ensuremath{m_{b}}] of the bottom pole at \ensuremath{M_{Z}}\par
\textit{Value:} 4.18248214152943\quad\textit{Unit:} $\mathrm{GeV}$\quad\textit{Comparison:} 4.1859 +0.0034/-0.0034 [PDG 2026]\quad\textit{Residual:} -1.005\par
\textit{Operation:} Declared conventional calculation record, with the inputs and reporting convention in Part III.\par
\textit{Calculation:} \path{programs/article_results.py}; \path{outputs/article_results.out}, line 474, identifier \path{SI.I05}.\par
{\footnotesize\textit{Derivation:} the finite inputs and the stated conventional calculation of Part III.}\\[7pt]
\hypertarget{row-T10}{}\textbf{T10}\label{row:T10} & \textbf{coloured/lepton C\_\allowbreak{}d}\par
\textit{Reading:} this closed world line's clock turning the leading zero \ensuremath{\div} the reference clock, the electron's phase on one directed edge or the world line the term names\par
\textit{Formula:} pole: C[the coloured/lepton statistic] of the down pole (the open levels carry 1/\ensuremath{\eta_0})\par
\textit{Value:} 10.4618\quad\textit{Comparison:} quark-mass statistic\quad\textit{Residual:} --\par
\textit{Operation:} Declared conventional calculation record, with the inputs and reporting convention in Part III.\par
\textit{Calculation:} \path{programs/article_results.py}; \path{outputs/article_results.out}, line 475, identifier \path{SI.T10}.\par
{\footnotesize\textit{Derivation:} the finite inputs and the stated conventional calculation of Part III.}\\[7pt]
\hypertarget{row-T11}{}\textbf{T11}\label{row:T11} & \textbf{coloured/lepton C\_\allowbreak{}s}\par
\textit{Reading:} this closed world line's clock turning the leading zero \ensuremath{\div} the reference clock, the electron's phase on one directed edge or the world line the term names\par
\textit{Formula:} pole: C[the coloured/lepton statistic] of the strange pole (the open levels carry 1/\ensuremath{\eta_0})\par
\textit{Value:} 10.4697\quad\textit{Comparison:} quark-mass statistic\quad\textit{Residual:} --\par
\textit{Operation:} Declared conventional calculation record, with the inputs and reporting convention in Part III.\par
\textit{Calculation:} \path{programs/article_results.py}; \path{outputs/article_results.out}, line 476, identifier \path{SI.T11}.\par
{\footnotesize\textit{Derivation:} the finite inputs and the stated conventional calculation of Part III.}\\[7pt]
\hypertarget{row-T12}{}\textbf{T12}\label{row:T12} & \textbf{coloured/lepton C\_\allowbreak{}b}\par
\textit{Reading:} this closed world line's clock turning the leading zero \ensuremath{\div} the reference clock, the electron's phase on one directed edge or the world line the term names\par
\textit{Formula:} pole: C[the coloured/lepton statistic] of the bottom pole\par
\textit{Value:} 10.3757734611098\quad\textit{Comparison:} quark-mass statistic\quad\textit{Residual:} --\par
\textit{Operation:} Declared conventional calculation record, with the inputs and reporting convention in Part III.\par
\textit{Calculation:} \path{programs/article_results.py}; \path{outputs/article_results.out}, line 477, identifier \path{SI.T12}.\par
{\footnotesize\textit{Derivation:} the finite inputs and the stated conventional calculation of Part III.}\\[7pt]
\hypertarget{row-S10}{}\textbf{S10}\label{row:S10} & \textbf{coloured/lepton Cd}\par
\textit{Reading:} this closed world line's clock turning the leading zero \ensuremath{\div} the reference clock, the electron's phase on one directed edge or the world line the term names\par
\textit{Formula:} pole: C[the coloured/lepton statistic] of the down pole (the open levels carry 1/\ensuremath{\eta_0})\par
\textit{Value:} 10.4618\quad\textit{Comparison:} 10.67 +0.16/-0.16 [quark-mass statistic]\quad\textit{Residual:} -1.30\par
\textit{Operation:} Declared conventional calculation record, with the inputs and reporting convention in Part III.\par
\textit{Calculation:} \path{programs/article_results.py}; \path{outputs/article_results.out}, line 478, identifier \path{SI.S10}.\par
{\footnotesize\textit{Derivation:} the finite inputs and the stated conventional calculation of Part III.}\\[7pt]
\hypertarget{row-S11}{}\textbf{S11}\label{row:S11} & \textbf{coloured/lepton Cs}\par
\textit{Reading:} this closed world line's clock turning the leading zero \ensuremath{\div} the reference clock, the electron's phase on one directed edge or the world line the term names\par
\textit{Formula:} pole: C[the coloured/lepton statistic] of the strange pole (the open levels carry 1/\ensuremath{\eta_0})\par
\textit{Value:} 10.4697\quad\textit{Comparison:} 10.57 +0.08/-0.08 [quark-mass statistic]\quad\textit{Residual:} -1.25\par
\textit{Operation:} Declared conventional calculation record, with the inputs and reporting convention in Part III.\par
\textit{Calculation:} \path{programs/article_results.py}; \path{outputs/article_results.out}, line 479, identifier \path{SI.S11}.\par
{\footnotesize\textit{Derivation:} the finite inputs and the stated conventional calculation of Part III.}\\[7pt]
\hypertarget{row-S12}{}\textbf{S12}\label{row:S12} & \textbf{coloured/lepton Cb}\par
\textit{Reading:} this closed world line's clock turning the leading zero \ensuremath{\div} the reference clock, the electron's phase on one directed edge or the world line the term names\par
\textit{Formula:} pole: C[the coloured/lepton statistic] of the bottom pole\par
\textit{Value:} 10.3757734611098\quad\textit{Comparison:} 10.38 +0.02/-0.02 [quark-mass statistic]\quad\textit{Residual:} -0.211\par
\textit{Operation:} Declared conventional calculation record, with the inputs and reporting convention in Part III.\par
\textit{Calculation:} \path{programs/article_results.py}; \path{outputs/article_results.out}, line 480, identifier \path{SI.S12}.\par
{\footnotesize\textit{Derivation:} the finite inputs and the stated conventional calculation of Part III.}\\[7pt]
\hypertarget{row-I06}{}\textbf{I06}\label{row:I06} & \textbf{strong coupling at the proton}\par
\textit{Reading:} the response of the rotations within the two sides to the curvature on a square cell \ensuremath{\div} its normalized form\par
\textit{Formula:} coefficient: C[four-loop running to \ensuremath{m_{p}}] of the colour two-source coefficient \ensuremath{\alpha_{s}}\par
\textit{Value:} 0.505523057661309\quad\textit{Comparison:} 0.50927 +0.0233/-0.0213 [four-loop evolution of the PDG 2024 alpha\_\allowbreak{}s(M\_\allowbreak{}Z) = 0.1180(9)]\quad\textit{Residual:} -0.176\par
\textit{Operation:} Declared conventional calculation record, with the inputs and reporting convention in Part III.\par
\textit{Calculation:} \path{programs/article_results.py}; \path{outputs/article_results.out}, line 481, identifier \path{SI.I06}.\par
{\footnotesize\textit{Derivation:} the finite inputs and the stated conventional calculation of Part III.}\\[7pt]
\hypertarget{row-S018}{}\textbf{S018}\label{row:S018} & \textbf{Lambda\_\allowbreak{}5, four-loop convention at M\_\allowbreak{}Z = 91.188 GeV}\par
\textit{Reading:} the response of the rotations within the two sides to the curvature on a square cell \ensuremath{\div} its normalized form\par
\textit{Formula:} coefficient: C[the four-loop length at \ensuremath{M_{Z}} = 91.188 GeV] of \ensuremath{\alpha_{s}}\par
\textit{Value:} 206.896190226441\quad\textit{Unit:} $\mathrm{MeV}$\quad\textit{Comparison:} four-loop convention: 206.90 MeV\quad\textit{Residual:} --\par
\textit{Operation:} Classical evolution of the stated finite inputs.\par
\textit{Calculation:} \path{programs/article_results.py}; \path{outputs/article_results.out}, line 482, identifier \path{SI.S018}.\par
{\footnotesize\textit{Derivation:} the finite inputs and the stated conventional calculation of Part III.}\\[7pt]
\hypertarget{row-Q01}{}\textbf{Q01}\label{row:Q01} & \textbf{electric running 1/alpha(M\_\allowbreak{}Z)}\par
\textit{Reading:} the coupling carried through the stated running \ensuremath{\div} its source value\par
\textit{Formula:} C[dispersion] of the computed vector current\par
\textit{Value:} 128.937243\quad\textit{Comparison:} 128.94 +0.014/-0.014 [on shell]\quad\textit{Residual:} -0.197\par
\textit{Operation:} Declared conventional calculation record, with the inputs and reporting convention in Part III.\par
\textit{Calculation:} \path{programs/article_results.py}; \path{outputs/article_results.out}, line 483, identifier \path{SI.Q01}.\par
{\footnotesize\textit{Derivation:} the finite inputs and the stated conventional calculation of Part III.}\\[7pt]
\hypertarget{row-Q11}{}\textbf{Q11}\label{row:Q11} & \textbf{proton decay matrix element}\par
\textit{Reading:} the proton-to-pion matrix element of the operator the rotations mixing the split carry \ensuremath{\div} its reference\par
\textit{Formula:} pole: C[continuous four-loop operator transport, computed bottom threshold] of the poles \ensuremath{m_{p}} and \ensuremath{m_{\pi}}0\par
\textit{Value:} 0.134483885087514\quad\textit{Unit:} $\mathrm{GeV}^{2}$\quad\textit{Comparison:} 0.131 +0.0136/-0.0136 [lattice 2017]\quad\textit{Residual:} +0.256\par
\textit{Operation:} Declared conventional calculation record, with the inputs and reporting convention in Part III.\par
\textit{Calculation:} \path{programs/article_results.py}; \path{outputs/article_results.out}, line 484, identifier \path{SI.Q11}.\par
{\footnotesize\textit{Derivation:} the finite inputs and the stated conventional calculation of Part III.}\\[7pt]
\hypertarget{row-T09}{}\textbf{T09}\label{row:T09} & \textbf{alpha\_\allowbreak{}s at bottom and charm thresholds}\par
\textit{Reading:} the response of the rotations within the two sides to the curvature on a square cell \ensuremath{\div} its normalized form\par
\textit{Formula:} coefficient: C[running to \ensuremath{m_{b}} and \ensuremath{m_{c}}] of \ensuremath{\alpha_{s}}\par
\textit{Value:} [0.224009369316102, 0.385356820505129]\quad\textit{Comparison:} running convention\quad\textit{Residual:} --\par
\textit{Operation:} Declared conventional calculation record, with the inputs and reporting convention in Part III.\par
\textit{Calculation:} \path{programs/article_results.py}; \path{outputs/article_results.out}, line 485, identifier \path{SI.T09}.\par
{\footnotesize\textit{Derivation:} the finite inputs and the stated conventional calculation of Part III.}\\[7pt]
\hypertarget{row-S14}{}\textbf{S14}\label{row:S14} & \textbf{five-flavour four-loop scale}\par
\textit{Reading:} the response of the rotations within the two sides to the curvature on a square cell \ensuremath{\div} its normalized form\par
\textit{Formula:} coefficient: C[the four-loop length at the constructed \ensuremath{M_{Z}}] of \ensuremath{\alpha_{s}}\par
\textit{Value:} 206.890203783025\quad\textit{Unit:} $\mathrm{MeV}$\quad\textit{Comparison:} four-loop convention\quad\textit{Residual:} --\par
\textit{Operation:} Classical evolution of the stated finite inputs.\par
\textit{Calculation:} \path{programs/article_results.py}; \path{outputs/article_results.out}, line 486, identifier \path{SI.S14}.\par
{\footnotesize\textit{Derivation:} the finite inputs and the stated conventional calculation of Part III.}\\[7pt]
\hypertarget{row-S017}{}\textbf{S017}\label{row:S017} & \textbf{five-flavour scale at the reporting M\_\allowbreak{}Z = 91.1876 GeV}\par
\textit{Reading:} the response of the rotations within the two sides to the curvature on a square cell \ensuremath{\div} its normalized form\par
\textit{Formula:} coefficient: C[the four-loop length at \ensuremath{M_{Z}} = 91.1876 GeV] of \ensuremath{\alpha_{s}}\par
\textit{Value:} 206.895282667594\quad\textit{Unit:} $\mathrm{MeV}$\quad\textit{Comparison:} four-loop convention at M\_\allowbreak{}Z = 91.1876 GeV: 206.9 MeV\quad\textit{Residual:} --\par
\textit{Operation:} Classical evolution of the stated finite inputs.\par
\textit{Calculation:} \path{programs/article_results.py}; \path{outputs/article_results.out}, line 487, identifier \path{SI.S017}.\par
{\footnotesize\textit{Derivation:} the finite inputs and the stated conventional calculation of Part III.}\\[7pt]
\hypertarget{row-I10}{}\textbf{I10}\label{row:I10} & \textbf{electron anomaly a\_\allowbreak{}e}\par
\textit{Reading:} the current of the four amplitudes at the pole \ensuremath{\div} the reference current: the nuclear magneton, the Compton length, the electron vertex\par
\textit{Formula:} residue: C[the electron vertex to five loops] of the residue \ensuremath{\alpha} and the charged poles; with the computed hadronic share\par
\textit{Value:} 0.00115965218050222\quad\textit{Comparison:} 0.00115965218059 +1.3e-13/-1.3e-13 [Fan et al., Phys. Rev. Lett. 130, 071801 (2023)]\quad\textit{Residual:} -0.675\par
\textit{Operation:} Declared conventional calculation record, with the inputs and reporting convention in Part III.\par
\textit{Calculation:} \path{programs/article_results.py}; \path{outputs/article_results.out}, line 488, identifier \path{SI.I10}.\par
{\footnotesize\textit{Derivation:} the finite inputs and the stated conventional calculation of Part III.}\\[7pt]
\hypertarget{row-I11}{}\textbf{I11}\label{row:I11} & \textbf{hydrogen binding}\par
\textit{Reading:} the current of the four amplitudes at the pole \ensuremath{\div} the reference current: the nuclear magneton, the Compton length, the electron vertex\par
\textit{Formula:} residue: C[the hydrogen series] of the residue \ensuremath{\alpha} and the pole \ensuremath{m_{p}}/\ensuremath{m_{e}}; with the computed hadronic term; compared at the spectral lines\par
\textit{Value:} 13.5984345996929\quad\textit{Unit:} $\mathrm{eV}$\quad\textit{Comparison:} 13.5984345997 (relative difference -6.710e-13)\quad\textit{Residual:} --\par
\textit{Operation:} Declared conventional calculation record, with the inputs and reporting convention in Part III.\par
\textit{Calculation:} \path{programs/article_results.py}; \path{outputs/article_results.out}, line 489, identifier \path{SI.I11}.\par
{\footnotesize\textit{Derivation:} the finite inputs and the stated conventional calculation of Part III.}\\[7pt]
\hypertarget{row-S06}{}\textbf{S06}\label{row:S06} & \textbf{muonic-hydrogen singlet capture}\par
\textit{Reading:} the amplitudes' current carried through the stated classical map \ensuremath{\div} its reference\par
\textit{Formula:} the whole axial shape w/(4 - 3w) with its Ward partner, through C[the capture map]\par
\textit{Value:} 717.110240723118\quad\textit{Unit:} $\mathrm{s}^{-1}$\quad\textit{Comparison:} 715.6 +7.4/-7.4 [MuCap]\quad\textit{Residual:} +0.204\par
\textit{Operation:} Declared conventional calculation record, with the inputs and reporting convention in Part III.\par
\textit{Calculation:} \path{programs/article_results.py}; \path{outputs/article_results.out}, line 490, identifier \path{SI.S06}.\par
{\footnotesize\textit{Derivation:} the finite inputs and the stated conventional calculation of Part III.}\\[7pt]
\hypertarget{row-S15}{}\textbf{S15}\label{row:S15} & \textbf{pion electronic/muonic decay ratio}\par
\textit{Reading:} the current of the four amplitudes at the pole \ensuremath{\div} the reference current: the nuclear magneton, the Compton length, the electron vertex\par
\textit{Formula:} residue: C[inclusive point-pseudoscalar QED through O(\ensuremath{\alpha})] of the charged poles \ensuremath{m_{e}}, \ensuremath{m_{\mu}}, \ensuremath{m_{\pi}}\ensuremath{\pm}\par
\textit{Value:} 0.000123292613490236\quad\textit{Comparison:} 0.00012327 +2.3e-07/-2.3e-07 [PDG 2024 inclusive ratio]\quad\textit{Residual:} +0.098\par
\textit{Operation:} Declared conventional calculation record, with the inputs and reporting convention in Part III.\par
\textit{Calculation:} \path{programs/article_results.py}; \path{outputs/article_results.out}, line 491, identifier \path{SI.S15}.\par
{\footnotesize\textit{Derivation:} the finite inputs and the stated conventional calculation of Part III.}\\[7pt]
\hypertarget{row-A001}{}\textbf{A001}\label{row:A001} & \textbf{induced pseudoscalar coupling at Q\textasciicircum{}2=0.88 m\_\allowbreak{}mu\textasciicircum{}2}\par
\textit{Reading:} the amplitudes' current carried through the stated classical map \ensuremath{\div} its reference\par
\textit{Formula:} the whole axial shape w/(4 - 3w) with its Ward partner at Q\ensuremath{{}^{2}} = 0.88 \ensuremath{m_{\mu}}\ensuremath{{}^{2}}\par
\textit{Value:} 8.36259039416189\quad\textit{Comparison:} 8.99 +0.626/-0.626 [ETMC 2024]\quad\textit{Residual:} -1.002\par
\textit{Operation:} Declared conventional calculation record, with the inputs and reporting convention in Part III.\par
\textit{Calculation:} \path{programs/article_results.py}; \path{outputs/article_results.out}, line 492, identifier \path{SI.A001}.\par
{\footnotesize\textit{Derivation:} the finite inputs and the stated conventional calculation of Part III.}\\[7pt]
\hypertarget{row-Q04}{}\textbf{Q04}\label{row:Q04} & \textbf{muon hadronic polarization}\par
\textit{Reading:} the amplitudes' current carried through the stated classical map \ensuremath{\div} its reference\par
\textit{Formula:} the computed vector current's leading polarization of the muon\par
\textit{Value:} 711.1437\quad\textit{Unit:} $10^{-10}$\quad\textit{Comparison:} 713.2 +6.1/-6.1 [2025 lattice]\quad\textit{Residual:} -0.337\par
\textit{Operation:} Declared conventional calculation record, with the inputs and reporting convention in Part III.\par
\textit{Calculation:} \path{programs/article_results.py}; \path{outputs/article_results.out}, line 493, identifier \path{SI.Q04}.\par
{\footnotesize\textit{Derivation:} the finite inputs and the stated conventional calculation of Part III.}\\[7pt]
\hypertarget{row-I07}{}\textbf{I07}\label{row:I07} & \textbf{top pole mass}\par
\textit{Reading:} the rotations that move h, its radial change or its heavy partner \ensuremath{\div} the length v of h\par
\textit{Formula:} pole: C[four-loop matching to the pole mass] of the top pole at \ensuremath{M_{Z}}\par
\textit{Value:} 172.76121248356\quad\textit{Unit:} $\mathrm{GeV}$\quad\textit{Comparison:} 172.1 +0.6/-0.6 [PDG 2026 top pole mass from cross sections]\quad\textit{Residual:} +1.102\par
\textit{Operation:} Declared conventional calculation record, with the inputs and reporting convention in Part III.\par
\textit{Calculation:} \path{programs/article_results.py}; \path{outputs/article_results.out}, line 494, identifier \path{SI.I07}.\par
{\footnotesize\textit{Derivation:} the finite inputs and the stated conventional calculation of Part III.}\\[7pt]
\hypertarget{row-T08}{}\textbf{T08}\label{row:T08} & \textbf{running top and Yukawa}\par
\textit{Reading:} the rotations that move h, its radial change or its heavy partner \ensuremath{\div} the length v of h\par
\textit{Formula:} pole: C[running to \ensuremath{m_{t}}(\ensuremath{m_{t}}), \ensuremath{y_{t}} = \ensuremath{\surd}2 \ensuremath{m_{t}}/v] of the top pole\par
\textit{Value:} [162.984094960692, 0.936132971038294]\quad\textit{Unit:} $\mathrm{GeV},\ 1$\quad\textit{Comparison:} running convention\quad\textit{Residual:} --\par
\textit{Operation:} Declared conventional calculation record, with the inputs and reporting convention in Part III.\par
\textit{Calculation:} \path{programs/article_results.py}; \path{outputs/article_results.out}, line 495, identifier \path{SI.T08}.\par
{\footnotesize\textit{Derivation:} the finite inputs and the stated conventional calculation of Part III.}\\[7pt]
\hypertarget{row-L21}{}\textbf{L21}\label{row:L21} & \textbf{clock-independent lower bound on the light neutral mass sum}\par
\textit{Reading:} the rotations that move h, its radial change or its heavy partner \ensuremath{\div} the length v of h\par
\textit{Formula:} pole: sum of the light poles at the clock endpoint \ensuremath{h_\nu} = 1\par
\textit{Value:} 0.0587890893281977\quad\textit{Unit:} $\mathrm{eV}$\quad\textit{Comparison:} bound for every clock\quad\textit{Residual:} --\par
\textit{Operation:} Classical evolution of the stated finite inputs.\par
\textit{Calculation:} \path{programs/article_results.py}; \path{outputs/article_results.out}, line 496, identifier \path{SI.L21}.\par
{\footnotesize\textit{Derivation:} the finite inputs and the stated conventional calculation of Part III.}\\[7pt]
\hypertarget{row-F02}{}\textbf{F02}\label{row:F02} & \textbf{cold/baryon omega\_\allowbreak{}c/omega\_\allowbreak{}b}\par
\textit{Reading:} the stationary weight on a type of state along the clock \ensuremath{\div} the reference weight\par
\textit{Formula:} stationary: stationary weight on all sixteen states with the heavy pole: a\ensuremath{{}^{16}}\ensuremath{M_{R}}/\ensuremath{m_{p}}, then C[\ensuremath{\kappa_{B}}] (the epoch with the nucleosynthesis network's \ensuremath{\kappa_{B}})\par
\textit{Value:} 5.34277169492503\quad\textit{Comparison:} 5.3643 +0.0646/-0.0646 [Planck 2018 ratio]\quad\textit{Residual:} -0.333\par
\textit{Operation:} Finite cold reading divided by the declared network composition of Sec.~\ref{sec:composition}.\par
\textit{Calculation:} \path{programs/article_results.py}; \path{outputs/article_results.out}, line 497, identifier \path{SI.F02}.\par
{\footnotesize\textit{Derivation:} the finite inputs and the stated conventional calculation of Part III.}\\[7pt]
\hypertarget{row-F05}{}\textbf{F05}\label{row:F05} & \textbf{photon temperature at the 24:11 epoch}\par
\textit{Reading:} the stationary weight on a type of state along the clock \ensuremath{\div} the reference weight\par
\textit{Formula:} stationary: the present, \ensuremath{\rho_{\Lambda}}/\ensuremath{\rho_{m}} = 24/11, then C[thermal history: T] (the epoch with the nucleosynthesis network's \ensuremath{\kappa_{B}})\par
\textit{Value:} 2.72549728670312\quad\textit{Unit:} $\mathrm{K}$\quad\textit{Comparison:} 2.72548 +0.00057/-0.00057 [FIRAS, Fixsen, Astrophys. J. 707, 916 (2009)]\quad\textit{Residual:} +0.030\par
\textit{Operation:} Declared conventional calculation record, with the inputs and reporting convention in Part III.\par
\textit{Calculation:} \path{programs/article_results.py}; \path{outputs/article_results.out}, line 498, identifier \path{SI.F05}.\par
{\footnotesize\textit{Derivation:} the finite inputs and the stated conventional calculation of Part III.}\\[7pt]
\hypertarget{row-F06}{}\textbf{F06}\label{row:F06} & \textbf{Hubble rate $H_{\rm comp}$}\par
\textit{Reading:} the stationary weight on a type of state along the clock \ensuremath{\div} the reference weight\par
\textit{Formula:} The thermal expansion $H_{\rm comp}$ at $\rho_\Lambda/\rho_m=24/11$, with the declared network composition.\par
\textit{Value:} 67.4316192077386\quad\textit{Unit:} $\mathrm{km}\,\mathrm{s}^{-1}\,\mathrm{Mpc}^{-1}$\quad\textit{Comparison:} 67.62 +0.5/-0.5 [P-ACT]\quad\textit{Residual:} -0.377\par
\textit{Operation:} Declared conventional calculation record, with the inputs and reporting convention in Part III.\par
\textit{Calculation:} \path{programs/article_results.py}; \path{outputs/article_results.out}, line 499, identifier \path{SI.F06}.\par
{\footnotesize\textit{Derivation:} the finite inputs and the stated conventional calculation of Part III.}\\[7pt]
\hypertarget{row-F07}{}\textbf{F07}\label{row:F07} & \textbf{age $t_{\rm age}$}\par
\textit{Reading:} the stationary weight on a type of state along the clock \ensuremath{\div} the reference weight\par
\textit{Formula:} The thermal age $t_{\rm age}$ at $\rho_\Lambda/\rho_m=24/11$, with the declared network composition.\par
\textit{Value:} 13.7944222597347\quad\textit{Unit:} $\mathrm{Gyr}$\quad\textit{Comparison:} 13.789 +0.018/-0.018 [P-ACT]\quad\textit{Residual:} +0.301\par
\textit{Operation:} Declared conventional calculation record, with the inputs and reporting convention in Part III.\par
\textit{Calculation:} \path{programs/article_results.py}; \path{outputs/article_results.out}, line 500, identifier \path{SI.F07}.\par
{\footnotesize\textit{Derivation:} the finite inputs and the stated conventional calculation of Part III.}\\[7pt]
\hypertarget{row-F08}{}\textbf{F08}\label{row:F08} & \textbf{matter share Omega\_\allowbreak{}m}\par
\textit{Reading:} the stationary weight on a type of state along the clock \ensuremath{\div} the reference weight\par
\textit{Formula:} stationary: the present, 24/11, then C[thermal history: \ensuremath{\Omega_{m}}] (the epoch with the nucleosynthesis network's \ensuremath{\kappa_{B}})\par
\textit{Value:} 0.314264682713255\quad\textit{Comparison:} 0.3116 +0.0071/-0.0071 [P-ACT]\quad\textit{Residual:} +0.375\par
\textit{Operation:} Declared conventional calculation record, with the inputs and reporting convention in Part III.\par
\textit{Calculation:} \path{programs/article_results.py}; \path{outputs/article_results.out}, line 501, identifier \path{SI.F08}.\par
{\footnotesize\textit{Derivation:} the finite inputs and the stated conventional calculation of Part III.}\\[7pt]
\hypertarget{row-F09}{}\textbf{F09}\label{row:F09} & \textbf{the \ensuremath{\Lambda} share Omega\_\allowbreak{}Lambda}\par
\textit{Reading:} the stationary weight on a type of state along the clock \ensuremath{\div} the reference weight\par
\textit{Formula:} stationary: the present, 24/11, then C[thermal history: \ensuremath{\Omega_{\Lambda}}] (the epoch with the nucleosynthesis network's \ensuremath{\kappa_{B}})\par
\textit{Value:} 0.685668398647101\quad\textit{Comparison:} 0.6847 +0.0073/-0.0073 [Planck 2018]\quad\textit{Residual:} +0.133\par
\textit{Operation:} Declared conventional calculation record, with the inputs and reporting convention in Part III.\par
\textit{Calculation:} \path{programs/article_results.py}; \path{outputs/article_results.out}, line 502, identifier \path{SI.F09}.\par
{\footnotesize\textit{Derivation:} the finite inputs and the stated conventional calculation of Part III.}\\[7pt]
\hypertarget{row-F10}{}\textbf{F10}\label{row:F10} & \textbf{helium mass fraction Y\_\allowbreak{}p}\par
\textit{Reading:} the stationary weight on a type of state along the clock \ensuremath{\div} the reference weight\par
\textit{Formula:} stationary: the nucleosynthesis network on the computed coefficients, \ensuremath{Y_{p}} = 4y/(1 + 4y) in the measurements' convention\par
\textit{Value:} 0.24694744653623\quad\textit{Comparison:} 0.2458 +0.0013/-0.0013 [LBT 2026, Y\_\allowbreak{}p = 4y/(1 + 4y) from He I/H I lines]\quad\textit{Residual:} +0.883\par
\textit{Operation:} Declared conventional calculation record, with the inputs and reporting convention in Part III.\par
\textit{Calculation:} \path{programs/article_results.py}; \path{outputs/article_results.out}, line 503, identifier \path{SI.F10}.\par
{\footnotesize\textit{Derivation:} the finite inputs and the stated conventional calculation of Part III.}\\[7pt]
\hypertarget{row-F11}{}\textbf{F11}\label{row:F11} & \textbf{baryon density omega\_\allowbreak{}b}\par
\textit{Reading:} the stationary weight on a type of state along the clock \ensuremath{\div} the reference weight\par
\textit{Formula:} stationary: the present, 24/11, then C[thermal history: \ensuremath{\omega_{b}}] (the epoch with the nucleosynthesis network's \ensuremath{\kappa_{B}})\par
\textit{Value:} 0.0224294823835265\quad\textit{Comparison:} 0.0225 +0.00011/-0.00011 [P-ACT]\quad\textit{Residual:} -0.641\par
\textit{Operation:} Declared conventional calculation record, with the inputs and reporting convention in Part III.\par
\textit{Calculation:} \path{programs/article_results.py}; \path{outputs/article_results.out}, line 504, identifier \path{SI.F11}.\par
{\footnotesize\textit{Derivation:} the finite inputs and the stated conventional calculation of Part III.}\\[7pt]
\hypertarget{row-F12}{}\textbf{F12}\label{row:F12} & \textbf{cold density omega\_\allowbreak{}c}\par
\textit{Reading:} the stationary weight on a type of state along the clock \ensuremath{\div} the reference weight\par
\textit{Formula:} stationary: the present, 24/11, then C[thermal history: \ensuremath{\omega_{c}}] (the epoch with the nucleosynthesis network's \ensuremath{\kappa_{B}})\par
\textit{Value:} 0.119835603610525\quad\textit{Comparison:} 0.1193 +0.0012/-0.0012 [P-ACT]\quad\textit{Residual:} +0.446\par
\textit{Operation:} Declared conventional calculation record, with the inputs and reporting convention in Part III.\par
\textit{Calculation:} \path{programs/article_results.py}; \path{outputs/article_results.out}, line 505, identifier \path{SI.F12}.\par
{\footnotesize\textit{Derivation:} the finite inputs and the stated conventional calculation of Part III.}\\[7pt]
\hypertarget{row-Q05}{}\textbf{Q05}\label{row:Q05} & \textbf{acoustic angle 100 theta\_\allowbreak{}* (common-area convention)}\par
\textit{Reading:} the stationary weight on a type of state along the clock \ensuremath{\div} the reference weight\par
\textit{Formula:} stationary: C[CAMB transfer, common-area convention] of the populations at the present\par
\textit{Value:} 1.04094999030065\quad\textit{Comparison:} 1.04094 +0.00025/-0.00025 [P-ACT]\quad\textit{Residual:} +0.040\par
\textit{Operation:} Declared conventional calculation record, with the inputs and reporting convention in Part III.\par
\textit{Calculation:} \path{programs/article_results.py}; \path{outputs/article_results.out}, line 506, identifier \path{SI.Q05}.\par
{\footnotesize\textit{Derivation:} the finite inputs and the stated conventional calculation of Part III.}\\[7pt]
\hypertarget{row-Q06}{}\textbf{Q06}\label{row:Q06} & \textbf{acoustic angle at optical depth one 100 theta}\par
\textit{Reading:} the same populations carried through the stated classical transport\par
\textit{Formula:} C[CAMB transport with the nucleosynthesis network's helium] of the populations\par
\textit{Value:} 1.04098061206\quad\textit{Comparison:} 1.04094 +0.00025/-0.00025 [P-ACT]\quad\textit{Residual:} +0.162\par
\textit{Operation:} Declared conventional calculation record, with the inputs and reporting convention in Part III.\par
\textit{Calculation:} \path{programs/article_results.py}; \path{outputs/article_results.out}, line 507, identifier \path{SI.Q06}.\par
{\footnotesize\textit{Derivation:} the finite inputs and the stated conventional calculation of Part III.}\\[7pt]
\hypertarget{row-Q07}{}\textbf{Q07}\label{row:Q07} & \textbf{acoustic angle at visibility peak 100 theta}\par
\textit{Reading:} the same populations carried through the stated classical transport\par
\textit{Formula:} C[CAMB transport with the nucleosynthesis network's helium] of the populations\par
\textit{Value:} 1.041686318178\quad\textit{Comparison:} 1.04161 +0.00023/-0.00023 [CMB-SPA]\quad\textit{Residual:} +0.332\par
\textit{Operation:} Declared conventional calculation record, with the inputs and reporting convention in Part III.\par
\textit{Calculation:} \path{programs/article_results.py}; \path{outputs/article_results.out}, line 508, identifier \path{SI.Q07}.\par
{\footnotesize\textit{Derivation:} the finite inputs and the stated conventional calculation of Part III.}\\[7pt]
\hypertarget{row-Q08}{}\textbf{Q08}\label{row:Q08} & \textbf{matter fluctuation sigma\_\allowbreak{}8}\par
\textit{Reading:} the same populations carried through the stated classical transport\par
\textit{Formula:} C[CAMB transport with the nucleosynthesis network's helium] of the populations\par
\textit{Value:} 0.8213353983338\quad\textit{Comparison:} 0.8149 +0.0063/-0.0063 [P-ACT]\quad\textit{Residual:} +1.021\par
\textit{Operation:} Declared conventional calculation record, with the inputs and reporting convention in Part III.\par
\textit{Calculation:} \path{programs/article_results.py}; \path{outputs/article_results.out}, line 509, identifier \path{SI.Q08}.\par
{\footnotesize\textit{Derivation:} the finite inputs and the stated conventional calculation of Part III.}\\[7pt]
\hypertarget{row-Q09}{}\textbf{Q09}\label{row:Q09} & \textbf{S\_\allowbreak{}8}\par
\textit{Reading:} the same populations carried through the stated classical transport\par
\textit{Formula:} \ensuremath{\sigma}\ensuremath{{}_{8}}(\ensuremath{\Omega_{m}}/0.3)\ensuremath{{}^{1/2}} of the carried \ensuremath{\sigma}\ensuremath{{}_{8}}\par
\textit{Value:} 0.840635453047808\quad\textit{Comparison:} 0.83 +0.014/-0.014 [P-ACT]\quad\textit{Residual:} +0.760\par
\textit{Operation:} Declared conventional calculation record, with the inputs and reporting convention in Part III.\par
\textit{Calculation:} \path{programs/article_results.py}; \path{outputs/article_results.out}, line 510, identifier \path{SI.Q09}.\par
{\footnotesize\textit{Derivation:} the finite inputs and the stated conventional calculation of Part III.}\\[7pt]
\hypertarget{row-Q10}{}\textbf{Q10}\label{row:Q10} & \textbf{drag ruler r\_\allowbreak{}d}\par
\textit{Reading:} the same populations carried through the stated classical transport\par
\textit{Formula:} C[CAMB transport with the nucleosynthesis network's helium] of the populations\par
\textit{Value:} 147.0826558961\quad\textit{Unit:} $\mathrm{Mpc}$\quad\textit{Comparison:} 147.14 +0.29/-0.29 [P-ACT]\quad\textit{Residual:} -0.198\par
\textit{Operation:} Declared conventional calculation record, with the inputs and reporting convention in Part III.\par
\textit{Calculation:} \path{programs/article_results.py}; \path{outputs/article_results.out}, line 511, identifier \path{SI.Q10}.\par
{\footnotesize\textit{Derivation:} the finite inputs and the stated conventional calculation of Part III.}\\[7pt]
\hypertarget{row-Q12}{}\textbf{Q12}\label{row:Q12} & \textbf{Muon magnetic-frequency ratio $R'_\mu$}\par
\textit{Formula:} $R'_\mu=a_\mu/[\mu_p(m_e/m_p)(1-\sigma'_p)(m_\mu/m_e)]$\par
\textit{Value:} 0.003707299080299\quad\textit{Comparison:} $0.00370730088(46)$ [Muon $g-2$ final result; shielded protons in spherical water at $25^\circ\mathrm C$, Ref.~\cite{MuonG2Final}]\quad\textit{Residual:} $-3.90$ with the recorded $4.62\times10^{-10}$ width\par
\textit{Operation:} Declared conventional calculation record, with the inputs and reporting convention in Part III.\par
\textit{Calculation:} \path{programs/article_results.py}; \path{outputs/article_results.out}, line 512, identifier \path{SI.Q12}.\par
{\footnotesize\textit{Derivation:} the finite inputs and the stated conventional calculation of Part III.}\\[7pt]
\end{longtable}
\endgroup

\subsection{Thirty-eight further composite comparison records}\label{sec:analyticregister}
The twelve explicitly evaluable meson and baryon expressions are calculated in Sec.~\ref{sec:literalcompositesSI}; their values and large residuals appear here. The other twenty-six entries carry the source expression and measurement in the comparison record.
\begingroup\small
\setlength{\tabcolsep}{3pt}\renewcommand{\arraystretch}{1.07}
\begin{longtable}{@{}>{\raggedright\arraybackslash}p{.065\textwidth}>{\raggedright\arraybackslash}p{.90\textwidth}@{}}
\toprule ID&Quantity, calculation and comparison\\\midrule
\endfirsthead\toprule ID&Quantity, calculation and comparison\\\midrule
\endhead\bottomrule\endfoot
\hypertarget{row-S21}{}\textbf{S21}\label{row:S21} & \textbf{B\_\allowbreak{}c* minus B\_\allowbreak{}c}\par
\textit{Reading:} the zero moved by the world lines bound together \ensuremath{\div} the same world lines apart\par
\textit{Formula:} pole difference: the colour kernel's contact with one mass for the motion and the spin, the rest energy\par
\textit{Value:} Analytic expression\quad\textit{Unit:} $\mathrm{MeV}$\quad\textit{Comparison:} 64.5 +1.7/-2.0 [ATLAS 2026, B\_\allowbreak{}c*\ensuremath{{}^{+}} \ensuremath{\to} B\_\allowbreak{}c\ensuremath{{}^{+}}\ensuremath{\gamma} with converted photons]\quad\textit{Residual:} --\par
\textit{Operation:} Analytic comparison record.\par
\textit{Calculation:} \path{programs/article_results.py}; \path{outputs/article_results.out}, line 513, identifier \path{SI.S21}.\par
{\footnotesize\textit{Derivation:} the several-line source expression stated in this row and Part III.}\\[7pt]
\hypertarget{row-N23}{}\textbf{N23}\label{row:N23} & \textbf{incompressibility K\_\allowbreak{}0}\par
\textit{Reading:} the pole of these bound world lines, carried through the stated classical map \ensuremath{\div} its reference\par
\textit{Formula:} C[one two-source response with both Pauli blockings at the saturation point]: 9n\ensuremath{{}_{0}}\ensuremath{{}^{2}}e\ensuremath{\prime\prime}(n\ensuremath{{}_{0}}), the scalar and vector strengths carrying \ensuremath{R_{\sigma}}(n)\ensuremath{{}^{2}}(1 - f(n))\par
\textit{Value:} Analytic expression\quad\textit{Unit:} $\mathrm{MeV}$\quad\textit{Comparison:} 230 +40/-40 [the isoscalar giant monopole resonance with RPA functionals]\quad\textit{Residual:} --\par
\textit{Operation:} Analytic comparison record.\par
\textit{Calculation:} \path{programs/article_results.py}; \path{outputs/article_results.out}, line 514, identifier \path{SI.N23}.\par
{\footnotesize\textit{Derivation:} the several-line source expression stated in this row and Part III.}\\[7pt]
\hypertarget{row-N24}{}\textbf{N24}\label{row:N24} & \textbf{symmetry energy J}\par
\textit{Reading:} the bulk response of many nucleons on one kernel \ensuremath{\div} its normalization\par
\textit{Formula:} coefficient: the kinetic part, the isospin-charge two-source response \ensuremath{C_{\tau}}n/2, the T = 0 iterated pion of neutron matter carried by the response's m*, and the one-pion Fock term, at n\ensuremath{{}_{0}}\par
\textit{Value:} Analytic expression\quad\textit{Unit:} $\mathrm{MeV}$\quad\textit{Comparison:} 31.6 +2.7/-2.7 [Li and Han 2013]\quad\textit{Residual:} --\par
\textit{Operation:} Analytic comparison record.\par
\textit{Calculation:} \path{programs/article_results.py}; \path{outputs/article_results.out}, line 515, identifier \path{SI.N24}.\par
{\footnotesize\textit{Derivation:} the several-line source expression stated in this row and Part III.}\\[7pt]
\hypertarget{row-N25}{}\textbf{N25}\label{row:N25} & \textbf{symmetry slope L}\par
\textit{Reading:} the pole of these bound world lines, carried through the stated classical map \ensuremath{\div} its reference\par
\textit{Formula:} 3n\ensuremath{{}_{0}}S\ensuremath{{}_{2}}\ensuremath{\prime}(n\ensuremath{{}_{0}}) of the same\par
\textit{Value:} Analytic expression\quad\textit{Unit:} $\mathrm{MeV}$\quad\textit{Comparison:} 58.7 +28.1/-28.1 [Oertel et al. 2017]\quad\textit{Residual:} --\par
\textit{Operation:} Analytic comparison record.\par
\textit{Calculation:} \path{programs/article_results.py}; \path{outputs/article_results.out}, line 516, identifier \path{SI.N25}.\par
{\footnotesize\textit{Derivation:} the several-line source expression stated in this row and Part III.}\\[7pt]
\hypertarget{row-J10}{}\textbf{J10}\label{row:J10} & \textbf{\ensuremath{\pi}\ensuremath{{}^{0}} by its world lines (MeV)}\par
\textit{Reading:} this closed world line's clock turning the leading zero \ensuremath{\div} the reference clock, the electron's phase on one directed edge or the world line the term names\par
\textit{Formula:} \ensuremath{\surd}[(P(\ensuremath{S_{u}},\ensuremath{S_{u}})\ensuremath{{}^{2}} + P(\ensuremath{S_{d}},\ensuremath{S_{d}})\ensuremath{{}^{2}})/2]\par
\textit{Value:} 134.956486146739\quad\textit{Unit:} $\mathrm{MeV}$\quad\textit{Comparison:} 134.9768 +0.0005 / -0.0005 [source not recorded]\quad\textit{Residual:} -40.627707\par
\textit{Operation:} Finite expression evaluated in Sec.~\ref{sec:literalcompositesSI}.\par
\textit{Calculation:} \path{programs/article_results.py}; \path{outputs/article_results.out}, line 517, identifier \path{SI.J10}.\par
{\footnotesize\textit{Derivation:} the finite inputs and the stated conventional calculation of Part III.}\\[7pt]
\hypertarget{row-J11}{}\textbf{J11}\label{row:J11} & \textbf{\ensuremath{\eta} (MeV)}\par
\textit{Reading:} this closed world line's clock turning the leading zero \ensuremath{\div} the reference clock, the electron's phase on one directed edge or the world line the term names\par
\textit{Formula:} P(\ensuremath{S_{d}},\ensuremath{S_{d}}); no join term\par
\textit{Value:} 157.925020972358\quad\textit{Unit:} $\mathrm{MeV}$\quad\textit{Comparison:} 547.862 +0.017 / -0.017 [source not recorded]\quad\textit{Residual:} -22937.469355\par
\textit{Operation:} Finite expression evaluated in Sec.~\ref{sec:literalcompositesSI}.\par
\textit{Calculation:} \path{programs/article_results.py}; \path{outputs/article_results.out}, line 518, identifier \path{SI.J11}.\par
{\footnotesize\textit{Derivation:} the finite inputs and the stated conventional calculation of Part III.}\\[7pt]
\hypertarget{row-J12}{}\textbf{J12}\label{row:J12} & \textbf{\ensuremath{\eta}\ensuremath{\prime} (MeV)}\par
\textit{Reading:} this closed world line's clock turning the leading zero \ensuremath{\div} the reference clock, the electron's phase on one directed edge or the world line the term names\par
\textit{Formula:} P(\ensuremath{S_{s}},\ensuremath{S_{s}}); no join term\par
\textit{Value:} 675.959220300687\quad\textit{Unit:} $\mathrm{MeV}$\quad\textit{Comparison:} 957.78 +0.06 / -0.06 [source not recorded]\quad\textit{Residual:} -4697.012995\par
\textit{Operation:} Finite expression evaluated in Sec.~\ref{sec:literalcompositesSI}.\par
\textit{Calculation:} \path{programs/article_results.py}; \path{outputs/article_results.out}, line 519, identifier \path{SI.J12}.\par
{\footnotesize\textit{Derivation:} the finite inputs and the stated conventional calculation of Part III.}\\[7pt]
\hypertarget{row-J13}{}\textbf{J13}\label{row:J13} & \textbf{D (mean) (MeV)}\par
\textit{Reading:} this closed world line's clock turning the leading zero \ensuremath{\div} the reference clock, the electron's phase on one directed edge or the world line the term names\par
\textit{Formula:} P(\ensuremath{S_{c}},(\ensuremath{S_{u}} + \ensuremath{S_{d}})/2)\par
\textit{Value:} Analytic expression\quad\textit{Unit:} $\mathrm{MeV}$\quad\textit{Comparison:} 1867.25 +0.03 / -0.03 [source not recorded]\quad\textit{Residual:} --\par
\textit{Operation:} Analytic comparison record.\par
\textit{Calculation:} \path{programs/article_results.py}; \path{outputs/article_results.out}, line 520, identifier \path{SI.J13}.\par
{\footnotesize\textit{Derivation:} the several-line source expression stated in this row and Part III.}\\[7pt]
\hypertarget{row-J14}{}\textbf{J14}\label{row:J14} & \textbf{D\_\allowbreak{}s (MeV)}\par
\textit{Reading:} this closed world line's clock turning the leading zero \ensuremath{\div} the reference clock, the electron's phase on one directed edge or the world line the term names\par
\textit{Formula:} P(\ensuremath{S_{c}},\ensuremath{S_{s}})\par
\textit{Value:} Analytic expression\quad\textit{Unit:} $\mathrm{MeV}$\quad\textit{Comparison:} 1968.35 +0.07 / -0.07 [source not recorded]\quad\textit{Residual:} --\par
\textit{Operation:} Analytic comparison record.\par
\textit{Calculation:} \path{programs/article_results.py}; \path{outputs/article_results.out}, line 521, identifier \path{SI.J14}.\par
{\footnotesize\textit{Derivation:} the several-line source expression stated in this row and Part III.}\\[7pt]
\hypertarget{row-J15}{}\textbf{J15}\label{row:J15} & \textbf{B (mean) (MeV)}\par
\textit{Reading:} this closed world line's clock turning the leading zero \ensuremath{\div} the reference clock, the electron's phase on one directed edge or the world line the term names\par
\textit{Formula:} P(\ensuremath{S_{b}},(\ensuremath{S_{u}} + \ensuremath{S_{d}})/2)\par
\textit{Value:} Analytic expression\quad\textit{Unit:} $\mathrm{MeV}$\quad\textit{Comparison:} 5279.565 +0.053 / -0.053 [source not recorded]\quad\textit{Residual:} --\par
\textit{Operation:} Analytic comparison record.\par
\textit{Calculation:} \path{programs/article_results.py}; \path{outputs/article_results.out}, line 522, identifier \path{SI.J15}.\par
{\footnotesize\textit{Derivation:} the several-line source expression stated in this row and Part III.}\\[7pt]
\hypertarget{row-J16}{}\textbf{J16}\label{row:J16} & \textbf{B\_\allowbreak{}s (MeV)}\par
\textit{Reading:} this closed world line's clock turning the leading zero \ensuremath{\div} the reference clock, the electron's phase on one directed edge or the world line the term names\par
\textit{Formula:} P(\ensuremath{S_{b}},\ensuremath{S_{s}})\par
\textit{Value:} Analytic expression\quad\textit{Unit:} $\mathrm{MeV}$\quad\textit{Comparison:} 5366.93 +0.1 / -0.1 [source not recorded]\quad\textit{Residual:} --\par
\textit{Operation:} Analytic comparison record.\par
\textit{Calculation:} \path{programs/article_results.py}; \path{outputs/article_results.out}, line 523, identifier \path{SI.J16}.\par
{\footnotesize\textit{Derivation:} the several-line source expression stated in this row and Part III.}\\[7pt]
\hypertarget{row-J19}{}\textbf{J19}\label{row:J19} & \textbf{\ensuremath{\omega} (MeV)}\par
\textit{Reading:} this closed world line's clock turning the leading zero \ensuremath{\div} the reference clock, the electron's phase on one directed edge or the world line the term names\par
\textit{Formula:} $\sqrt{s_R}$, with $s_R=\Re s_\rho$ in Eq.~\eqref{eq:realsquaredpoleSI}.\par
\textit{Value:} 759.215750293876\quad\textit{Unit:} $\mathrm{MeV}$\quad\textit{Comparison:} 782.66 +0.13 / -0.13 [source not recorded]\quad\textit{Residual:} -180.340382\par
\textit{Operation:} Finite expression evaluated in Sec.~\ref{sec:literalcompositesSI}.\par
\textit{Calculation:} \path{programs/article_results.py}; \path{outputs/article_results.out}, line 524, identifier \path{SI.J19}.\par
{\footnotesize\textit{Derivation:} the finite inputs and the stated conventional calculation of Part III.}\\[7pt]
\hypertarget{row-J20}{}\textbf{J20}\label{row:J20} & \textbf{K*\ensuremath{\pm} (MeV)}\par
\textit{Reading:} this closed world line's clock turning the leading zero \ensuremath{\div} the reference clock, the electron's phase on one directed edge or the world line the term names\par
\textit{Formula:} $\sqrt{s_R+m_{K^\pm}^2}$.\par
\textit{Value:} 905.601302988073\quad\textit{Unit:} $\mathrm{MeV}$\quad\textit{Comparison:} 891.88 +0.23 / -0.23 [source not recorded]\quad\textit{Residual:} +59.657839\par
\textit{Operation:} Finite expression evaluated in Sec.~\ref{sec:literalcompositesSI}.\par
\textit{Calculation:} \path{programs/article_results.py}; \path{outputs/article_results.out}, line 525, identifier \path{SI.J20}.\par
{\footnotesize\textit{Derivation:} the finite inputs and the stated conventional calculation of Part III.}\\[7pt]
\hypertarget{row-J21}{}\textbf{J21}\label{row:J21} & \textbf{K*\ensuremath{{}^{0}} (MeV)}\par
\textit{Reading:} this closed world line's clock turning the leading zero \ensuremath{\div} the reference clock, the electron's phase on one directed edge or the world line the term names\par
\textit{Formula:} $\sqrt{s_R+m_{K^0}^2}$.\par
\textit{Value:} 907.759789219815\quad\textit{Unit:} $\mathrm{MeV}$\quad\textit{Comparison:} 895.56 +0.2 / -0.2 [source not recorded]\quad\textit{Residual:} +60.998946\par
\textit{Operation:} Finite expression evaluated in Sec.~\ref{sec:literalcompositesSI}.\par
\textit{Calculation:} \path{programs/article_results.py}; \path{outputs/article_results.out}, line 526, identifier \path{SI.J21}.\par
{\footnotesize\textit{Derivation:} the finite inputs and the stated conventional calculation of Part III.}\\[7pt]
\hypertarget{row-J22}{}\textbf{J22}\label{row:J22} & \textbf{\ensuremath{\varphi} (MeV)}\par
\textit{Reading:} this closed world line's clock turning the leading zero \ensuremath{\div} the reference clock, the electron's phase on one directed edge or the world line the term names\par
\textit{Formula:} $\sqrt{s_R+P(S_s,S_s)^2}$.\par
\textit{Value:} 1016.52812209196\quad\textit{Unit:} $\mathrm{MeV}$\quad\textit{Comparison:} 1019.46 +0.016 / -0.016 [source not recorded]\quad\textit{Residual:} -183.242369\par
\textit{Operation:} Finite expression evaluated in Sec.~\ref{sec:literalcompositesSI}.\par
\textit{Calculation:} \path{programs/article_results.py}; \path{outputs/article_results.out}, line 527, identifier \path{SI.J22}.\par
{\footnotesize\textit{Derivation:} the finite inputs and the stated conventional calculation of Part III.}\\[7pt]
\hypertarget{row-J23}{}\textbf{J23}\label{row:J23} & \textbf{J/\ensuremath{\psi} (MeV)}\par
\textit{Reading:} this closed world line's clock turning the leading zero \ensuremath{\div} the reference clock, the electron's phase on one directed edge or the world line the term names\par
\textit{Formula:} Source to construct: the c c\ensuremath{\overline{\vphantom{x}}} source\par
\textit{Value:} Analytic expression\quad\textit{Unit:} $\mathrm{MeV}$\quad\textit{Comparison:} 3096.9 +0.006 / -0.006 [source not recorded]\quad\textit{Residual:} --\par
\textit{Operation:} Analytic comparison record.\par
\textit{Calculation:} \path{programs/article_results.py}; \path{outputs/article_results.out}, line 528, identifier \path{SI.J23}.\par
{\footnotesize\textit{Derivation:} the several-line source expression stated in this row and Part III.}\\[7pt]
\hypertarget{row-J24}{}\textbf{J24}\label{row:J24} & \textbf{\ensuremath{\Upsilon}(1S) (MeV)}\par
\textit{Reading:} this closed world line's clock turning the leading zero \ensuremath{\div} the reference clock, the electron's phase on one directed edge or the world line the term names\par
\textit{Formula:} Source to construct: the b b\ensuremath{\overline{\vphantom{x}}} source\par
\textit{Value:} Analytic expression\quad\textit{Unit:} $\mathrm{MeV}$\quad\textit{Comparison:} 9460.4 +0.1 / -0.1 [source not recorded]\quad\textit{Residual:} --\par
\textit{Operation:} Analytic comparison record.\par
\textit{Calculation:} \path{programs/article_results.py}; \path{outputs/article_results.out}, line 529, identifier \path{SI.J24}.\par
{\footnotesize\textit{Derivation:} the several-line source expression stated in this row and Part III.}\\[7pt]
\hypertarget{row-J25}{}\textbf{J25}\label{row:J25} & \textbf{\ensuremath{\Sigma}\ensuremath{{}^{+}} (MeV)}\par
\textit{Reading:} this closed world line's clock turning the leading zero \ensuremath{\div} the reference clock, the electron's phase on one directed edge or the world line the term names\par
\textit{Formula:} Source to construct: the charge assignment\par
\textit{Value:} Analytic expression\quad\textit{Unit:} $\mathrm{MeV}$\quad\textit{Comparison:} 1189.37 +0.07 / -0.07 [source not recorded]\quad\textit{Residual:} --\par
\textit{Operation:} Analytic comparison record.\par
\textit{Calculation:} \path{programs/article_results.py}; \path{outputs/article_results.out}, line 530, identifier \path{SI.J25}.\par
{\footnotesize\textit{Derivation:} the several-line source expression stated in this row and Part III.}\\[7pt]
\hypertarget{row-J26}{}\textbf{J26}\label{row:J26} & \textbf{\ensuremath{\Sigma}\ensuremath{{}^{0}} (MeV)}\par
\textit{Reading:} this closed world line's clock turning the leading zero \ensuremath{\div} the reference clock, the electron's phase on one directed edge or the world line the term names\par
\textit{Formula:} \ensuremath{m_{n}}(5/2)\ensuremath{{}^{59/320}}\par
\textit{Value:} 1112.492814404643\quad\textit{Unit:} $\mathrm{MeV}$\quad\textit{Comparison:} 1192.642 +0.024 / -0.024 [source not recorded]\quad\textit{Residual:} -3339.549400\par
\textit{Operation:} Finite expression evaluated in Sec.~\ref{sec:literalcompositesSI}.\par
\textit{Calculation:} \path{programs/article_results.py}; \path{outputs/article_results.out}, line 531, identifier \path{SI.J26}.\par
{\footnotesize\textit{Derivation:} the finite inputs and the stated conventional calculation of Part III.}\\[7pt]
\hypertarget{row-J27}{}\textbf{J27}\label{row:J27} & \textbf{\ensuremath{\Sigma}\ensuremath{{}^{-}} (MeV)}\par
\textit{Reading:} this closed world line's clock turning the leading zero \ensuremath{\div} the reference clock, the electron's phase on one directed edge or the world line the term names\par
\textit{Formula:} Source to construct: the charge assignment\par
\textit{Value:} Analytic expression\quad\textit{Unit:} $\mathrm{MeV}$\quad\textit{Comparison:} 1197.449 +0.029 / -0.029 [source not recorded]\quad\textit{Residual:} --\par
\textit{Operation:} Analytic comparison record.\par
\textit{Calculation:} \path{programs/article_results.py}; \path{outputs/article_results.out}, line 532, identifier \path{SI.J27}.\par
{\footnotesize\textit{Derivation:} the several-line source expression stated in this row and Part III.}\\[7pt]
\hypertarget{row-J28}{}\textbf{J28}\label{row:J28} & \textbf{\ensuremath{\Xi}\ensuremath{{}^{0}} (MeV)}\par
\textit{Reading:} this closed world line's clock turning the leading zero \ensuremath{\div} the reference clock, the electron's phase on one directed edge or the world line the term names\par
\textit{Formula:} \ensuremath{m_{n}}(5/2)\ensuremath{{}^{118/320}}\par
\textit{Value:} 1317.247562638359\quad\textit{Unit:} $\mathrm{MeV}$\quad\textit{Comparison:} 1314.86 +0.2 / -0.2 [source not recorded]\quad\textit{Residual:} +11.937813\par
\textit{Operation:} Finite expression evaluated in Sec.~\ref{sec:literalcompositesSI}.\par
\textit{Calculation:} \path{programs/article_results.py}; \path{outputs/article_results.out}, line 533, identifier \path{SI.J28}.\par
{\footnotesize\textit{Derivation:} the finite inputs and the stated conventional calculation of Part III.}\\[7pt]
\hypertarget{row-J29}{}\textbf{J29}\label{row:J29} & \textbf{\ensuremath{\Xi}\ensuremath{{}^{-}} (MeV)}\par
\textit{Reading:} this closed world line's clock turning the leading zero \ensuremath{\div} the reference clock, the electron's phase on one directed edge or the world line the term names\par
\textit{Formula:} Source to construct: the charge assignment\par
\textit{Value:} Analytic expression\quad\textit{Unit:} $\mathrm{MeV}$\quad\textit{Comparison:} 1321.71 +0.07 / -0.07 [source not recorded]\quad\textit{Residual:} --\par
\textit{Operation:} Analytic comparison record.\par
\textit{Calculation:} \path{programs/article_results.py}; \path{outputs/article_results.out}, line 534, identifier \path{SI.J29}.\par
{\footnotesize\textit{Derivation:} the several-line source expression stated in this row and Part III.}\\[7pt]
\hypertarget{row-J30}{}\textbf{J30}\label{row:J30} & \textbf{\ensuremath{\Sigma}* (mean) (MeV)}\par
\textit{Reading:} this closed world line's clock turning the leading zero \ensuremath{\div} the reference clock, the electron's phase on one directed edge or the world line the term names\par
\textit{Formula:} 3b\ensuremath{m_{p}}(5/2)\ensuremath{{}^{20/320}}\par
\textit{Value:} 1304.065364234984\quad\textit{Unit:} $\mathrm{MeV}$\quad\textit{Comparison:} 1384.61 +0.36 / -0.36 [source not recorded]\quad\textit{Residual:} -223.735099\par
\textit{Operation:} Finite expression evaluated in Sec.~\ref{sec:literalcompositesSI}.\par
\textit{Calculation:} \path{programs/article_results.py}; \path{outputs/article_results.out}, line 535, identifier \path{SI.J30}.\par
{\footnotesize\textit{Derivation:} the finite inputs and the stated conventional calculation of Part III.}\\[7pt]
\hypertarget{row-J31}{}\textbf{J31}\label{row:J31} & \textbf{\ensuremath{\Xi}* (mean) (MeV)}\par
\textit{Reading:} this closed world line's clock turning the leading zero \ensuremath{\div} the reference clock, the electron's phase on one directed edge or the world line the term names\par
\textit{Formula:} 3b\ensuremath{m_{p}}(5/2)\ensuremath{{}^{40/320}}\par
\textit{Value:} 1380.926649456699\quad\textit{Unit:} $\mathrm{MeV}$\quad\textit{Comparison:} 1533.4 +0.34 / -0.34 [source not recorded]\quad\textit{Residual:} -448.451031\par
\textit{Operation:} Finite expression evaluated in Sec.~\ref{sec:literalcompositesSI}.\par
\textit{Calculation:} \path{programs/article_results.py}; \path{outputs/article_results.out}, line 536, identifier \path{SI.J31}.\par
{\footnotesize\textit{Derivation:} the finite inputs and the stated conventional calculation of Part III.}\\[7pt]
\hypertarget{row-J32}{}\textbf{J32}\label{row:J32} & \textbf{\ensuremath{\Omega}\ensuremath{{}^{-}} (MeV)}\par
\textit{Reading:} this closed world line's clock turning the leading zero \ensuremath{\div} the reference clock, the electron's phase on one directed edge or the world line the term names\par
\textit{Formula:} 3b\ensuremath{m_{p}}(5/2)\ensuremath{{}^{60/320}}\par
\textit{Value:} 1462.318119535674\quad\textit{Unit:} $\mathrm{MeV}$\quad\textit{Comparison:} 1672.45 +0.29 / -0.29 [source not recorded]\quad\textit{Residual:} -724.592691\par
\textit{Operation:} Finite expression evaluated in Sec.~\ref{sec:literalcompositesSI}.\par
\textit{Calculation:} \path{programs/article_results.py}; \path{outputs/article_results.out}, line 537, identifier \path{SI.J32}.\par
{\footnotesize\textit{Derivation:} the finite inputs and the stated conventional calculation of Part III.}\\[7pt]
\hypertarget{row-J33}{}\textbf{J33}\label{row:J33} & \textbf{\ensuremath{\Lambda}\_\allowbreak{}c (MeV)}\par
\textit{Reading:} this closed world line's clock turning the leading zero \ensuremath{\div} the reference clock, the electron's phase on one directed edge or the world line the term names\par
\textit{Formula:} Source to construct: the c source\par
\textit{Value:} Analytic expression\quad\textit{Unit:} $\mathrm{MeV}$\quad\textit{Comparison:} 2286.46 +0.14 / -0.14 [source not recorded]\quad\textit{Residual:} --\par
\textit{Operation:} Analytic comparison record.\par
\textit{Calculation:} \path{programs/article_results.py}; \path{outputs/article_results.out}, line 538, identifier \path{SI.J33}.\par
{\footnotesize\textit{Derivation:} the several-line source expression stated in this row and Part III.}\\[7pt]
\hypertarget{row-J34}{}\textbf{J34}\label{row:J34} & \textbf{\ensuremath{\Lambda}\_\allowbreak{}b (MeV)}\par
\textit{Reading:} this closed world line's clock turning the leading zero \ensuremath{\div} the reference clock, the electron's phase on one directed edge or the world line the term names\par
\textit{Formula:} Source to construct: the b source\par
\textit{Value:} Analytic expression\quad\textit{Unit:} $\mathrm{MeV}$\quad\textit{Comparison:} 5619.57 +0.16 / -0.16 [source not recorded]\quad\textit{Residual:} --\par
\textit{Operation:} Analytic comparison record.\par
\textit{Calculation:} \path{programs/article_results.py}; \path{outputs/article_results.out}, line 539, identifier \path{SI.J34}.\par
{\footnotesize\textit{Derivation:} the several-line source expression stated in this row and Part III.}\\[7pt]
\hypertarget{row-S19}{}\textbf{S19}\label{row:S19} & \textbf{nucleon scalar residue sigma\_\allowbreak{}piN}\par
\textit{Reading:} the current of the four amplitudes at the pole \ensuremath{\div} the reference current: the nuclear magneton, the Compton length, the electron vertex\par
\textit{Formula:} residue: the scalar insertion at the nucleon pole: its \ensuremath{\pi}, K, \ensuremath{\eta_8} and \ensuremath{\Delta} cloud, \ensuremath{\Sigma_{\varphi}} (m\ensuremath{\widehat{\vphantom{x}}}\ensuremath{\partial}\ensuremath{M_{\varphi}}\ensuremath{{}^{2}}/\ensuremath{\partial}m\ensuremath{\widehat{\vphantom{x}}})\ensuremath{\partial}\ensuremath{\Sigma_{\varphi}}/\ensuremath{\partial}\ensuremath{M_{\varphi}}\ensuremath{{}^{2}} = 45.315 MeV, and its core, the levels of the three light world lines at one depth, 3\ensuremath{Z_{l}} m\ensuremath{\widehat{\vphantom{x}}}(1 + 2A/u), with the light world line's depth \ensuremath{Z_{l}} = 1.678 from its share \ensuremath{\lambda} = \ensuremath{m_{p}}/3 moving around a world line at rest\par
\textit{Value:} Analytic expression\quad\textit{Unit:} $\mathrm{MeV}$\quad\textit{Comparison:} 59.1 +3.5/-3.5 [Roy--Steiner, Hoferichter et al. 2015]\quad\textit{Residual:} --\par
\textit{Operation:} Analytic comparison record.\par
\textit{Calculation:} \path{programs/article_results.py}; \path{outputs/article_results.out}, line 540, identifier \path{SI.S19}.\par
{\footnotesize\textit{Derivation:} the several-line source expression stated in this row and Part III.}\\[7pt]
\hypertarget{row-S20}{}\textbf{S20}\label{row:S20} & \textbf{nucleon strange scalar residue sigma\_\allowbreak{}s}\par
\textit{Reading:} the current of the four amplitudes at the pole \ensuremath{\div} the reference current: the nuclear magneton, the Compton length, the electron vertex\par
\textit{Formula:} residue: the strange scalar insertion at the nucleon pole through its K and \ensuremath{\eta_8} cloud\par
\textit{Value:} Analytic expression\quad\textit{Unit:} $\mathrm{MeV}$\quad\textit{Comparison:} 44.9 +6.4/-6.4 [FLAG 2024, N\_\allowbreak{}f = 2+1 lattice average]\quad\textit{Residual:} --\par
\textit{Operation:} Analytic comparison record.\par
\textit{Calculation:} \path{programs/article_results.py}; \path{outputs/article_results.out}, line 541, identifier \path{SI.S20}.\par
{\footnotesize\textit{Derivation:} the several-line source expression stated in this row and Part III.}\\[7pt]
\hypertarget{row-G14}{}\textbf{G14}\label{row:G14} & \textbf{deuteron asymptotic S normalization A\_\allowbreak{}S}\par
\textit{Reading:} the current of the four amplitudes at the pole \ensuremath{\div} the reference current: the nuclear magneton, the Compton length, the electron vertex\par
\textit{Formula:} residue: the S normalization of the two-nucleon pole's residue, the two-nucleon kernel\par
\textit{Value:} Analytic expression\quad\textit{Unit:} $\mathrm{fm}^{-1/2}$\quad\textit{Comparison:} 0.8845 +0.0008/-0.0008 [de Swart, Terheggen and Stoks 1995, Nijmegen]\quad\textit{Residual:} --\par
\textit{Operation:} Analytic comparison record.\par
\textit{Calculation:} \path{programs/article_results.py}; \path{outputs/article_results.out}, line 542, identifier \path{SI.G14}.\par
{\footnotesize\textit{Derivation:} the several-line source expression stated in this row and Part III.}\\[7pt]
\hypertarget{row-G15}{}\textbf{G15}\label{row:G15} & \textbf{deuteron D/S ratio eta}\par
\textit{Reading:} the bound pole's eigenvector \ensuremath{\div} its S part\par
\textit{Formula:} eigenvector: the D/S direction of the two-nucleon pole's residue\par
\textit{Value:} Analytic expression\quad\textit{Comparison:} 0.0253 +0.0002/-0.0002 [de Swart et al. 1995, recommended]\quad\textit{Residual:} --\par
\textit{Operation:} Analytic comparison record.\par
\textit{Calculation:} \path{programs/article_results.py}; \path{outputs/article_results.out}, line 543, identifier \path{SI.G15}.\par
{\footnotesize\textit{Derivation:} the several-line source expression stated in this row and Part III.}\\[7pt]
\hypertarget{row-G16}{}\textbf{G16}\label{row:G16} & \textbf{deuteron effective range at the pole rho\_\allowbreak{}d}\par
\textit{Reading:} the current of the four amplitudes at the pole \ensuremath{\div} the reference current: the nuclear magneton, the Compton length, the electron vertex\par
\textit{Formula:} residue: 1/\ensuremath{\gamma} - 2/\ensuremath{N_{p}}\ensuremath{{}^{2}} of the two-nucleon pole\par
\textit{Value:} Analytic expression\quad\textit{Unit:} $\mathrm{fm}$\quad\textit{Comparison:} 1.765 +0.004/-0.004 [de Swart et al. 1995]\quad\textit{Residual:} --\par
\textit{Operation:} Analytic comparison record.\par
\textit{Calculation:} \path{programs/article_results.py}; \path{outputs/article_results.out}, line 544, identifier \path{SI.G16}.\par
{\footnotesize\textit{Derivation:} the several-line source expression stated in this row and Part III.}\\[7pt]
\hypertarget{row-G17}{}\textbf{G17}\label{row:G17} & \textbf{triplet scattering length a\_\allowbreak{}t}\par
\textit{Reading:} the amplitudes' current carried through the stated classical map \ensuremath{\div} its reference\par
\textit{Formula:} the two-nucleon kernel at zero energy\par
\textit{Value:} Analytic expression\quad\textit{Unit:} $\mathrm{fm}$\quad\textit{Comparison:} 5.4194 +0.002/-0.002 [de Swart et al. 1995, Nijmegen]\quad\textit{Residual:} --\par
\textit{Operation:} Analytic comparison record.\par
\textit{Calculation:} \path{programs/article_results.py}; \path{outputs/article_results.out}, line 545, identifier \path{SI.G17}.\par
{\footnotesize\textit{Derivation:} the several-line source expression stated in this row and Part III.}\\[7pt]
\hypertarget{row-G18}{}\textbf{G18}\label{row:G18} & \textbf{triplet effective range r\_\allowbreak{}t}\par
\textit{Reading:} the amplitudes' current carried through the stated classical map \ensuremath{\div} its reference\par
\textit{Formula:} the two-nucleon kernel at zero energy\par
\textit{Value:} Analytic expression\quad\textit{Unit:} $\mathrm{fm}$\quad\textit{Comparison:} 1.7536 +0.0025/-0.0025 [de Swart et al. 1995, Nijmegen]\quad\textit{Residual:} --\par
\textit{Operation:} Analytic comparison record.\par
\textit{Calculation:} \path{programs/article_results.py}; \path{outputs/article_results.out}, line 546, identifier \path{SI.G18}.\par
{\footnotesize\textit{Derivation:} the several-line source expression stated in this row and Part III.}\\[7pt]
\hypertarget{row-G19}{}\textbf{G19}\label{row:G19} & \textbf{deuteron point matter radius r\_\allowbreak{}m}\par
\textit{Reading:} the current of the four amplitudes at the pole \ensuremath{\div} the reference current: the nuclear magneton, the Compton length, the electron vertex\par
\textit{Formula:} residue: (1/2)\ensuremath{\langle}r\ensuremath{{}^{2}}\ensuremath{\rangle}\ensuremath{{}^{1/2}} of the two-nucleon pole's wave\par
\textit{Value:} Analytic expression\quad\textit{Unit:} $\mathrm{fm}$\quad\textit{Comparison:} 1.9676 +0.001/-0.001 [Nijmegen high-quality potentials, de Swart et al. 1995]; the isotope shift's 1.97507 is the structure radius r\_\allowbreak{}str (G19a)\quad\textit{Residual:} --\par
\textit{Operation:} Analytic comparison record.\par
\textit{Calculation:} \path{programs/article_results.py}; \path{outputs/article_results.out}, line 547, identifier \path{SI.G19}.\par
{\footnotesize\textit{Derivation:} the several-line source expression stated in this row and Part III.}\\[7pt]
\hypertarget{row-G19a}{}\textbf{G19a}\label{row:G19a} & \textbf{deuteron structure radius r\_\allowbreak{}str (the 1S--2S number)}\par
\textit{Reading:} the current of the four amplitudes at the pole \ensuremath{\div} the reference current: the nuclear magneton, the Compton length, the electron vertex\par
\textit{Formula:} residue: \ensuremath{\surd}(\ensuremath{r_{m}}\ensuremath{{}^{2}} + the two-body and spin--orbit monopoles) of the two-nucleon pole\par
\textit{Value:} Analytic expression\quad\textit{Unit:} $\mathrm{fm}$\quad\textit{Comparison:} 1.97507 +0.00078/-0.00078 [H--D 1S--2S isotope shift with bound-state QED, Jentschura et al. 2011]\quad\textit{Residual:} --\par
\textit{Operation:} Analytic comparison record.\par
\textit{Calculation:} \path{programs/article_results.py}; \path{outputs/article_results.out}, line 548, identifier \path{SI.G19a}.\par
{\footnotesize\textit{Derivation:} the several-line source expression stated in this row and Part III.}\\[7pt]
\hypertarget{row-G20}{}\textbf{G20}\label{row:G20} & \textbf{deuteron quadrupole moment, impulse part Q\_\allowbreak{}d}\par
\textit{Reading:} the current of the four amplitudes at the pole \ensuremath{\div} the reference current: the nuclear magneton, the Compton length, the electron vertex\par
\textit{Formula:} residue: (1/20)\ensuremath{\int}r\ensuremath{{}^{2}}w(\ensuremath{\surd}8u - w) of the two-nucleon pole\par
\textit{Value:} Analytic expression\quad\textit{Unit:} $\mathrm{fm}^{2}$\quad\textit{Comparison:} 0.271 +0.001/-0.001 [the impulse Q\ensuremath{{}_{0}} of the Nijmegen high-quality potentials, a model value]\quad\textit{Residual:} --\par
\textit{Operation:} Analytic comparison record.\par
\textit{Calculation:} \path{programs/article_results.py}; \path{outputs/article_results.out}, line 549, identifier \path{SI.G20}.\par
{\footnotesize\textit{Derivation:} the several-line source expression stated in this row and Part III.}\\[7pt]
\hypertarget{row-G21}{}\textbf{G21}\label{row:G21} & \textbf{deuteron electric polarizability alpha\_\allowbreak{}E}\par
\textit{Reading:} the amplitudes' current carried through the stated classical map \ensuremath{\div} its reference\par
\textit{Formula:} the two-nucleon pole's E1 response with the free P and F waves\par
\textit{Value:} Analytic expression\quad\textit{Unit:} $\mathrm{fm}^{3}$\quad\textit{Comparison:} 0.637 +0.029/-0.029 [SLEGS 2026: the Baldin \ensuremath{\alpha}\_\allowbreak{}E + \ensuremath{\beta}\_\allowbreak{}M less a pionless \ensuremath{\beta}\_\allowbreak{}M]\quad\textit{Residual:} --\par
\textit{Operation:} Analytic comparison record.\par
\textit{Calculation:} \path{programs/article_results.py}; \path{outputs/article_results.out}, line 550, identifier \path{SI.G21}.\par
{\footnotesize\textit{Derivation:} the several-line source expression stated in this row and Part III.}\\[7pt]
\end{longtable}
\endgroup

\subsection{Further current and nuclear comparisons}
The following comparisons accompany the current, radial-kernel and nuclear
calculations of Part III. They retain the same evaluations and
reference conventions.\cite{Repository}
\begingroup\small
\setlength{\tabcolsep}{3pt}
\renewcommand{\arraystretch}{1.12}
\begin{longtable}{@{}>{\raggedright\arraybackslash}p{.23\textwidth}>{\raggedright\arraybackslash}p{.22\textwidth}>{\raggedright\arraybackslash}p{.36\textwidth}>{\raggedright\arraybackslash}p{.13\textwidth}@{}}
\caption{Further comparisons of the carried current and nuclear calculations. A quoted residual is reproduced in its stated unit.}\label{tab:furthercomparisons}\\
\toprule Reading&Evaluation&Reference&Residual\\\midrule
\endfirsthead
\multicolumn{4}{@{}l}{\textit{Table \thetable\ continued}}\\
\toprule Reading&Evaluation&Reference&Residual\\\midrule
\endhead
\bottomrule\endlastfoot
Second pion lifetime&$26.09718624891\,\mathrm{ns}$&$26.033(5)\,\mathrm{ns}$, retained pion-decay record&$+12.84\,\sigma$\\[5pt]
Deuteron $A_S$&$0.88221\,\mathrm{fm}^{-1/2}$&$0.8845(8)\,\mathrm{fm}^{-1/2}$, de Swart et al. (1995), G14&---\\[5pt]
Deuteron $A_D/A_S$&$0.025208$&$0.0253(2)$, de Swart et al. (1995), G15&---\\[5pt]
Deuteron $\alpha_E$&$0.6309\,\mathrm{fm}^{3}$&$0.637(29)\,\mathrm{fm}^{3}$, SLEGS 2026 convention in G21&---\\[5pt]
Triplet scattering length&Retained radial kernel&$5.4194(20)\fm$, de Swart et al. (1995), G17&$-6.58\,\sigma$\\[5pt]
Triplet effective range&Retained radial kernel&$1.7536(25)\fm$, de Swart et al. (1995), G18&$-7.39\,\sigma$\\[5pt]
Direct zero-energy test&Retained radial kernel&Direct test in the nuclear calculation&$6.8\,\sigma$ (magnitude)\\[5pt]
Singlet isospin contact&Shift from $C_\tau$&Singlet pole requires a contact 7.5\% above the one-channel contact&About 97\% of required shift\\[5pt]
Triton--helion mirror splitting&$\alpha m_p/9$&$0.763807\MeV$, retained nuclear record&$-0.398\%$\\[5pt]
Triton nuclear mass&$2808.921773830\MeV$&CODATA 2022 comparison of Proposition 186&About $+724\,\sigma$\\[5pt]
Helion nuclear mass&$2808.389208140\MeV$&CODATA 2022 comparison of Proposition 186&About $-2731\,\sigma$\\[5pt]
Helium-4 nuclear mass&$3727.377390382\MeV$&CODATA 2022 comparison of Proposition 186&About $-1685\,\sigma$\\[5pt]
Neon-20 atomic binding&$160.772109354\MeV$&$160.6448\MeV$, AME2020 in the original N07/S13 record&$+0.07925\%$\\[5pt]
Silicon-28 and sulfur-32 clusters&Rigid-count continuation&Retained cluster comparisons&About one further effective bond\\[5pt]
Bulk nuclear masses&$B_{\rm bulk}$&2542 experimental nuclei with $A\ge5$&$15.188\MeV$ RMS\\[5pt]
Bulk masses with shell correction&Specified shell correction&Same 2542 experimental nuclei&$11.427\MeV$ RMS\\[5pt]
Folded charge radii&Eq.~\eqref{eq:attempt-S12-5}&Calcium-48 through uranium-238&$0.68\%$ RMS\\[5pt]
Carbon-12 charge radius&Eq.~\eqref{eq:attempt-S12-5}&Retained nuclear-radius record&$+5.2\%$\\[5pt]
Calcium-40 charge radius&Eq.~\eqref{eq:attempt-S12-5}&Retained nuclear-radius record&$-5.9\%$\\[5pt]
Helium-4 charge radius&$1.672542\fm$&$1.67824(83)\fm$, retained nuclear-radius record&$-0.34\%$; $6.9\,\sigma$ in magnitude\\[5pt]
Positive two-mode mass bound&$m_p\sqrt{12\log(4/3)}$&Measured $D,D_s,B,B_s$ masses in the full register&All exceed the bound\\
\end{longtable}
\endgroup

\section{Calculation and formal proofs}
\label{sec:provenance}

\subsection{Programs and printed results}
\label{sec:programrecord}

The repository at \url{https://github.com/yongzic1-lgtm/physics-without-an-outside-observer} contains the directed-step
calculations, their comparison inputs and their printed results. Each of the
252 register entries above gives a physical line of
\path{outputs/article_results.out}. That line prints its identifier, expression,
value, unit and numerical origin. The conventional calculation records are
inputs in \path{data}; their reporting maps and the explicit recalculations
are in Part III. The twelve literal composite expressions are evaluated
from their displayed formulas.

The six evaluations give the following run record. The last column counts
physical output lines, including the printed headings. Together they contain
186 successful checks.
\begin{center}
\begin{tabular}{@{}lrr@{}}
\toprule Program in \path{programs/} & Checks & Output lines\\
\midrule
\path{main_line_formulas.py} & 29/29 & 32\\
\path{classical_experiments.py} & 25/25 & 709\\
\path{table_from_the_step.py} & 16/16 & 159\\
\path{geometry_law_and_rows.py} & 28/28 & 308\\
\path{whole_picture.py} & 73/73 & 254\\
\path{article_results.py} & 15/15 & 553\\
\bottomrule
\end{tabular}
\end{center}

The run uses Python 3.12.14, NumPy 2.3.5, SciPy 1.17.0 and the
\path{rundec} 0.7 interface to CRunDec 3.1, with
\path{OPENBLAS_NUM_THREADS=1}. The exact commands, versions and output hashes
are in \path{outputs/verification.json}. From the repository root, each
program runs with \texttt{python -I programs/NAME.py}; its output is the
same-named \path{.out} file under \path{outputs/}.

The full numerical register is \path{docs/NUMBERS.md}. Its structured record
is \path{outputs/supplemental_numbers.json}. The measurement values, widths,
source names and conversion inputs are in \path{docs/INPUTS.md} and
\path{data/}. The output index gives the exact address of each value.

\subsection{Statements and their proofs}
\label{sec:proofrecord}

The Lean development contains 334 source modules and 6,894 theorem
declarations: 6,863 explicit theorem or lemma statements and 31 aliases.
Every original mathematical definition, theorem statement and proof retains
its code tokens. The common
entry point \path{Comparison.System} imports a closure of 282 project
modules with 6,188 declarations; the System modules themselves contain
2,456 declarations. The remaining source libraries and their independent
entry points are present in \path{lean/}.

The pinned Lean 4.35.0 run compiled 79 modules containing 1,797 source
declarations. The two unchanged-source runs of
\path{Comparison.Existence.PlaquetteData} ended with exit code 137 under
the 8 GiB memory limit. The exact commands, completed modules and memory
controller records are in \path{lean/verification/BUILD.md} and
\path{lean/verification/build_scope.json}.

Every declaration is indexed by its module and name in
\path{lean/theorems.json}, with its exact statement, file and source line.
The index also records the active variables, typeclass hypotheses and
local notation surrounding each statement.
The import graph and its dependency order are in \path{lean/modules.json}.
The compiled metadata records 84 source declarations exported from nine
modules: their elaborated types, dependencies and axiom sets. \path{docs/PROOFS.md}
follows the geometry and eight physical steps and states the scope of each formal result.
The index places stable declaration names beside the article's current equation numbers.

The exact finite statements supply the matrix, determinant, series and
source identities used in the SI. The SI writes out the continuum
variations, normalizations and reporting operations at their places in
the derivation. The article concordance below identifies both parts of
each calculation.

The build uses the retained \path{lean-toolchain}, \path{lakefile.toml}
and \path{lake-manifest.json}. From \path{lean/},
\texttt{lake exe cache get} obtains the pinned Mathlib objects and
\texttt{LEAN\_NUM\_THREADS=1 lake build Comparison.System} builds the common entry point.
The complete build commands and their logs are in
\path{lean/verification/}.


\subsection{Every numbered result and table entry of the article}
\label{sec:articleconcordance}

The article contains 96 numbered main-text equations and 57 equations in
Appendices A--E. Table I has 34 rows. Table II has 22 matter readings,
18 entries of the two overlap matrices, and ten mixing-parameter rows.
The following 237 addresses follow their printed order.

The column AR gives physical lines of \path{outputs/article_results.out},
produced by \path{programs/article_results.py}. Each such line also prints
its result identifier and expression. Table-II addresses distinguish the
native mass, the conventional conversion, the ratio and the individual
matrix entry. The conventional-conversion lines print the article record,
the stated CRunDec calculation and their difference.

The Lean column names the formal statements used for the result. Where an
SI calculation completes the displayed physical result, that calculation
is named in italics. The full statement, its hypotheses, its source file
and line, the entire printed program text, and the scope of each route are
in \path{docs/ARTICLE_CONCORDANCE.md} and
\path{docs/article_concordance.json}.

Every abbreviated Lean name begins with \path{Comparison.}.
The further abbreviations are
\path{Law.} = \path{System.Law.};
\path{Num.}, \path{Fam.}, \path{Neu.}, and \path{Comp.} =
\path{System.Completion.Num.}, \path{System.Completion.Families.},
\path{System.Completion.Neutral.}, and \path{System.Completion.Comp.};
\path{Form.} = \path{Formulas.};
\path{Whole.} = \path{WholePicture.};
\path{Walk.} = \path{Law.}; and
\path{Series.} = \path{ArticleStatements.Series.}.
\path{System.} and \path{Article.} keep their full relative namespaces.

\begingroup
\footnotesize
\setlength{\tabcolsep}{3pt}
\renewcommand{\arraystretch}{1.10}

\begin{longtable}{@{}>{\raggedright\arraybackslash}p{.57in}>{\raggedright\arraybackslash}p{1.16in}>{\raggedright\arraybackslash}p{.75in}>{\raggedright\arraybackslash}p{4.05in}@{}}
\caption{All numbered article equations.}\label{tab:articleequations}\\
\toprule Article & Supplemental Material & Printed line & Lean statement and SI calculation\\
\midrule\endfirsthead
\multicolumn{4}{l}{\textit{All numbered article equations (continued)}}\\
\toprule Article & Supplemental Material & Printed line & Lean statement and SI calculation\\
\midrule\endhead
\bottomrule\endfoot
(1) & Sec.~\ref{sec:geometry} & AR 194 & \path{System.P1.chance_sum}\newline\emph{P1 probability pairing}.\\
(2) & \eqref{eq:symplectictranspose} & AR 195 & \path{System.comparison_P2}, \path{partnerTwo_kron}\\
(3) & \eqref{eq:fullspacepostulate} & AR 194 & \path{distinction_sq}, \path{System.uniqueness}\newline\emph{P3 scalar-square condition}.\\
(4) & \eqref{eq:reciprocaltransport} & AR 202 & \path{System.prop_32}\\
(5) & \eqref{eq:postulatestep}, \eqref{eq:B} & AR 202 & \path{System.stateAt_continues}, \path{System.B_stateAt}\newline\emph{P5 continuation}.\\
(6) & \eqref{eq:quaternionic} & AR 194 & \path{System.quaternionic_plane}, \path{distinction_sq}\\
(7) & \eqref{eq:quaternionic} & AR 16, 194 & \path{System.uniqueness}, \path{System.systems_isomorphic}\\
(8) & \eqref{eq:explicitclifford} & AR 194 & \path{System.prop_4}\\
(9) & \eqref{eq:agreementstate} & AR 195 & \path{System.prop_11}\\
(10) & \eqref{eq:paircomparison} & AR 195 & \path{System.prop_11}, \path{System.comparison_P2}\\
(11) & \eqref{eq:paircomparison}, \eqref{eq:onekernel} & AR 17, 195,196 & \path{System.prop_10}\\
(12) & \eqref{eq:paircomparison} & AR 195 & \path{System.prop_12}\\
(13) & \eqref{eq:pfaffiansi} & AR 194 & \path{System.prop_16}, \path{System.prop_18}\\
(14) & \eqref{eq:divisionroot} & AR 197 & \path{System.prop_19}\\
(15) & \eqref{eq:divisionroot}, \eqref{eq:nullquotient}, \eqref{eq:localcone} & AR 197, 201 & \path{System.prop_21}\\
(16) & \eqref{eq:adj}, \eqref{eq:B} & AR 18, 202 & \path{System.prop_26}\\
(17) & \eqref{eq:closedlaw}, \eqref{eq:eulerproduct}, \eqref{eq:lawinverse} & AR 203 & \path{System.prop_54}\newline\emph{SI: history expansion}.\\
(18) & \eqref{eq:bass} & AR 203 & \path{System.prop_54}, \path{Form.law_bass_form}\\
(19) & \eqref{eq:onekernel} & AR 203 & \path{Form.law_gap}, \path{Form.kN_gap}\\
(20) & \eqref{eq:onekernel} & AR 9, 204 & \path{Form.lawPoly_eq}, \path{System.prop_54}\\
(21) & \eqref{eq:encounterfactor}, \eqref{eq:encounterproduct} & AR 2--8, 204 & \path{Walk.pattern_factor}, \path{Walk.singlet_factor}\\
(22) & Sec.~\ref{sec:step-two} & AR 198 & \path{System.prop_27}\\
(23) & \eqref{eq:commonphase} & AR 198 & \path{Form.halfStep_carries_clock}, \path{Form.quaternion_triple}\\
(24) & \eqref{eq:nullquotient}, \eqref{eq:localcone} & AR 201 & \path{GroupsContinuumLaw.normalized_ray_iff}, \path{GroupsContinuumLaw.frame_orthonormal}\newline\emph{SI: tangent-metric sign and units}.\\
(25) & \eqref{eq:verticalparallel} & AR 197 & \path{System.prop_33}, \path{System.prop_36}\\
(26) & \eqref{eq:spinreturn} & AR 197 & \path{System.prop_34}\\
(27) & \eqref{eq:verticalfourspinor} & AR 194 & \path{System.prop_37}\\
(28) & \eqref{eq:support}, \eqref{eq:fourquarters} & AR 205, 290 & \path{System.prop_35}\\
(29) & \eqref{eq:fourquarters} & AR 19, 205 & \path{Form.chi_quarters}\\
(30) & \eqref{eq:returns}, \eqref{eq:signedjointreturn} & AR 20, 215 & \path{System.prop_30}, \path{Num.Phi_closed}\\
(31) & \eqref{eq:support} & AR 290--292 & \path{System.prop_44}, \path{System.prop_69}\\
(32) & \eqref{eq:stepreflection}, \eqref{eq:stepnorm}, \eqref{eq:scattering} & AR 21, 206 & \path{System.prop_28}, \path{Form.transmission_echo}, \path{Form.kN_plane_gap}\\
(33) & \eqref{eq:returns}, \eqref{eq:jointreturn}, \eqref{eq:visitmoments} & AR 214, 216 & \path{System.prop_29}\\
(34) & \eqref{eq:returns} & AR 22, 214, 219 & \path{System.prop_29}, \path{Num.a_eq}\\
(35) & \eqref{eq:unitaryclocksharing} & AR 198, 208 & \path{clock_turns_singlet}, \path{Form.one_kron_expI}\ensuremath{\gamma}\path{_singlet}\\
(36) & \eqref{eq:finiteclockunitary}, \eqref{eq:finiteclockdispersion} & AR 208 & \path{System.prop_60}\\
(37) & \eqref{eq:packetclock}, \eqref{eq:clockexpansion}, \eqref{eq:clocklimit} & AR 208 & \path{System.prop_60}\newline\emph{SI: physical units}.\\
(38) & \eqref{eq:exactgap}, \eqref{eq:finite_mass_conversion} & AR 207 & \path{Form.schur_symbol}, \path{System.prop_59}\\
(39) & \eqref{eq:clocklimit}, \eqref{eq:continuumgreen} & AR 207,208 & \path{System.prop_60}, \path{Article.dirac_mass_shell}\newline\emph{SI: Lorentz continuation}.\\
(40) & \eqref{eq:continuumgreen} & AR 209 & \path{Form.schur_symbol}\newline\emph{SI: Fourier integral}.\\
(41) & \eqref{eq:schur}, \eqref{eq:schurderivative}, \eqref{eq:twoinsertions} & AR 203 & \path{Law.schur_source_deriv}, \path{Law.schur_source_mixed}\\
(42) & \eqref{eq:fullrestform}, \eqref{eq:fullstatenorm}, \eqref{eq:positiveclock} & AR 213 & \path{Law.schur_source_deriv}\newline\emph{SI: complete clock norm}.\\
(43) & \eqref{eq:residue}, \eqref{eq:sharedamplitude} & AR 213 & \path{Law.residue_at_simple_zero}, \path{Law.zero_shift_pencil}\\
(44) & Sec.~\ref{sec:common}, \eqref{eq:sourceoperations} & AR 288 & \path{System.prop_58}, \path{bornM_trace}\\
(45) & \eqref{eq:positivemassform}, Sec.~\ref{sec:step-five} & AR 210 & \path{System.prop_39}\newline\emph{SI: continuum derivative}.\\
(46) & \eqref{eq:curlsi}, \eqref{eq:orientedtwoform} & AR 210 & \path{System.prop_38}\\
(47) & \eqref{eq:tenclifford}, \eqref{eq:paircomplexstructure} & AR 199 & \path{System.prop_45}\\
(48) & \eqref{eq:electromagneticcentre} & AR 199 & \path{System.prop_45}, \path{System.S3_3}\\
(49) & \eqref{eq:retainedschur}, \eqref{eq:tracechain}, \eqref{eq:stationaryweights}, \eqref{eq:EH} & AR 23,24, 217 & \path{System.prop_31}, \path{Num.EH_eq}\\
(50) & \eqref{eq:green} & AR 25, 218 & \path{Num.b_eq}, \path{System.S3_11}\\
(51) & \eqref{eq:fnu} & AR 26, 223, 291 & \path{Num.f_eq}, \path{Num.nu_eq}\\
(52) & \eqref{eq:clockmarks}, \eqref{eq:MQ}, \eqref{eq:h} & AR 27, 224, 293 & \path{System.eq_33}, \path{System.prop_206}\\
(53) & \eqref{eq:Nh} & AR 28, 225 & \path{System.eq_34}, \path{System.prop_72}\\
(54) & \eqref{eq:electricblock}, \eqref{eq:alpha} & AR 29, 226, 230 & \path{System.eq_32}, \path{System.row_D02}\\
(55) & \eqref{eq:gaugecomponents}, \eqref{eq:canonicalgauge}, \eqref{eq:proof-SCa-8} & AR 210 & \path{System.prop_62}, \path{Article.gauge_stationary_iff}\newline\emph{SI: continuum Maxwell variation}.\\
(56) & \eqref{eq:crossing}, \eqref{eq:crossingchi}, \eqref{eq:alphas} & AR 30, 220, 228 & \path{Whole.mixed_moment_central_subspaces}, \path{Num.alphaS_sq}\newline\emph{SI: singlet-extended clock}.\\
(57) & \eqref{eq:canonicalgauge}, \eqref{eq:alphas} & AR 210, 274 & \path{System.prop_62}, \path{Article.gauge_stationary_iff}\newline\emph{SI: continuum Yang--Mills variation}.\\
(58) & \eqref{eq:positiveconnection}, \eqref{eq:thetasectors}, \eqref{eq:unitsectornumber} & AR 272 & \path{System.prop_54}, \path{Series.theta_bar_series}\newline\emph{SI: positive unit sector}.\\
(59) & \eqref{eq:restpencil}, \eqref{eq:chargednecessity}, \eqref{eq:endpointclockhalf} & AR 264 & \path{Law.schur_source_deriv}\newline\emph{SI: rest pencil and equal shares}.\\
(60) & \eqref{eq:marked}, \eqref{eq:markedsourceclock}, \eqref{eq:meunit} & AR 31, 222, 236 & \path{Form.one_arc_explicit}, \path{System.prop_203}\\
(61) & \eqref{eq:protonsourcecount}, \eqref{eq:completeprotonsource}, \eqref{eq:proton} & AR 32, 237, 294 & \path{Num.protonPole_closed}, \path{System.row_B12}\\
(62) & \eqref{eq:neutronsourcefunctional}, \eqref{eq:neutroncontrasts}, \eqref{eq:neutron} & AR 33, 239, 295 & \path{Num.neutronExponent_closed}, \path{System.row_B13}\\
(63) & \eqref{eq:pairingcompression}, \eqref{eq:gradesourcefunctional}, \eqref{eq:layers}, \eqref{eq:threecumulativegrades} & AR 34, 241, 292 & \path{Fam.pairing_response}, \path{System.prop_69}\newline\emph{SI: completed family source}.\\
(64) & \eqref{eq:chargedsource}, \eqref{eq:leptons} & AR 35, 242 & \path{System.row_B01}, \path{System.row_B02}\newline\emph{SI: diagonal source assembly}.\\
(65) & \eqref{eq:quarks}, \eqref{eq:chargedsource} & AR 36,37, 244, 246 & \path{System.row_B04}, \path{System.row_B05}, \path{System.row_B09}, \path{System.row_B10}, \path{System.row_B11}\newline\emph{SI: diagonal source assembly}.\\
(66) & \eqref{eq:ladderinvariant} & AR 38, 244, 246 & \path{System.row_B06}, \path{System.row_B07}\\
(67) & \eqref{eq:weakcontactmatrix}, \eqref{eq:GF} & AR 39--41, 227, 232, 248 & \path{Comp.weak_higgs_top}, \path{System.row_D03}, \path{System.row_D15}\\
(68) & \eqref{eq:Nh}, \eqref{eq:Ncon}, \eqref{eq:realvector} & AR 42, 250 & \path{System.row_D08}, \path{System.row_D10}, \path{System.row_D09}\newline\emph{SI: kinetic normalization}.\\
(69) & \eqref{eq:weak_coefficients}, \eqref{eq:weakpencil}, \eqref{eq:actionpolecoefficients} & AR 43, 252 & \path{System.row_D11}, \path{System.row_D12}\newline\emph{SI: kinetic normalization}.\\
(70) & \eqref{eq:weakpencil}, \eqref{eq:weakmasses} & AR 44, 200, 252 & \path{System.prop_47}, \path{System.row_D11}, \path{System.row_D12}, \path{System.row_D13}\\
(71) & \eqref{eq:radialsusceptibility}, \eqref{eq:radialfermionsource}, \eqref{eq:radial} & AR 45,46, 219, 254, 256 & \path{Comp.weak_higgs_top}, \path{System.row_D16}, \path{System.row_D17}\\
(72) & \eqref{eq:incoming}, \eqref{eq:typicallogtransfer}, \eqref{eq:W} & AR 47, 234, 296 & \path{Form.W_minor}, \path{Num.W_closed}\newline\emph{SI: stationary incoming source}.\\
(73) & \eqref{eq:proof-SE-1}, \eqref{eq:proof-SE-2}, \eqref{eq:proof-SE-3} & AR 212 & \path{Law.curvature_jet}, \path{Law.trace_four_form}\\
(74) & \eqref{eq:proof-SE-4}, \eqref{eq:proof-SE-3} & AR 48, 212 & \path{System.eq_29}, \path{Law.trace_four_form}\newline\emph{SI: continuum action}.\\
(75) & \eqref{eq:proof-SE-5}, \eqref{eq:proof-SE-6}, \eqref{eq:proof-SE-7} & AR 212, 288 & \path{Law.curvature_jet}, \path{Law.trace_four_form}\newline\emph{SI: continuum variation}.\\
(76) & \eqref{eq:gravitystrength} & AR 49--51, 234 & \path{System.row_D04}, \path{System.row_D05}\newline\emph{electron calibration}.\\
(77) & \eqref{eq:weakfieldgravity}, \eqref{eq:continuumgreen} & AR 209, 288 & \path{Law.trace_four_form}, \path{Form.schur_symbol}\newline\emph{SI: weak-field limit}.\\
(78) & \eqref{eq:lapse}, \eqref{eq:fullclock} & AR 213 & \path{Form.lapse_partials}, \path{Form.lapse_displacement}\\
(79) & \eqref{eq:neutralalgebra}, \eqref{eq:neutralbasis} & AR 262 & \path{Neu.finrank_kept}, \path{Neu.kept_su5}\\
(80) & \eqref{eq:coherentcomparison}, \eqref{eq:retainedcomparison}, \eqref{eq:neutralpullback} & AR 262 & \path{Neu.Dcoh_eq}, \path{Neu.Dcoh_mul}, \path{Neu.keptH_eq}\newline\emph{SI: source normalization}.\\
(81) & \eqref{eq:neutralgram}, \eqref{eq:neutraldeterminant}, \eqref{eq:neutralthresholds} & AR 52, 262 & \path{Neu.retained_det}, \path{System.row_T01}\\
(82) & \eqref{eq:neutralhalfphase}, \eqref{eq:neutralblock}, \eqref{eq:neutralfull} & AR 53,54, 262 & \path{ArticleStatements.Radar.radar_quarter}, \path{Article.neutrino_seesaw}\\
(83) & \eqref{eq:neutralroots}, \eqref{eq:neutrinovalues} & AR 55, 262 & \path{Article.neutrino_seesaw}, \path{Article.neutrinos}\\
(84) & \eqref{eq:V}, \eqref{eq:U} & AR 258 & \path{System.row_C01}, \path{System.row_C02}, \path{System.row_C03}\newline\emph{rotation definitions}.\\
(85) & \eqref{eq:V} & AR 56, 258 & \path{System.row_C01}, \path{System.row_C02}, \path{System.row_C03}, \path{System.row_C04}\\
(86) & \eqref{eq:mixingcharacters}, \eqref{eq:mixingcharactervalues}, \eqref{eq:tiltidentity} & AR 13, 57, 258, 282 & \path{System.prop_73}, \path{Num.ns_eq}\\
(87) & \eqref{eq:V}, \eqref{eq:U}, \eqref{eq:taurow} & AR 69, 258, 260 & \path{System.row_C05}, \path{System.row_C10}, \path{System.row_C11}, \path{System.row_C13}, \path{System.row_C14}, \path{System.row_T06}\newline\emph{SI: rotation products}.\\
(88) & \eqref{eq:whole_source}, \eqref{eq:curvaturewhole} & AR 58, 284 & \path{Num.tau_eq_q}, \path{Num.s_eq}\\
(89) & \eqref{eq:curvaturewhole} & AR 59,60, 284 & \path{System.row_D06}, \path{System.row_T03}\newline\emph{SI: scale conversion}.\\
(90) & \eqref{eq:connectedreturn}, \eqref{eq:positivevolume}, \eqref{eq:singlepopulation} & AR 61, 280 & \path{Num.excursionF_eq_GK}, \path{Num.sourceC_eq}, \path{Num.positive_states}\\
(91) & \eqref{eq:threeandtwelve}, \eqref{eq:primordial} & AR 62--65, 280, 282 & \path{Num.sourceC_deriv}, \path{Num.sourceC_second}, \path{System.prop_127}\newline\emph{SI: spectral observables}.\\
(92) & \eqref{eq:neutrinotemperatureratio}, \eqref{eq:epoch} & AR 66, 286 & \path{System.row_F04}, \path{System.prop_127}\\
(93) & \eqref{eq:coldnative}, Sec.~\ref{sec:evolution} & AR 67, 286 & \path{Num.a_eq}, \path{System.row_T01}, \path{Num.W_closed}\newline\emph{recorded baryonic conversion}.\\
(94) & \eqref{eq:friedmann}, \eqref{eq:proof-SE-8} & AR 286 & \path{Law.curvature_jet}, \path{Law.trace_four_form}\newline\emph{SI: Friedmann evolution}.\\
(95) & \eqref{eq:epochcount}, \eqref{eq:epoch} & AR 68, 286 & \path{System.prop_127}, \path{Article.one_surface}\newline\emph{SI: epoch conversion}.\\
(96) & \eqref{eq:completeaction}, Table~\ref{tab:actioncoefficients}, \eqref{eq:actiongaugevariation}, \eqref{eq:actionscalarvariation} & AR 212, 228, 230, 232, 234, 242, 244, 246, 250, 252, 254, 256, 258, 260, 262, 272, 284 & \path{Article.local_form}, \path{Law.curvature_jet}, \path{Law.trace_four_form}, \path{Comp.weak_higgs_top}\newline\emph{SI: common action and variation}.\\
(A.1) & \eqref{eq:explicitclifford} & AR 194 & \path{System.prop_4}, \path{System.Geometry.J_eq_one_kron_i}\ensuremath{\sigma}\path{y}\\
(A.2) & \eqref{eq:sideoverlap} & AR 194 & \path{System.prop_6}\\
(A.3) & \eqref{eq:crossingidentity} & AR 10, 195 & \path{partnerTwo_kron}, \path{one_kron_line_singlet}\\
(A.4) & \eqref{eq:unitaryclocksharing} & AR 198 & \path{clock_turns_singlet}, \path{Form.one_kron_expI}\ensuremath{\gamma}\path{_singlet}\\
(A.5) & \eqref{eq:paircomparison}, \eqref{eq:returnvertex} & AR 195 & \path{System.prop_10}\\
(A.6) & \eqref{eq:pairmarginal} & AR 195 & \path{Article.partial_trace_singlet}\newline\emph{SI: arbitrary local unitary}.\\
(B.1) & \eqref{eq:bass} & AR 202 & \path{System.prop_54}\\
(B.2) & \eqref{eq:bass} & AR 203 & \path{System.prop_54}\\
(B.3) & \eqref{eq:onekernel} & AR 204 & \path{System.prop_54}, \path{Form.lawPoly_eq}\\
(B.4) & \eqref{eq:signeddet} & AR 215 & \path{System.prop_54}\\
(B.5) & \eqref{eq:lawinverse} & AR 203 & \path{System.prop_55}\\
(B.6) & \eqref{eq:returns}, \eqref{eq:signedjointreturn} & AR 214,215 & \path{System.prop_55}\\
(B.7) & \eqref{eq:returns}, \eqref{eq:signedjointreturn} & AR 214,215 & \path{System.prop_29}, \path{System.prop_30}\\
(B.8) & \eqref{eq:retainedschur}, \eqref{eq:tracechain} & AR 217 & \path{System.prop_31}, \path{schur_det}\\
(B.9) & \eqref{eq:marked}, \eqref{eq:markedsourceclock} & AR 222 & \path{Form.one_arc_explicit}, \path{System.prop_203}\\
(B.10) & \eqref{eq:green} & AR 218 & \path{System.S3_11}\\
(B.11) & \eqref{eq:MQ}, \eqref{eq:h} & AR 224, 293 & \path{System.eq_33}, \path{System.prop_206}\\
(B.12) & \eqref{eq:Nh} & AR 225 & \path{System.eq_34}, \path{System.prop_72}\\
(B.13) & \eqref{eq:electricblock} & AR 226 & \path{System.eq_32}, \path{System.row_D02}\\
(C.1) & \eqref{eq:schur} & AR 203 & \path{schur_LKR}, \path{schur_inv}\\
(C.2) & \eqref{eq:schur}, \eqref{eq:schurderivative}, \eqref{eq:twoinsertions} & AR 203 & \path{schur_inv}, \path{Law.schur_source_deriv}, \path{Law.schur_source_mixed}\\
(C.3) & \eqref{eq:residue}, \eqref{eq:sharedamplitude} & AR 213 & \path{Law.residue_at_simple_zero}, \path{Law.zero_shift_pencil}\\
(C.4) & \eqref{eq:fullstatenorm}, \eqref{eq:fullrestform} & AR 213 & \path{Law.schur_source_deriv}\newline\emph{SI: complete clock norm}.\\
(C.5) & \eqref{eq:selfenergyderivative} & AR 213 & \path{Law.schur_source_deriv}\newline\emph{SI: endpoint derivative terms}.\\
(C.6) & \eqref{eq:endpointclockhalf} & AR 264 & \path{Law.schur_source_deriv}\newline\emph{SI: equal chiral shares}.\\
(C.7) & \eqref{eq:groundedpetersen} & AR 218, 239 & \path{Fam.rooted_plane_block}, \path{Fam.GammaP_mul_inv}\\
(C.8) & \eqref{eq:pairingisometry}, \eqref{eq:pairingcompression} & AR 241, 292 & \path{Fam.pairCols_gram}, \path{Fam.pairing_response}\\
(D.1) & \eqref{eq:positivemassform} & AR 210 & \path{System.prop_39}\newline\emph{SI: off-diagonal expansion}.\\
(D.2) & \eqref{eq:circulationperturbation}, \eqref{eq:circulationperron} & AR 211 & \path{System.S3_17}\\
(D.3) & \eqref{eq:circulation48}, \eqref{eq:circulationbound} & AR 211 & \path{System.eq_35}, \path{System.S3_17}\\
(D.4) & \eqref{eq:proof-SCa-3}, \eqref{eq:proof-SCa-4} & AR 210 & \path{System.prop_62}\newline\emph{SI: trace cancellation}.\\
(D.5) & \eqref{eq:proof-SCa-8}, \eqref{eq:proof-SCa-9} & AR 210 & \path{System.prop_62}, \path{Form.kappa_gap}\\
(D.6) & \eqref{eq:exactgap}, \eqref{eq:continuumgreen} & AR 207 & \path{Form.schur_symbol}\\
(D.7) & \eqref{eq:exactgap}, \eqref{eq:finite_mass_conversion} & AR 207 & \path{Form.schur_symbol}\\
(D.8) & \eqref{eq:exactgap}, \eqref{eq:finite_mass_conversion} & AR 207 & \path{Form.chi_mass_shell}, \path{System.prop_59}\\
(D.9) & \eqref{eq:continuumgreen} & AR 209 & \path{Form.schur_symbol}\newline\emph{SI: Fourier integral}.\\
(D.10) & \eqref{eq:proof-SE-1}, \eqref{eq:proof-SE-2} & AR 212 & \path{Law.curvature_jet}\\
(D.11) & \eqref{eq:proof-SE-3} & AR 212 & \path{Law.trace_four_form}, \path{System.prop_61}\\
(D.12) & \eqref{eq:proof-SE-3}, \eqref{eq:proof-SE-4} & AR 212 & \path{System.eq_29}, \path{Law.trace_four_form}\newline\emph{SI: curvature contraction}.\\
(D.13) & \eqref{eq:proof-SE-6}, \eqref{eq:proof-SE-7} & AR 212 & \path{Law.curvature_jet}, \path{Law.trace_four_form}\newline\emph{SI: continuum variation}.\\
(E.1) & \eqref{eq:coherentcomparison} & AR 262 & \path{Neu.Dcoh_eq}\\
(E.2) & \eqref{eq:neutralalgebra}, \eqref{eq:neutralbasis} & AR 262 & \path{Neu.finrank_kept}, \path{Neu.kept_su5}\\
(E.3) & \eqref{eq:neutralpullback} & AR 262 & \path{Neu.Dcoh_mul}, \path{Neu.keptH_eq}\\
(E.4) & \eqref{eq:neutralgram} & AR 262 & \path{Neu.keptN_gram}, \path{Neu.keptH_eq}\newline\emph{SI: Gram determinant}.\\
(E.5) & \eqref{eq:neutraldeterminant} & AR 262 & \path{Neu.retained_det}, \path{System.row_T01}\\
(E.6) & \eqref{eq:classsolutions}, \eqref{eq:jointreturn} & AR 216 & \path{System.prop_29}\newline\emph{SI: class elimination}.\\
(E.7) & \eqref{eq:jointreturn} & AR 216 & \path{Num.excursionF_vector}, \path{System.prop_29}\\
(E.8) & \eqref{eq:connectedreturn} & AR 216, 280 & \path{Num.excursionF_eq_GK}, \path{Num.GK_law}\\
(E.9) & \eqref{eq:singlepopulation} & AR 280 & \path{Num.sourceC_eq}, \path{Num.positive_states}\\
(E.10) & \eqref{eq:positivestatisticalsource}, \eqref{eq:singlepopulation} & AR 14,15, 280 & \path{Num.sourceC_eq}, \path{Num.GK_law}\newline\emph{SI: positive coefficient expansion}.\\
(E.11) & \eqref{eq:threeandtwelve} & AR 15, 280, 282 & \path{Num.sourceC_deriv}, \path{Num.sourceC_second}\\
(E.12) & \eqref{eq:protonpairgram} & AR 237, 294 & \path{Article.gramTheta_eq}\\
(E.13) & \eqref{eq:gramheight}, \eqref{eq:protonpairgram} & AR 294 & \path{Form.detF0_is_transmission}\newline\emph{SI: all three eigenvalues}.\\
(E.14) & \eqref{eq:positiveconnection} & AR 272 & \path{System.prop_54}\\
(E.15) & \eqref{eq:thetasectors} & AR 272 & \path{Series.theta_bar_series}\newline\emph{positive symmetric sector law}.\\
(E.16) & \eqref{eq:projectorcurvature}, \eqref{eq:unitsectorcurvature}, \eqref{eq:unitsectornumber} & AR 11,12, 194, 212 & \path{System.prop_4}, \path{System.prop_5}, \path{Law.trace_four_form}\newline\emph{SI: sphere integral}.\\
(E.17) & \eqref{eq:thetasectors}, \eqref{eq:unitsectornumber} & AR 272 & \path{Series.theta_bar_series}\newline\emph{SI: positive unit sector}.\\
\end{longtable}
\begin{longtable}{@{}>{\raggedright\arraybackslash}p{.57in}>{\raggedright\arraybackslash}p{1.16in}>{\raggedright\arraybackslash}p{.75in}>{\raggedright\arraybackslash}p{4.05in}@{}}
\caption{Every row of article Table I.}\label{tab:articletableone}\\
\toprule Article & Supplemental Material & Printed line & Lean statement and SI calculation\\
\midrule\endfirsthead
\multicolumn{4}{l}{\textit{Every row of article Table I (continued)}}\\
\toprule Article & Supplemental Material & Printed line & Lean statement and SI calculation\\
\midrule\endhead
\bottomrule\endfoot
\path{T1.01} & \eqref{eq:gravitystrength} & AR 91 & \path{System.row_D04}\\
\path{T1.02} & \eqref{eq:gravitystrength} & AR 92 & \path{System.row_D05}\\
\path{T1.03} & \eqref{eq:proof-SE-4}, \eqref{eq:proof-SE-3} & AR 93 & \path{System.eq_29}, \path{Law.trace_four_form}\\
\path{T1.04} & \eqref{eq:curvaturewhole} & AR 94 & \path{System.row_D06}\\
\path{T1.05} & \eqref{eq:crossing}, \eqref{eq:crossingchi}, \eqref{eq:alphas} & AR 95 & \path{System.row_D01}\\
\path{T1.06} & \eqref{eq:positiveconnection}, \eqref{eq:thetasectors}, \eqref{eq:unitsectornumber} & AR 96 & \path{System.row_K01}\\
\path{T1.07} & \eqref{eq:electricblock}, \eqref{eq:alpha} & AR 97 & \path{System.row_D02}\\
\path{T1.08} & \eqref{eq:weak_coefficients}, \eqref{eq:weakpencil}, \eqref{eq:actionpolecoefficients} & AR 98 & \path{System.row_D11}, \path{System.row_D12}\\
\path{T1.09} & \eqref{eq:weak_coefficients}, \eqref{eq:weakpencil}, \eqref{eq:actionpolecoefficients} & AR 99 & \path{System.row_D11}, \path{System.row_D12}\\
\path{T1.10} & \eqref{eq:weakcontactmatrix}, \eqref{eq:GF} & AR 100 & \path{System.row_D03}\\
\path{T1.11} & \eqref{eq:weakcontactmatrix}, \eqref{eq:GF} & AR 101 & \path{System.row_D15}\\
\path{T1.12} & \eqref{eq:Nh}, \eqref{eq:Ncon}, \eqref{eq:realvector} & AR 102 & \path{System.row_D08}\\
\path{T1.13} & \eqref{eq:Nh}, \eqref{eq:Ncon}, \eqref{eq:realvector} & AR 103 & \path{System.row_D10}\\
\path{T1.14} & \eqref{eq:Nh}, \eqref{eq:Ncon}, \eqref{eq:realvector} & AR 104 & \path{System.row_D09}\\
\path{T1.15} & \eqref{eq:fnu} & AR 105 & \path{System.row_D14}\\
\path{T1.16} & \eqref{eq:weakpencil}, \eqref{eq:weakmasses} & AR 106 & \path{System.row_D11}\\
\path{T1.17} & \eqref{eq:weakpencil}, \eqref{eq:weakmasses} & AR 107 & \path{System.row_D12}\\
\path{T1.18} & \eqref{eq:weakpencil}, \eqref{eq:weakmasses} & AR 108 & \path{System.row_D13}\\
\path{T1.19} & \eqref{eq:radialsusceptibility}, \eqref{eq:radialfermionsource}, \eqref{eq:radial} & AR 109 & \path{Comp.weak_higgs_top}, \path{System.row_D16}, \path{System.row_D17}\\
\path{T1.20} & \eqref{eq:radialsusceptibility}, \eqref{eq:radialfermionsource}, \eqref{eq:radial} & AR 110 & \path{System.row_D16}\\
\path{T1.21} & \eqref{eq:radialsusceptibility}, \eqref{eq:radialfermionsource}, \eqref{eq:radial} & AR 111 & \path{Comp.weak_higgs_top}, \path{System.row_D16}, \path{System.row_D17}\\
\path{T1.22} & \eqref{eq:radialsusceptibility}, \eqref{eq:radialfermionsource}, \eqref{eq:radial} & AR 112 & \path{System.row_D17}\\
\path{T1.23} & \eqref{eq:protonsourcecount}, \eqref{eq:completeprotonsource}, \eqref{eq:proton} & AR 113 & \path{System.row_B12}\\
\path{T1.24} & \eqref{eq:neutronsourcefunctional}, \eqref{eq:neutroncontrasts}, \eqref{eq:neutron} & AR 114 & \path{System.row_B13}\\
\path{T1.25} & \eqref{eq:marked}, \eqref{eq:markedsourceclock}, \eqref{eq:meunit} & AR 115 & \path{Form.one_arc_explicit}, \path{System.prop_203}\newline\emph{Electron calibration.}\\
\path{T1.26} & \eqref{eq:neutralgram}, \eqref{eq:neutraldeterminant}, \eqref{eq:neutralthresholds} & AR 116 & \path{System.row_T01}\\
\path{T1.27} & \eqref{eq:threeandtwelve}, \eqref{eq:primordial} & AR 117 & \path{System.row_F01}\\
\path{T1.28} & \eqref{eq:threeandtwelve}, \eqref{eq:primordial} & AR 118, 125 & \path{System.row_F03}\\
\path{T1.29} & \eqref{eq:threeandtwelve}, \eqref{eq:primordial} & AR 119 & \path{System.row_C15}\\
\path{T1.30} & \eqref{eq:neutrinotemperatureratio}, \eqref{eq:epoch} & AR 120 & \path{System.row_F04}\\
\path{T1.31} & \eqref{eq:coldnative}, Sec.~\ref{sec:evolution} & AR 121 & \path{Num.a_eq}, \path{System.row_T01}, \path{Num.W_closed}\newline\emph{Program: $a^{16}M_R/m_p$.}\\
\path{T1.32} & \eqref{eq:coldnative}, Sec.~\ref{sec:evolution} & AR 122 & \path{Num.a_eq}, \path{System.row_T01}, \path{Num.W_closed}\newline\emph{Recorded baryonic conversion.}\\
\path{T1.33} & \eqref{eq:epochcount}, \eqref{eq:epoch} & AR 123 & \path{System.row_F13}\\
\path{T1.34} & \eqref{eq:epochcount}, \eqref{eq:epoch} & AR 124 & \path{System.prop_127}, \path{Article.one_surface}\\
\end{longtable}
\begin{longtable}{@{}>{\raggedright\arraybackslash}p{.57in}>{\raggedright\arraybackslash}p{1.16in}>{\raggedright\arraybackslash}p{.75in}>{\raggedright\arraybackslash}p{4.05in}@{}}
\caption{Every displayed entry of article Table II.}\label{tab:articletabletwo}\\
\toprule Article & Supplemental Material & Printed line & Lean statement and SI calculation\\
\midrule\endfirsthead
\multicolumn{4}{l}{\textit{Every displayed entry of article Table II (continued)}}\\
\toprule Article & Supplemental Material & Printed line & Lean statement and SI calculation\\
\midrule\endhead
\bottomrule\endfoot
\path{T2.F1.nu} & \eqref{eq:neutralroots}, \eqref{eq:neutrinovalues} & AR 126 & \path{Article.neutrino_seesaw}, \path{Article.neutrinos}\\
\path{T2.F1.e} & \eqref{eq:marked}, \eqref{eq:markedsourceclock}, \eqref{eq:meunit} & AR 127 & \path{Form.one_arc_explicit}, \path{System.prop_203}\newline\emph{Electron calibration.}\\
\path{T2.F1.u.native} & \eqref{eq:quarks}, \eqref{eq:chargedsource} & AR 128 & \path{System.row_B09}\newline\emph{Calibrated mass; finite ratio in Lean.}\\
\path{T2.F1.u.carried} & \eqref{eq:quarks}, \eqref{eq:chargedsource}, Sec.~\ref{sec:evolution}, Sec.~\ref{sec:quarkconversions} & AR 129--131 & \path{System.row_B04}, \path{System.row_B05}, \path{System.row_B09}, \path{System.row_B10}, \path{System.row_B11}\newline\emph{CRunDec conversion; finite inputs in Lean.}\\
\path{T2.F1.d.native} & \eqref{eq:quarks}, \eqref{eq:chargedsource} & AR 132 & \path{System.row_B04}, \path{System.row_B05}, \path{System.row_B09}, \path{System.row_B10}, \path{System.row_B11}\newline\emph{Calibrated mass; finite ratio in Lean.}\\
\path{T2.F1.d.carried} & \eqref{eq:quarks}, \eqref{eq:chargedsource}, Sec.~\ref{sec:evolution}, Sec.~\ref{sec:quarkconversions} & AR 133--135 & \path{System.row_B04}, \path{System.row_B05}, \path{System.row_B09}, \path{System.row_B10}, \path{System.row_B11}\newline\emph{CRunDec conversion; finite inputs in Lean.}\\
\path{T2.F2.nu.mass} & \eqref{eq:neutralroots}, \eqref{eq:neutrinovalues} & AR 136 & \path{Article.neutrino_seesaw}, \path{Article.neutrinos}\\
\path{T2.F2.nu.split} & \eqref{eq:neutralroots}, \eqref{eq:neutrinovalues} & AR 137 & \path{System.row_E01}\\
\path{T2.F2.mu.mass} & \eqref{eq:chargedsource}, \eqref{eq:leptons} & AR 138 & \path{System.row_B01}\newline\emph{Calibrated mass; finite ratio in Lean.}\\
\path{T2.F2.mu.ratio} & \eqref{eq:chargedsource}, \eqref{eq:leptons} & AR 139 & \path{System.row_B01}\\
\path{T2.F2.c.native} & \eqref{eq:quarks}, \eqref{eq:chargedsource} & AR 140 & \path{System.row_B11}\newline\emph{Calibrated mass; finite ratio in Lean.}\\
\path{T2.F2.c.carried} & \eqref{eq:quarks}, \eqref{eq:chargedsource}, Sec.~\ref{sec:evolution}, Sec.~\ref{sec:quarkconversions} & AR 141--143 & \path{System.row_B04}, \path{System.row_B05}, \path{System.row_B09}, \path{System.row_B10}, \path{System.row_B11}\newline\emph{CRunDec conversion; finite inputs in Lean.}\\
\path{T2.F2.s.native} & \eqref{eq:quarks}, \eqref{eq:chargedsource} & AR 144 & \path{System.row_B04}\newline\emph{Calibrated mass; finite ratio in Lean.}\\
\path{T2.F2.s.carried} & \eqref{eq:quarks}, \eqref{eq:chargedsource}, Sec.~\ref{sec:evolution}, Sec.~\ref{sec:quarkconversions} & AR 145--147 & \path{System.row_B04}, \path{System.row_B05}, \path{System.row_B09}, \path{System.row_B10}, \path{System.row_B11}\newline\emph{CRunDec conversion; finite inputs in Lean.}\\
\path{T2.F3.nu.mass} & \eqref{eq:neutralroots}, \eqref{eq:neutrinovalues} & AR 148 & \path{System.row_E06}\\
\path{T2.F3.nu.split} & \eqref{eq:neutralroots}, \eqref{eq:neutrinovalues} & AR 149 & \path{System.row_E02}\\
\path{T2.F3.tau.mass} & \eqref{eq:chargedsource}, \eqref{eq:leptons} & AR 150 & \path{System.row_L01}\newline\emph{Calibrated mass; finite ratio in Lean.}\\
\path{T2.F3.tau.ratio} & \eqref{eq:chargedsource}, \eqref{eq:leptons} & AR 151 & \path{System.row_B02}\\
\path{T2.F3.t.native} & \eqref{eq:quarks}, \eqref{eq:chargedsource} & AR 152 & \path{System.row_B07}\newline\emph{Calibrated mass; finite ratio in Lean.}\\
\path{T2.F3.t.carried} & \eqref{eq:quarks}, \eqref{eq:chargedsource}, Sec.~\ref{sec:evolution}, Sec.~\ref{sec:quarkconversions} & AR 153--155 & \path{System.row_B04}, \path{System.row_B05}, \path{System.row_B09}, \path{System.row_B10}, \path{System.row_B11}\newline\emph{CRunDec conversion; finite inputs in Lean.}\\
\path{T2.F3.b.native} & \eqref{eq:quarks}, \eqref{eq:chargedsource} & AR 156 & \path{System.row_B05}\newline\emph{Calibrated mass; finite ratio in Lean.}\\
\path{T2.F3.b.carried} & \eqref{eq:quarks}, \eqref{eq:chargedsource}, Sec.~\ref{sec:evolution}, Sec.~\ref{sec:quarkconversions} & AR 157--159 & \path{System.row_B04}, \path{System.row_B05}, \path{System.row_B09}, \path{System.row_B10}, \path{System.row_B11}\newline\emph{CRunDec conversion; finite inputs in Lean.}\\
\path{T2.V11} & \eqref{eq:V} & AR 160 & \path{System.row_C01}, \path{System.row_C02}, \path{System.row_C03}, \path{System.row_C04}\newline\emph{Program: rotation-product entry.}\\
\path{T2.V12} & \eqref{eq:V} & AR 161 & \path{System.row_C01}, \path{System.row_C02}, \path{System.row_C03}, \path{System.row_C04}\newline\emph{Program: rotation-product entry.}\\
\path{T2.V13} & \eqref{eq:V} & AR 162 & \path{System.row_C01}, \path{System.row_C02}, \path{System.row_C03}, \path{System.row_C04}\newline\emph{Program: rotation-product entry.}\\
\path{T2.V21} & \eqref{eq:V} & AR 163 & \path{System.row_C01}, \path{System.row_C02}, \path{System.row_C03}, \path{System.row_C04}\newline\emph{Program: rotation-product entry.}\\
\path{T2.V22} & \eqref{eq:V} & AR 164 & \path{System.row_C01}, \path{System.row_C02}, \path{System.row_C03}, \path{System.row_C04}\newline\emph{Program: rotation-product entry.}\\
\path{T2.V23} & \eqref{eq:V} & AR 165 & \path{System.row_C01}, \path{System.row_C02}, \path{System.row_C03}, \path{System.row_C04}\newline\emph{Program: rotation-product entry.}\\
\path{T2.V31} & \eqref{eq:V} & AR 166 & \path{System.row_C01}, \path{System.row_C02}, \path{System.row_C03}, \path{System.row_C04}\newline\emph{Program: rotation-product entry.}\\
\path{T2.V32} & \eqref{eq:V} & AR 167 & \path{System.row_C01}, \path{System.row_C02}, \path{System.row_C03}, \path{System.row_C04}\newline\emph{Program: rotation-product entry.}\\
\path{T2.V33} & \eqref{eq:V} & AR 168 & \path{System.row_C01}, \path{System.row_C02}, \path{System.row_C03}, \path{System.row_C04}\newline\emph{Program: rotation-product entry.}\\
\path{T2.U11} & \eqref{eq:U} & AR 169 & \path{System.row_C10}, \path{System.row_C11}, \path{System.row_C13}, \path{System.row_C14}, \path{System.row_T06}\newline\emph{Program: rotation-product entry.}\\
\path{T2.U12} & \eqref{eq:U} & AR 170 & \path{System.row_C10}, \path{System.row_C11}, \path{System.row_C13}, \path{System.row_C14}, \path{System.row_T06}\newline\emph{Program: rotation-product entry.}\\
\path{T2.U13} & \eqref{eq:U} & AR 171 & \path{System.row_C10}, \path{System.row_C11}, \path{System.row_C13}, \path{System.row_C14}, \path{System.row_T06}\newline\emph{Program: rotation-product entry.}\\
\path{T2.U21} & \eqref{eq:U} & AR 172 & \path{System.row_C10}, \path{System.row_C11}, \path{System.row_C13}, \path{System.row_C14}, \path{System.row_T06}\newline\emph{Program: rotation-product entry.}\\
\path{T2.U22} & \eqref{eq:U} & AR 173 & \path{System.row_C10}, \path{System.row_C11}, \path{System.row_C13}, \path{System.row_C14}, \path{System.row_T06}\newline\emph{Program: rotation-product entry.}\\
\path{T2.U23} & \eqref{eq:U} & AR 174 & \path{System.row_C10}, \path{System.row_C11}, \path{System.row_C13}, \path{System.row_C14}, \path{System.row_T06}\newline\emph{Program: rotation-product entry.}\\
\path{T2.U31} & \eqref{eq:U} & AR 175 & \path{System.row_C10}, \path{System.row_C11}, \path{System.row_C13}, \path{System.row_C14}, \path{System.row_T06}\newline\emph{Program: rotation-product entry.}\\
\path{T2.U32} & \eqref{eq:U} & AR 176 & \path{System.row_C10}, \path{System.row_C11}, \path{System.row_C13}, \path{System.row_C14}, \path{System.row_T06}\newline\emph{Program: rotation-product entry.}\\
\path{T2.U33} & \eqref{eq:U} & AR 177 & \path{System.row_C10}, \path{System.row_C11}, \path{System.row_C13}, \path{System.row_C14}, \path{System.row_T06}\newline\emph{Program: rotation-product entry.}\\
\path{T2.CKM.s12} & \eqref{eq:V} & AR 178 & \path{System.row_C01}\\
\path{T2.CKM.s23} & \eqref{eq:V} & AR 179 & \path{System.row_C02}\\
\path{T2.CKM.s13} & \eqref{eq:V} & AR 180 & \path{System.row_C03}\\
\path{T2.CKM.delta} & \eqref{eq:V} & AR 181 & \path{System.row_C04}\\
\path{T2.CKM.J} & \eqref{eq:V}, \eqref{eq:U}, \eqref{eq:taurow} & AR 182 & \path{System.row_C05}\\
\path{T2.PMNS.s12} & \eqref{eq:V}, \eqref{eq:U}, \eqref{eq:taurow} & AR 183 & \path{System.row_C10}\\
\path{T2.PMNS.s13} & \eqref{eq:V}, \eqref{eq:U}, \eqref{eq:taurow} & AR 184 & \path{System.row_C11}\\
\path{T2.PMNS.s23} & \eqref{eq:V}, \eqref{eq:U}, \eqref{eq:taurow} & AR 185 & \path{System.row_C13}\\
\path{T2.PMNS.delta} & \eqref{eq:V}, \eqref{eq:U}, \eqref{eq:taurow} & AR 186 & \path{System.row_C14}\\
\path{T2.PMNS.J} & \eqref{eq:V}, \eqref{eq:U}, \eqref{eq:taurow} & AR 187 & \path{System.row_T06}\\
\end{longtable}
\endgroup


\newpage
\begin{thebibliography}{99}\raggedright
\setlength{\itemsep}{2pt}

\bibitem{Camphuis}
E. Camphuis et al. (SPT-3G Collaboration),
SPT-3G D1: CMB temperature and polarization power spectra and cosmology from
2019 and 2020 observations of the SPT-3G Main field,
Phys. Rev. D \textbf{113}, 083504 (2026),
\href{https://doi.org/10.1103/7wt3-9v2y}{doi:10.1103/7wt3-9v2y};
\url{https://arxiv.org/abs/2506.20707}.

\bibitem{JUNO2026}
JUNO Collaboration, Measurement of Neutrino Oscillation Parameters at JUNO
and Indication of the Neutrino Mass Ordering (2026),
\url{https://arxiv.org/abs/2609.37895}.

\bibitem{ACT2025}
T. Louis et al. (ACT Collaboration), The Atacama Cosmology Telescope:
DR6 Power Spectra, Likelihoods and $\Lambda$CDM Parameters,
JCAP \textbf{11} (2025), 062,
\href{https://doi.org/10.1088/1475-7516/2025/11/062}{doi:10.1088/1475-7516/2025/11/062}.

\bibitem{PDG2026}
F. Takahashi et al. (Particle Data Group), \emph{Review of Particle Physics},
Int. J. Mod. Phys. A \textbf{41}, 2630011 (2026),
gauge- and Higgs-boson and quark summary tables,
\url{https://pdg.lbl.gov/2026/tables/contents_tables.html};\
quark mass summary (90\% confidence intervals),\
\url{https://pdg.lbl.gov/2026/tables/rpp2026-sum-quarks.pdf}.

\bibitem{Ihara}
Y.~Ihara, On discrete subgroups of the two by two projective linear group over $p$-adic fields,
J. Math. Soc. Japan \textbf{18}, 219--235 (1966).

\bibitem{Bass}
H.~Bass, The Ihara--Selberg zeta function of a tree lattice,
Int. J. Math. \textbf{3}, 717--797 (1992).

\bibitem{Kac}
M.~Kac, On the notion of recurrence in discrete stochastic processes,
Bull. Amer. Math. Soc. \textbf{53}, 1002--1010 (1947).

\bibitem{ChauKeung}
L.-L.~Chau and W.-Y.~Keung, Comments on the parametrization of the Kobayashi--Maskawa matrix,
Phys. Rev. Lett. \textbf{53}, 1802--1805 (1984).

\bibitem{Jarlskog}
C.~Jarlskog, Commutator of the quark mass matrices in the standard electroweak model and a measure of maximal CP nonconservation,
Phys. Rev. Lett. \textbf{55}, 1039--1042 (1985).

\bibitem{MacDowellMansouri}
S. W. MacDowell and F. Mansouri, Unified geometric theory of gravity and supergravity,
Phys. Rev. Lett. \textbf{38}, 739--742 (1977);
erratum \textbf{38}, 1376 (1977),
\href{https://doi.org/10.1103/PhysRevLett.38.739}{doi:10.1103/PhysRevLett.38.739}.

\bibitem{Minkowski}
P.~Minkowski, $\mu\to e\gamma$ at a rate of one out of $10^9$ muon decays?,
Phys. Lett. B \textbf{67}, 421--428 (1977).

\bibitem{Procopio}
L.~M.~Procopio et al., Experimental superposition of orders of quantum gates,
Nat. Commun. \textbf{6}, 7913 (2015),
\url{https://doi.org/10.1038/ncomms8913}.

\bibitem{GibbonsHawking}
G.~W.~Gibbons and S.~W.~Hawking, Cosmological event horizons, thermodynamics, and particle creation,
Phys. Rev. D \textbf{15}, 2738--2751 (1977).

\bibitem{Beta4}
T.~van Ritbergen, J.~A.~M.~Vermaseren, and S.~A.~Larin,
The four-loop $\beta$-function in quantum chromodynamics,
Phys. Lett. B \textbf{400}, 379--384 (1997),
\url{https://arxiv.org/abs/hep-ph/9701390}.

\bibitem{CODATA}
P.~J.~Mohr, D.~B.~Newell, B.~N.~Taylor, and E.~Tiesinga,
CODATA recommended values of the fundamental physical constants: 2022,
Rev. Mod. Phys. \textbf{97}, 025002 (2025).
NIST 2022 reference chart and numerical listing:
\url{https://physics.nist.gov/cuu/pdf/wall_2022.pdf};
\url{https://physics.nist.gov/cuu/Constants/Table/allascii.txt}.

\bibitem{SI}
Bureau International des Poids et Mesures,
\emph{The International System of Units (SI)}, 9th ed.,
\url{https://www.bipm.org/en/publications/si-brochure}.

\bibitem{RunDec}
B. Schmidt and M. Steinhauser, CRunDec: a C++ package for running and decoupling of the strong coupling and quark masses,
Comput. Phys. Commun. \textbf{183}, 1845--1848 (2012),
\url{https://arxiv.org/abs/1201.6149}.

\bibitem{CRunDec}
F. Herren and M. Steinhauser, Version 3 of RunDec and CRunDec,
Comput. Phys. Commun. \textbf{224}, 333--345 (2018),
\url{https://arxiv.org/abs/1703.03751}.

\bibitem{NuFIT6}
I. Esteban et al., NuFit-6.0: Updated global analysis of three-flavor neutrino
oscillations, JHEP \textbf{12} (2024), 216,
\href{https://doi.org/10.1007/JHEP12(2024)216}{doi:10.1007/JHEP12(2024)216};
\url{https://arxiv.org/abs/2410.05380}, Table 1.

\bibitem{MuonG2Final}
D. P. Aguillard et al. (Muon $g-2$ Collaboration), Measurement of the Positive Muon Anomalous Magnetic Moment to 127 ppb,
Phys. Rev. Lett. \textbf{135}, 101802 (2025),
\href{https://doi.org/10.1103/7clf-sm2v}{doi:10.1103/7clf-sm2v};
\url{https://arxiv.org/abs/2506.03069}, Table II.
\bibitem{Repository}
\emph{Physics Without an Outside Observer}, release 1.0, \url{https://github.com/yongzic1-lgtm/physics-without-an-outside-observer}, accompanying programs, data, calculation outputs and Lean proofs; see \texttt{README.md}, \texttt{docs/NUMBERS.md}, \texttt{docs/PROOFS.md} and \path{docs/ARTICLE_CONCORDANCE.md}.
\end{thebibliography}

\end{document}
